% !TEX root = print.tex

\part{Domains, traces and manifolds}
\label{part:domains}

This part carries the function spaces from $\R^d$ to domains and manifolds.
\Cref{ch:lipschitz-domains} describes the Besov spaces of Lipschitz graph domains and of bounded
Lipschitz domains by intrinsic differences and as restrictions, and builds the surface measure
from the area formula; \Cref{ch:traces-lipschitz} proves the trace theorems on Lipschitz
boundaries for $0 < s < 1$ and on the half-space for every $s > 0$; \Cref{ch:sobolev-domains}
builds the Sobolev spaces on open sets as Hilbert spaces, with their duality, the trace space
$H^{1/2}(\Gamma)$, the extension theorem for $\Sob{1}{p}$, Green's formula and the kernel of the
trace; and \Cref{ch:manifolds} defines Besov spaces on compact manifolds by charts and proves the
trace theorem on their boundary.

\chapter{Lipschitz domains}
\label{ch:lipschitz-domains}

This chapter works on a finite-dimensional real inner product space $V$ of dimension $n$ with
Lebesgue measure, and with $L_p = L_p(V)$ for $1 \leq p < \infty$; the moduli of smoothness on
subsets, the local means and the pointwise multipliers are those of \Cref{ch:rd}. A Lipschitz
graph domain lies above the graph of a Lipschitz function, and a bounded Lipschitz domain agrees
near each boundary point with a Lipschitz graph domain. On Lipschitz graph domains the
restrictions of the Besov spaces of $V$ are described by intrinsic local means and by intrinsic
differences, with Rychkov's universal extension (\Cref{thm:graph-extension},
\Cref{thm:restriction-local-means} and \Cref{thm:graph-besov}); on bounded Lipschitz domains
they are described by intrinsic differences, by pasting local extensions
(\Cref{thm:local-extension} and \Cref{thm:bounded-lipschitz-besov}). The area formula for Lipschitz graphs
(\Cref{thm:lipschitz-graph-area}) expresses the Euclidean Hausdorff measure of the boundary in a
chart, and \Cref{sec:surface-measure} builds from it the surface measure, $L_2(\Gamma)$ and the
traces into it. \Cref{sec:local-lipschitz-graphs} shows that the bounded Lipschitz domains are
the sets of the usual definition by local Lipschitz graphs, with the open balls as examples, and
\Cref{sec:convex-domains} adds the bounded open convex sets and the open unit cube.

\section{Lipschitz graph domains}

\begin{definition}[Lipschitz graph domains]
\label{def:lipschitz-graph}
\lean{Rd.LipschitzGraph, Rd.LipschitzGraph.domain, Rd.LipschitzGraph.cone,
  Rd.LipschitzGraph.continuous_phi, Rd.LipschitzGraph.isOpen_domain,
  Rd.LipschitzGraph.measurableSet_domain}
\leanok
Let $e \in V$ be a unit vector, $M \geq 0$, and $\varphi : V \to \R$ with
$|\varphi(x) - \varphi(y)| \leq M \|x - y\|$ and $\varphi(x + t e) = \varphi(x)$ for all
$x, y \in V$ and $t \in \R$. The Lipschitz graph domain and its cone are
\[
  \Omega = \{x \in V : \varphi(x) < \langle x, e \rangle\}, \qquad
  K = \{y \in V : M \|y - \langle y, e \rangle e\| < \langle y, e \rangle\} .
\]
The set $\Omega$ is open.
\end{definition}

For $V = \R^d$ and $e = e_d$ this is $\Omega = \{x : x_d > \varphi(x_1, \dots, x_{d-1})\}$.

\begin{lemma}[The cone]
\label{lem:graph-cone}
\lean{Rd.LipschitzGraph.add_mem, Rd.LipschitzGraph.smul_mem_cone,
  Rd.LipschitzGraph.add_nsmul_mem, Rd.LipschitzGraph.coneRadius,
  Rd.LipschitzGraph.coneRadius_pos, Rd.LipschitzGraph.coneRadius_le,
  Rd.LipschitzGraph.two_mul_add_one_mul_coneRadius_le, Rd.LipschitzGraph.ball_subset_cone,
  Rd.LipschitzGraph.closedCone, Rd.LipschitzGraph.cone_subset_closedCone,
  Rd.LipschitzGraph.isClosed_closedCone, Rd.LipschitzGraph.smul_mem_closedCone,
  Rd.LipschitzGraph.phi_add_le, Rd.LipschitzGraph.add_mem_of_mem_closedCone}
\leanok
\uses{def:lipschitz-graph}
$\Omega + K \subseteq \Omega$ and $c K \subseteq K$ for $c > 0$, so $x + l h \in \Omega$ for
$x \in \Omega$, $h \in K$ and $l \in \N$. The cone $K$ contains the ball of radius
$\rho = 1/(8(M + 1))$ around $e/4$. For all $x, y \in V$,
$\varphi(x + y) \leq \varphi(x) + M \|y - \langle y, e \rangle e\|$. The closed cone
$\bar K = \{y : M \|y - \langle y, e \rangle e\| \leq \langle y, e \rangle\}$ is closed, contains
$K$, satisfies $c \bar K \subseteq \bar K$ for $c \geq 0$, and $\Omega + \bar K \subseteq \Omega$.
\end{lemma}

\begin{proof}
\leanok
Write $y = d + \langle y, e \rangle e$ with $d \perp e$. Since $\varphi$ is constant along $e$,
$\varphi(x + y) = \varphi(x + d) \leq \varphi(x) + M \|d\|$, and
$\varphi(x) + M \|d\| < \langle x, e \rangle + \langle y, e \rangle$
for $x \in \Omega$ and $y \in \bar K$, strict since $\varphi(x) < \langle x, e \rangle$; this gives
$\Omega + \bar K \subseteq \Omega$, and $\Omega + K \subseteq \Omega$ follows from
$K \subseteq \bar K$. For $\|y - e/4\| < \rho$, $\langle y, e \rangle > 1/4 - \rho$ and
$\|y - \langle y, e \rangle e\| < 2\rho$, and $2 M \rho \leq 1/4 - \rho$.
\end{proof}

\begin{definition}[Kernels supported in the reflected cone]
\label{def:cone-supported}
\lean{Rd.LipschitzGraph.ConeSupported, Rd.LipschitzGraph.ConeSupported.add,
  Rd.LipschitzGraph.ConeSupported.sub, Rd.LipschitzGraph.ConeSupported.dilate,
  Rd.LipschitzGraph.ConeSupported.finsetSum}
\leanok
\uses{lem:graph-cone, lem:dilate}
A function $k : V \to \C$ is supported in $-K$ if $k(y) \neq 0$ implies $-y \in K$. Linear
combinations of such functions and their dilations $k_c$, $c > 0$, are supported in $-K$.
\end{definition}

\begin{lemma}[The closure of a graph domain]
\label{lem:graph-domain-closure-cone}
\lean{Rd.LipschitzGraph.add_mem_of_mem_closure, Rd.LipschitzGraph.mem_cone_of_mem_tsupport_dilate}
\leanok
\uses{def:lipschitz-graph, lem:graph-cone, lem:dilate}
For every Lipschitz graph domain $\Omega_0 = \{x : \varphi(x) < \langle x, e \rangle\}$ with cone
$K$, every $x \in \overline{\Omega_0}$ and every $y \in K$, $x + y \in \Omega_0$. For every
$\eta \colon V \to \C$ with $\operatorname{supp} \eta \subseteq K$, every $c > 0$ and every
$x \in \operatorname{supp} \eta_c$, $x \in K$.
\end{lemma}

\begin{proof}
\leanok
On $\overline{\Omega_0}$, $\varphi(x) \leq \langle x, e \rangle$ by continuity of $\varphi$, and
$\varphi(x + y) \leq \varphi(x) + M \|y - \langle y, e \rangle e\| < \langle x + y, e \rangle$ by
\Cref{lem:graph-cone} and the strict inequality of the cone. If
$x \in \operatorname{supp} \eta_c$, then $c x \in K$ (\Cref{lem:dilate}), and
$c^{-1} K \subseteq K$.
\end{proof}

\begin{lemma}[The cone bump]
\label{lem:cone-bump}
\lean{Rd.LipschitzGraph.coneBumpBase, Rd.LipschitzGraph.coneBump,
  Rd.LipschitzGraph.coneBump_radius, Rd.LipschitzGraph.coneSupported_coneBump,
  Rd.LipschitzGraph.coneSupported_momentKernel, Rd.LipschitzGraph.coneSupported_scaled,
  Rd.LipschitzGraph.coneSupported_kernel, Rd.LipschitzGraph.coneSupported_dual}
\leanok
\uses{lem:graph-cone, def:cone-supported, def:local-means}
There is a smooth $g : V \to \C$ with $\int_V g = 1$ that vanishes outside the ball of radius
$1/2$ and is supported in $-K$. For every $L \geq 1$ the kernels $m_{2^j}$, $k_j$ and $\eta_j$
of \Cref{def:local-means} built from $g$ are supported in $-K$.
\end{lemma}

\begin{proof}
\leanok
Take the normalised smooth bump of radius $\rho$ around $-e/4$ (\Cref{lem:graph-cone}). It
vanishes outside the ball of radius $1/4 + \rho \leq 1/2$, and $g(y) \neq 0$ gives
$\|y + e/4\| < \rho$, so $-y$ lies in the ball of radius $\rho$ around $e/4$ and in $K$. The kernels are linear combinations of dilations of $g$.
\end{proof}

\begin{lemma}[Multiplication by an indicator]
\label{lem:indicator-lp}
\lean{Rd.memLp_indicator_coe, Rd.norm_toLp_indicator_le, Rd.indicatorLp,
  Rd.indicatorLp_apply_ae, Rd.setIntegral_mul_eq_lpPairing, Rd.enorm_indicatorLp_le}
\leanok
Let $\mu$ be a measure on $X$, $S \subseteq X$ measurable and $1 \leq p \leq \infty$. Then
$h \mapsto 1_S h$ is a bounded linear operator on $L_p(\mu)$ of norm at most $1$. For all
$a, b \in L_2(\mu)$, $\int_S a b \, d\mu = \langle 1_S a, b \rangle$ with the bilinear pairing
$\langle a, b \rangle = \int a b \, d\mu$ of $L_2(\mu)$ with itself.
\end{lemma}

\begin{proof}
\leanok
$|1_S h| \leq |h|$ pointwise, and $1_S a = a$ on $S$, $1_S a = 0$ off $S$ (Mathlib's
\texttt{ContinuousLinearMap.\allowbreak lpPairing\_\allowbreak eq\_\allowbreak integral}).
\end{proof}

\begin{lemma}[Restrictions in $L_p$]
\label{lem:indicator-lp-eq-iff}
\lean{Rd.indicatorLp_eq_iff_ae_eq_restrict}
\leanok
\uses{lem:indicator-lp}
For every measure $\mu$ on $X$, every measurable $S \subseteq X$, every $1 \leq p \leq \infty$
and all $F, f \in L_p(\mu)$, $1_S F = 1_S f$ in $L_p(\mu)$ if and only if $F = f$
$\mu$-almost everywhere on $S$.
\end{lemma}

\begin{proof}
\leanok
Both statements say that $1_S F = 1_S f$ $\mu$-almost everywhere (Mathlib's
\texttt{MeasureTheory.\allowbreak ae\_\allowbreak eq\_\allowbreak restrict\_\allowbreak
iff\_\allowbreak indicator\_\allowbreak ae\_\allowbreak eq}), and the class $1_S h$ of
\Cref{lem:indicator-lp} agrees almost everywhere with the function $1_S h$.
\end{proof}

\begin{lemma}[The support principle]
\label{lem:support-principle}
\lean{Rd.LipschitzGraph.conv_indicator_eq, Rd.LipschitzGraph.cutoff,
  Rd.LipschitzGraph.cutoff_apply_ae, Rd.LipschitzGraph.enorm_cutoff,
  Rd.LipschitzGraph.enorm_cutoff_le, Rd.LipschitzGraph.indicator_ae_eq_of_cutoff_eq,
  Rd.LipschitzGraph.cutoff_convLp_cutoff, Rd.LipschitzGraph.cutoff_convLp_eq}
\leanok
\uses{def:cone-supported, lem:graph-cone, lem:indicator-lp, lem:conv-op}
Let $k \in L_1$ be supported in $-K$ and $1 \leq p < \infty$. For $x \in \Omega$ and
$f : V \to \C$, $(k * (1_\Omega f))(x) = (k * f)(x)$. Hence
$1_\Omega (k * (1_\Omega h)) = 1_\Omega (k * h)$ for $h \in L_p$, and
$1_\Omega (k * F) = 1_\Omega (k * f)$ whenever $F, f \in L_p$ agree almost everywhere on
$\Omega$.
\end{lemma}

\begin{proof}
\leanok
If $k(y) \neq 0$, then $-y \in K$ and $x - y \in x + K \subseteq \Omega$, so
$k(y) (1_\Omega f)(x - y) = k(y) f(x - y)$.
\end{proof}

\section{The extension operator for Lipschitz graph domains}

In this section $\Omega$ is a Lipschitz graph domain, $g$ the bump of \Cref{lem:cone-bump},
$L \geq 1$, and $m$, $k_j$, $\eta_j$ are the kernels of \Cref{def:local-means} built from $g$.
The extension operator is Rychkov's universal extension \cite{Rychkov}, written with these
kernels.

\begin{definition}[Intrinsic local means]
\label{def:intrinsic-local-means}
\lean{Rd.LipschitzGraph.intrinsicNorm, Rd.LipschitzGraph.intrinsicNorm_congr,
  Rd.LipschitzGraph.intrinsicNorm_le_localMeansNorm}
\leanok
\uses{lem:cone-bump, lem:support-principle}
For $1 \leq p < \infty$, $s \in \R$, $1 \leq q \leq \infty$ and $f \in L_p$, the intrinsic local
means quasi-norm is $\|(2^{js} \|k_j * f\|_{L_p(\Omega)})_j\|_{\ell_q}$. It depends only on the
class of $f|_\Omega$ and is at most the local means quasi-norm
$\|(2^{js} \|k_j * f\|_p)_j\|_{\ell_q}$.
\end{definition}

\begin{theorem}[The extension operator]
\label{thm:graph-extension}
\lean{Rd.LipschitzGraph.extTerm, Rd.LipschitzGraph.summable_extTerm,
  Rd.LipschitzGraph.extension, Rd.LipschitzGraph.tendsto_sum_extTerm,
  Rd.LipschitzGraph.cutoff_extension, Rd.LipschitzGraph.modNorm_extension_le}
\leanok
\uses{def:intrinsic-local-means, lem:support-principle, lem:reproducing-formula,
  lem:local-means-synthesis}
Let $1 \leq p < \infty$, $s > 0$, $1 \leq q \leq \infty$, and let $f \in L_p$ have a finite
intrinsic local means quasi-norm. Then the series
\[
  \operatorname{Ext} f = \sum_{j=0}^\infty \eta_j * (1_\Omega (k_j * f))
\]
converges absolutely in $L_p$, and $\operatorname{Ext} f = f$ almost everywhere on $\Omega$. For
$s < L$ there is $C < \infty$, independent of $f$, with
\[
  \|\operatorname{Ext} f\|_p + \|(2^{is} \omega_L(\operatorname{Ext} f, 2^{-i})_p)_i\|_{\ell_q}
  \leq C \, \|(2^{js} \|k_j * f\|_{L_p(\Omega)})_j\|_{\ell_q} .
\]
\end{theorem}

\begin{proof}
\leanok
$\|\eta_j * (1_\Omega (k_j * f))\|_p \leq 2 \|m\|_1 \, 2^{-js}
\|(2^{is} \|k_i * f\|_{L_p(\Omega)})_i\|_{\ell_q}$, which is summable for $s > 0$. By
\Cref{lem:support-principle},
$1_\Omega (\eta_j * (1_\Omega (k_j * f))) = 1_\Omega (\eta_j * k_j * f)$, so on $\Omega$ the
partial sums are those of \Cref{lem:reproducing-formula}, which converge to $f$. The bound is
\Cref{lem:local-means-synthesis} with $h_j = 1_\Omega (k_j * f)$.
\end{proof}

\begin{theorem}[Restrictions by intrinsic local means]
\label{thm:restriction-local-means}
\lean{Rd.LipschitzGraph.exists_besov_extension_iff}
\leanok
\uses{thm:graph-extension, def:intrinsic-local-means, lem:local-means-upper}
Let $1 \leq p < \infty$, $0 < s < L$, $1 \leq q \leq \infty$ and $f \in L_p$. There is
$F \in L_p$ with $F = f$ almost everywhere on $\Omega$ and
$\|F\|_p + \|(2^{ks} \omega_L(F, 2^{-k})_p)_k\|_{\ell_q} < \infty$ if and only if
$\|(2^{js} \|k_j * f\|_{L_p(\Omega)})_j\|_{\ell_q} < \infty$, and then
$F = \operatorname{Ext} f$ is such a function.
\end{theorem}

\begin{proof}
\leanok
If $F$ exists, the intrinsic quasi-norm of $f$ equals that of $F$ and is at most the local means
quasi-norm of $F$, which is finite by \Cref{lem:local-means-upper}. The converse is
\Cref{thm:graph-extension}.
\end{proof}

\section{Intrinsic differences on Lipschitz graph domains}

The notation is that of the previous section.

\begin{lemma}[Steps in the cone]
\label{lem:cone-steps}
\lean{Rd.LipschitzGraph.subset_stepSet, Rd.LipschitzGraph.subset_stepSet_of_coneSupported,
  Rd.LipschitzGraph.enorm_cutoff_sub_convLp_scaled_le}
\leanok
\uses{lem:graph-cone, def:rd-modulus-on, thm:moment-approx, lem:cone-bump}
For $h \in K$ and $r \in \N$, $\Omega \subseteq \Omega_{rh}$. Consequently, for
$1 \leq p < \infty$, $j \in \N$ and $f \in L_p$,
\[
  \|f - m_{2^j} * f\|_{L_p(\Omega)} \leq \|g\|_1 \, \omega_L(f, 2^{-(j+1)}, \Omega)_p .
\]
\end{lemma}

\begin{proof}
\leanok
The first statement is \Cref{lem:graph-cone}. For $g(z) \neq 0$ the step
$-2^{-j} z = 2^{-j} (-z)$ lies in $K$ and has length at most $2^{-(j+1)}$; apply
\Cref{thm:moment-approx} with $U = T = \Omega$, $R = 1/2$ and $c = 2^j$.
\end{proof}

\begin{lemma}[Intrinsic local means and moduli on $\Omega$]
\label{lem:intrinsic-upper}
\lean{Rd.LipschitzGraph.enorm_cutoff_localMean_zero_le,
  Rd.LipschitzGraph.enorm_cutoff_localMean_succ_le, Rd.LipschitzGraph.intrinsicConst,
  Rd.LipschitzGraph.intrinsicConst_ne_top, Rd.LipschitzGraph.intrinsicNorm_le,
  Rd.LipschitzGraph.intrinsicNorm_lt_top}
\leanok
\uses{lem:cone-steps, lem:support-principle, lem:rd-modulus-on, def:intrinsic-local-means}
Let $1 \leq p < \infty$ and $f \in L_p$. Then
$\|k_0 * f\|_{L_p(\Omega)} \leq \|m\|_1 \|f\|_{L_p(\Omega)}$ and
$\|k_{j+1} * f\|_{L_p(\Omega)} \leq 2 \|g\|_1 \, \omega_L(f, 2^{-(j+1)}, \Omega)_p$ for
$j \in \N$. Hence, for $s \in \R$ and $1 \leq q \leq \infty$,
\begin{multline*}
  \|(2^{js} \|k_j * f\|_{L_p(\Omega)})_j\|_{\ell_q} \\
  \leq (\|m\|_1 + 2 \|g\|_1) \bigl( \|f\|_{L_p(\Omega)}
  + \|(2^{js} \omega_L(f, 2^{-j}, \Omega)_p)_j\|_{\ell_q} \bigr) .
\end{multline*}
\end{lemma}

\begin{proof}
\leanok
$1_\Omega (k_0 * f) = 1_\Omega (k_0 * (1_\Omega f))$ by \Cref{lem:support-principle}. For
$j \in \N$, $k_{j+1} * f = (f - m_{2^j} * f) - (f - m_{2^{j+1}} * f)$, and \Cref{lem:cone-steps}
bounds both terms by $\|g\|_1 \, \omega_L(f, 2^{-(j+1)}, \Omega)_p$, since the modulus
increases in $t$.
\end{proof}

\begin{theorem}[Besov spaces on Lipschitz graph domains by intrinsic differences]
\label{thm:graph-besov}
\lean{Rd.LipschitzGraph.modNormOn_le_modNorm, Rd.LipschitzGraph.cutoff_extension_of_modNormOn,
  Rd.LipschitzGraph.modNorm_extension_le_modNormOn,
  Rd.LipschitzGraph.exists_besov_extension_iff_modulusOn,
  Rd.LipschitzGraph.exists_memApprox_extension_iff}
\leanok
\uses{lem:intrinsic-upper, thm:graph-extension, lem:rd-modulus-on, thm:rd-besov}
Let $1 \leq p < \infty$, $0 < s < r$, $1 \leq q \leq \infty$ and $f \in L_p$. The following are
equivalent:
\begin{enumerate}
  \item $\|f\|_{L_p(\Omega)} + \|(2^{ks} \omega_r(f, 2^{-k}, \Omega)_p)_k\|_{\ell_q} < \infty$;
  \item there is $F \in L_p$ with $F = f$ almost everywhere on $\Omega$ and
    $\|F\|_p + \|(2^{ks} \omega_r(F, 2^{-k})_p)_k\|_{\ell_q} < \infty$;
  \item there is $F$ in the approximation space $A^s_q$ of the band-limited scheme with
    $F = f$ almost everywhere on $\Omega$.
\end{enumerate}
For every $F$ as in (2), the quantity in (1) is at most
$\|F\|_p + \|(2^{ks} \omega_r(F, 2^{-k})_p)_k\|_{\ell_q}$. Under (1), $F = \operatorname{Ext} f$
with $L = r$ satisfies (2), and there is $C < \infty$, independent of $f$, with
\begin{multline*}
  \|\operatorname{Ext} f\|_p + \|(2^{ks} \omega_r(\operatorname{Ext} f, 2^{-k})_p)_k\|_{\ell_q} \\
  \leq C \bigl( \|f\|_{L_p(\Omega)} + \|(2^{ks} \omega_r(f, 2^{-k}, \Omega)_p)_k\|_{\ell_q}
  \bigr) .
\end{multline*}
\end{theorem}

\begin{proof}
\leanok
If $F = f$ on $\Omega$, then $\|f\|_{L_p(\Omega)} = \|F\|_{L_p(\Omega)} \leq \|F\|_p$ and
$\omega_r(f, t, \Omega)_p = \omega_r(F, t, \Omega)_p \leq \omega_r(F, t)_p$ by
\Cref{lem:rd-modulus-on}. Conversely, (1) bounds the intrinsic local means by
\Cref{lem:intrinsic-upper}, and \Cref{thm:graph-extension} gives (2) with the bound for
$\operatorname{Ext} f$. The equivalence of (2) and (3) is \Cref{thm:rd-besov}.
\end{proof}

\section{Local extensions}

\begin{lemma}[The support principle for two sets]
\label{lem:support-principle-two}
\lean{Rd.conv_indicator_eq_of_forall, Rd.indicatorLp_convLp_indicatorLp}
\leanok
\uses{lem:conv-op, lem:indicator-lp}
Let $k \in L_1$, $1 \leq p < \infty$ and $U, W \subseteq V$ measurable with $x - y \in W$ whenever
$x \in U$ and $k(y) \neq 0$. Then $1_U (k * (1_W h)) = 1_U (k * h)$ for every $h \in L_p$.
\end{lemma}

\begin{proof}
\leanok
For $x \in U$, $k(y) (1_W h)(x - y) = k(y) h(x - y)$ for every $y$.
\end{proof}

\begin{lemma}[Bumps of every radius]
\label{lem:bumps-every-radius}
\lean{Rd.LocalMeansBump.rescale, Rd.LocalMeansBump.rescale_radius,
  Rd.LocalMeansBump.scaled_eq_zero, Rd.LocalMeansBump.dual_eq_zero,
  Rd.LocalMeansBump.kernel_eq_zero, Rd.centeredBumpBase, Rd.LocalMeansBump.centered,
  Rd.LocalMeansBump.centered_radius, Rd.LipschitzGraph.coneSupported_momentKernel_of,
  Rd.LipschitzGraph.coneSupported_scaled_of, Rd.LipschitzGraph.coneSupported_dual_of,
  Rd.LipschitzGraph.coneSupported_rescale}
\leanok
\uses{def:local-means, def:cone-supported, lem:dilate}
If $g$ is a bump for local means of radius $R$ (\Cref{def:local-means}) and $c > 0$, then $g_c$ is
one of radius $R/c$; for every $R > 0$ the normalised smooth bump around $0$ is one of radius $R$.
The kernels $m_{2^j}$, $k_j$ and $\eta_j$ of a bump of radius $R$ vanish outside the ball of radius
$L R$, and they are supported in $-K$ if $g$ is.
\end{lemma}

\begin{proof}
\leanok
$\int g_c = \int g$ (\Cref{lem:dilate}); the kernels are linear combinations of dilations
$g_{c'}$ with $c' \geq 1/L$.
\end{proof}

\begin{lemma}[Patches]
\label{lem:patches}
\lean{Rd.LipschitzGraph.sub_mem_patch_of_mem_closedCone, Rd.LipschitzGraph.patch_support,
  Rd.LipschitzGraph.patch_step, Rd.setOf_lt_infDist_compl_subset, Rd.interior_patch_support,
  Rd.interior_patch_step, Rd.measurableSet_infDist_gt}
\leanok
\uses{lem:graph-cone, lem:bumps-every-radius, def:rd-modulus-on}
Let $g$ be a bump of radius $R$, $L \in \N$ and $\varepsilon, \delta, \rho > 0$.
\begin{enumerate}
  \item If $\Omega$ is a Lipschitz graph domain, $g$ is supported in $-K$ and $L R \leq \varepsilon$,
    then for $x \in \Omega \cap B(c, \rho - 2\varepsilon)$ and $\eta_j(y) \neq 0$,
    $x - y \in \Omega \cap B(c, \rho - \varepsilon)$, and the steps $h = -c'^{-1} z$, $c' \geq 1$,
    $g(z) \neq 0$, satisfy $\Omega \cap B(c, \rho - \varepsilon) \subseteq T_{Lh}$ for
    $T = \Omega \cap B(c, \rho)$.
  \item If $O \subseteq V$ and $L R \leq \delta$, then for $\operatorname{dist}(x, O^c) > 2\delta$ and
    $\eta_j(y) \neq 0$, $\operatorname{dist}(x - y, O^c) > \delta$, and the same steps satisfy
    $\{\operatorname{dist}(\cdot, O^c) > \delta\} \subseteq O_{Lh}$.
\end{enumerate}
For a Lipschitz graph domain $\Omega$ with the closed cone $\bar K$, every
$x \in \Omega \cap B(c, \rho - \varepsilon)$ and every $y$ with $-y \in \bar K$ and
$\|y\| \leq \varepsilon$, $x - y \in \Omega \cap B(c, \rho)$. For every $O \subseteq V$ and every
$r \geq 0$, $\{\operatorname{dist}(\cdot, O^c) > r\} \subseteq O$.
\end{lemma}

\begin{proof}
\leanok
The supports of the kernels lie in the ball of radius $L R$ and, in (1), in $-K$; steps in the
cone stay in $\Omega$ (\Cref{lem:graph-cone}), and points within distance $L R$ stay in the larger
balls.
\end{proof}

\begin{lemma}[Local means on a set]
\label{lem:local-means-on-sets}
\lean{Rd.LocalMeansBump.localNormOn, Rd.LocalMeansBump.enorm_indicatorLp_sub_convLp_scaled_le,
  Rd.LocalMeansBump.enorm_indicatorLp_localMean_zero_le,
  Rd.LocalMeansBump.enorm_indicatorLp_localMean_succ_le, Rd.LocalMeansBump.localNormOn_le}
\leanok
\uses{thm:moment-approx, def:local-means, lem:rd-modulus-on}
Let $g$ be a bump of radius $R \leq 1/2$, $1 \leq p < \infty$, and let $W, T \subseteq V$ with $W$
measurable and $W \subseteq T_{Lh}$ for the steps $h = -c^{-1} z$, $c \geq 1$, $g(z) \neq 0$. For
$f \in L_p$, $s \in \R$ and $1 \leq q \leq \infty$,
\[
  \|(2^{js} \|k_j * f\|_{L_p(W)})_j\|_{\ell_q}
  \leq (\|m\|_1 + 2 \|g\|_1) \bigl( \|f\|_p + \|(2^{js} \omega_L(f, 2^{-j}, T)_p)_j\|_{\ell_q}
  \bigr) .
\]
\end{lemma}

\begin{proof}
\leanok
As in \Cref{lem:intrinsic-upper}, with \Cref{thm:moment-approx} for the two sets $W$ and $T$.
\end{proof}

\begin{theorem}[Local extensions]
\label{thm:local-extension}
\lean{Rd.LocalMeansBump.localExtTerm, Rd.LocalMeansBump.localExtension,
  Rd.LocalMeansBump.summable_localExtTerm, Rd.LocalMeansBump.tendsto_sum_localExtTerm,
  Rd.LocalMeansBump.modNorm_localExtension_le, Rd.LocalMeansBump.indicatorLp_localExtension,
  Rd.LocalMeansBump.exists_local_extension, Rd.LocalMeansBump.tendsto_sum_localExtTerm_of_summable,
  Rd.LocalMeansBump.indicatorLp_localExtension_of_summable,
  Rd.LocalMeansBump.enorm_localExtTerm_le_of_forall_le,
  Rd.LocalMeansBump.tsum_enorm_localExtTerm_le,
  Rd.LocalMeansBump.summable_localExtTerm_of_forall_le}
\leanok
\uses{lem:support-principle-two, lem:local-means-on-sets, lem:reproducing-formula,
  lem:local-means-synthesis}
Let $g$ be a bump of radius $R \leq 1/2$, $L \geq 1$, $1 \leq p < \infty$, $0 < s < L$ and
$1 \leq q \leq \infty$. Let $U \subseteq W$ and $U \subseteq T$ be measurable with $x - y \in W$
for $x \in U$ and $\eta_j(y) \neq 0$, and with $W \subseteq T_{Lh}$ for the steps
$h = -c^{-1} z$, $c \geq 1$, $g(z) \neq 0$. If
$\|f\|_{L_p(T)} + \|(2^{ks} \omega_L(f, 2^{-k}, T)_p)_k\|_{\ell_q} < \infty$, then there is
$F \in L_p$ with $\|F\|_p + \|(2^{ks} \omega_L(F, 2^{-k})_p)_k\|_{\ell_q} < \infty$ and $F = f$
almost everywhere on $U$. For every measurable $W$, every $s > 0$ and every $f \in L_p$ with
$2^{js} \|\mathbf 1_W (k_j * f)\|_p \leq X < \infty$ for every $j$,
$\sum_j \|\eta_j * \mathbf 1_W (k_j * f)\|_p \leq 2 \|m\|_1 (1 - 2^{-s})^{-1} X$.
\end{theorem}

\begin{proof}
\leanok
Take $F = \sum_j \eta_j * 1_W (k_j * 1_T f)$, which converges by \Cref{lem:local-means-on-sets}
and has the bound of \Cref{lem:local-means-synthesis}. On $U$, \Cref{lem:support-principle-two}
removes the cut-off $1_W$, and the reproducing formula (\Cref{lem:reproducing-formula}) gives
$1_T f = f$. The agreement on $U$ uses only the convergence of the series.
\end{proof}

The extension operator of \Cref{thm:graph-extension} is the case $U = W = T = \Omega$.

\section{Bounded Lipschitz domains}

\begin{definition}[Bounded Lipschitz domains]
\label{def:bounded-lipschitz}
\lean{Rd.BoundedLipschitzDomain, Rd.BoundedLipschitzDomain.inter_ball_eq,
  Rd.BoundedLipschitzDomain.measurableSet, Rd.BoundedLipschitzDomain.isCompact_closure,
  Rd.BoundedLipschitzDomain.isCompact_frontier, Rd.isCompact_frontier_of_isBounded}
\leanok
\uses{def:lipschitz-graph}
A bounded Lipschitz domain is a bounded open set $\Omega \subseteq V$ with finitely many open balls
$B_j = B(c_j, \rho_j)$ with $c_j \in \partial\Omega$, $\partial \Omega \subseteq \bigcup_j B_j$, and
Lipschitz graph domains $\Omega_j$ with $B_j \cap \Omega = B_j \cap \Omega_j$.
\end{definition}

This is the definition of Stein and Triebel, as in \cite[Remark 3.2]{SchneiderVybiral}; the
direction $e_j$ of each graph domain replaces the rotations. In Lean the balls and the graph
domains are data of the structure. In finite dimension the carriers of these structures are
exactly the sets of the usual definition, and the open balls and the bounded open convex sets are
among them (\Cref{thm:bounded-lipschitz-sets}, \Cref{thm:ball-lipschitz} and
\Cref{thm:convex-lipschitz}).

\begin{lemma}[Shrinking and the inner margin]
\label{lem:bounded-lipschitz-geometry}
\lean{Rd.BoundedLipschitzDomain.exists_shrink, Rd.BoundedLipschitzDomain.nonempty_compl,
  Rd.BoundedLipschitzDomain.exists_interior_margin, Rd.BoundedLipschitzDomain.patch,
  Rd.BoundedLipschitzDomain.isOpen_patch, Rd.BoundedLipschitzDomain.isBounded_patch,
  Rd.BoundedLipschitzDomain.closure_subset_iUnion_patch}
\leanok
\uses{def:bounded-lipschitz, lem:patches}
There is $\varepsilon > 0$ with $\partial\Omega \subseteq \bigcup_j B(c_j, \rho_j - 2\varepsilon)$.
In dimension $n \geq 1$, $\Omega \neq V$. If $\Omega \neq V$, there is $\delta > 0$ with
$\operatorname{dist}(x, \Omega^c) > 2\delta$ for every $x \in \Omega$ outside these balls. For
such $\varepsilon$ and $\delta \geq 0$ the patches $B(c_j, \rho_j - 2\varepsilon)$ and
$\{\operatorname{dist}(\cdot, \Omega^c) > 2\delta\}$ are bounded open sets that cover
$\overline\Omega$.
\end{lemma}

\begin{proof}
\leanok
Both by compactness: $\partial\Omega$ is covered by the increasing open sets
$\bigcup_j B(c_j, \rho_j - 1/(n+1))$, and $\operatorname{dist}(\cdot, \Omega^c)$ attains a positive
minimum on the compact set $\overline{\Omega} \setminus \bigcup_j B(c_j, \rho_j - 2\varepsilon)
\subseteq \Omega$. In dimension $n \geq 1$ the space $V$ is unbounded and $\Omega$ is bounded.
The inner patch lies in $\Omega$ (\Cref{lem:patches}), and a point of $\overline\Omega$ lies on
$\partial\Omega$, in one of the balls, or in $\Omega$ outside them.
\end{proof}

\begin{lemma}[Resolutions of unity]
\label{lem:resolution-of-unity}
\lean{Rd.exists_smooth_partition, Rd.exists_smooth_partition_nhds, Rd.exists_iteratedFDeriv_bound}
\leanok
For a closed set $S \subseteq V$ and a finite open cover $(U_i)$ of $S$ there are smooth
$\psi_i \geq 0$ with $\operatorname{supp} \psi_i \subseteq U_i$, $\sum_i \psi_i = 1$ on $S$ and
$\sum_i \psi_i \leq 1$ on $V$; if $S$ is compact, they can be chosen with $\sum_i \psi_i = 1$ on a
neighbourhood of $S$. Smooth functions with compact support have bounded derivatives of every
order.
\end{lemma}

\begin{proof}
\leanok
Mathlib's smooth partitions of unity on manifolds, for the model space $V$; the derivatives are
continuous with compact support. For compact $S$, apply the first statement to a compact $L$ with
$S \subseteq \operatorname{int} L \subseteq L \subseteq \bigcup_i U_i$ (Mathlib's
\texttt{exists\_\allowbreak compact\_\allowbreak between}).
\end{proof}

\begin{lemma}[Pasting local extensions]
\label{lem:pasting}
\lean{Rd.lpModulus_finsetSum_le, Rd.modNorm_finsetSum_lt_top, Rd.exists_pasting}
\leanok
\uses{lem:resolution-of-unity, thm:multipliers}
Let $O \subseteq V$ be measurable, $(U_i)$ a finite cover of $\overline O$ by bounded open sets, and
$F_i \in L_p$ with $\|F_i\|_p + \|(2^{ks} \omega_r(F_i, 2^{-k})_p)_k\|_{\ell_q} < \infty$ and
$F_i = f$ almost everywhere on $O \cap U_i$, where $s < r$. Then $f$ agrees on $O$ with some $F$ with
$\|F\|_p + \|(2^{ks} \omega_r(F, 2^{-k})_p)_k\|_{\ell_q} < \infty$.
\end{lemma}

\begin{proof}
\leanok
$F = \sum_i \psi_i F_i$ for a resolution of unity subordinate to $(U_i)$
(\Cref{lem:resolution-of-unity}): each $\psi_i F_i$ has a finite norm by \Cref{thm:multipliers},
and $\psi_i F_i = \psi_i f$ on $O$.
\end{proof}

\begin{lemma}[Restrictions and smaller sets]
\label{lem:restriction-moduli}
\lean{Rd.modNormOn_le_modNorm_of_indicatorLp, Rd.BoundedLipschitzDomain.modNormOn_mono_set}
\leanok
\uses{lem:rd-modulus-on}
Let $O \subseteq V$ be measurable. If $F = f$ almost everywhere on $O$, then
$\|f\|_{L_p(O)} + \|(2^{ks} \omega_r(f, 2^{-k}, O)_p)_k\|_{\ell_q} \leq
\|F\|_p + \|(2^{ks} \omega_r(F, 2^{-k})_p)_k\|_{\ell_q}$, and the left side decreases when $O$ is
replaced by a measurable subset.
\end{lemma}

\begin{proof}
\leanok
\Cref{lem:rd-modulus-on}.
\end{proof}

\begin{lemma}[The patches of a bounded Lipschitz domain]
\label{lem:bounded-patches}
\lean{Rd.BoundedLipschitzDomain.exists_local_extension_boundary,
  Rd.BoundedLipschitzDomain.exists_local_extension_interior}
\leanok
\uses{thm:local-extension, lem:patches, lem:bumps-every-radius, lem:restriction-moduli}
Let $1 \leq p < \infty$, $0 < s < r$, $1 \leq q \leq \infty$, $\varepsilon, \delta > 0$, and
$f \in L_p$ with $\|f\|_{L_p(\Omega)} + \|(2^{ks} \omega_r(f, 2^{-k}, \Omega)_p)_k\|_{\ell_q} < \infty$.
For every $j$, $f$ agrees on $\Omega \cap B(c_j, \rho_j - 2\varepsilon)$, and on
$\{\operatorname{dist}(\cdot, \Omega^c) > 2\delta\}$, with functions with finite Besov quasi-norms
on $V$.
\end{lemma}

\begin{proof}
\leanok
\Cref{thm:local-extension} with the patches of \Cref{lem:patches}: for the boundary balls the cone
bump of $\Omega_j$ rescaled to radius $\leq \varepsilon / r$, inside a centred bump of radius
$\leq \delta / r$; the moduli on the smaller sets are bounded by those on $\Omega$
(\Cref{lem:restriction-moduli}).
\end{proof}

\begin{theorem}[Besov spaces on bounded Lipschitz domains by intrinsic differences]
\label{thm:bounded-lipschitz-besov}
\lean{Rd.BoundedLipschitzDomain.exists_besov_extension_of_modulusOn,
  Rd.BoundedLipschitzDomain.exists_besov_extension_iff_modulusOn,
  Rd.BoundedLipschitzDomain.exists_memApprox_extension_iff}
\leanok
\uses{lem:bounded-lipschitz-geometry, lem:bounded-patches, lem:pasting, lem:restriction-moduli,
  thm:rd-besov}
Let $\Omega$ be a bounded Lipschitz domain, $1 \leq p < \infty$, $0 < s < r$,
$1 \leq q \leq \infty$ and $f \in L_p$. The following are equivalent:
\begin{enumerate}
  \item $\|f\|_{L_p(\Omega)} + \|(2^{ks} \omega_r(f, 2^{-k}, \Omega)_p)_k\|_{\ell_q} < \infty$;
  \item there is $F \in L_p$ with $F = f$ almost everywhere on $\Omega$ and
    $\|F\|_p + \|(2^{ks} \omega_r(F, 2^{-k})_p)_k\|_{\ell_q} < \infty$;
  \item there is $F$ in the approximation space $A^s_q$ of the band-limited scheme with $F = f$
    almost everywhere on $\Omega$.
\end{enumerate}
\end{theorem}

\begin{proof}
\leanok
(2) gives (1) by \Cref{lem:restriction-moduli}. For the converse, if $\Omega = V$ take $F = f$.
Otherwise \Cref{lem:bounded-lipschitz-geometry} gives $\varepsilon$ and $\delta$, the balls
$B(c_j, \rho_j - 2\varepsilon)$ and the set $\{\operatorname{dist}(\cdot, \Omega^c) > 2\delta\}$
cover $\overline\Omega$, \Cref{lem:bounded-patches} gives the local extensions, and
\Cref{lem:pasting} pastes them. (2) and (3) agree by \Cref{thm:rd-besov}.
\end{proof}

\section{The area formula for Lipschitz graphs}
\label{sec:lipschitz-graph-area}

Throughout this section $\Omega = \{x : \varphi(x) < \langle x, e \rangle\}$ is a Lipschitz graph
domain (\Cref{def:lipschitz-graph}) with Lipschitz constant $M$, hyperplane $H = e^\perp$ of
dimension $n - 1$, projection $P$ and graph map $\gamma(x') = x' + \varphi(x') e$
(\Cref{def:graph-map}). For $d \in \N$, $\mathcal H^d$ is Mathlib's $d$-dimensional Euclidean
Hausdorff measure (\texttt{MeasureTheory.\allowbreak Measure.\allowbreak
euclideanHausdorffMeasure}): the Hausdorff measure $\mu_H^d$ times the constant that makes it the
Lebesgue measure on every $d$-dimensional inner product space (Mathlib's
\texttt{InnerProductSpace.\allowbreak euclideanHausdorffMeasure\_\allowbreak eq\_\allowbreak
volume}); isometries preserve it. The area formula expresses $\mathcal H^{n-1}$ on the boundary
of $\Omega$ through the graph map. The proof is the one of Federer
\cite[3.2.3]{Federer} and Evans and Gariepy \cite[Section 3.3]{EvansGariepy}, for a graph, which
is injective and antilipschitz with constant $1$, so that every point of the graph has exactly one
preimage.

\begin{lemma}[Lipschitz maps and Euclidean Hausdorff measures]
\label{lem:hausdorff-lipschitz-image}
\lean{LipschitzOnWith.euclideanHausdorffMeasure_image_le,
  Rd.euclideanHausdorffMeasure_image_le_of_near}
\leanok
For all metric spaces $X$, $Y$, every $K \geq 0$, every $s \subseteq X$, every $f : X \to Y$ that
is $K$-Lipschitz on $s$ and every $d \in \N$,
$\mathcal H^d(f(s)) \leq K^d \mathcal H^d(s)$. For all normed groups $X$, $Y$, all
$f_1, f_2 : X \to Y$, $P \subseteq X$, $\delta \geq 0$ and $d \in \N$ with
$\|x - y\| \leq \|f_2(x) - f_2(y)\|$ and
$\|(f_1(x) - f_1(y)) - (f_2(x) - f_2(y))\| \leq \delta \|x - y\|$ for all $x, y \in P$,
\[
  \mathcal H^d(f_1(P)) \leq (1 + \delta)^d \, \mathcal H^d(f_2(P)).
\]
\end{lemma}

\begin{proof}
\leanok
The first bound is Mathlib's \texttt{LipschitzOnWith.\allowbreak hausdorffMeasure\_\allowbreak
image\_\allowbreak le} for $\mu_H^d$, multiplied by the constant. For the second, $f_2$ is
injective on $P$, so it has a left inverse $q$ on $P$ (Mathlib's
\texttt{Function.\allowbreak invFunOn}); $f_1 \circ q$ maps $f_2(P)$ onto $f_1(P)$, and for
$x, y \in P$,
$\|f_1(x) - f_1(y)\| \leq \|f_2(x) - f_2(y)\| + \delta \|x - y\|
\leq (1 + \delta) \|f_2(x) - f_2(y)\|$.
\end{proof}

\begin{definition}[The area factor]
\label{def:area-factor}
\lean{Rd.areaFactor}
\leanok
For $m \in \N$ and $\varepsilon \in \R$, the area factor is
$C_\varepsilon = (1 + 2\varepsilon)^m (1 + \varepsilon)$.
\end{definition}

In the area formula $m = n - 1$: comparing a graph with a linear graph costs the factor
$(1 + 2\varepsilon)^m$, and comparing the Gram density with a constant costs $1 + \varepsilon$.

\begin{lemma}[Three elementary bounds]
\label{lem:area-elementary}
\lean{Rd.abs_sqrt_one_add_sq_sub_le, Rd.tendsto_areaFactor,
  Rd.le_of_forall_pos_le_ofReal_mul_of_tendsto}
\leanok
\uses{def:area-factor}
For all $a, b \in \R$, $|\sqrt{1 + a^2} - \sqrt{1 + b^2}| \leq |a - b|$. For every $m \in \N$,
$C_\varepsilon \to 1$ as $\varepsilon \to 0+$. For all $a, b \in [0, \infty]$ and every
$c : \R \to \R$ with $c(\varepsilon) \to 1$ as $\varepsilon \to 0+$ and
$a \leq \max(c(\varepsilon), 0) \, b$ for every $\varepsilon > 0$, $a \leq b$.
\end{lemma}

\begin{proof}
\leanok
$|\sqrt{1 + a^2} - \sqrt{1 + b^2}| (\sqrt{1 + a^2} + \sqrt{1 + b^2}) = |a - b| |a + b|$ and
$|a + b| \leq \sqrt{1 + a^2} + \sqrt{1 + b^2}$. The area factor is continuous in $\varepsilon$
with value $1$ at $0$. In the last statement $\max(c(\varepsilon), 0) \, b \to b$ as
$\varepsilon \to 0+$.
\end{proof}

\begin{definition}[Graphs over the hyperplane]
\label{def:graph-over}
\lean{Rd.LipschitzGraph.graphOver, Rd.LipschitzGraph.linearGraph,
  Rd.LipschitzGraph.coe_linearGraph}
\leanok
\uses{def:graph-map}
For $h : H \to \R$, the graph of $h$ over $H$ is the map $\gamma_h(x') = x' + h(x') e$ from $H$ to
$V$. For $g \in H$, the linear graph of $g$ is the linear map
$S_g = \gamma_{\langle g, \cdot \rangle}$, $S_g x' = x' + \langle g, x' \rangle e$, from $H$ to
$V$.
\end{definition}

\begin{lemma}[Increments of graphs over the hyperplane]
\label{lem:graph-over}
\lean{Rd.LipschitzGraph.proj_graphOver, Rd.LipschitzGraph.graphOver_injective,
  Rd.LipschitzGraph.norm_coe_add_smul_e_sq, Rd.LipschitzGraph.norm_le_norm_coe_add_smul_e,
  Rd.LipschitzGraph.graphOver_sub, Rd.LipschitzGraph.norm_sub_le_norm_graphOver_sub,
  Rd.LipschitzGraph.graphOver_sub_sub}
\leanok
\uses{def:graph-over}
For all $h, h_1, h_2 : H \to \R$, $x', y' \in H$ and $a \in \R$: $P \gamma_h(x') = x'$, so
$\gamma_h$ is injective; $\|x' + a e\|^2 = \|x'\|^2 + a^2$;
the increment $\Delta_h = \gamma_h(x') - \gamma_h(y')$ is $(x' - y') + (h(x') - h(y')) e$, so
$\|x' - y'\| \leq \|\Delta_h\|$; and
\[
  \Delta_{h_1} - \Delta_{h_2} = \bigl((h_1(x') - h_1(y')) - (h_2(x') - h_2(y'))\bigr) e .
\]
\end{lemma}

\begin{proof}
\leanok
$P e = 0$ and $P x' = x'$ on $H$ (\Cref{def:graph-map}); Pythagoras, since $x' \perp e$. The
bound $\|x' - y'\| \leq \|\Delta_h\|$ holds since $x' - y' = P \Delta_h$ and $\|P z\| \leq \|z\|$.
\end{proof}

\begin{lemma}[Hausdorff measures of nearby graphs]
\label{lem:graph-over-hausdorff}
\lean{Rd.LipschitzGraph.euclideanHausdorffMeasure_graphOver_image_le}
\leanok
\uses{lem:graph-over, lem:hausdorff-lipschitz-image}
For all $h_1, h_2 : H \to \R$, $P \subseteq H$, $\delta \geq 0$ and $d \in \N$ with
$|(h_1(x') - h_1(y')) - (h_2(x') - h_2(y'))| \leq \delta \|x' - y'\|$ for all $x', y' \in P$,
\[
  \mathcal H^d(\gamma_{h_1}(P)) \leq (1 + \delta)^d \, \mathcal H^d(\gamma_{h_2}(P)).
\]
\end{lemma}

\begin{proof}
\leanok
\Cref{lem:hausdorff-lipschitz-image} with $f_2 = \gamma_{h_2}$, by \Cref{lem:graph-over}.
\end{proof}

For finite-dimensional real inner product spaces $U$, $W$, the norm determinant of a linear map
$T : U \to W$ is $\sqrt{\det(T^* T)}$ (Mathlib's \texttt{LinearMap.\allowbreak normDet} and
\texttt{LinearMap.\allowbreak normDet\_\allowbreak sq}), and
$\mathcal H^{\dim U}(T(P)) = \sqrt{\det(T^* T)} \, |P|$ for every $P \subseteq U$ (Mathlib's
\texttt{LinearMap.\allowbreak euclideanHausdorffMeasure\_\allowbreak image\_\allowbreak
eq\_\allowbreak normDet\_\allowbreak mul\_\allowbreak volume}).

\begin{lemma}[The norm determinant of a linear graph]
\label{lem:linear-graph-normdet}
\lean{Rd.LipschitzGraph.inner_linearGraph, Rd.LipschitzGraph.adjoint_comp_linearGraph,
  Rd.LipschitzGraph.normDet_linearGraph}
\leanok
\uses{def:graph-over}
For every $g \in H$ and all $x', y' \in H$,
$\langle S_g x', S_g y' \rangle
= \langle x', y' \rangle + \langle g, x' \rangle \langle g, y' \rangle$,
$S_g^* S_g x' = x' + \langle g, x' \rangle g$, and $\sqrt{\det(S_g^* S_g)} = \sqrt{1 + \|g\|^2}$.
\end{lemma}

\begin{proof}
\leanok
$x', y' \perp e$ and $\|e\| = 1$. Hence $\langle S_g^* S_g x', y' \rangle = \langle S_g x', S_g y'
\rangle = \langle x' + \langle g, x' \rangle g, y' \rangle$ for all $y' \in H$. The map
$x' \mapsto x' + \langle g, x' \rangle g$ is a transvection, of determinant
$1 + \langle g, g \rangle$ (Mathlib's \texttt{LinearMap.\allowbreak transvection.\allowbreak det}).
\end{proof}

\begin{theorem}[The area formula for linear graphs]
\label{thm:linear-graph-area}
\lean{Rd.LipschitzGraph.euclideanHausdorffMeasure_linearGraph_image}
\leanok
\uses{lem:linear-graph-normdet}
For every $g \in H$ and every $P \subseteq H$,
\[
  \mathcal H^{n-1}(S_g(P)) = \sqrt{1 + \|g\|^2} \, |P| ,
\]
with the Lebesgue measure $|P|$ of $H$.
\end{theorem}

\begin{proof}
\leanok
Mathlib's area formula for linear maps, with $\dim H = n - 1$ and the norm determinant of
\Cref{lem:linear-graph-normdet}.
\end{proof}

\paragraph{The linear graph as a tilted stretch.} Explicitly, $S_g$ for $g \neq 0$ is a stretch of
$H$ along $g/\|g\|$ followed by an isometry of $H$ into $V$ (\Cref{lem:tilt}); then
$\mathcal H^{n-1}(S_g(P)) = |\det D| \, |P|$ for the stretch $D$, of determinant
$\sqrt{1 + \|g\|^2}$ (\Cref{lem:det-stretch}, Mathlib's
\texttt{MeasureTheory.\allowbreak Measure.\allowbreak addHaar\_\allowbreak image\_\allowbreak
linearMap}).

\begin{definition}[Stretches]
\label{def:stretch}
\lean{Rd.stretch, Rd.stretch_apply}
\leanok
Let $E$ be a real inner product space, $u \in E$ and $c \in \R$. The stretch by $c$ along $u$ is
the linear map $D x = x + (c - 1) \langle u, x \rangle u$ of $E$.
\end{definition}

For a unit vector $u$, $D$ multiplies $u$ by $c$ and fixes the vectors orthogonal to $u$; it is
Mathlib's transvection (\texttt{LinearMap.\allowbreak transvection}) of the linear form
$(c - 1)\langle u, \cdot \rangle$ and the vector $u$.

\begin{lemma}[The determinant of a stretch]
\label{lem:det-stretch}
\lean{Rd.det_stretch}
\leanok
\uses{def:stretch}
For every finite-dimensional real inner product space $E$, every unit vector $u \in E$ and every
$c \in \R$, the stretch by $c$ along $u$ has determinant $c$.
\end{lemma}

\begin{proof}
\leanok
A transvection $x \mapsto x + f(x) u$ has determinant $1 + f(u)$ (Mathlib's
\texttt{LinearMap.\allowbreak transvection.\allowbreak det}), here
$1 + (c - 1) \langle u, u \rangle = c$.
\end{proof}

\begin{definition}[The tilt]
\label{def:tilt}
\lean{Rd.LipschitzGraph.tilt}
\leanok
\uses{def:graph-map}
For $u \in H$ and $b \in \R$, with $c = \sqrt{1 + b^2}$, the tilt is the map
$J x' = x' + \langle u, x' \rangle \bigl(c^{-1}(u + b e) - u\bigr)$ from $H$ to $V$.
\end{definition}

For a unit vector $u$ the tilt maps $u$ to the unit vector $(u + b e)/c$ and fixes the vectors of
$H$ orthogonal to $u$.

\begin{lemma}[Linear graphs are tilted stretches]
\label{lem:tilt}
\lean{Rd.LipschitzGraph.tilt_sub, Rd.LipschitzGraph.norm_tilt, Rd.LipschitzGraph.isometry_tilt,
  Rd.LipschitzGraph.tilt_stretch, Rd.LipschitzGraph.linearGraph_eq_tilt_comp_stretch}
\leanok
\uses{def:tilt, def:stretch, def:graph-over}
For every unit vector $u \in H$ and every $b \in \R$, the tilt $J$ is an isometry of $H$ into $V$,
and $J(D x') = x' + b \langle u, x' \rangle e$ for every $x' \in H$, where $D$ is the stretch by
$\sqrt{1 + b^2}$ along $u$. For every $g \in H \setminus \{0\}$, $S_g = J \circ D$ for
$u = g/\|g\|$ and $b = \|g\|$.
\end{lemma}

\begin{proof}
\leanok
$J$ is linear. With $a = \langle u, x' \rangle$ and $w = c^{-1}(u + b e) - u$,
$\langle x', w \rangle = a (c^{-1} - 1)$ and $\|w\|^2 = 2 - 2c^{-1}$, since $c^2 = 1 + b^2$, so
$\|x' + a w\|^2 = \|x'\|^2$. Further $\langle u, D x' \rangle = c a$, and
$D x' + c a w = x' + (c - 1) a u + a (u + b e) - c a u = x' + a b e$.
\end{proof}

\paragraph{The graph map and the Gram density.}

\begin{lemma}[The graph map]
\label{lem:graph-map-area}
\lean{Rd.LipschitzGraph.phiH, Rd.LipschitzGraph.graphMap_eq_graphOver,
  Rd.LipschitzGraph.graphMap_injective, Rd.LipschitzGraph.lipschitzWith_phiH,
  Rd.LipschitzGraph.continuous_phiH, Rd.LipschitzGraph.lipschitzWith_graphMap,
  Rd.LipschitzGraph.antilipschitzWith_graphMap, Rd.LipschitzGraph.range_graphMap,
  Rd.LipschitzGraph.measurableEmbedding_graphMap}
\leanok
\uses{def:graph-over, lem:graph-over, lem:boundary-charts}
The graph map is the graph $\gamma = \gamma_{\varphi|_H}$ of the $M$-Lipschitz function
$\varphi|_H$ over $H$. It is injective, $(1 + M)$-Lipschitz, satisfies
$\|x' - y'\| \leq \|\gamma(x') - \gamma(y')\|$ for all $x', y' \in H$, is a measurable embedding
(measurable, injective, and it maps measurable sets to measurable sets), and its range is the
boundary of $\Omega$.
\end{lemma}

\begin{proof}
\leanok
\Cref{lem:graph-over} and \Cref{def:graph-map}; an antilipschitz and uniformly continuous map of
the complete space $H$ is a closed embedding, hence a measurable embedding. The boundary of
$\Omega$ is $\{y : \gamma(Py) = y\}$ (\Cref{lem:boundary-charts}).
\end{proof}

\begin{lemma}[The graph map keeps null sets]
\label{lem:graph-map-null}
\lean{Rd.LipschitzGraph.euclideanHausdorffMeasure_graphMap_image_eq_zero}
\leanok
\uses{lem:graph-map-area, lem:hausdorff-lipschitz-image}
For every null set $N \subseteq H$, $\mathcal H^{n-1}(\gamma(N)) = 0$.
\end{lemma}

\begin{proof}
\leanok
$\mathcal H^{n-1}(\gamma(N)) \leq (1 + M)^{n-1} \mathcal H^{n-1}(N)$ by
\Cref{lem:hausdorff-lipschitz-image}, and $\mathcal H^{n-1}$ is the Lebesgue measure on $H$.
\end{proof}

\begin{definition}[The gradient and the Gram density]
\label{def:graph-gradient}
\lean{Rd.LipschitzGraph.gradPhi, Rd.LipschitzGraph.gram}
\leanok
\uses{lem:graph-map-area}
For $x' \in H$, $\nabla\varphi(x') \in H$ is the gradient of $\varphi|_H$ at $x'$ (Mathlib's
\texttt{gradient}, which is $0$ where $\varphi|_H$ is not differentiable), and
$\sqrt{1 + \|\nabla\varphi(x')\|^2}$ is the Gram density of the graph map.
\end{definition}

\begin{lemma}[Properties of the gradient]
\label{lem:graph-gradient}
\lean{Rd.LipschitzGraph.inner_gradPhi, Rd.LipschitzGraph.norm_gradPhi_le,
  Rd.LipschitzGraph.one_le_gram, Rd.LipschitzGraph.gram_pos, Rd.LipschitzGraph.gram_le,
  Rd.LipschitzGraph.abs_gram_sub_le, Rd.LipschitzGraph.measurable_gradPhi,
  Rd.LipschitzGraph.measurable_gram, Rd.LipschitzGraph.ae_differentiableAt_phiH}
\leanok
\uses{def:graph-gradient, lem:area-elementary}
For all $x', v, g \in H$: $\langle \nabla\varphi(x'), v \rangle$ is the derivative of
$\varphi|_H$ at $x'$ in the direction $v$, $\|\nabla\varphi(x')\| \leq M$,
$1 \leq \sqrt{1 + \|\nabla\varphi(x')\|^2} \leq \sqrt{1 + M^2}$ and
$|\sqrt{1 + \|\nabla\varphi(x')\|^2} - \sqrt{1 + \|g\|^2}| \leq \|\nabla\varphi(x') - g\|$. The
gradient and the Gram density are measurable, and $\varphi|_H$ is differentiable almost
everywhere on $H$.
\end{lemma}

\begin{proof}
\leanok
The derivative of an $M$-Lipschitz function has norm at most $M$; \Cref{lem:area-elementary} and
the reverse triangle inequality. The derivative is measurable (Mathlib's
\texttt{measurable\_\allowbreak fderiv}), and Rademacher's theorem (Mathlib's
\texttt{LipschitzWith.\allowbreak ae\_\allowbreak differentiableAt}).
\end{proof}

\begin{definition}[Good pieces]
\label{def:good-piece}
\lean{Rd.LipschitzGraph.IsGoodPiece}
\leanok
\uses{def:graph-gradient}
For $\varepsilon \in \R$ and $g \in H$, a set $P \subseteq H$ is a good piece if
$|(\varphi(x') - \varphi(y')) - \langle g, x' - y' \rangle| \leq 2 \varepsilon \|x' - y'\|$ and
$\|\nabla\varphi(x') - g\| \leq \varepsilon$ for all $x', y' \in P$.
\end{definition}

\begin{lemma}[The Gram density on good pieces]
\label{lem:good-piece-gram}
\lean{Rd.LipschitzGraph.IsGoodPiece.mono, Rd.LipschitzGraph.IsGoodPiece.sqrt_le_mul_gram,
  Rd.LipschitzGraph.IsGoodPiece.gram_le_mul_sqrt}
\leanok
\uses{def:good-piece, lem:graph-gradient}
For every $\varepsilon \geq 0$, every $g \in H$ and every good piece $P$ for $\varepsilon$ and $g$,
every subset of $P$ is a good piece for $\varepsilon$ and $g$, and for every $x' \in P$,
$\sqrt{1 + \|g\|^2} \leq (1 + \varepsilon) \sqrt{1 + \|\nabla\varphi(x')\|^2}$ and
$\sqrt{1 + \|\nabla\varphi(x')\|^2} \leq (1 + \varepsilon) \sqrt{1 + \|g\|^2}$.
\end{lemma}

\begin{proof}
\leanok
The two square roots differ by at most $\varepsilon$ (\Cref{lem:graph-gradient}), and both are at
least $1$.
\end{proof}

\begin{lemma}[The area formula on good pieces, up to a factor]
\label{lem:good-piece-bounds}
\lean{Rd.LipschitzGraph.IsGoodPiece.euclideanHausdorffMeasure_image_le,
  Rd.LipschitzGraph.IsGoodPiece.lintegral_gram_le}
\leanok
\uses{lem:good-piece-gram, lem:graph-over-hausdorff, thm:linear-graph-area, def:area-factor}
For every $\varepsilon \geq 0$, every $g \in H$ and every good piece $P$ for $\varepsilon$ and $g$,
with the area factor $C_\varepsilon$ for $m = n - 1$,
\[
  \mathcal H^{n-1}(\gamma(P)) \leq C_\varepsilon \int_P \sqrt{1 + \|\nabla\varphi(x')\|^2} \, dx'
\]
and
\[
  \int_P \sqrt{1 + \|\nabla\varphi(x')\|^2} \, dx' \leq C_\varepsilon \mathcal H^{n-1}(\gamma(P)).
\]
\end{lemma}

\begin{proof}
\leanok
By \Cref{lem:graph-over-hausdorff} with $\delta = 2\varepsilon$, $\mathcal H^{n-1}(\gamma(P))$
and $\mathcal H^{n-1}(S_g(P)) = \sqrt{1 + \|g\|^2} \, |P|$ (\Cref{thm:linear-graph-area}) differ
at most by the factor $(1 + 2\varepsilon)^{n-1}$, and $\sqrt{1 + \|g\|^2} \, |P|$ and the integral
differ at most by the factor $1 + \varepsilon$ (\Cref{lem:good-piece-gram}).
\end{proof}

\begin{definition}[Good sets]
\label{def:good-set}
\lean{Rd.LipschitzGraph.goodSet}
\leanok
\uses{def:graph-gradient}
Let $(u_l)_l$ be a dense sequence in $H$. For $\varepsilon \in \R$ and $i, k \in \N$, the good set
consists of the $x' \in H$ at which $\varphi|_H$ is differentiable, with
$\|\nabla\varphi(x') - u_i\| < \varepsilon$ and
$|\varphi(u_l) - \varphi(x') - \langle \nabla\varphi(x'), u_l - x' \rangle| \leq
\varepsilon \|u_l - x'\|$ for every $l$ with $\|u_l - x'\| < 1/(k+1)$.
\end{definition}

\begin{lemma}[Good sets]
\label{lem:good-set}
\lean{Rd.LipschitzGraph.abs_taylor_le_of_mem_goodSet,
  Rd.LipschitzGraph.isGoodPiece_goodSet_inter_ball,
  Rd.LipschitzGraph.exists_mem_goodSet_inter_ball, Rd.LipschitzGraph.measurableSet_goodSet}
\leanok
\uses{def:good-set, def:good-piece, lem:graph-gradient}
For every $\varepsilon \in \R$, all $i, k \in \N$ and every $x'$ in the good set,
$|\varphi(y') - \varphi(x') - \langle \nabla\varphi(x'), y' - x' \rangle| \leq
\varepsilon \|y' - x'\|$ for every $y' \in H$ with $\|y' - x'\| < 1/(k+1)$. For every
$c \in H$, the intersection of the good set with the ball of radius $1/(2(k+1))$ around $c$ is a
good piece for $\varepsilon$ and $u_i$. For every $\varepsilon > 0$ and every $x'$ at which
$\varphi|_H$ is differentiable there are $i, k, j \in \N$ such that $x'$ lies in the good set of
$i$ and $k$ and in the ball of radius $1/(2(k+1))$ around $u_j$. The good sets are measurable.
\end{lemma}

\begin{proof}
\leanok
The inequality is closed in $y'$ and holds on the dense sequence. On the intersection with the
ball, $\|x' - y'\| < 1/(k+1)$, and the bound at $y'$ together with
$\|\nabla\varphi(y') - u_i\| < \varepsilon$ gives the good piece. For the covering, $k$ comes from
the definition of the derivative, $i$ and $j$ from density. Measurability: the set of
differentiability points is measurable (Mathlib's
\texttt{measurableSet\_\allowbreak of\_\allowbreak differentiableAt}), the gradient is measurable,
and the Taylor condition quantifies over countably many $l$.
\end{proof}

\begin{lemma}[A countable partition into good pieces]
\label{lem:good-partition}
\lean{Rd.LipschitzGraph.exists_goodPartition}
\leanok
\uses{lem:good-set, lem:good-piece-gram, lem:graph-gradient}
For every $\varepsilon > 0$ there are pairwise disjoint measurable sets $Q_0, Q_1, \ldots
\subseteq H$ whose union has a null complement, and $g_0, g_1, \ldots \in H$, such that every
$Q_m$ is a good piece for $\varepsilon$ and $g_m$.
\end{lemma}

\begin{proof}
\leanok
Enumerate the intersections of \Cref{lem:good-set} by $\N$ and make them disjoint (Mathlib's
\texttt{disjointed}); their union contains every point of differentiability, and the others form
a null set by Rademacher's theorem.
\end{proof}

\begin{lemma}[The area formula up to a factor]
\label{lem:area-fixed-accuracy}
\lean{Rd.LipschitzGraph.euclideanHausdorffMeasure_graphMap_image_le_mul,
  Rd.LipschitzGraph.lintegral_gram_le_mul}
\leanok
\uses{lem:good-partition, lem:good-piece-bounds, lem:graph-map-null, lem:graph-map-area}
For every $\varepsilon > 0$ and every measurable $A \subseteq H$, with the area factor
$C_\varepsilon$ for $m = n - 1$,
\[
  \mathcal H^{n-1}(\gamma(A)) \leq C_\varepsilon \int_A \sqrt{1 + \|\nabla\varphi(x')\|^2} \, dx'
\]
and
\[
  \int_A \sqrt{1 + \|\nabla\varphi(x')\|^2} \, dx' \leq C_\varepsilon \mathcal H^{n-1}(\gamma(A)).
\]
\end{lemma}

\begin{proof}
\leanok
With the partition of \Cref{lem:good-partition}, $A \cap Q_m$ is a good piece for every $m$
(\Cref{lem:good-piece-gram}), and the rest of $A$ is a null set whose image is
$\mathcal H^{n-1}$-null (\Cref{lem:graph-map-null}). Sum the bounds of \Cref{lem:good-piece-bounds}:
$\mathcal H^{n-1}$ is countably subadditive, and the images $\gamma(A \cap Q_m)$ are pairwise
disjoint and measurable (\Cref{lem:graph-map-area}).
\end{proof}

\begin{theorem}[The area formula for Lipschitz graphs]
\label{thm:lipschitz-graph-area}
\lean{Rd.LipschitzGraph.euclideanHausdorffMeasure_graphMap_image,
  Rd.LipschitzGraph.map_graphMap_withDensity, Rd.LipschitzGraph.lintegral_graphMap,
  Rd.LipschitzGraph.integral_graphMap}
\leanok
\uses{lem:area-fixed-accuracy, lem:area-elementary, lem:graph-map-area}
For every measurable $A \subseteq H$,
\[
  \mathcal H^{n-1}(\gamma(A)) = \int_A \sqrt{1 + \|\nabla\varphi(x')\|^2} \, dx' .
\]
The image under $\gamma$ of the measure on $H$ with density $\sqrt{1 + \|\nabla\varphi\|^2}$ is
$\mathcal H^{n-1}$ restricted to $\gamma(H)$, and for every $F : V \to [0, \infty]$,
\[
  \int_{\gamma(H)} F \, d\mathcal H^{n-1}
  = \int_H F(\gamma(x')) \sqrt{1 + \|\nabla\varphi(x')\|^2} \, dx' ,
\]
and the same identity holds for every $F : V \to \C$ that is $\mathcal H^{n-1}$-integrable on
$\gamma(H)$.
\end{theorem}

\begin{proof}
\leanok
\Cref{lem:area-fixed-accuracy} and the last two parts of \Cref{lem:area-elementary}, in both
directions. The measure form follows on measurable sets $B \subseteq V$ with
$\gamma(\gamma^{-1}(B)) = B \cap \gamma(H)$, and the integral form from it, since $\gamma$ is a
measurable embedding and the density is finite.
\end{proof}

By \Cref{lem:graph-map-area}, $\gamma(H)$ is the boundary of $\Omega$, so the theorem expresses
$\mathcal H^{n-1}$ on the boundary in the chart. The integral form holds for every $F$,
measurable or not; in Lean the complex form holds for every $F : V \to \C$, both sides being $0$
by Mathlib's convention when $F$ is not integrable.

\begin{definition}[The outward unit normal]
\label{def:graph-normal}
\lean{Rd.LipschitzGraph.normalAt}
\leanok
\uses{def:graph-gradient}
For $x' \in H$, the outward unit normal at $\gamma(x')$ is
$\nu(x') = (\nabla\varphi(x') - e)/\sqrt{1 + \|\nabla\varphi(x')\|^2}$.
\end{definition}

Where $\varphi|_H$ is differentiable, $\nu(x')$ is the normalised gradient at $\gamma(x')$ of
$x \mapsto \varphi(Px) - \langle x, e \rangle$, a function that is negative on $\Omega$.

\begin{lemma}[Properties of the normal]
\label{lem:graph-normal}
\lean{Rd.LipschitzGraph.norm_gradPhi_sub_e, Rd.LipschitzGraph.norm_normalAt,
  Rd.LipschitzGraph.inner_normalAt_e, Rd.LipschitzGraph.inner_normalAt_neg_of_mem_cone,
  Rd.LipschitzGraph.measurable_normalAt}
\leanok
\uses{def:graph-normal, lem:graph-gradient, lem:graph-over, lem:graph-cone}
For every $x' \in H$, $\|\nabla\varphi(x') - e\| = \sqrt{1 + \|\nabla\varphi(x')\|^2}$,
$\|\nu(x')\| = 1$, $\langle \nu(x'), e \rangle = -1/\sqrt{1 + \|\nabla\varphi(x')\|^2}$, and
$\langle \nu(x'), y \rangle < 0$ for every $y$ in the cone $K$ of $\Omega$. The normal is
measurable.
\end{lemma}

\begin{proof}
\leanok
Pythagoras, since $\nabla\varphi(x') \perp e$. For $y \in K$,
$\langle \nabla\varphi(x'), y \rangle
= \langle \nabla\varphi(x'), y - \langle y, e \rangle e \rangle
\leq M \|y - \langle y, e \rangle e\| < \langle y, e \rangle$.
\end{proof}

\section{The surface measure of a bounded Lipschitz domain}
\label{sec:surface-measure}

Throughout this section $\Omega$ is a bounded Lipschitz domain in $V$ with boundary $\Gamma$,
boundary balls $B_j$, graph domains $\Omega_j$ with Lipschitz constants $M_j$ over the
hyperplanes $H_j$, graph maps $\gamma_j(x') = x' + \varphi_j(x') e_j$ and gradients
$\nabla\varphi_j$ (\Cref{def:bounded-lipschitz}, \Cref{def:graph-gradient}), and $\mathcal H^d$
is the $d$-dimensional Euclidean Hausdorff measure of \Cref{sec:lipschitz-graph-area}. As in
\Cref{sec:boundary-trace-space}, $\psi^0$ is the fixed resolution of unity on $\Gamma$,
$\operatorname{Tr}_\Gamma$ the boundary trace, $N_\Gamma$ the trace kernel, $\gamma$ and $E$ the
trace and the extension of the trace space $H^{1/2}(\Gamma)$, and boundary traces of $H^1$
functions are taken with $L = 2$, $p = q = 2$ and $s = 1/2$. The surface measure of $\Gamma$ is
$\mathcal H^{n-1}$ on $\Gamma$, defined without charts; the area formula for Lipschitz graphs
(\Cref{thm:lipschitz-graph-area}) gives its form in the charts.

\begin{lemma}[The boundary in a boundary ball]
\label{lem:boundary-ball-graph}
\lean{Rd.BoundedLipschitzDomain.frontier_inter_ball_eq}
\leanok
\uses{lem:boundary-charts}
For every $j$, $\Gamma \cap B_j = \gamma_j(H_j) \cap B_j$.
\end{lemma}

\begin{proof}
\leanok
For $y \in B_j$, $y \in \Gamma$ if and only if $\gamma_j(P_j y) = y$ (\Cref{lem:boundary-charts}),
and $\gamma_j(x') \in \Gamma$ for every $x'$ with $\gamma_j(x') \in B_j$.
\end{proof}

\begin{definition}[The surface measure]
\label{def:surface-measure}
\lean{Rd.BoundedLipschitzDomain.surfaceMeasure, Rd.BoundedLipschitzDomain.ae_mem_frontier,
  Rd.BoundedLipschitzDomain.surfaceMeasure_compl_frontier}
\leanok
\uses{def:bounded-lipschitz}
The surface measure of $\Gamma$ is the measure $\sigma$ on $V$ with
$\sigma(B) = \mathcal H^{n-1}(B \cap \Gamma)$ for every measurable $B \subseteq V$. Almost every
point of $V$ with respect to $\sigma$ lies on $\Gamma$, and $\sigma(V \setminus \Gamma) = 0$.
\end{definition}

On the part of $\Gamma$ that lies in a hyperplane $H$, $\sigma$ is the Lebesgue measure of $H$.

\begin{theorem}[The surface measure in the charts]
\label{thm:surface-measure-charts}
\lean{Rd.BoundedLipschitzDomain.surfaceMeasure_restrict_ball,
  Rd.BoundedLipschitzDomain.lintegral_surfaceMeasure_of_support,
  Rd.BoundedLipschitzDomain.lintegral_chart_eq,
  Rd.BoundedLipschitzDomain.surfaceMeasure_eq_sum_withDensity,
  Rd.BoundedLipschitzDomain.lintegral_surfaceMeasure,
  Rd.BoundedLipschitzDomain.lintegral_surfaceMeasure_le_sum}
\leanok
\uses{def:surface-measure, lem:boundary-ball-graph, thm:lipschitz-graph-area,
  def:boundary-partition}
For every $j$, $\sigma$ restricted to $B_j$ is $\mathcal H^{n-1}$ restricted to
$\gamma_j(H_j) \cap B_j$. For every $j$ and every $F : V \to [0, \infty]$ with $F = 0$ outside
$B_j$,
\[
  \int_\Gamma F \, d\sigma
  = \int_{H_j} F(\gamma_j(x')) \sqrt{1 + \|\nabla\varphi_j(x')\|^2} \, dx' .
\]
For all $j$, $k$ and every $F : V \to [0, \infty]$ with $F = 0$ outside $B_j \cap B_k$,
\[
  \int_{H_j} F(\gamma_j(x')) \sqrt{1 + \|\nabla\varphi_j(x')\|^2} \, dx'
  = \int_{H_k} F(\gamma_k(y')) \sqrt{1 + \|\nabla\varphi_k(y')\|^2} \, dy' .
\]
For every resolution of unity $\psi$ on $\Gamma$ with $\psi_j \geq 0$ on $\Gamma$ for every $j$,
$\sigma = \sum_j \psi_j \sigma$, and for every $F : V \to [0, \infty]$,
\[
  \int_\Gamma F \, d\sigma = \sum_j \int_{H_j} \psi_j(\gamma_j(x')) F(\gamma_j(x'))
  \sqrt{1 + \|\nabla\varphi_j(x')\|^2} \, dx' .
\]
For every $F : V \to [0, \infty]$,
\[
  \int_\Gamma F \, d\sigma \leq \sum_j \int_{H_j} (\mathbf 1_{B_j} F)(\gamma_j(x'))
  \sqrt{1 + \|\nabla\varphi_j(x')\|^2} \, dx' .
\]
\end{theorem}

\begin{proof}
\leanok
$\Gamma \cap B_j = \gamma_j(H_j) \cap B_j$ (\Cref{lem:boundary-ball-graph}) gives the first
statement. For $F = 0$ outside $B_j$, $\int_\Gamma F \, d\sigma$ is the integral of $F$ over
$\gamma_j(H_j)$ with respect to $\mathcal H^{n-1}$, which is the chart integral by the area
formula (\Cref{thm:lipschitz-graph-area}); the second statement follows by applying this in both
charts. For the third, $\sum_j \psi_j = 1$ on $\Gamma$ and $\sigma$ is carried by $\Gamma$, so
$\sigma = \sum_j \psi_j \sigma$; then
$\int_\Gamma F \, d\sigma = \sum_j \int_\Gamma \psi_j F \, d\sigma$ and $\psi_j F = 0$ outside
$B_j$. For the fourth, the balls $B_j$ cover $\Gamma$, so $\sigma$ is
carried by $\bigcup_j B_j$ and $\int_\Gamma F \, d\sigma \leq \sum_j \int_\Gamma \mathbf 1_{B_j} F
\, d\sigma$.
\end{proof}

Both sides of the second statement equal $\int_\Gamma F \, d\sigma$, so the chart integrals agree
on the overlaps of the charts. None of the statements asks for $F$ to be measurable.

\begin{theorem}[The surface measure in the charts, for complex integrals]
\label{thm:surface-measure-charts-complex}
\lean{Rd.BoundedLipschitzDomain.integral_surfaceMeasure_of_support,
  Rd.BoundedLipschitzDomain.integral_surfaceMeasure}
\leanok
\uses{thm:surface-measure-charts, thm:lipschitz-graph-area, def:boundary-partition}
For every $j$ and every $\sigma$-integrable $f : V \to \C$ with $f = 0$ outside $B_j$,
\[
  \int_\Gamma f \, d\sigma
  = \int_{H_j} f(\gamma_j(x')) \sqrt{1 + \|\nabla\varphi_j(x')\|^2} \, dx' .
\]
For every resolution of unity $\psi$ on $\Gamma$ and every $\sigma$-integrable $f : V \to \C$,
\[
  \int_\Gamma f \, d\sigma = \sum_j \int_{H_j} \psi_j(\gamma_j(x')) f(\gamma_j(x'))
  \sqrt{1 + \|\nabla\varphi_j(x')\|^2} \, dx' .
\]
\end{theorem}

\begin{proof}
\leanok
The first statement is the complex form of the area formula (\Cref{thm:lipschitz-graph-area}) on
$\gamma_j(H_j) \cap B_j = \Gamma \cap B_j$ (\Cref{thm:surface-measure-charts}). For the second,
$f = \sum_j \psi_j f$ on $\Gamma$, each $\psi_j f$ is $\sigma$-integrable since $\psi_j$ is
bounded, and $\psi_j f = 0$ outside $B_j$.
\end{proof}

The resolution of unity in the second statement may take negative values.

\begin{theorem}[Finiteness and null sets of the surface measure]
\label{thm:surface-measure-null-sets}
\lean{Rd.BoundedLipschitzDomain.isFiniteMeasure_surfaceMeasure,
  Rd.BoundedLipschitzDomain.ae_surfaceMeasure_ball_iff,
  Rd.BoundedLipschitzDomain.ae_graphMap_of_ae_surfaceMeasure,
  Rd.BoundedLipschitzDomain.ae_surfaceMeasure_iff_aeBoundary}
\leanok
\uses{thm:surface-measure-charts, thm:lipschitz-graph-area, lem:graph-gradient,
  lem:graph-map-area, def:ae-boundary}
$\sigma(V) < \infty$. For every $j$ and every property $P$ of the points of $V$, $P(y)$ holds for
$\sigma$-almost every $y \in B_j$ if and only if $P(\gamma_j(x'))$ holds for almost every
$x' \in H_j$ with $\gamma_j(x') \in B_j$; in particular, for every $j$, a property that holds
$\sigma$-almost everywhere holds at $\gamma_j(x')$ for almost every $x' \in H_j$ with
$\gamma_j(x') \in B_j$. A property holds $\sigma$-almost everywhere if and only if it holds
almost everywhere on $\Gamma$ in the sense of \Cref{def:ae-boundary}.
\end{theorem}

\begin{proof}
\leanok
By the fourth statement of \Cref{thm:surface-measure-charts} with $F = 1$ and
$\sqrt{1 + \|\nabla\varphi_j\|^2} \leq \sqrt{1 + M_j^2}$ (\Cref{lem:graph-gradient}), we get
\[
  \sigma(V) \leq \sum_j \sqrt{1 + M_j^2} \, |\gamma_j^{-1}(B_j)| ,
\]
and $\gamma_j^{-1}(B_j)$ is
bounded since $\|x' - y'\| \leq \|\gamma_j(x') - \gamma_j(y')\|$ (\Cref{lem:graph-map-area}). On
$B_j$, $\sigma$ is the image under $\gamma_j$ of the measure on $H_j$ with density
$\sqrt{1 + \|\nabla\varphi_j\|^2} \geq 1$ (\Cref{thm:lipschitz-graph-area},
\Cref{thm:surface-measure-charts}), which has the null sets of the Lebesgue measure of $H_j$;
this holds for every property, measurable or not, since $\gamma_j$ is a measurable embedding. The
last statement follows since the balls $B_j$ cover $\Gamma$ and $\sigma$-almost every point lies
on $\Gamma$.
\end{proof}

\begin{corollary}[The boundary is a null set]
\label{cor:boundary-null}
\lean{Rd.BoundedLipschitzDomain.volume_frontier}
\leanok
\uses{thm:surface-measure-null-sets, def:surface-measure}
$|\Gamma| = 0$.
\end{corollary}

\begin{proof}
\leanok
If there is no boundary ball, $\Gamma = \emptyset$. Otherwise $V$ contains a unit vector, so
$n \geq 1$, and $\mathcal H^{n-1}(\Gamma) = \sigma(V) < \infty$
(\Cref{thm:surface-measure-null-sets}). For $n - 1 < n$, every set has $\mathcal H^n$-measure
zero or $\mathcal H^{n-1}$-measure infinity (Mathlib's \texttt{MeasureTheory.\allowbreak
Measure.\allowbreak euclideanHausdorffMeasure\_\allowbreak zero\_\allowbreak or\_\allowbreak
top}), and $\mathcal H^n$ is the Lebesgue measure of $V$.
\end{proof}

\paragraph{The trace in $L_2(\Gamma)$.}

\begin{theorem}[The boundary trace in $L_p(\Gamma)$]
\label{thm:boundary-trace-lp}
\lean{Rd.BoundedLipschitzDomain.stronglyMeasurable_boundaryTrace,
  Rd.BoundedLipschitzDomain.lintegral_enorm_boundaryTrace_le,
  Rd.BoundedLipschitzDomain.memLp_boundaryTrace}
\leanok
\uses{def:boundary-trace, thm:boundary-trace-charts, thm:surface-measure-charts,
  lem:graph-gradient}
Let $1 \leq p < \infty$ and $\psi$ a resolution of unity on $\Gamma$. For every $L \in \N$ and
every $F \in L_p(V)$, $\operatorname{Tr}_\Gamma F$ is measurable. Let moreover $L \geq 1$, $s > 0$,
$s + 1/p < L$, $1 \leq q \leq \infty$ and
$\|F\|_p + \|(2^{k(s + 1/p)} \omega_L(F, 2^{-k})_p)_k\|_{\ell_q} < \infty$. Then
\[
  \int_\Gamma |\operatorname{Tr}_\Gamma F|^p \, d\sigma
  \leq \sum_j \sqrt{1 + M_j^2} \, \|\operatorname{Tr}_j F\|_{L_p(H_j)}^p ,
\]
and $\operatorname{Tr}_\Gamma F \in L_p(\sigma)$.
\end{theorem}

\begin{proof}
\leanok
$\operatorname{Tr}_\Gamma F = \sum_k \psi_k (\operatorname{Tr}_k F \circ P_k)$ is a finite sum of
products of continuous functions with measurable functions composed with the continuous
projections $P_k$. By the fourth statement of \Cref{thm:surface-measure-charts} the integral is at
most $\sum_j \int_{H_j} (\mathbf 1_{B_j} |\operatorname{Tr}_\Gamma F|^p)(\gamma_j(x'))
\sqrt{1 + \|\nabla\varphi_j(x')\|^2} \, dx'$, and on $\gamma_j^{-1}(B_j)$,
$\operatorname{Tr}_\Gamma F \circ \gamma_j = \operatorname{Tr}_j F$ almost everywhere
(\Cref{thm:boundary-trace-charts}) and the Gram density is at most $\sqrt{1 + M_j^2}$
(\Cref{lem:graph-gradient}). Each $\operatorname{Tr}_j F$ lies in $L_p(H_j)$.
\end{proof}

\begin{definition}[The space $L_2(\Gamma)$]
\label{def:boundary-l2}
\lean{Rd.BoundedLipschitzDomain.BoundaryL2}
\leanok
\uses{def:surface-measure}
$L_2(\Gamma)$ is the space $L_2(\sigma)$ of square integrable complex functions with respect to
the surface measure, modulo $\sigma$-null sets.
\end{definition}

\begin{theorem}[The trace of $H^1(V)$ in $L_2(\Gamma)$]
\label{thm:ambient-trace-l2}
\lean{Rd.BoundedLipschitzDomain.memLp_boundaryTrace_toL2,
  Rd.BoundedLipschitzDomain.exists_eLpNorm_boundaryTrace_toL2_le,
  Rd.BoundedLipschitzDomain.boundaryTrace_toL2_add,
  Rd.BoundedLipschitzDomain.boundaryTrace_toL2_smul,
  Rd.BoundedLipschitzDomain.ambientTraceLinear, Rd.BoundedLipschitzDomain.ambientTraceL2,
  Rd.BoundedLipschitzDomain.ambientTraceL2_apply_ae}
\leanok
\uses{thm:boundary-trace-lp, lem:graph-trace-sobolev-one, thm:sobolev-one-besov-constant,
  thm:trace-kernel-boundary, thm:boundary-trace-charts, lem:trace-linear,
  thm:surface-measure-null-sets, def:boundary-l2}
For every resolution of unity $\psi$ on $\Gamma$ and every $U \in H^1(V)$,
$\operatorname{Tr}_\Gamma \iota U \in L_2(\Gamma)$. There is $C < \infty$ with
$\|\operatorname{Tr}_\Gamma \iota U\|_{L_2(\Gamma)} \leq C \|U\|$ for every resolution of unity
$\psi$ and every $U \in H^1(V)$. For every resolution of unity $\psi$, all $U, W \in H^1(V)$ and
every $c \in \C$, $\operatorname{Tr}_\Gamma \iota (U + W) = \operatorname{Tr}_\Gamma \iota U +
\operatorname{Tr}_\Gamma \iota W$ and $\operatorname{Tr}_\Gamma \iota (c U) = c
\operatorname{Tr}_\Gamma \iota U$ almost everywhere on $\Gamma$. The map
$T_\Gamma U = \operatorname{Tr}_\Gamma \iota U$, with the boundary trace of $\psi^0$, is a bounded
linear operator $T_\Gamma : H^1(V) \to L_2(\Gamma)$, and $T_\Gamma U = \operatorname{Tr}_\Gamma
\iota U$ $\sigma$-almost everywhere for every resolution of unity $\psi$ and every
$U \in H^1(V)$.
\end{theorem}

\begin{proof}
\leanok
\Cref{thm:boundary-trace-lp} applies with $s = 1/2$, $p = q = 2$, $L = 2$
(\Cref{thm:sobolev-one-besov-constant}). With the constants $C_j$ of
\Cref{lem:graph-trace-sobolev-one} for the graph domains $\Omega_j$,
\[
  \int_\Gamma |\operatorname{Tr}_\Gamma \iota U|^2 \, d\sigma
  \leq \sum_j \sqrt{1 + M_j^2} \, C_j^2 \|U\|^2 ,
\]
so $C = (\sum_j \sqrt{1 + M_j^2} \, C_j^2)^{1/2}$. Additivity is \Cref{thm:trace-kernel-boundary}
for $U + W$ and $W$. For homogeneity, in the chart $j$ the boundary trace is
$\operatorname{Tr}_j$ (\Cref{thm:boundary-trace-charts}), which commutes with complex scalars
(\Cref{lem:trace-linear}). Almost everywhere on $\Gamma$ is $\sigma$-almost everywhere
(\Cref{thm:surface-measure-null-sets}), and the boundary traces of two resolutions of unity agree
almost everywhere on $\Gamma$ (\Cref{thm:boundary-trace-charts}).
\end{proof}

\begin{definition}[The traces into $L_2(\Gamma)$]
\label{def:trace-l2}
\lean{Rd.BoundedLipschitzDomain.traceL2, Rd.BoundedLipschitzDomain.boundaryTraceL2}
\leanok
\uses{thm:ambient-trace-l2, def:trace-space}
The trace $\tau : H^{1/2}(\Gamma) \to L_2(\Gamma)$ is $\tau t = T_\Gamma E t$, reading
$E t \in H^1(\Omega)$ as an element of $H^1(V)$, and the trace
$\gamma_{L_2} : H^1(\Omega) \to L_2(\Gamma)$ is $\gamma_{L_2} = \tau \circ \gamma$. Both are
bounded linear operators.
\end{definition}

\begin{theorem}[The traces in $L_2(\Gamma)$]
\label{thm:trace-l2}
\lean{Rd.BoundedLipschitzDomain.traceL2_apply_ae, Rd.BoundedLipschitzDomain.traceL2_injective,
  Rd.BoundedLipschitzDomain.boundaryTraceL2_apply_ae,
  Rd.BoundedLipschitzDomain.boundaryTraceL2_eq_ambientTraceL2,
  Rd.BoundedLipschitzDomain.ker_boundaryTraceL2,
  Rd.BoundedLipschitzDomain.exists_traceL2_ae_eq,
  Rd.BoundedLipschitzDomain.memLp_of_memBoundaryBesov}
\leanok
\uses{def:trace-l2, thm:ambient-trace-l2, thm:trace-space-functions, thm:trace-kernel-boundary,
  thm:surface-measure-null-sets}
Let $\psi$ be a resolution of unity on $\Gamma$, with $B^{1/2}_{2,2}(\Gamma)$ and the boundary
trace taken with respect to $\psi$. For every $t \in H^{1/2}(\Gamma)$,
$\tau t = \operatorname{Tr}_\Gamma \iota E t$ $\sigma$-almost everywhere, and $\tau$ is injective.
For every $u \in H^1(\Omega)$, $\gamma_{L_2} u = \operatorname{Tr}_\Gamma \iota u$
$\sigma$-almost everywhere, $\gamma_{L_2} u = T_\Gamma u$ with $u$ read as an element of $H^1(V)$,
and $\gamma_{L_2} u = 0$ if and only if $u \in N_\Gamma$. For every
$g \in B^{1/2}_{2,2}(\Gamma)$ there is $t \in H^{1/2}(\Gamma)$ with $\tau t = g$
$\sigma$-almost everywhere; in particular $g \in L_2(\Gamma)$.
\end{theorem}

\begin{proof}
\leanok
The first statement is \Cref{thm:ambient-trace-l2}. If $\tau t = \tau t'$, then
$\operatorname{Tr}_\Gamma \iota E t = \operatorname{Tr}_\Gamma \iota E t'$ $\sigma$-almost
everywhere, hence almost everywhere on $\Gamma$ (\Cref{thm:surface-measure-null-sets}), so
$t = t'$ (\Cref{thm:trace-space-functions}). By \Cref{thm:trace-space-functions},
$\operatorname{Tr}_\Gamma \iota E \gamma u = \operatorname{Tr}_\Gamma \iota u$ almost everywhere
on $\Gamma$, which gives the formula for $\gamma_{L_2}$; with the formula for $T_\Gamma$
(\Cref{thm:ambient-trace-l2}), $\gamma_{L_2} u = T_\Gamma u$. Then $\gamma_{L_2} u = 0$ if and
only if
$\operatorname{Tr}_\Gamma \iota u = 0$ almost everywhere on $\Gamma$, that is $u \in N_\Gamma$
(\Cref{thm:trace-kernel-boundary}). The last statement is \Cref{thm:trace-space-functions}.
\end{proof}

So $\gamma_{L_2}$ is the trace $H^1(\Omega) \to L_2(\Gamma)$ as a bounded operator, with kernel
$N_\Gamma$, and $\tau$ identifies $H^{1/2}(\Gamma)$ with $B^{1/2}_{2,2}(\Gamma)$ inside
$L_2(\Gamma)$.

\paragraph{The unit normal.}

\begin{definition}[The outward unit normal field]
\label{def:boundary-normal}
\lean{Rd.BoundedLipschitzDomain.normal}
\leanok
\uses{def:graph-normal, def:trace-space-setting}
For $y \in V$, $\nu(y) = \sum_j \psi^0_j(y) \, \nu_j(P_j y)$, with the outward unit normals
$\nu_j$ of the graph domains $\Omega_j$ (\Cref{def:graph-normal}).
\end{definition}

\begin{lemma}[Properties of the normal field]
\label{lem:boundary-normal}
\lean{Rd.BoundedLipschitzDomain.measurable_normal, Rd.BoundedLipschitzDomain.norm_normal_le}
\leanok
\uses{def:boundary-normal, lem:graph-normal}
$\nu$ is measurable, and $\|\nu(y)\| \leq 1$ for every $y \in V$.
\end{lemma}

\begin{proof}
\leanok
The $\psi^0_j$ are continuous and the $\nu_j$ measurable (\Cref{lem:graph-normal}). For the bound,
$\|\nu(y)\| \leq \sum_j \psi^0_j(y) \|\nu_j(P_j y)\| = \sum_j \psi^0_j(y) \leq 1$.
\end{proof}

For $y \in \Gamma$, $\sum_j \psi^0_j(y) = 1$, so $\nu(y)$ is a convex combination of the chart
normals $\nu_j(P_j y)$ of the charts with $\psi^0_j(y) > 0$. These agree $\sigma$-almost
everywhere by the Gauss--Green theorem in two charts (\Cref{thm:chart-normals-agree}), so that
$\nu = \nu_j \circ P_j$ $\sigma$-almost everywhere on $B_j$ (\Cref{thm:boundary-normal-chart}).

\section{Bounded Lipschitz domains from local graphs}
\label{sec:local-lipschitz-graphs}

The bounded Lipschitz domains of \Cref{def:bounded-lipschitz} carry their balls and graph domains
as data. This section starts from the usual definition, a condition on a set, and shows that in
finite dimension it describes exactly the carriers of these structures; the open balls are
examples. In this section $V$ is any real inner product space; finite dimension enters only
through the compactness of $\partial\Omega$. For $\nu \in V$, $P_\nu$ is the orthogonal
projection onto $\nu^\perp$, a hyperplane for a unit vector $\nu$.

\begin{definition}[Sets that are Lipschitz near a point]
\label{def:lipschitz-near-point}
\lean{Rd.IsLipschitzDomainAt, Rd.IsBoundedLipschitzDomain}
\leanok
Let $\Omega \subseteq V$ and $z \in V$. The set $\Omega$ is Lipschitz near $z$ if there are
$r > 0$, a unit vector $\nu \in V$, $L \geq 0$ and an $L$-Lipschitz function
$\varphi : \nu^\perp \to \R$ with
\[
  \Omega \cap B(z, r)
  = \{x \in V : \varphi(P_\nu(x - z)) < \langle \nu, x - z \rangle\} \cap B(z, r) .
\]
The set $\Omega$ is a bounded Lipschitz set if it is open and bounded and Lipschitz near every
$z \in \partial\Omega$.
\end{definition}

Near $z$, $\Omega$ lies above the graph of $\varphi$ in the direction $\nu$. This is the local
condition of \cite[Definition 1.2.1.1]{Grisvard} for open sets with Lipschitz boundary, with
boundedness added, the rotation of coordinates replaced by the direction $\nu$, and balls in place
of hypercubes. At a boundary point $\varphi(0) = 0$, so a small cylinder around $z$ lies in the
ball and a small ball lies in each cylinder; both forms define the same sets.

\begin{lemma}[The projection onto a hyperplane]
\label{lem:hyperplane-projection}
\lean{Rd.coe_orthogonalProjectionOnto_orthogonalComplement_singleton,
  Rd.norm_sub_smul_sq_eq_norm_orthogonalProjectionOnto_sq_add}
\leanok
For every unit vector $\nu \in V$, every $w \in V$ and every $a \in \R$,
$P_\nu w = w - \langle \nu, w \rangle \nu$ and
\[
  \|w - a \nu\|^2 = \|P_\nu w\|^2 + (\langle \nu, w \rangle - a)^2 .
\]
\end{lemma}

\begin{proof}
\leanok
The first formula is Mathlib's projection onto a line
(\texttt{Submodule.\allowbreak starProjection\_\allowbreak unit\_\allowbreak singleton}),
subtracted from the identity. Then
$w - a \nu = P_\nu w + (\langle \nu, w \rangle - a) \nu$ with $P_\nu w \perp \nu$, and Pythagoras
gives the identity.
\end{proof}

\begin{lemma}[Local graphs and Lipschitz graph domains]
\label{lem:local-graph-domain}
\lean{Rd.LipschitzGraph.ofLocal, Rd.LipschitzGraph.domain_ofLocal,
  Rd.LipschitzGraph.contDiff_phi_ofLocal, Rd.LipschitzGraph.localPhi,
  Rd.LipschitzGraph.lipschitzWith_localPhi, Rd.LipschitzGraph.contDiff_localPhi,
  Rd.LipschitzGraph.domain_eq_localPhi}
\leanok
\uses{def:lipschitz-graph, def:graph-map, lem:hyperplane-projection}
For every $z \in V$, every unit vector $\nu \in V$, every $L \geq 0$ and every $L$-Lipschitz
$\varphi : \nu^\perp \to \R$, the function
$x \mapsto \langle z, \nu \rangle + \varphi(P_\nu(x - z))$ is $L$-Lipschitz and constant along
$\nu$, and its Lipschitz graph domain with direction $\nu$ is
$\{x \in V : \varphi(P_\nu(x - z)) < \langle \nu, x - z \rangle\}$; for every
$n \in \N \cup \{\infty, \omega\}$ (Mathlib's orders of smoothness, $\omega$ meaning analytic), if
$\varphi$ is of class $C^n$, then so is this function. For every Lipschitz graph domain
$\Omega = \{x : \psi(x) < \langle x, e \rangle\}$ with graph function $\psi$ and constant $M$, and
every $z \in V$, the function $\psi_z(y) = \psi(z + y) - \langle z, e \rangle$ on $H = e^\perp$ is
$M$-Lipschitz and
\[
  \Omega = \{x \in V : \psi_z(P_e(x - z)) < \langle e, x - z \rangle\} ,
\]
with $P_e$ the orthogonal projection onto $H$; for every $n \in \N \cup \{\infty, \omega\}$, if
$\psi$ is of class $C^n$, then so is $\psi_z$.
\end{lemma}

\begin{proof}
\leanok
$P_\nu \nu = 0$ (Mathlib's \texttt{Submodule.\allowbreak orthogonalProjectionOnto\_\allowbreak
orthogonalComplement\_\allowbreak singleton\_\allowbreak eq\_\allowbreak zero}) gives the
constancy along $\nu$, and $\|P_\nu y\| \leq \|y\|$ the Lipschitz bound; the domain is
$\langle z, \nu \rangle + \varphi(P_\nu(x - z)) < \langle x, \nu \rangle$, which is the stated
inequality. For the second statement,
$z + P_e(x - z) = x - \langle x - z, e \rangle e$, so $\psi(z + P_e(x - z)) = \psi(x)$ since
$\psi$ is constant along $e$. Both functions are compositions of $\varphi$ or $\psi$ with affine
maps, so they keep the class $C^n$.
\end{proof}

\begin{theorem}[Lipschitz near a point by graph domains]
\label{thm:lipschitz-near-point}
\lean{Rd.isLipschitzDomainAt_of_ball_inter_eq, Rd.IsLipschitzDomainAt.exists_ball_inter_eq,
  Rd.isLipschitzDomainAt_iff, Rd.IsLipschitzDomainAt.congr_ball,
  Rd.localGraph_inter_ball_eq_of_eqOn, Rd.isLipschitzDomainAt_of_lipschitzOnWith}
\leanok
\uses{def:lipschitz-near-point, lem:local-graph-domain}
For every $z \in V$, every $\nu \in V$, every $\rho \in \R$ and all
$\varphi, \psi : \nu^\perp \to \R$ that agree on $\{y \in \nu^\perp : \|y\| < \rho\}$, the sets
$\{x : \varphi(P_\nu(x - z)) < \langle \nu, x - z \rangle\}$ and
$\{x : \psi(P_\nu(x - z)) < \langle \nu, x - z \rangle\}$ have the same intersection with
$B(z, \rho)$. For every $\Omega \subseteq V$ and every $z \in V$, $\Omega$ is Lipschitz near $z$
if and only if there are $r > 0$ and a Lipschitz graph domain $\Omega'$ with
$B(z, r) \cap \Omega = B(z, r) \cap \Omega'$. For all $\Omega, \Omega' \subseteq V$, every
$z \in V$ and every $\rho > 0$ with $B(z, \rho) \cap \Omega = B(z, \rho) \cap \Omega'$, if
$\Omega'$ is Lipschitz near $z$, then so is $\Omega$. For every $\Omega \subseteq V$, every
$z \in V$,
every $r > 0$, every unit vector $\nu \in V$, every $L \geq 0$ and every
$\varphi : \nu^\perp \to \R$ that is $L$-Lipschitz on $\{y \in \nu^\perp : \|y\| < r\}$ with
$\Omega \cap B(z, r) = \{x : \varphi(P_\nu(x - z)) < \langle \nu, x - z \rangle\} \cap B(z, r)$, the
set $\Omega$ is Lipschitz near $z$.
\end{theorem}

\begin{proof}
\leanok
The first statement holds since $\|P_\nu(x - z)\| \leq \|x - z\| < \rho$ for $x \in B(z, \rho)$.
Both directions of the equivalence are \Cref{lem:local-graph-domain}. For the third statement,
shrink the radius $r$ of a ball on which $\Omega'$ agrees with a Lipschitz graph domain to
$\min(r, \rho)$. For the last statement, McShane's extension (Mathlib's
\texttt{LipschitzOnWith.\allowbreak extend\_\allowbreak real}) gives an $L$-Lipschitz
$\psi : \nu^\perp \to \R$ equal to $\varphi$ on the ball of radius $r$, and the first statement
applies.
\end{proof}

So the condition can be checked with a function that is Lipschitz only on the part of the
hyperplane that the ball sees.

\begin{theorem}[Bounded Lipschitz sets and bounded Lipschitz domains]
\label{thm:bounded-lipschitz-sets}
\lean{Rd.exists_finset_cover_frontier, Rd.frontierCenters, Rd.frontier_subset_frontierCenters,
  Rd.frontierCenterAt, Rd.BoundedLipschitzDomain.ofCharts,
  Rd.BoundedLipschitzDomain.carrier_ofCharts, Rd.BoundedLipschitzDomain.exists_graph_ofCharts_eq,
  Rd.IsBoundedLipschitzDomain.exists_chart, Rd.IsBoundedLipschitzDomain.chart,
  Rd.IsBoundedLipschitzDomain.chart_fst_pos, Rd.IsBoundedLipschitzDomain.ball_inter_chart,
  Rd.BoundedLipschitzDomain.ofIsBoundedLipschitzDomain,
  Rd.BoundedLipschitzDomain.carrier_ofIsBoundedLipschitzDomain,
  Rd.BoundedLipschitzDomain.exists_ball_inter_eq_graph,
  Rd.BoundedLipschitzDomain.isLipschitzDomainAt,
  Rd.BoundedLipschitzDomain.isBoundedLipschitzDomain, Rd.isBoundedLipschitzDomain_iff_exists}
\leanok
\uses{def:bounded-lipschitz, thm:lipschitz-near-point}
For every bounded Lipschitz domain $\Omega$ with balls $B_j$ and graph domains $\Omega_j$ and every
$z \in \partial\Omega$, there are $j$ and $\rho > 0$ with
$B(z, \rho) \cap \Omega = B(z, \rho) \cap \Omega_j$; in particular $\Omega$ is Lipschitz near $z$,
and $\Omega$ is a bounded Lipschitz set. Let $V$ be finite-dimensional. For every bounded open
$\Omega \subseteq V$, every family of radii $r_z > 0$ and every family of Lipschitz graph domains
$\Omega'_z$, $z \in \partial\Omega$, with $B(z, r_z) \cap \Omega = B(z, r_z) \cap \Omega'_z$ for
every $z \in \partial\Omega$, there is a bounded Lipschitz domain with carrier $\Omega$ whose
graph domains are among the $\Omega'_z$. For every $\Omega \subseteq V$, $\Omega$ is a bounded
Lipschitz set if and only if it is the carrier of a bounded Lipschitz domain.
\end{theorem}

\begin{proof}
\leanok
Let $\Omega$ be a bounded Lipschitz domain and $z \in \partial\Omega$. Then $z \in B(c_j, \rho_j)$
for some $j$, the ball $B(z, \rho_j - \|z - c_j\|)$ lies in $B(c_j, \rho_j)$, and on it $\Omega$
agrees with $\Omega_j$; \Cref{thm:lipschitz-near-point} applies. For a family of charts, the
boundary of the bounded set $\Omega$ is compact in finite dimension, so finitely many of the balls
$B(z, r_z)$ cover it (Mathlib's \texttt{IsCompact.\allowbreak elim\_\allowbreak
finite\_\allowbreak subcover}); their centres, radii and graph domains form a bounded Lipschitz
domain with carrier $\Omega$. For a bounded Lipschitz set $\Omega$,
\Cref{thm:lipschitz-near-point} gives such a family of charts.
\end{proof}

\begin{definition}[The sphere profile]
\label{def:sphere-profile}
\lean{Rd.sphereProfile}
\leanok
For $R, s \in \R$ let $g_R(s) = R - \sqrt{R^2 - \min(s, R/2)^2}$.
\end{definition}

For $R > 0$, near a point $z$ of the sphere $\|x - c\| = R$, the ball $B(c, R)$ is the region
above the graph of $y \mapsto g_R(\|y\|)$ over the tangent hyperplane (\Cref{thm:ball-lipschitz}).

\begin{lemma}[The sphere profile is $1$-Lipschitz]
\label{lem:sphere-profile}
\lean{Rd.abs_sqrt_sq_sub_sq_sub_le, Rd.abs_sphereProfile_sub_le,
  Rd.lipschitzWith_sphereProfile_norm}
\leanok
\uses{def:sphere-profile}
For every $R > 0$ and all $s_1, s_2 \in [0, R/2]$,
$|\sqrt{R^2 - s_1^2} - \sqrt{R^2 - s_2^2}| \leq |s_1 - s_2|$. For every $R > 0$ and all
$s_1, s_2 \geq 0$, $|g_R(s_1) - g_R(s_2)| \leq |s_1 - s_2|$. For every $R > 0$ and every normed
group $E$, the function $y \mapsto g_R(\|y\|)$ on $E$ is $1$-Lipschitz.
\end{lemma}

\begin{proof}
\leanok
Both square roots are at least $R/2$, and
\[
  |\sqrt{R^2 - s_1^2} - \sqrt{R^2 - s_2^2}| \, (\sqrt{R^2 - s_1^2} + \sqrt{R^2 - s_2^2})
  = |s_1 - s_2| \, (s_1 + s_2)
\]
with $s_1 + s_2 \leq R$. The function $s \mapsto \min(s, R/2)$ is $1$-Lipschitz, and so is the
norm. The bound for the square roots is the argument of the Lean project
\texttt{quadratic-\allowbreak forms-\allowbreak sobolev-\allowbreak lean}, available at
\url{https://github.com/dbenbenn/quadratic-forms-sobolev-lean}.
\end{proof}

\begin{theorem}[The ball]
\label{thm:ball-lipschitz}
\lean{Rd.mem_ball_iff_sphereProfile_lt, Rd.isLipschitzDomainAt_ball,
  Rd.norm_inv_smul_sub_eq_one, Rd.inter_ball_eq_sphereProfile,
  Rd.isBoundedLipschitzDomain_ball, Rd.BoundedLipschitzDomain.ball,
  Rd.BoundedLipschitzDomain.carrier_ball}
\leanok
\uses{lem:hyperplane-projection, def:sphere-profile, lem:sphere-profile,
  def:lipschitz-near-point, thm:bounded-lipschitz-sets}
Let $c \in V$.
\begin{enumerate}
  \item For every $R > 0$, every unit vector $\nu \in V$ and all $z, x \in V$ with $c - z = R \nu$
    and $\|x - z\| < R/2$, $x \in B(c, R)$ if and only if
    $g_R(\|P_\nu(x - z)\|) < \langle \nu, x - z \rangle$.
  \item For every $R > 0$ and every $z \in V$ with $\|z - c\| = R$, the vector
    $\nu = (c - z)/R$ is a unit vector,
    \[
      B(c, R) \cap B(z, R/2)
      = \{x : g_R(\|P_\nu(x - z)\|) < \langle \nu, x - z \rangle\} \cap B(z, R/2) ,
    \]
    and the ball $B(c, R)$ is Lipschitz near $z$.
  \item For every $R \in \R$, the ball $B(c, R)$ is a bounded Lipschitz set; if $V$ is
    finite-dimensional, it is the carrier of a bounded Lipschitz domain.
\end{enumerate}
\end{theorem}

\begin{proof}
\leanok
Let $t = \langle \nu, x - z \rangle$ and $p = P_\nu(x - z)$. By \Cref{lem:hyperplane-projection},
$\|x - z\|^2 = \|p\|^2 + t^2$ and $\|x - c\|^2 = \|x - z - R\nu\|^2 = \|p\|^2 + (t - R)^2$. From
$\|x - z\| < R/2$, $\|p\| < R/2$ and $t < R/2$, so $g_R(\|p\|) = R - \sqrt{R^2 - \|p\|^2}$ and
$R - t > 0$; then $\|x - c\| < R$ if and only if $(R - t)^2 < R^2 - \|p\|^2$, that is
$R - t < \sqrt{R^2 - \|p\|^2}$. For $z$ on the sphere take radius $R/2$, $\nu = (c - z)/R$, so
that $c - z = R\nu$ and $\|\nu\| = \|z - c\|/R = 1$, and $\varphi(y) = g_R(\|y\|)$, which is
$1$-Lipschitz by \Cref{lem:sphere-profile}; the first statement gives the set identity and the
condition of \Cref{def:lipschitz-near-point}. For $R > 0$ the boundary of the ball lies in the
sphere $\|z - c\| = R$; for $R \leq 0$ the ball is empty. The last statement is
\Cref{thm:bounded-lipschitz-sets}.
\end{proof}

\section{Bounded convex open sets}
\label{sec:convex-domains}

Every bounded open convex set is a bounded Lipschitz set \cite{Grisvard}. Mathlib's gauge of a
convex set describes it as a sublevel set; the graph form of \Cref{def:lipschitz-near-point}
comes from the lengths of parallel rays instead. In this section $V$ is any real inner product
space; for a unit vector $\nu \in V$, $P_\nu$ is the orthogonal projection onto $\nu^\perp$.

\begin{definition}[The ray length]
\label{def:ray-length}
\lean{Rd.rayLength}
\leanok
For every $\Omega \subseteq V$ and all $a, \nu, y \in V$, the ray length is
\[
  \ell(y) = \ell_{\Omega, a, \nu}(y) = \sup\, \{s \in \R : a + y - s \nu \in \Omega\} ,
\]
with the convention of Lean's real numbers that the supremum of a set that is empty or unbounded
above is $0$.
\end{definition}

For a bounded convex open set $\Omega$, a unit vector $\nu$ and $a + y \in \Omega$, the points of
the ray from $a + y$ in the direction $-\nu$ that lie in $\Omega$ are exactly those with parameter
$s < \ell(y)$ (\Cref{lem:ray-length}).

\begin{lemma}[The ray in a convex open set]
\label{lem:ray-length}
\lean{Rd.nonempty_raySet, Rd.le_norm_add_of_sub_smul_mem, Rd.bddAbove_rayLength, Rd.le_rayLength,
  Rd.rayLength_le, Rd.mem_iff_lt_rayLength}
\leanok
\uses{def:ray-length}
For every $\Omega \subseteq V$ and all $a, \nu, y \in V$ with $a + y \in \Omega$, the set of
\Cref{def:ray-length} contains $0$. Let $\Omega \subseteq V$ and $C \in \R$ with $\|x\| \leq C$ for
every $x \in \Omega$. For every unit vector $\nu \in V$, all $a, y \in V$ and every $s \in \R$ with
$a + y - s\nu \in \Omega$, $s \leq \|a + y\| + C$ and $s \leq \ell(y)$. For every unit vector
$\nu \in V$ and all $a, y \in V$ with $a + y \in \Omega$, $\ell(y) \leq 2C$. If $\Omega$ is
moreover convex and open, then for every unit vector $\nu \in V$, all $a, y \in V$ with
$a + y \in \Omega$ and every $s \geq 0$, $a + y - s\nu \in \Omega$ if and only if $s < \ell(y)$.
\end{lemma}

\begin{proof}
\leanok
The triangle inequality gives $|s| = \|s\nu\| \leq \|a + y\| + \|a + y - s\nu\|$, so the set of
\Cref{def:ray-length} is bounded above and $\ell(y)$ is its least upper bound; for
$a + y \in \Omega$ it contains $0$ and is bounded by $2C$. If $a + y - s\nu \in \Omega$, openness
gives $\varepsilon > 0$ with $a + y - (s + \varepsilon/2)\nu \in \Omega$, so $s < \ell(y)$. If
$s < \ell(y)$, there is $b > s$ with $a + y - b\nu \in \Omega$ (Mathlib's
\texttt{exists\_\allowbreak lt\_\allowbreak of\_\allowbreak lt\_\allowbreak csSup}), and
$a + y - s\nu = (1 - s/b)(a + y) + (s/b)(a + y - b\nu) \in \Omega$ by convexity.
\end{proof}

The convex hull of a point of $\Omega$ and the ball $B(a, \delta)$ is a cone around the segment
between them. It contains the shortened rays from nearby starting points, which makes the ray
length Lipschitz near $a$.

\begin{lemma}[The cone estimate]
\label{lem:ray-length-lipschitz}
\lean{Rd.sub_smul_mem_of_norm_lt, Rd.rayLength_le_rayLength_add, Rd.cone_const_nonneg,
  Rd.lipschitzOnWith_rayLength, Rd.lipschitzOnWith_sub_rayLength}
\leanok
\uses{lem:ray-length}
Let $\Omega \subseteq V$ be convex and $C \in \R$ with $\|x\| \leq C$ for every $x \in \Omega$, and
let $a \in V$ and $\delta > 0$ with $B(a, \delta) \subseteq \Omega$. Then $4C/\delta \geq 0$. For
every $\nu \in V$, all $y_1, y_2 \in V$ with $\|y_1\|, \|y_2\| < \delta/4$ and every $s \in \R$
with $a + y_1 - s\nu \in \Omega$, $a + y_2 - (1 - \mu) s \nu \in \Omega$, where
$\mu = 2\|y_1 - y_2\|/\delta$. For every unit vector $\nu \in V$ and all $y_1, y_2 \in V$ with
$\|y_1\|, \|y_2\| < \delta/4$,
\[
  \ell(y_1) \leq \ell(y_2) + \frac{4C}{\delta} \|y_1 - y_2\| .
\]
For every unit vector $\nu \in V$, the ray length is $4C/\delta$-Lipschitz on $B(0, \delta/4)$,
and for every $d \in \R$ the function $y \mapsto d - \ell(y)$ is $4C/\delta$-Lipschitz on
$\{y \in \nu^\perp : \|y\| < \delta/4\}$.
\end{lemma}

\begin{proof}
\leanok
The constant is nonnegative since $C \geq \|a\| \geq 0$. For $y_1 = y_2$ the membership statement
is the hypothesis. Otherwise $0 < \mu \leq 1$, since $\|y_1 - y_2\| < \delta/2$, and
$b = y_1 + \mu^{-1}(y_2 - y_1)$ has $\|b\| \leq \|y_1\| + \delta/2 < \delta$, so
$a + b \in \Omega$; then
\[
  a + y_2 - (1 - \mu) s \nu = (1 - \mu)(a + y_1 - s\nu) + \mu (a + b) \in \Omega .
\]
Hence $(1 - \mu) s \leq \ell(y_2)$ by \Cref{lem:ray-length}, and $s \leq 2C$ gives
$s \leq \ell(y_2) + 2\mu C = \ell(y_2) + (4C/\delta)\|y_1 - y_2\|$; the supremum over $s$ is the
bound for $\ell(y_1)$. The Lipschitz bounds follow, with the roles of $y_1$ and $y_2$ exchanged for
$d - \ell$ (Mathlib's
\texttt{LipschitzOnWith.\allowbreak of\_\allowbreak le\_\allowbreak add\_\allowbreak mul}).
\end{proof}

\begin{theorem}[Bounded convex open sets]
\label{thm:convex-lipschitz}
\lean{Rd.mem_iff_sub_rayLength_lt, Rd.isLipschitzDomainAt_of_convex,
  Rd.isBoundedLipschitzDomain_of_convex, Rd.BoundedLipschitzDomain.ofConvex,
  Rd.BoundedLipschitzDomain.carrier_ofConvex}
\leanok
\uses{lem:hyperplane-projection, lem:ray-length, lem:ray-length-lipschitz,
  def:lipschitz-near-point, thm:lipschitz-near-point, thm:bounded-lipschitz-sets}
Let $\Omega \subseteq V$ be convex, open and bounded.
\begin{enumerate}
  \item For every $C \in \R$ with $\|x\| \leq C$ for every $x \in \Omega$, all $a \in V$ and
    $\delta > 0$ with $B(a, \delta) \subseteq \Omega$, every unit vector $\nu \in V$, all $z \in V$
    and $d \geq \delta$ with $a - z = d\nu$, and every $x \in V$ with $\|x - z\| < \delta/4$,
    $x \in \Omega$ if and only if $d - \ell(P_\nu(x - z)) < \langle \nu, x - z \rangle$.
  \item For every $z \in \partial\Omega$, $\Omega$ is Lipschitz near $z$.
  \item $\Omega$ is a bounded Lipschitz set; if $V$ is finite-dimensional, it is the carrier of a
    bounded Lipschitz domain.
\end{enumerate}
\end{theorem}

\begin{proof}
\leanok
Let $t = \langle \nu, x - z \rangle$. By \Cref{lem:hyperplane-projection},
$x = a + P_\nu(x - z) - (d - t)\nu$; moreover $\|P_\nu(x - z)\| \leq \|x - z\| < \delta/4$, so
$a + P_\nu(x - z) \in \Omega$, and $t \leq \|x - z\| < \delta \leq d$. \Cref{lem:ray-length} with
$s = d - t$ gives the first statement. For $z \in \partial\Omega$, the set $\Omega$ is nonempty
since $z$ lies in its closure; choose $a \in \Omega$ and $\delta > 0$ with $B(a, \delta) \subseteq
\Omega$. Since $\Omega$ is open, $z \notin \Omega$, so $d = \|a - z\| \geq \delta$. With
$\nu = (a - z)/d$, the first statement and \Cref{lem:ray-length-lipschitz} give the hypotheses of
the second statement of \Cref{thm:lipschitz-near-point} for $r = \delta/4$ and
$\varphi(y) = d - \ell(y)$. The last statement follows by \Cref{thm:bounded-lipschitz-sets}.
\end{proof}

Near $z$ the graph function $y \mapsto \|a - z\| - \ell(y)$ is $4C/\delta$-Lipschitz: the
Lipschitz constant of the boundary is at most four times the ratio of a bound $C$ for the norms on
$\Omega$ to the radius $\delta$ of a ball inside $\Omega$. The local statement needs no
boundedness: half-spaces and epigraphs of convex functions are examples.

\begin{corollary}[Open convex sets near a frontier point]
\label{cor:open-convex-lipschitz-near-point}
\lean{Rd.isLipschitzDomainAt_of_convex_of_isOpen}
\leanok
\uses{thm:convex-lipschitz, thm:lipschitz-near-point}
For every convex open set $\Omega \subseteq V$ and every $z \in \partial\Omega$, $\Omega$ is
Lipschitz near $z$.
\end{corollary}

\begin{proof}
\leanok
The set $\Omega' = \Omega \cap B(z, 1)$ is convex, open and bounded, and agrees with $\Omega$ on
$B(z, 1)$. It does not contain $z$, and $z \in \overline{\Omega} \cap B(z, 1) \subseteq
\overline{\Omega'}$ (Mathlib's \texttt{IsOpen.\allowbreak closure\_\allowbreak inter}), so
$z \in \partial\Omega'$. By \Cref{thm:convex-lipschitz}, $\Omega'$ is Lipschitz near $z$, and by
\Cref{thm:lipschitz-near-point} so is $\Omega$.
\end{proof}

\begin{corollary}[The open unit cube]
\label{cor:unit-cube-lipschitz}
\lean{Rd.isBoundedLipschitzDomain_unitCube, Rd.BoundedLipschitzDomain.unitCube,
  Rd.BoundedLipschitzDomain.carrier_unitCube}
\leanok
\uses{thm:convex-lipschitz}
For every $n \in \N$, the open unit cube $(0, 1)^n \subseteq \R^n$ is a bounded Lipschitz
set and the carrier of a bounded Lipschitz domain.
\end{corollary}

\begin{proof}
\leanok
The cube is convex, open as the intersection of the $n$ open slabs $\{x : 0 < x_i < 1\}$, and
bounded, since $\|x\| \leq \sqrt{n}$ on it (Mathlib's
\texttt{EuclideanSpace.\allowbreak norm\_\allowbreak eq}); \Cref{thm:convex-lipschitz} applies.
\end{proof}

\chapter{Traces on Lipschitz domains}
\label{ch:traces-lipschitz}

In this chapter $V$ is a finite-dimensional real inner product space of dimension $n$ with
Lebesgue measure, and $1 \leq p < \infty$. The trace on the boundary of a Lipschitz graph domain
is built from the atomic decomposition of \Cref{ch:rd}: the cell atoms, composed with the graph
map $\gamma$ over the hyperplane $H = e^\perp$, are atoms on $H$ (\Cref{lem:atom-traces}). This
gives the trace theorem on Lipschitz graph domains for $0 < s < 1$
(\Cref{thm:trace-graph} and \Cref{thm:graph-trace-space}), on the half-space for every $s > 0$
(\Cref{thm:half-space-trace-space}), and, through the change of chart
(\Cref{thm:change-of-chart}), on bounded Lipschitz domains, whose boundary Besov spaces are
defined in charts by first differences (\Cref{thm:boundary-trace-space}).

\section{Traces on Lipschitz graph domains}

Let $\Omega = \{x : \phi(x) < \langle x, e \rangle\}$ be a Lipschitz graph domain
(\Cref{def:lipschitz-graph}) with Lipschitz constant $M$. Its boundary is the graph of $\phi$ over
the hyperplane $H = e^\perp$, of dimension $n - 1$ with its Lebesgue measure, and the trace is a
function on $H$.

\begin{definition}[The graph map]
\label{def:graph-map}
\lean{Rd.LipschitzGraph.hyperplane, Rd.LipschitzGraph.e_ne_zero, Rd.LipschitzGraph.inner_e_self,
  Rd.LipschitzGraph.inner_coe_e, Rd.LipschitzGraph.sub_inner_smul_mem,
  Rd.LipschitzGraph.finrank_hyperplane_add_one, Rd.LipschitzGraph.finrank_hyperplane_eq_sub_one,
  Rd.LipschitzGraph.finrank_hyperplane_eq,
  Rd.LipschitzGraph.orthogonal_hyperplane, Rd.LipschitzGraph.proj, Rd.LipschitzGraph.coe_proj,
  Rd.LipschitzGraph.inner_coe_eq_inner_proj, Rd.LipschitzGraph.norm_proj_le,
  Rd.LipschitzGraph.proj_coe, Rd.LipschitzGraph.proj_e, Rd.LipschitzGraph.graphMap,
  Rd.LipschitzGraph.proj_graphMap, Rd.LipschitzGraph.inner_graphMap,
  Rd.LipschitzGraph.phi_graphMap, Rd.LipschitzGraph.norm_graphMap_sub_le,
  Rd.LipschitzGraph.continuous_graphMap, Rd.LipschitzGraph.norm_sub_proj_le,
  Rd.LipschitzGraph.graphMap_add_mem}
\leanok
\uses{def:lipschitz-graph, lem:graph-cone}
Let $P x = x - \langle x, e \rangle e$ be the orthogonal projection onto $H$ and
$\gamma(x') = x' + \phi(x') e$ for $x' \in H$. Then $H^\perp = \R e$,
$\langle a, v \rangle = \langle a, P v \rangle$ for all $a \in H$ and $v \in V$,
$P \gamma(x') = x'$,
$\langle \gamma(x'), e \rangle = \phi(x') = \phi(\gamma(x'))$,
$\|\gamma(x') - \gamma(y')\| \leq (1 + M) \|x' - y'\|$, $\|x' - Pc\| \leq \|\gamma(x') - c\|$, and
$\gamma(x') + K \subseteq \Omega$ for the cone $K$ of \Cref{lem:graph-cone}.
\end{definition}

The last property holds since a step $y \in K$ raises $\langle \cdot, e \rangle$ by more than
$M \|Py\|$ and $\phi$ by at most $M \|Py\|$.

\begin{lemma}[Traces of atoms]
\label{lem:atom-traces}
\lean{Rd.min_one_mul_le, Rd.LipschitzGraph.isRegularAtom_comp_graphMap}
\leanok
\uses{def:graph-map, lem:pointwise-atoms}
Let $a$ be continuous with $a = 0$ outside $\overline{B}(c, d 2^{-j})$, $|a| \leq A$ and
$|a(x + h) - a(x)| \leq A \min(1, 2^j \|h\|)$. Then $a \circ \gamma$ is a regular atom on $H$ of
level $j$, order and exponent $1$, centre $Pc$, radius factor $d$ and constant
$(1 + M) A (2 d^{n-1} |B_1|)^{1/p}$.
\end{lemma}

\begin{proof}
\leanok
As in the proof of \cite[Theorem 4.2]{SchneiderVybiral}: $\|\gamma(x' + h') - \gamma(x')\| \leq
(1 + M) \|h'\|$, and $\min(1, K t) \leq K \min(1, t)$ for $K \geq 1$; \Cref{lem:pointwise-atoms} on
$H$.
\end{proof}

\begin{definition}[Traces of the terms]
\label{def:trace-terms}
\lean{Rd.LipschitzGraph.traceFun, Rd.LipschitzGraph.traceCoef,
  Rd.LipschitzGraph.continuous_traceFun, Rd.LipschitzGraph.traceFun_eq_tsum,
  Rd.LipschitzGraph.enorm_traceCoef_le, Rd.LipschitzGraph.levelNorm_traceCoef_le}
\leanok
\uses{def:graph-map, lem:cell-decomposition, lem:cell-coefficients}
For a bump, $L \in \N$ and $h \in L_p$, the trace of the term $\eta_j * h$ is the continuous
function $x' \mapsto (\eta_j * h)(\gamma(x'))$ on $H$. With the coefficients $\lambda'_{j,m}$ equal
to $\lambda_{j,m}$ for the cells whose atoms meet the graph and $0$ otherwise,
$(\eta_j * h)(\gamma(x')) = \sum_m \lambda'_{j,m} a_{j,m}(\gamma(x'))$ and
$\|\lambda'_{j,\cdot}\|_{\ell_p} \leq 2^{jn/p} \|h\|_p$.
\end{definition}

\begin{lemma}[The trace atoms]
\label{lem:trace-atoms}
\lean{Rd.LipschitzGraph.norm_graphMap_sub_cellCenter_le, Rd.LipschitzGraph.exists_trace_atoms}
\leanok
\uses{def:trace-terms, lem:cell-atoms-regular, lem:dual-kernel-bounds, lem:atom-traces,
  lem:cells-partition}
There are $d \geq 0$, $C < \infty$ and $N$ such that for every level $j$ and $h \in L_p$ the
restrictions $a_{j,m} \circ \gamma$ of the cell atoms of $\eta_j * h$ are regular atoms on $H$ of
order and exponent $1$ with centres $P c_{j,m}$, and the centres $P c_{j,m}$ with
$\lambda'_{j,m} \neq 0$ have overlap at most $N$ on balls of radius $(d + 1) 2^{-j}$ in $H$.
\end{lemma}

\begin{proof}
\leanok
The cell atoms satisfy the pointwise bounds of \Cref{lem:atom-traces} by
\Cref{lem:cell-atoms-regular} and \Cref{lem:dual-kernel-bounds}. If $\lambda'_{j,m} \neq 0$, the
atom meets the graph at some $\gamma(y')$ with $\|\gamma(y') - c_{j,m}\| \leq d 2^{-j}$; if moreover
$\|x' - P c_{j,m}\| \leq (d + 1) 2^{-j}$, then
$\|\gamma(x') - c_{j,m}\| \leq (d + (1 + M)(2d + 1)) 2^{-j}$, and \Cref{lem:cells-partition} counts
these centres.
\end{proof}

\begin{lemma}[The traces of the terms in $L_p(H)$]
\label{lem:trace-terms-lp}
\lean{Rd.LipschitzGraph.eLpNorm_traceFun_le, Rd.LipschitzGraph.memLp_traceFun,
  Rd.LipschitzGraph.traceTerm, Rd.LipschitzGraph.traceTerm_ae}
\leanok
\uses{lem:trace-atoms, lem:level-sums}
There is $K < \infty$ with $\|(\eta_j * h) \circ \gamma\|_{L_p(H)} \leq K 2^{j/p} \|h\|_p$ for all
$j$ and $h \in L_p$.
\end{lemma}

\begin{proof}
\leanok
\Cref{lem:level-sums} on $H$, of dimension $n - 1$, with \Cref{lem:trace-atoms}:
$2^{-j(n-1)/p} \|\lambda'_{j,\cdot}\|_{\ell_p} \leq 2^{-j(n-1)/p} 2^{jn/p} \|h\|_p$.
\end{proof}

\begin{theorem}[The trace of a series]
\label{thm:trace-series}
\lean{Rd.LipschitzGraph.traceSeries, Rd.LipschitzGraph.summable_traceTerm,
  Rd.LipschitzGraph.tendsto_sum_traceTerm, Rd.LipschitzGraph.modNorm_traceSeries_le}
\leanok
\uses{lem:trace-terms-lp, lem:trace-atoms, thm:rd-atomic-synthesis}
Let $1 \leq p < \infty$, $0 < s < 1$ and $1 \leq q \leq \infty$. There is $K < \infty$ such that for
$h_j \in L_p$ with $X = \|(2^{j(s + 1/p)} \|h_j\|_p)_j\|_{\ell_q} < \infty$ the series
$\operatorname{Tr} = \sum_j (\eta_j * h_j) \circ \gamma$ converges absolutely in $L_p(H)$ and
\[
  \|\operatorname{Tr}\|_{L_p(H)} + \|(2^{ks} \omega_1(\operatorname{Tr}, 2^{-k})_p)_k\|_{\ell_q}
  \leq K X.
\]
\end{theorem}

\begin{proof}
\leanok
The terms have $\|(\eta_j * h_j) \circ \gamma\|_{L_p(H)} \leq K 2^{j/p} \|h_j\|_p \leq K X 2^{-js}$
(\Cref{lem:trace-terms-lp}). \Cref{thm:rd-atomic-synthesis} on $H$ with the atoms of
\Cref{lem:trace-atoms} and $\sigma = r = 1 > s$ bounds the left side by the coefficient
quasi-norm on $H$, and since $s - (n - 1)/p = (s + 1/p) - n/p$,
$2^{j(s - (n-1)/p)} \|\lambda'_{j,\cdot}\|_{\ell_p} \leq 2^{j(s + 1/p)} \|h_j\|_p$.
\end{proof}

\begin{theorem}[The trace theorem on Lipschitz graph domains]
\label{thm:trace-graph}
\lean{Rd.LipschitzGraph.trace, Rd.LipschitzGraph.modNorm_trace_le}
\leanok
\uses{thm:trace-series, lem:intrinsic-upper}
The trace of $f \in L_p$ is
$\operatorname{Tr} f = \sum_j (\eta_j * 1_\Omega (k_j * f)) \circ \gamma$, with the kernels of the
cone bump. Let $1 \leq p < \infty$, $0 < s < 1$, $1 \leq q \leq \infty$ and $L \in \N$. There is
$K < \infty$ such that
\begin{multline*}
  \|\operatorname{Tr} f\|_{L_p(H)} + \|(2^{ks} \omega_1(\operatorname{Tr} f, 2^{-k})_p)_k\|_{\ell_q}
  \\ \leq K \bigl( \|f\|_{L_p(\Omega)} +
    \|(2^{k(s + 1/p)} \omega_L(f, 2^{-k}, \Omega)_p)_k\|_{\ell_q} \bigr)
\end{multline*}
whenever the right side is finite.
\end{theorem}

\begin{proof}
\leanok
This is \cite[Theorem 4.2]{SchneiderVybiral} for $p \geq 1$. \Cref{thm:trace-series} with
$h_j = 1_\Omega (k_j * f)$, and \Cref{lem:intrinsic-upper} with smoothness $s + 1/p$.
\end{proof}

\begin{lemma}[Linearity and continuity of the trace]
\label{lem:trace-linear}
\lean{Rd.LipschitzGraph.traceTerm_sub, Rd.LipschitzGraph.trace_sub,
  Rd.LipschitzGraph.traceTerm_smul, Rd.LipschitzGraph.trace_smul,
  Rd.LipschitzGraph.enorm_trace_le, Rd.LipschitzGraph.tendsto_trace_of_modNorm_sub_le,
  Rd.LipschitzGraph.enorm_trace_le_modNormOn}
\leanok
\uses{thm:trace-graph, lem:trace-terms-lp, lem:local-means-upper, lem:intrinsic-upper, def:conv}
For every $c \in \C$ and every $f \in L_p(V)$, $\operatorname{Tr}(c f) = c \operatorname{Tr} f$.
The trace is additive on the functions with a finite intrinsic norm of smoothness $s + 1/p$,
$s > 0$, and there is $K < \infty$ with
$\|\operatorname{Tr} f\|_{L_p(H)} \leq K (\|f\|_p + \sum_k 2^{k/p} \omega_L(f, 2^{-k})_p)$ and
$\|\operatorname{Tr} f\|_{L_p(H)} \leq
K (\|f\|_{L_p(\Omega)} + \sum_k 2^{k/p} \omega_L(f, 2^{-k}, \Omega)_p)$.
If $f$ and the $f_J$ have finite intrinsic norms of smoothness $s + 1/p$ and
$\|f - f_J\|_p + \sum_k 2^{k/p} \omega_L(f - f_J, 2^{-k})_p \to 0$, then
$\operatorname{Tr} f_J \to \operatorname{Tr} f$ in $L_p(H)$.
\end{lemma}

\begin{proof}
\leanok
The terms are linear in $h_j$, since convolution is (\Cref{def:conv}), and the series converge;
a series commutes with complex scalars (Mathlib's \texttt{tsum\_\allowbreak const\_\allowbreak
smul''}). The bound:
$\|\operatorname{Tr} f\|_{L_p(H)} \leq \sum_j K 2^{j/p} \|k_j * f\|_p$
(\Cref{lem:trace-terms-lp}), and the local means are bounded by the moduli
(\Cref{lem:local-means-upper}); with the cut-offs $1_\Omega (k_j * f)$ kept, by the moduli on
$\Omega$ (\Cref{lem:intrinsic-upper}). The convergence follows from the linearity and the bound for
$f - f_J$.
\end{proof}

\begin{lemma}[Approximate identities at points of continuity]
\label{lem:approx-identity-points}
\lean{Rd.conv_dilate_near}
\leanok
\uses{lem:dilate}
Let $\Phi$ be integrable with $\int \Phi = 1$ and $\Phi = 0$ outside $\overline{B}(0, R)$, and let
$F_J(z) \to z_0$ as $J \to \infty$ and $z \to y$. Then $(\Phi_{2^J} * F_J)(z) \to z_0$ as
$J \to \infty$ and $z \to y$.
\end{lemma}

\begin{proof}
\leanok
$(\Phi_{2^J} * F_J)(z) - z_0 = \int \Phi_{2^J}(w) (F_J(z - w) - z_0) \, dw$, only
$\|w\| \leq R 2^{-J}$ enters, and $\|\Phi_{2^J}\|_1 = \|\Phi\|_1$.
\end{proof}

\begin{lemma}[The reproducing formula at every point]
\label{lem:reproducing-pointwise}
\lean{Rd.LocalMeansBump.hasCompactSupport_scaled, Rd.LocalMeansBump.continuous_scaled,
  Rd.LocalMeansBump.hasCompactSupport_kernel, Rd.LocalMeansBump.continuous_kernel,
  Rd.LocalMeansBump.sum_conv_dual_localMean_apply}
\leanok
\uses{lem:reproducing-formula, lem:dual-kernel-bounds}
For $f \in L_p$ and every $y \in V$,
$\sum_{j \leq J} (\eta_j * (k_j * f))(y) = (m_{2^J} * (m_{2^J} * f))(y)$.
\end{lemma}

\begin{proof}
\leanok
Both sides are continuous and agree almost everywhere by \Cref{lem:reproducing-formula}.
\end{proof}

\begin{lemma}[The trace as a limit of mollified restrictions]
\label{lem:trace-mollified}
\lean{Rd.LipschitzGraph.sum_traceFun_eq_conv, Rd.LipschitzGraph.exists_subseq_tendsto_trace}
\leanok
\uses{thm:trace-series, lem:reproducing-pointwise, lem:support-principle-two, def:graph-map}
Let $s > 0$ and let $f \in L_p$ have a finite intrinsic norm of smoothness $s + 1/p$. The partial
sums of the trace at $x' \in H$ are $(m_{2^J} * (m_{2^J} * f))(\gamma(x'))$, and a subsequence of
them converges to $\operatorname{Tr} f(x')$ for almost every $x' \in H$.
\end{lemma}

\begin{proof}
\leanok
At $\gamma(x')$ the cut-off does not matter, since the kernels are supported in $-K$ and
$\gamma(x') + K \subseteq \Omega$ (\Cref{def:graph-map}, \Cref{lem:support-principle-two}), so
\Cref{lem:reproducing-pointwise} gives the partial sums. They converge to $\operatorname{Tr} f$ in
$L_p(H)$ (\Cref{thm:trace-series}), and a subsequence converges almost everywhere.
\end{proof}

\begin{theorem}[The trace of continuous functions]
\label{thm:trace-continuous}
\lean{Rd.LipschitzGraph.sum_traceFun_eq, Rd.LipschitzGraph.trace_ae_eq_comp_graphMap}
\leanok
\uses{lem:trace-mollified, lem:approx-identity-points}
Let $L \geq 1$, $s > 0$, and let $f \in L_p$ have a finite norm as in \Cref{thm:trace-graph} and a
continuous representative $G$. Then $\operatorname{Tr} f = G \circ \gamma$ almost everywhere on $H$.
\end{theorem}

\begin{proof}
\leanok
The partial sums of the trace at $x'$ are $(m_{2^J} * (m_{2^J} * G))(\gamma(x'))$, and a subsequence
converges to $\operatorname{Tr} f$ almost everywhere (\Cref{lem:trace-mollified}). By
\Cref{lem:approx-identity-points}, applied twice, the whole sequence converges to $G(\gamma(x'))$
for every $x'$.
\end{proof}

\section{The trace space of a Lipschitz graph domain}

To extend functions from $H$ into $V$ with smoothness above $1$, the graph is smoothed at each
scale: the extended atoms need second differences of order $(2^j \|h\|)^2$, which a merely
Lipschitz graph does not give.

\begin{lemma}[Convolution with Lipschitz functions]
\label{lem:conv-lipschitz}
\lean{Rd.higherDiff_one_apply, Rd.higherDiff_two_apply, Rd.integral_higherDiff_eq_zero,
  Rd.norm_conv_sub_conv_le,
  Rd.norm_conv_sub_self_le, Rd.norm_higherDiff_two_conv_le,
  Rd.LocalMeansBump.integral_norm_scaled, Rd.LocalMeansBump.integral_norm_higherDiff_scaled_le}
\leanok
\uses{lem:conv-pointwise, lem:dual-kernel-bounds}
Let $u$ satisfy $|u(x) - u(y)| \leq M \|x - y\|$, and let $K$ be integrable, continuous, with
compact support and $K = 0$ outside $\overline{B}(0, R)$. Then
$|(K * u)(x) - (K * u)(y)| \leq M \|K\|_1 \|x - y\|$; if $\int K = 1$, then
$|(K * u)(x) - u(x)| \leq M R \|K\|_1$; and
$|\Delta_h^2 (K * u)(x)| \leq M (R + 2\|h\|) \|\Delta_h^2 K\|_1$. For the dilated moment kernels,
$\|\Delta_v^L m_{2^j}\|_1 \leq (2^j \|v\|)^L \|D^L m\|_1$.
\end{lemma}

\begin{proof}
\leanok
Differences integrate to $0$, and $\Delta_h^2 (K * u) = (\Delta_h^2 K) * u$
(\Cref{lem:conv-pointwise}), so
$\Delta_h^2 (K * u)(x) = \int \Delta_h^2 K(z) (u(x - z) - u(x)) \, dz$ with
$\|z\| \leq R + 2\|h\|$ on the support of $\Delta_h^2 K$.
\end{proof}

\begin{definition}[The mollified graph]
\label{def:mollified-graph}
\lean{Rd.LipschitzGraph.graphFun, Rd.LipschitzGraph.norm_graphFun_sub_le,
  Rd.LipschitzGraph.continuous_graphFun, Rd.LipschitzGraph.graphBump,
  Rd.LipschitzGraph.graphBump_radius, Rd.LipschitzGraph.graphBumpNorm,
  Rd.LipschitzGraph.graphBumpNorm_nonneg, Rd.LipschitzGraph.graphBumpDeriv,
  Rd.LipschitzGraph.graphBumpDeriv_nonneg, Rd.LipschitzGraph.graphBump_scaled_eq_zero,
  Rd.LipschitzGraph.integrable_graphBump_mul, Rd.LipschitzGraph.smoothGraph,
  Rd.LipschitzGraph.continuous_smoothGraph}
\leanok
\uses{def:graph-map, def:moment-kernel}
Let $\rho$ be the moment kernel of order $2$ of a bump of radius $1$ on $H$. The mollified graph at
scale $2^{-j}$ is $\phi_j = \operatorname{Re} (\rho_{2^j} * \phi)$ on $H$.
\end{definition}

\begin{lemma}[Bounds of the mollified graph]
\label{lem:mollified-graph}
\lean{Rd.LipschitzGraph.abs_smoothGraph_sub_le, Rd.LipschitzGraph.abs_smoothGraph_sub_phi_le,
  Rd.LipschitzGraph.abs_second_smoothGraph_le}
\leanok
\uses{def:mollified-graph, lem:conv-lipschitz}
$\phi_j$ is Lipschitz with constant $M \|\rho\|_1$, $|\phi_j - \phi| \leq 2 M \|\rho\|_1 2^{-j}$, and
$|\phi_j(x' + 2h') - 2\phi_j(x' + h') + \phi_j(x')| \leq 4 M \|D^2 \rho\|_1 2^j \|h'\|^2$ for
$\|h'\| \leq 2^{-j}$.
\end{lemma}

\begin{proof}
\leanok
\Cref{lem:conv-lipschitz} with $R = 2 \cdot 2^{-j}$.
\end{proof}

\begin{definition}[The vertical cut-off]
\label{def:vertical-cutoff}
\lean{Rd.LipschitzGraph.cutoffLevel, Rd.LipschitzGraph.cutoffLevel_pos,
  Rd.LipschitzGraph.profileBump, Rd.LipschitzGraph.profile, Rd.LipschitzGraph.contDiff_profile,
  Rd.LipschitzGraph.hasCompactSupport_profile, Rd.LipschitzGraph.norm_profile_le,
  Rd.LipschitzGraph.profile_eq_one, Rd.LipschitzGraph.profile_eq_zero,
  Rd.LipschitzGraph.exists_profile_bound, Rd.LipschitzGraph.profileConst,
  Rd.LipschitzGraph.norm_iteratedFDeriv_profile_le, Rd.LipschitzGraph.profileConst_nonneg,
  Rd.LipschitzGraph.vertCoord, Rd.LipschitzGraph.vertCutoff, Rd.LipschitzGraph.norm_vertCutoff_le,
  Rd.LipschitzGraph.vertCutoff_graphMap, Rd.LipschitzGraph.vertCutoff_eq_zero,
  Rd.LipschitzGraph.continuous_vertCutoff}
\leanok
\uses{lem:mollified-graph}
Let $c = 2 M \|\rho\|_1 + 1$ and let $\chi$ be a smooth bump on $\R$ with $0 \leq \chi \leq 1$,
$\chi = 1$ on $[-c, c]$ and $\chi = 0$ outside $(-(c + 1), c + 1)$. The vertical cut-off is
$v_j(x) = \chi(\psi_j(x))$ with $\psi_j(x) = 2^j (\langle x, e \rangle - \phi_j(Px))$. Then
$v_j(\gamma(x')) = 1$, since $|\psi_j(\gamma(x'))| = 2^j |\phi(x') - \phi_j(x')| < c$, and
$v_j(x) = 0$ if $|\langle x, e \rangle - \phi_j(Px)| \geq (c + 1) 2^{-j}$.
\end{definition}

\begin{lemma}[Differences of the vertical cut-off]
\label{lem:vertical-cutoff}
\lean{Rd.LipschitzGraph.norm_profile_sub_le, Rd.LipschitzGraph.norm_profile_second_le,
  Rd.LipschitzGraph.abs_vertCoord_sub_le, Rd.LipschitzGraph.abs_second_vertCoord_le,
  Rd.LipschitzGraph.vertCutoffConst, Rd.LipschitzGraph.four_le_vertCutoffConst,
  Rd.LipschitzGraph.norm_vertCutoff_add_sub_le, Rd.LipschitzGraph.norm_vertCutoff_second_le,
  Rd.LipschitzGraph.norm_higherDiff_vertCutoff_le, Rd.LipschitzGraph.enorm_higherDiff_vertCutoff_le}
\leanok
\uses{def:vertical-cutoff, lem:mollified-graph, lem:smooth-differences}
There is $B < \infty$ with $|\Delta_h^m v_j| \leq B \min(1, (2^j \|h\|)^m)$ for $m \leq 2$.
\end{lemma}

\begin{proof}
\leanok
$\psi_j$ is Lipschitz with constant $2^j (1 + M \|\rho\|_1)$ and has second differences at most
$4 M \|D^2 \rho\|_1 (2^j \|h\|)^2$ for $\|h\| \leq 2^{-j}$ (\Cref{lem:mollified-graph}); the linear
part $\langle x, e \rangle$ has none. With $\psi_l = \psi_j(x + lh)$ and $t = \psi_1 - \psi_0$,
\[
  \Delta_h^2 v_j(x) = [\chi(\psi_2) - \chi(2\psi_1 - \psi_0)] + \Delta_t^2 \chi(\psi_0),
\]
and \Cref{lem:smooth-differences} bounds both terms by the derivatives of $\chi$. For
$\|h\| > 2^{-j}$, $|\Delta_h^m v_j| \leq 4$.
\end{proof}

\begin{lemma}[Extended atoms]
\label{lem:extended-atoms}
\lean{Rd.min_one_rpow_mul_le, Rd.LipschitzGraph.extendAtom, Rd.LipschitzGraph.extendAtom_graphMap,
  Rd.LipschitzGraph.continuous_extendAtom, Rd.LipschitzGraph.higherDiff_comp_proj,
  Rd.LipschitzGraph.enorm_extendAtom_le, Rd.LipschitzGraph.enorm_higherDiff_two_extendAtom_le,
  Rd.LipschitzGraph.extendRadius, Rd.LipschitzGraph.extendCenter,
  Rd.LipschitzGraph.proj_extendCenter, Rd.LipschitzGraph.extendAtom_eq_zero,
  Rd.LipschitzGraph.isRegularAtom_extendAtom}
\leanok
\uses{lem:vertical-cutoff, lem:leibniz-differences, lem:pointwise-atoms}
Let $a$ be continuous on $H$, $a = 0$ outside $\overline{B}(c', d 2^{-j})$, and
$|\Delta_h^l a| \leq A_0 \min(1, (2^j \|h\|)^l)$ for $l \leq 2$. Then $A(x) = a(Px) v_j(x)$ has
$A \circ \gamma = a$, vanishes outside
$\overline{B}(c' + \phi_j(c') e, (d (1 + M \|\rho\|_1) + c + 1) 2^{-j})$, and
$|\Delta_h^2 A| \leq 4 A_0 B \min(1, (2^j \|h\|)^2)$; so $A$ is a regular atom of order and exponent
$2$ on $V$.
\end{lemma}

\begin{proof}
\leanok
By Leibniz's rule for differences (\Cref{lem:leibniz-differences}),
$\Delta_h^2 A(x) = \sum_{l \leq 2} \binom{2}{l} \Delta_{Ph}^l a(Px) \, \Delta_h^{2 - l} v_j(x + lh)$,
and $\min(1, t^l) \min(1, t^{2-l}) \leq \min(1, t^2)$. \Cref{lem:pointwise-atoms} concludes.
\end{proof}

\begin{theorem}[The extension from the boundary]
\label{thm:graph-trace-extension}
\lean{Rd.LipschitzGraph.extLevelH, Rd.LipschitzGraph.continuous_extLevelH,
  Rd.LipschitzGraph.extRadiusH, Rd.LipschitzGraph.extRadiusH_nonneg,
  Rd.LipschitzGraph.exists_extension_atoms, Rd.LipschitzGraph.extLevel,
  Rd.LipschitzGraph.extLevel_eq_tsum, Rd.LipschitzGraph.continuous_extLevel,
  Rd.LipschitzGraph.eLpNorm_extLevel_le, Rd.LipschitzGraph.memLp_extLevel,
  Rd.LipschitzGraph.traceExtTerm, Rd.LipschitzGraph.traceExtTerm_ae,
  Rd.LipschitzGraph.graphExtension, Rd.LipschitzGraph.summable_traceExtTerm,
  Rd.LipschitzGraph.modNorm_graphExtension_le}
\leanok
\uses{lem:extended-atoms, thm:atomic-decomposition, lem:cell-atoms-dual, thm:rd-atomic-synthesis,
  thm:local-means}
Let $1 \leq p < \infty$, $0 < s < 1$ and $1 \leq q \leq \infty$. For $g \in L_p(H)$ let
$g_j = \eta_j * (k_j * g)$ be the terms of the reproducing formula on $H$, with moment order $1$,
and $F_j(x) = g_j(Px) v_j(x)$. There is $K < \infty$ such that $E g = \sum_j F_j$ satisfies
\begin{multline*}
  \|E g\|_p + \|(2^{k(s + 1/p)} \omega_2(E g, 2^{-k})_p)_k\|_{\ell_q}
  \\ \leq K \bigl( \|g\|_{L_p(H)} + \|(2^{ks} \omega_1(g, 2^{-k})_p)_k\|_{\ell_q} \bigr).
\end{multline*}
\end{theorem}

\begin{proof}
\leanok
This is \cite[Theorem 4.10]{SchneiderVybiral} for $p \geq 1$. $F_j = \sum_m \lambda_{j,m} A_{j,m}$
with the extensions of the cell atoms of $g_j$ (\Cref{lem:extended-atoms}), whose centres project to
the cell centres of $H$ and so have bounded overlap. The atomic synthesis on $V$
(\Cref{thm:rd-atomic-synthesis}) with exponent $2 > s + 1/p$ bounds the left side by the coefficient
quasi-norm with smoothness $s + 1/p$ on $V$, which is the one with smoothness $s$ on $H$ since
$(s + 1/p) - n/p = s - (n-1)/p$, that is, the local means quasi-norm of $g$
(\Cref{lem:cell-coefficients}, \Cref{thm:local-means}).
\end{proof}

\begin{lemma}[Moduli of different orders]
\label{lem:moduli-orders}
\lean{Rd.lpModulus_add_order_le, Rd.modNorm_order_le,
  Rd.LipschitzGraph.intrinsicNorm_lt_top_of_modNorm}
\leanok
\uses{lem:rd-modulus, lem:intrinsic-upper}
$\omega_{r + l}(f, t)_p \leq 2^l \omega_r(f, t)_p$. Hence a finite norm with moduli of order $2$ on
$V$ gives a finite intrinsic norm of every order $L \geq 2$ on $\Omega$.
\end{lemma}

\begin{proof}
\leanok
$\Delta_h^{r + l} f = \Delta_h^l \Delta_h^r f$ and $\|\Delta_h^l g\|_p \leq 2^l \|g\|_p$.
\end{proof}

\begin{theorem}[The trace of the extension]
\label{thm:trace-extension}
\lean{Rd.modNorm_one_tsum_le, Rd.LipschitzGraph.exists_modNorm_traceExtTerm_le,
  Rd.LipschitzGraph.trace_sum_traceExtTerm, Rd.LipschitzGraph.trace_graphExtension}
\leanok
\uses{thm:graph-trace-extension, thm:trace-continuous, lem:trace-linear, lem:moduli-orders,
  lem:reproducing-formula}
Let $1 \leq p < \infty$, $0 < s < 1$, $1 \leq q \leq \infty$, $L \geq 2$, and $g \in L_p(H)$ with
$\|g\|_p + \|(2^{ks} \omega_1(g, 2^{-k})_p)_k\|_{\ell_q} < \infty$. Then
$\operatorname{Tr}(E g) = g$.
\end{theorem}

\begin{proof}
\leanok
The partial sums $\sum_{j \leq J} F_j$ are continuous, so their traces are
$\sum_{j \leq J} F_j \circ \gamma = \sum_{j \leq J} g_j$ (\Cref{thm:trace-continuous}), which converge
to $g$ (\Cref{lem:reproducing-formula}). They converge to $E g$ in the norm
$\|\cdot\|_p + \sum_k 2^{k/p} \omega_L(\cdot, 2^{-k})_p$, since each level has norm at most a constant
times $\|k_j * g\|_p$ by \Cref{lem:single-level} and \Cref{lem:moduli-orders}, and
$\sum_j \|k_j * g\|_p < \infty$; there the trace is linear and continuous (\Cref{lem:trace-linear}).
\end{proof}

\begin{theorem}[The trace space of a Lipschitz graph domain]
\label{thm:graph-trace-space}
\lean{Rd.LipschitzGraph.exists_trace_eq_iff}
\leanok
\uses{thm:trace-graph, thm:trace-extension, lem:moduli-orders}
Let $1 \leq p < \infty$, $0 < s < 1$, $1 \leq q \leq \infty$ and $L \geq 2$. A function
$g \in L_p(H)$ is the trace of some $f$ with
$\|f\|_{L_p(\Omega)} + \|(2^{k(s + 1/p)} \omega_L(f, 2^{-k}, \Omega)_p)_k\|_{\ell_q} < \infty$ if and
only if $\|g\|_p + \|(2^{ks} \omega_1(g, 2^{-k})_p)_k\|_{\ell_q} < \infty$.
\end{theorem}

\begin{proof}
\leanok
This is \cite[Theorem 4.11]{SchneiderVybiral} for one Lipschitz graph domain and $p \geq 1$. The
trace is bounded (\Cref{thm:trace-graph}), and $E g$ has the trace $g$
(\Cref{thm:trace-extension}) and a finite intrinsic norm (\Cref{thm:graph-trace-extension},
\Cref{lem:moduli-orders}).
\end{proof}

\section{Traces of every order on the half-space}

For the half-space the boundary is flat, and the restrictions and extensions of smooth atoms stay
smooth. This removes the restriction $s < 1$ of the Lipschitz case. Let $e$ be a unit vector,
$\Omega = \{x : \langle x, e \rangle > 0\}$ the Lipschitz graph domain with graph function $0$,
and $H = e^\perp$.

\begin{lemma}[Restrictions of smooth atoms to the hyperplane]
\label{lem:smooth-atoms-hyperplane}
\lean{Rd.norm_iteratedFDeriv_comp_clm_le, Rd.LipschitzGraph.flat_graphMap,
  Rd.LipschitzGraph.isSmoothAtom_comp_graphMap_flat}
\leanok
\uses{def:smooth-atoms, def:graph-map}
The graph map of $\Omega$ is the inclusion of $H$. A smooth atom on $V$ of order $K$ with centre
$c$ restricts to a smooth atom on $H$ of order $K$ with centre $Pc$, the same radius factor and the
same constant.
\end{lemma}

\begin{proof}
\leanok
$\|D^i (f \circ L)(x)\| \leq \|D^i f(Lx)\| \|L\|^i$ for a continuous linear map $L$, and the
inclusion has norm at most $1$. For $x' \in H$, $\|x' - Pc\| \leq \|x' - c\|$.
\end{proof}

\begin{theorem}[The trace theorem of every order]
\label{thm:trace-every-order}
\lean{Rd.LipschitzGraph.exists_trace_atoms_flat, Rd.LipschitzGraph.modNorm_traceSeries_le_flat,
  Rd.LipschitzGraph.modNorm_trace_le_flat}
\leanok
\uses{lem:smooth-atoms-hyperplane, lem:cell-atoms-smooth, lem:smooth-atoms-regular,
  thm:trace-graph, thm:rd-atomic-synthesis}
Let $1 \leq p < \infty$, $0 < s < r$, $1 \leq q \leq \infty$ and $L \in \N$. There is
$K < \infty$ such that
\begin{multline*}
  \|\operatorname{Tr} f\|_{L_p(H)} + \|(2^{ks} \omega_r(\operatorname{Tr} f, 2^{-k})_p)_k\|_{\ell_q}
  \\ \leq K \bigl(\|f\|_{L_p(\Omega)} +
    \|(2^{k(s + 1/p)} \omega_L(f, 2^{-k}, \Omega)_p)_k\|_{\ell_q}\bigr)
\end{multline*}
for all $f \in L_p(V)$ for which the right side is finite.
\end{theorem}

\begin{proof}
\leanok
As \Cref{thm:trace-graph}: the restrictions of the cell atoms are smooth atoms of order $r$ on $H$
(\Cref{lem:cell-atoms-smooth}, \Cref{lem:smooth-atoms-hyperplane}), hence regular atoms of order
and exponent $r$ (\Cref{lem:smooth-atoms-regular}). The centres of the cells meeting $H$ lie within
$3d 2^{-j} + r 2^{-j}$ of the point, and the synthesis on $H$ with exponent $r > s$ concludes.
\end{proof}

\begin{lemma}[Extended atoms of every order]
\label{lem:extended-atoms-every-order}
\lean{Rd.LipschitzGraph.smoothGraph_flat, Rd.LipschitzGraph.vertCutoff_flat,
  Rd.LipschitzGraph.exists_isSmoothAtom_extendAtom_flat}
\leanok
\uses{def:vertical-cutoff, lem:extended-atoms, def:smooth-atoms}
For $\Omega$ the vertical cut-off is $v_j(x) = \chi(2^j \langle x, e \rangle)$. For every order $K$
there is $B \geq 0$ such that the extension $A(x) = a(Px) \chi(2^j \langle x, e \rangle)$ of a
smooth atom $a$ on $H$ of order $K$, level $j$ and constant $A$ is a smooth atom on $V$ of order
$K$ with constant $2^K B A$.
\end{lemma}

\begin{proof}
\leanok
The mollified graph of the graph function $0$ vanishes. Leibniz's rule bounds
$\|D^l A\| \leq \sum_i \binom{l}{i} A 2^{ji} B 2^{j(l-i)} = 2^l A B 2^{jl}$, since
$\|D^i (a \circ P)\| \leq A 2^{ji}$ and the derivatives of $x \mapsto \chi(2^j \langle x, e \rangle)$
of order $i$ are at most $2^{ji} \sup |\chi^{(i)}|$.
\end{proof}

\begin{definition}[The extension of every order]
\label{def:flat-extension}
\lean{Rd.LipschitzGraph.flatLevelH, Rd.LipschitzGraph.continuous_flatLevelH,
  Rd.LipschitzGraph.flatLevel, Rd.LipschitzGraph.flatLevel_eq_tsum,
  Rd.LipschitzGraph.continuous_flatLevel, Rd.LipschitzGraph.flatTerm,
  Rd.LipschitzGraph.flatTerm_ae, Rd.LipschitzGraph.flatExtension}
\leanok
\uses{def:local-means, lem:cell-decomposition}
For $g \in L_p(H)$ and a moment order $L_H$ let $g_j = \eta_j * (k_j * g)$ be the terms of the
reproducing formula on $H$. The extension is $E g = \sum_j F_j$ with the continuous levels
$F_j(x) = g_j(Px) \chi(2^j \langle x, e \rangle) = \sum_m \lambda_{j,m} A_{j,m}(x)$, the
extensions of the cell atoms of $g_j$.
\end{definition}

\begin{theorem}[The extension of every order is bounded]
\label{thm:flat-extension-bounded}
\lean{Rd.LipschitzGraph.exists_flat_extension_atoms, Rd.LipschitzGraph.eLpNorm_flatLevel_le,
  Rd.LipschitzGraph.memLp_flatLevel, Rd.LipschitzGraph.summable_flatTerm,
  Rd.LipschitzGraph.coeffNorm_flat_eq, Rd.LipschitzGraph.modNorm_flatExtension_le_localMeansNorm,
  Rd.LipschitzGraph.modNorm_flatExtension_le}
\leanok
\uses{def:flat-extension, lem:extended-atoms-every-order, lem:cell-atoms-smooth,
  thm:rd-atomic-synthesis, thm:local-means}
Let $1 \leq p < \infty$, $s > 0$ and $s + 1/p < r$. There is $K < \infty$ with
\begin{multline*}
  \|E g\|_p + \|(2^{k(s + 1/p)} \omega_r(E g, 2^{-k})_p)_k\|_{\ell_q} \\ \leq
  K \bigl(\|g\|_{L_p(H)} + \|(2^{ks} \omega_{L_H}(g, 2^{-k})_p)_k\|_{\ell_q}\bigr) .
\end{multline*}
\end{theorem}

\begin{proof}
\leanok
The extended atoms are smooth of order $r$ (\Cref{lem:extended-atoms-every-order}), hence regular
atoms of order and exponent $r$ on $V$; they project to the cell centres of $H$, so they have
bounded overlap. The coefficient quasi-norm with smoothness $s + 1/p$ in dimension $n$ is the local
means quasi-norm of $g$ with smoothness $s$, since $(s + 1/p) - n/p = s - (n-1)/p$, and
\Cref{thm:local-means} bounds it.
\end{proof}

\begin{theorem}[The trace of the extension of every order]
\label{thm:trace-flat-extension}
\lean{Rd.LipschitzGraph.exists_modNorm_flatTerm_le, Rd.LipschitzGraph.trace_sum_flatTerm,
  Rd.LipschitzGraph.trace_flatExtension}
\leanok
\uses{thm:flat-extension-bounded, thm:trace-continuous, thm:trace-extension}
Let $L \geq 2$, $L_H \geq 1$, $s > 0$ with $s + 1/p < L$, and $g \in L_p(H)$ with
$\|(2^{js} \|k_j * g\|_p)_j\|_{\ell_q} < \infty$. Then $\operatorname{Tr}(E g) = g$.
\end{theorem}

\begin{proof}
\leanok
As \Cref{thm:trace-extension}: the partial sums $\sum_{j \leq J} F_j$ are continuous with traces
$\sum_{j \leq J} g_j \to g$ (\Cref{thm:trace-continuous}), and they converge to $E g$ in the norm of
$\Besov{1/p}{p}{1}$ with order $L$, where the trace is continuous.
\end{proof}

\begin{theorem}[The trace space of the half-space]
\label{thm:half-space-trace-space}
\lean{Rd.LipschitzGraph.exists_trace_eq_iff_flat}
\leanok
\uses{thm:trace-every-order, thm:flat-extension-bounded, thm:trace-flat-extension}
Let $1 \leq p < \infty$, $s > 0$, $1 \leq q \leq \infty$, $L \geq 2$ with $s + 1/p < L$, and
$r > s$. Then $g \in L_p(H)$ is the trace of some $f$ with
$\|f\|_{L_p(\Omega)} + \|(2^{k(s+1/p)} \omega_L(f, 2^{-k}, \Omega)_p)_k\|_{\ell_q} < \infty$ if and
only if $\|g\|_p + \|(2^{ks} \omega_r(g, 2^{-k})_p)_k\|_{\ell_q} < \infty$. So
$\Besov{s}{p}{q}(H)$ is the trace space of $\Besov{s + 1/p}{p}{q}$ of the half-space for every
$s > 0$.
\end{theorem}

\begin{proof}
\leanok
\Cref{thm:trace-every-order} bounds the trace. Conversely, the extension with moment order
$L_H = r$ has finite norm (\Cref{thm:flat-extension-bounded}) and trace $g$
(\Cref{thm:trace-flat-extension}).
\end{proof}

\begin{lemma}[Traces of functions continuous on the closed half-space]
\label{lem:trace-continuous-half-space}
\lean{Rd.foldHalfSpace, Rd.continuous_foldHalfSpace, Rd.foldHalfSpace_of_nonneg,
  Rd.inner_foldHalfSpace, Rd.foldHalfSpace_mem, Rd.norm_foldHalfSpace,
  Rd.trace_ae_eq_of_continuousOn}
\leanok
\uses{thm:trace-continuous, lem:trace-locality}
Let $f \in L_p(V)$ have a finite intrinsic norm of smoothness $s + 1/p$, $s > 0$, on $\Omega$, and
let $G$ be continuous on the closed half-space $\{x : \langle x, e \rangle \geq 0\}$, vanish there
outside a compact set, and agree with $f$ almost everywhere on $\Omega$. Then
$\operatorname{Tr} f = G$ almost everywhere on $H$.
\end{lemma}

\begin{proof}
\leanok
The folding $F(x) = x + (|\langle x, e \rangle| - \langle x, e \rangle) e$ is continuous, preserves
the norm, maps $V$ into the closed half-space and fixes it. So $G \circ F$ is continuous with
compact support and agrees with $f$ on $\Omega$. The trace depends only on $1_\Omega f$
(\Cref{lem:trace-locality}), and the trace of a continuous function is its restriction
(\Cref{thm:trace-continuous}).
\end{proof}

\section{Traces in charts}

Three properties make the trace a function on the boundary, independent of the graph domain that
describes the boundary near a point: locality, approximation by continuous functions, and the
change of chart. In this section $\Omega$, $\Omega_1$ and $\Omega_2$ are Lipschitz graph domains,
with hyperplanes $H$, $H_1$, $H_2$, graph maps $\gamma$, $\gamma_1$, $\gamma_2$ and projections $P$,
$P_1$, $P_2$.

\begin{lemma}[Locality of the trace]
\label{lem:trace-locality}
\lean{Rd.LipschitzGraph.conv_scaled_convLp_eq_zero,
  Rd.LipschitzGraph.trace_ae_eq_zero_of_ae_eq_zero, Rd.LipschitzGraph.intrinsicNorm_sub_le,
  Rd.LipschitzGraph.trace_ae_eq_of_ae_eq, Rd.LipschitzGraph.trace_congr}
\leanok
\uses{lem:trace-mollified, lem:trace-linear, lem:support-principle-two}
Let $s > 0$, let $f_1, f_2 \in L_p$ have finite intrinsic norms of smoothness $s + 1/p$, and let
$U \subseteq V$ be open with $f_1 = f_2$ almost everywhere on $\Omega \cap U$. Then
$\operatorname{Tr} f_1 = \operatorname{Tr} f_2$ almost everywhere on $\gamma^{-1}(U)$. For
$U = V$ this holds for all $f_1, f_2 \in L_p$: the trace depends only on $1_\Omega f$.
\end{lemma}

\begin{proof}
\leanok
By linearity (\Cref{lem:trace-linear}) it suffices to treat $f = f_1 - f_2$, whose intrinsic norm is
at most the sum of the two. The mollified restriction $(m_{2^J} * (m_{2^J} * f))(\gamma(x'))$ only
sees $f$ on $\Omega \cap B(\gamma(x'), 2 L R 2^{-J})$, since the kernels are supported in the cone
(\Cref{lem:support-principle-two}); so it vanishes for large $J$ if $\gamma(x') \in U$, and a
subsequence converges to $\operatorname{Tr} f(x')$ almost everywhere (\Cref{lem:trace-mollified}).
For $U = V$, the trace is built from the cut-offs $1_\Omega (k_j * f)$, which depend only on
$1_\Omega f$ (\Cref{lem:support-principle-two}).
\end{proof}

\begin{lemma}[Approximation by the partial sums of the reproducing formula]
\label{lem:partial-sums-approx}
\lean{Rd.LocalMeansBump.partialSum, Rd.LocalMeansBump.partialSum_ae,
  Rd.LocalMeansBump.continuous_partialSum_rep, Rd.LocalMeansBump.sum_range_ite_lt,
  Rd.LocalMeansBump.tendsto_modNorm_sub_partialSum, Rd.LocalMeansBump.modNorm_partialSum_lt_top}
\leanok
\uses{lem:reproducing-formula, lem:local-means-synthesis, lem:intrinsic-upper}
The partial sums $G_J = \sum_{j \leq J} \eta_j * (k_j * f)$ are continuous and, for $0 < \sigma < L$,
have finite norms $\|G_J\|_p + \|(2^{k\sigma} \omega_L(G_J, 2^{-k})_p)_k\|_{\ell_q}$. If
$\|f\|_p + \|(2^{k\sigma} \omega_L(f, 2^{-k})_p)_k\|_{\ell_q} < \infty$ and $0 < \sigma' < \sigma$,
$\sigma' < L$, then
\[
  \|f - G_J\|_p + \sum_k 2^{k\sigma'} \omega_L(f - G_J, 2^{-k})_p \to 0 .
\]
\end{lemma}

\begin{proof}
\leanok
The tail $f - G_J = \sum_{j > J} \eta_j * (k_j * f)$ is bounded by the synthesis estimate
(\Cref{lem:local-means-synthesis}) by
$\sum_{j > J} 2^{j\sigma'} \|k_j * f\|_p \leq
\sum_{j > J} 2^{-j(\sigma - \sigma')} \|(2^{i\sigma} \|k_i * f\|_p)_i\|_{\ell_q}$, and the local
means quasi-norm is finite (\Cref{lem:intrinsic-upper}).
\end{proof}

\begin{lemma}[The traces of the partial sums]
\label{lem:trace-partial-sums}
\lean{Rd.LipschitzGraph.intrinsicNorm_lt_top_of_modNorm',
  Rd.LipschitzGraph.tendsto_trace_partialSum}
\leanok
\uses{lem:partial-sums-approx, lem:trace-linear, lem:restriction-moduli}
Let $s > 0$, $s + 1/p < L$ and
$\|F\|_p + \|(2^{k(s + 1/p)} \omega_L(F, 2^{-k})_p)_k\|_{\ell_q} < \infty$. Then $F$ and the $G_J$
have finite intrinsic norms on every Lipschitz graph domain, and
$\operatorname{Tr} G_J \to \operatorname{Tr} F$ in $L_p(H)$.
\end{lemma}

\begin{proof}
\leanok
The trace is linear and bounded by $\|\cdot\|_p + \sum_k 2^{k/p} \omega_L(\cdot, 2^{-k})_p$
(\Cref{lem:trace-linear}), and \Cref{lem:partial-sums-approx} with $\sigma' = 1/p$.
\end{proof}

\begin{lemma}[Lipschitz images of null sets]
\label{lem:lipschitz-null}
\lean{Rd.volume_image_eq_zero_of_lipschitzWith}
\leanok
A Lipschitz map between finite-dimensional inner product spaces of the same dimension maps null
sets to null sets.
\end{lemma}

\begin{proof}
\leanok
In dimension $k$ the Hausdorff measure $\mu_H^k$ is a Haar measure, so it has the null sets of
Lebesgue measure, and a $K$-Lipschitz map $f$ has $\mu_H^k(f(N)) \leq K^k \mu_H^k(N)$.
\end{proof}

\begin{theorem}[Change of chart]
\label{thm:change-of-chart}
\lean{Rd.LipschitzGraph.lipschitzWith_proj_graphMap,
  Rd.LipschitzGraph.finrank_hyperplane_eq_finrank, Rd.LipschitzGraph.ae_proj_graphMap,
  Rd.LipschitzGraph.trace_ae_eq_trace_comp}
\leanok
\uses{lem:trace-partial-sums, thm:trace-continuous, lem:lipschitz-null}
Let $s > 0$, $s + 1/p < L$ and
$\|F\|_p + \|(2^{k(s + 1/p)} \omega_L(F, 2^{-k})_p)_k\|_{\ell_q} < \infty$. For almost every
$x' \in H_1$ such that $\gamma_1(x')$ lies on the graph of $\Omega_2$, that is
$\gamma_2(P_2 \gamma_1(x')) = \gamma_1(x')$,
\[
  \operatorname{Tr}_1 F(x') = \operatorname{Tr}_2 F(P_2 \gamma_1(x')) .
\]
\end{theorem}

\begin{proof}
\leanok
The traces of the continuous $G_J$ are $G_J \circ \gamma_1$ and $G_J \circ \gamma_2$
(\Cref{thm:trace-continuous}), and they converge to $\operatorname{Tr}_1 F$ and $\operatorname{Tr}_2 F$
(\Cref{lem:trace-partial-sums}); a common subsequence converges almost everywhere on $H_1$ and on
$H_2$. The map $P_1 \circ \gamma_2$ is Lipschitz with constant $1 + M_2$, and $H_1$, $H_2$ both have
dimension $n - 1$, so it maps the exceptional null set of $H_2$ to a null set of $H_1$
(\Cref{lem:lipschitz-null}). Outside of it, $G_J(\gamma_1(x')) = G_J(\gamma_2(P_2 \gamma_1(x')))$
converges to both values.
\end{proof}

\begin{lemma}[Traces of products with smooth functions]
\label{lem:trace-products}
\lean{Rd.LipschitzGraph.trace_mulLp_ae}
\leanok
\uses{thm:change-of-chart, thm:multipliers, thm:trace-continuous, lem:trace-partial-sums}
Let $s > 0$, $s + 1/p < L$, $F$ as in \Cref{thm:change-of-chart}, and $\psi$ smooth with bounded
derivatives up to order $L$. Then $\operatorname{Tr}(\psi F) = (\psi \circ \gamma) \operatorname{Tr} F$
almost everywhere on $H$.
\end{lemma}

\begin{proof}
\leanok
$\operatorname{Tr}(\psi G_J) = (\psi G_J) \circ \gamma$ and $\operatorname{Tr} G_J = G_J \circ \gamma$
(\Cref{thm:trace-continuous}); $\psi G_J \to \psi F$ in the norm with smoothness $1/p$ by
\Cref{thm:multipliers} and \Cref{lem:partial-sums-approx}, and subsequences converge almost
everywhere.
\end{proof}

\begin{lemma}[Traces of finite sums]
\label{lem:trace-sums}
\lean{Rd.LipschitzGraph.trace_zero, Rd.LipschitzGraph.trace_finsetSum,
  Rd.LipschitzGraph.intrinsicNorm_add_le, Rd.LipschitzGraph.intrinsicNorm_finsetSum_lt_top,
  Rd.LipschitzGraph.trace_finsetSum_of_intrinsicNorm}
\leanok
\uses{lem:trace-linear, lem:pasting}
For $s > 0$ and summands $F_i$ with finite norms
$\|F_i\|_p + \|(2^{k(s + 1/p)} \omega_L(F_i, 2^{-k})_p)_k\|_{\ell_q}$, or with finite intrinsic
norms of smoothness $s + 1/p$,
$\operatorname{Tr} \sum_i F_i = \sum_i \operatorname{Tr} F_i$.
\end{lemma}

\begin{proof}
\leanok
Induction on the number of summands with \Cref{lem:trace-linear}; finite sums keep the norms finite
(\Cref{lem:pasting}), and the intrinsic norms are subadditive.
\end{proof}

\section{Traces on bounded Lipschitz domains}

In this section $\Omega$ is a bounded Lipschitz domain with boundary $\Gamma$, boundary balls
$B_j$ and graph domains $\Omega_j$ (\Cref{def:bounded-lipschitz}), with hyperplanes $H_j$, graph
maps $\gamma_j$, projections $P_j$ and traces $\operatorname{Tr}_j$.

\begin{lemma}[The boundary in the charts]
\label{lem:boundary-charts}
\lean{Rd.LipschitzGraph.phi_proj, Rd.LipschitzGraph.graphMap_proj_eq_iff,
  Rd.LipschitzGraph.mem_frontier_domain_iff, Rd.BoundedLipschitzDomain.mem_frontier_iff,
  Rd.BoundedLipschitzDomain.graphMap_mem_frontier, Rd.BoundedLipschitzDomain.graphMap_proj_graphMap}
\leanok
\uses{def:bounded-lipschitz, def:graph-map}
The boundary of a Lipschitz graph domain is its graph
$\{y : \langle y, e \rangle = \phi(y)\} = \{y : \gamma(Py) = y\}$. For $y \in B_j$,
$y \in \Gamma$ if and only if $\gamma_j(P_j y) = y$. Hence a boundary point $\gamma_j(x')$ in
$B_j \cap B_k$ satisfies $\gamma_k(P_k \gamma_j(x')) = \gamma_j(x')$.
\end{lemma}

\begin{proof}
\leanok
A point of the graph is the limit of the points $y + te \in \Omega$, $t > 0$, since
$\phi(y + te) = \phi(y)$, and it is not in the open set $\Omega$. The boundary is local in open sets,
and $\Omega \cap B_j = \Omega_j \cap B_j$.
\end{proof}

\begin{definition}[Resolutions of unity on the boundary]
\label{def:boundary-partition}
\lean{Rd.BoundedLipschitzDomain.BoundaryPartition,
  Rd.BoundedLipschitzDomain.BoundaryPartition.toComplex,
  Rd.BoundedLipschitzDomain.BoundaryPartition.contDiff_toComplex,
  Rd.BoundedLipschitzDomain.BoundaryPartition.mem_ball_of_ne_zero,
  Rd.BoundedLipschitzDomain.BoundaryPartition.toComplex_of_eq_zero,
  Rd.BoundedLipschitzDomain.BoundaryPartition.sum_toComplex_eq_one,
  Rd.BoundedLipschitzDomain.BoundaryPartition.hasCompactSupport_toComplex,
  Rd.BoundedLipschitzDomain.nonempty_boundaryPartition}
\leanok
\uses{lem:resolution-of-unity}
A resolution of unity on the boundary consists of smooth $\psi_j$ with
$\operatorname{supp} \psi_j \subseteq B_j$ and $\sum_j \psi_j = 1$ on $\Gamma$. It exists by
\Cref{lem:resolution-of-unity}, and the $\psi_j$ have compact support.
\end{definition}

\begin{lemma}[Cut-offs]
\label{lem:boundary-cutoffs}
\lean{Rd.BoundedLipschitzDomain.BoundaryPartition.exists_cutoff}
\leanok
\uses{def:boundary-partition}
For each $j$ there is a smooth $\zeta_j$ with compact support, $\zeta_j \neq 0$ only in $B_j$, and
$\zeta_j = 1$ where $\psi_j \neq 0$.
\end{lemma}

\begin{proof}
\leanok
$\operatorname{supp} \psi_j$ is closed in $B_j = B(c_j, \rho_j)$, so it lies in some $B(c_j, r)$ with
$r < \rho_j$; take a bump function equal to $1$ on $\overline{B}(c_j, r)$ with support $B_j$.
\end{proof}

\begin{definition}[Almost everywhere on the boundary]
\label{def:ae-boundary}
\lean{Rd.BoundedLipschitzDomain.aeBoundary, Rd.BoundedLipschitzDomain.eventually_aeBoundary_iff,
  Rd.BoundedLipschitzDomain.eventually_aeBoundary_of_chart}
\leanok
\uses{lem:boundary-charts, thm:change-of-chart}
A property holds almost everywhere on $\Gamma$ if for every $j$ it holds at $\gamma_j(x')$ for almost
every $x' \in H_j$ with $\gamma_j(x') \in B_j$. A property of $\gamma_k(x')$ for almost every
$x' \in H_k$ with $\gamma_k(x') \in B_k$ holds almost everywhere on $\Gamma \cap B_k$.
\end{definition}

\begin{definition}[Besov spaces on the boundary]
\label{def:boundary-besov}
\lean{Rd.BoundedLipschitzDomain.MemBoundaryBesov, Rd.BoundedLipschitzDomain.MemBoundaryBesov.congr}
\leanok
\uses{def:boundary-partition, def:ae-boundary}
A function $g$ belongs to $B^s_{p,q}(\Gamma)$ if for every $j$ the chart piece
$(\psi_j g) \circ \gamma_j$ agrees almost everywhere on $H_j$ with some $G_j \in L_p(H_j)$ with
$\|G_j\|_p + \|(2^{ks} \omega_1(G_j, 2^{-k})_p)_k\|_{\ell_q} < \infty$. This only depends on $g$
almost everywhere on $\Gamma$.
\end{definition}

This is \cite[Definition 3.3]{SchneiderVybiral} with Lebesgue measure on the hyperplanes and first
differences, which describe $B^s_{p,q}(H_j)$ for $0 < s < 1$.

\begin{definition}[The boundary trace]
\label{def:boundary-trace}
\lean{Rd.BoundedLipschitzDomain.boundaryTrace}
\leanok
\uses{def:boundary-partition, thm:trace-graph}
The boundary trace of $F \in L_p(V)$ is
$\operatorname{Tr}_\Gamma F(y) = \sum_k \psi_k(y) \operatorname{Tr}_k F(P_k y)$.
\end{definition}

\begin{theorem}[The boundary trace in the charts]
\label{thm:boundary-trace-charts}
\lean{Rd.BoundedLipschitzDomain.boundaryTrace_graphMap,
  Rd.BoundedLipschitzDomain.boundaryTrace_eventuallyEq,
  Rd.BoundedLipschitzDomain.boundaryTrace_eventuallyEq_of_continuous}
\leanok
\uses{def:boundary-trace, thm:change-of-chart, lem:boundary-charts, def:ae-boundary,
  thm:trace-continuous}
Let $s > 0$, $s + 1/p < L$ and
$\|F\|_p + \|(2^{k(s + 1/p)} \omega_L(F, 2^{-k})_p)_k\|_{\ell_q} < \infty$. Then
$\operatorname{Tr}_\Gamma F(\gamma_j(x')) = \operatorname{Tr}_j F(x')$ for almost every $x' \in H_j$
with $\gamma_j(x') \in B_j$. Hence $\operatorname{Tr}_\Gamma F$ does not depend on the resolution of
unity up to null sets of $\Gamma$, and $\operatorname{Tr}_\Gamma F = G$ almost everywhere on $\Gamma$
if $F$ has a continuous representative $G$.
\end{theorem}

\begin{proof}
\leanok
If $\psi_k(\gamma_j(x')) \neq 0$, then $\gamma_j(x') \in B_k$ lies on the graph of $\Omega_k$
(\Cref{lem:boundary-charts}), and
$\operatorname{Tr}_k F(P_k \gamma_j(x')) = \operatorname{Tr}_j F(x')$ for almost every such $x'$
(\Cref{thm:change-of-chart}); and $\sum_k \psi_k = 1$ on $\Gamma$. For continuous $F$,
\Cref{thm:trace-continuous}.
\end{proof}

\begin{lemma}[Locality of the boundary trace]
\label{lem:boundary-trace-locality}
\lean{Rd.BoundedLipschitzDomain.boundaryTrace_congr}
\leanok
\uses{thm:boundary-trace-charts, lem:trace-locality}
For $F_1$, $F_2$ as in \Cref{thm:boundary-trace-charts} with $F_1 = F_2$ almost everywhere on
$\Omega$, $\operatorname{Tr}_\Gamma F_1 = \operatorname{Tr}_\Gamma F_2$ almost everywhere on $\Gamma$.
\end{lemma}

\begin{proof}
\leanok
In the chart $j$ both are the traces on $\Omega_j$ (\Cref{thm:boundary-trace-charts}), and
$F_1 = F_2$ on $\Omega_j \cap B_j = \Omega \cap B_j$, so \Cref{lem:trace-locality} with $U = B_j$.
\end{proof}

\begin{theorem}[The boundary trace is bounded]
\label{thm:boundary-trace-bounded}
\lean{Rd.BoundedLipschitzDomain.trace_mulLp_ae_boundaryTrace,
  Rd.BoundedLipschitzDomain.memBoundaryBesov_boundaryTrace}
\leanok
\uses{thm:boundary-trace-charts, lem:trace-products, thm:trace-graph, thm:multipliers,
  def:boundary-besov}
Let $1 \leq p < \infty$, $0 < s < 1$, $1 \leq q \leq \infty$, $s + 1/p < L$ and
$\|F\|_p + \|(2^{k(s + 1/p)} \omega_L(F, 2^{-k})_p)_k\|_{\ell_q} < \infty$. Then
$\operatorname{Tr}_\Gamma F \in B^s_{p,q}(\Gamma)$.
\end{theorem}

\begin{proof}
\leanok
The chart piece $(\psi_j \operatorname{Tr}_\Gamma F) \circ \gamma_j$ is
$(\psi_j \circ \gamma_j) \operatorname{Tr}_j F = \operatorname{Tr}_j (\psi_j F)$ almost everywhere
(\Cref{thm:boundary-trace-charts}, \Cref{lem:trace-products}), and $\psi_j F$ has a finite norm
(\Cref{thm:multipliers}), so \Cref{thm:trace-graph} applies.
\end{proof}

\begin{theorem}[The extension from the boundary]
\label{thm:boundary-trace-extension}
\lean{Rd.BoundedLipschitzDomain.exists_boundaryTrace_eq}
\leanok
\uses{def:boundary-besov, lem:boundary-cutoffs, thm:graph-trace-extension, thm:trace-extension,
  lem:moduli-orders, lem:trace-products, thm:change-of-chart, lem:trace-sums,
  thm:boundary-trace-charts}
Let $1 \leq p < \infty$, $0 < s < 1$, $1 \leq q \leq \infty$ and $L \geq 2$. Every
$g \in B^s_{p,q}(\Gamma)$ is $\operatorname{Tr}_\Gamma F$ almost everywhere on $\Gamma$ for some $F$
with $\|F\|_p + \|(2^{k(s + 1/p)} \omega_L(F, 2^{-k})_p)_k\|_{\ell_q} < \infty$.
\end{theorem}

\begin{proof}
\leanok
Let $G_j$ be the chart pieces of $g$, $E_j$ the extensions of the graph domains
(\Cref{thm:graph-trace-extension}, \Cref{lem:moduli-orders}) and $\zeta_j$ the cut-offs of
\Cref{lem:boundary-cutoffs}, and $F = \sum_j \zeta_j E_j G_j$. In the chart $k$, at almost every
$x'$ with $y = \gamma_k(x') \in B_k$,
$\operatorname{Tr}_k(\zeta_j E_j G_j)(x') = \zeta_j(y) \operatorname{Tr}_k(E_j G_j)(x')$
(\Cref{lem:trace-products}). If $\zeta_j(y) \neq 0$, then $y \in B_j$, and by the change of chart
(\Cref{thm:change-of-chart}) and \Cref{thm:trace-extension} this is
$\zeta_j(y) G_j(P_j y) = \zeta_j(y) \psi_j(y) g(y) = \psi_j(y) g(y)$; otherwise $\psi_j(y) = 0$.
Summing (\Cref{lem:trace-sums}), $\operatorname{Tr}_k F(x') = g(y)$, and
\Cref{thm:boundary-trace-charts} concludes.
\end{proof}

\begin{theorem}[The trace space of a bounded Lipschitz domain]
\label{thm:boundary-trace-space}
\lean{Rd.add_one_div_toReal_lt, Rd.BoundedLipschitzDomain.exists_boundaryTrace_eq_iff,
  Rd.BoundedLipschitzDomain.memBoundaryBesov_iff_of_boundaryPartition,
  Rd.BoundedLipschitzDomain.IsBoundaryTrace, Rd.ae_eq_of_indicatorLp_eq,
  Rd.BoundedLipschitzDomain.memBoundaryBesov_iff_exists_isBoundaryTrace,
  Rd.BoundedLipschitzDomain.exists_isBoundaryTrace_of_modNormOn_lt_top,
  Rd.BoundedLipschitzDomain.IsBoundaryTrace.unique}
\leanok
\uses{thm:boundary-trace-bounded, thm:boundary-trace-extension, lem:boundary-trace-locality,
  thm:bounded-lipschitz-besov, lem:restriction-moduli}
Let $1 \leq p < \infty$, $0 < s < 1$, $1 \leq q \leq \infty$ and $L \geq 2$. Then
$g \in B^s_{p,q}(\Gamma)$ if and only if $g = \operatorname{Tr}_\Gamma F$ almost everywhere on $\Gamma$
for some $F$ with $\|F\|_p + \|(2^{k(s + 1/p)} \omega_L(F, 2^{-k})_p)_k\|_{\ell_q} < \infty$, if and
only if $g$ is the boundary trace of some $f$ with
$\|f\|_{L_p(\Omega)} + \|(2^{k(s + 1/p)} \omega_L(f, 2^{-k}, \Omega)_p)_k\|_{\ell_q} < \infty$, that
is, $g = \operatorname{Tr}_\Gamma F$ almost everywhere on $\Gamma$ for every such $F$ that agrees with
$f$ on $\Omega$. Every such $f$ has a boundary trace in $B^s_{p,q}(\Gamma)$, and every two
boundary traces of $f$ agree almost everywhere on $\Gamma$. In particular $B^s_{p,q}(\Gamma)$ does
not depend on the resolution of unity.
\end{theorem}

\begin{proof}
\leanok
This is \cite[Theorem 4.11]{SchneiderVybiral} for $p \geq 1$. The first equivalence is
\Cref{thm:boundary-trace-bounded} and \Cref{thm:boundary-trace-extension}. For the second, the
functions with a finite intrinsic norm on $\Omega$ are the restrictions of such $F$
(\Cref{thm:bounded-lipschitz-besov}), and all of them have the same boundary trace
(\Cref{lem:boundary-trace-locality}); for one such $F$, $\operatorname{Tr}_\Gamma F$ is a boundary
trace of $f$ in $B^s_{p,q}(\Gamma)$ (\Cref{thm:boundary-trace-bounded}), and every boundary trace
of $f$ agrees with it. Both resolutions of unity give the space of boundary traces,
and the boundary traces agree almost everywhere on $\Gamma$ (\Cref{thm:boundary-trace-charts}).
\end{proof}

A boundary trace of $f$ in the sense of \Cref{thm:boundary-trace-space} is a condition on all
$F$ that agree with $f$ on $\Omega$ and have a finite norm; when $0 < \sigma < L$, $q \geq 1$ and
$f$ has a finite intrinsic norm on $\Omega$, at least one such $F$ exists
(\Cref{thm:bounded-lipschitz-besov}). The following notion contains this finiteness.

\begin{definition}[Boundary traces of functions with a finite intrinsic norm]
\label{def:boundary-trace-finite-norm}
\lean{Rd.BoundedLipschitzDomain.HasBoundaryTrace, Rd.BoundedLipschitzDomain.hasBoundaryTrace_iff}
\leanok
\uses{def:boundary-trace, def:ae-boundary, def:rd-modulus-on}
Let $1 \leq p < \infty$, $\sigma \in \R$, $L \in \N$ and $1 \leq q \leq \infty$. A function
$f \in L_p(V)$ \emph{has the boundary trace} $g$ if
$\|f\|_{L_p(\Omega)} + \|(2^{k\sigma} \omega_L(f, 2^{-k}, \Omega)_p)_k\|_{\ell_q} < \infty$ and
$g = \operatorname{Tr}_\Gamma F$ almost everywhere on $\Gamma$ for every $F \in L_p(V)$ with
$F = f$ almost everywhere on $\Omega$ and
$\|F\|_p + \|(2^{k\sigma} \omega_L(F, 2^{-k})_p)_k\|_{\ell_q} < \infty$.
\end{definition}

For $0 < \sigma < L$ the trace condition concerns at least one $F$.

\begin{theorem}[The trace space with finite intrinsic norms]
\label{thm:boundary-trace-space-finite-norm}
\lean{Rd.BoundedLipschitzDomain.memBoundaryBesov_iff_exists_hasBoundaryTrace,
  Rd.BoundedLipschitzDomain.HasBoundaryTrace.memBoundaryBesov,
  Rd.BoundedLipschitzDomain.exists_hasBoundaryTrace_of_modNormOn_lt_top,
  Rd.BoundedLipschitzDomain.HasBoundaryTrace.unique}
\leanok
\uses{def:boundary-trace-finite-norm, thm:boundary-trace-space, def:boundary-besov,
  thm:bounded-lipschitz-besov}
Let $1 \leq p < \infty$ and $1 \leq q \leq \infty$.
\begin{enumerate}
  \item For every $0 < s < 1$, every $L \geq 2$ and every function $g$ on $V$,
    $g \in B^s_{p,q}(\Gamma)$ if and only if some $f \in L_p(V)$ has the boundary trace $g$
    for $\sigma = s + 1/p$; every $f$ with
    $\|f\|_{L_p(\Omega)} + \|(2^{k(s + 1/p)} \omega_L(f, 2^{-k}, \Omega)_p)_k\|_{\ell_q} < \infty$
    has a boundary trace $g \in B^s_{p,q}(\Gamma)$ for $\sigma = s + 1/p$.
  \item For every $L \in \N$ and every $0 < \sigma < L$, any two boundary traces $g, g'$ of the
    same $f$ for $\sigma$ agree almost everywhere on $\Gamma$.
\end{enumerate}
\end{theorem}

\begin{proof}
\leanok
The first statement is the corresponding one of \Cref{thm:boundary-trace-space}, whose hypothesis
of a finite intrinsic norm is part of \Cref{def:boundary-trace-finite-norm}. For the second, by
\Cref{thm:bounded-lipschitz-besov} with $r = L$ some admissible $F$ exists, and
$g = \operatorname{Tr}_\Gamma F = g'$ almost everywhere on $\Gamma$.
\end{proof}

\chapter{Sobolev spaces on domains}
\label{ch:sobolev-domains}

In this chapter $V$ is a finite-dimensional real inner product space of dimension $n$ with
Lebesgue measure, and $\Omega \subseteq V$ is open or a bounded Lipschitz domain with boundary
$\Gamma$, as each section states; weak derivatives, Mathlib's spaces $H^s(V)$ and test functions
are those of \Cref{ch:sobolev-rd}, the boundary traces those of \Cref{ch:traces-lipschitz}.
\Cref{sec:lipschitz-domain-h1} restricts $H^1(V)$ to bounded Lipschitz domains, describes the
restrictions by intrinsic differences, shows that they lie in $\Sob{1}{2}(\Omega)$, and
identifies their boundary traces with $B^{1/2}_{2,2}(\Gamma)$.
\Cref{sec:sobolev-restriction-spaces} builds the Hilbert spaces $H^s(\Omega)$,
$\tilde H^s(\Omega)$ and $H^s_0(\Omega)$ inside $H^s(V)$, for every open $\Omega$ and every real
$s$, and \Cref{sec:sobolev-duality} proves the duality of $\tilde H^s(\Omega)$ and
$H^{-s}(\Omega)$. \Cref{sec:boundary-trace-space} defines the trace space $H^{1/2}(\Gamma)$ of a
bounded Lipschitz domain as the orthogonal complement of the trace kernel in $H^1(\Omega)$.
\Cref{sec:divergence-kernels} and \Cref{sec:sobolev-extension} prove the extension theorem for
$\Sob{1}{p}$ on bounded Lipschitz domains, $1 < p < \infty$, of first order and with an
extension for each function. \Cref{sec:lipschitz-graph-green} and
\Cref{sec:lipschitz-domain-green} prove Gauss--Green on Lipschitz graph domains and Green's
formula on bounded Lipschitz domains, and \Cref{sec:trace-kernel} proves that the kernel of the
trace is $H^1_0(\Omega)$.

On a bounded Lipschitz domain $\Omega$, the space $H^1(\Omega)$ of this chapter
(\Cref{sec:lipschitz-domain-h1}) is the space of restrictions of $H^1(V)$, the definition of
Chandler-Wilde, Hewett and Moiola \cite{ChandlerWildeHewettMoiola}; for restrictions of $H^1(V)$,
Rellich's theorem holds on every bounded set (\Cref{sec:rellich-compactness}). By this definition
every restriction lies in the space $\Sob{1}{2}(\Omega)$ of weak derivatives, on every open set;
the converse inclusion is the extension theorem. On bounded Lipschitz domains the two spaces
agree by the extension theorem of Calder\'on and Stein \cite[Chapter VI]{SteinSingular}
(\Cref{thm:sobolev-extension}); on some bounded open sets $\Sob{1}{2}(\Omega)$ is strictly
larger, as for a disc with a slit removed, and its embedding into $L_2(\Omega)$ can fail to be
compact, as for the rooms-and-passages domains \cite[Chapter 6]{AdamsFournier}. The norm bounds
on $V$ are explicit or named (\Cref{ch:sobolev-rd}); on bounded Lipschitz domains the extension
and trace theorems carry existential constants, and the statements are equivalences of
membership. On bounded Lipschitz domains the functions of
$H^1(\Omega)$ with boundary trace zero form the closure of the test functions in $H^1(\Omega)$
\cite[Theorem 3.40]{McLean} (\Cref{thm:trace-kernel-zero-boundary}); the spaces $H^1(\Omega)$ and
$H^1_0(\Omega)$ of that theorem are the Hilbert spaces of \Cref{sec:sobolev-restriction-spaces}.

\section{\texorpdfstring{$H^1$}{H\^{}1} on bounded Lipschitz domains}
\label{sec:lipschitz-domain-h1}

Throughout this section $\Omega$ is a bounded Lipschitz domain (\Cref{def:bounded-lipschitz})
with boundary $\Gamma$, and $L_2 = L_2(V)$. For $f \in L_2$, $\|f\|_{L_2(\Omega)}$ is the norm of
$1_\Omega f$, and $\omega_r(f, t, \Omega)_2$ is the modulus of smoothness of $f$ inside $\Omega$
(\Cref{def:rd-modulus-on}). The space $H^1(\Omega)$ below consists of the restrictions of
$\Sob{1}{2}(V) = H^1(V)$; on $\Omega$ its functions have weak derivatives in $L_2(\Omega)$, and
their boundary traces form $B^{1/2}_{2,2}(\Gamma)$.

\begin{definition}[$H^1$ on a bounded Lipschitz domain]
\label{def:sobolev-lipschitz-domain}
\lean{Rd.BoundedLipschitzDomain.toOpens, Rd.BoundedLipschitzDomain.coe_toOpens,
  Rd.BoundedLipschitzDomain.MemSobolevOn}
\leanok
\uses{def:bounded-lipschitz, def:weak-line-derivative}
A function $f \in L_2(V)$ lies in $H^1(\Omega)$ if there is $F \in \Sob{1}{2}(V)$ with $F = f$
almost everywhere on $\Omega$.
\end{definition}

Only the values of $f$ on $\Omega$ matter, and by \Cref{thm:sobolev-one-two} the functions $F$
are those of Mathlib's $H^1(V)$. The open set $\Omega$ carries the distributions of
\Cref{def:weak-line-derivative}. Every $f \in H^1(\Omega)$ lies in $\Sob{1}{2}(\Omega)$
(\Cref{thm:sobolev-lipschitz-domain-weak}), and the two spaces agree by the extension theorem
of Calder\'on and Stein \cite[Chapter VI]{SteinSingular}.

\begin{lemma}[$H^1(\Omega)$ is a linear space]
\label{lem:sobolev-lipschitz-domain-basic}
\lean{Rd.BoundedLipschitzDomain.memSobolevOn_of_memSobolev,
  Rd.BoundedLipschitzDomain.zero_memSobolevOn, Rd.BoundedLipschitzDomain.MemSobolevOn.congr_ae,
  Rd.BoundedLipschitzDomain.MemSobolevOn.add, Rd.BoundedLipschitzDomain.MemSobolevOn.smul}
\leanok
\uses{def:sobolev-lipschitz-domain, lem:weak-derivative-basic}
Every $f \in \Sob{1}{2}(V)$ lies in $H^1(\Omega)$. For all $f, f' \in L_2(V)$ with
$f \in H^1(\Omega)$ and $f = f'$ almost everywhere on $\Omega$, $f' \in H^1(\Omega)$. The space
$H^1(\Omega)$ is a complex linear subspace of $L_2(V)$.
\end{lemma}

\begin{proof}
\leanok
For $f \in \Sob{1}{2}(V)$ take $F = f$; in particular $0 \in H^1(\Omega)$
(\Cref{lem:weak-derivative-basic}). For $f'$ with $f' = f$ on $\Omega$ take the same $F$. For
sums and multiples take the sums and multiples of the functions $F$, which lie in
$\Sob{1}{2}(V)$ by \Cref{lem:weak-derivative-basic}.
\end{proof}

\begin{theorem}[$H^1(\Omega)$ by intrinsic differences]
\label{thm:sobolev-lipschitz-domain-intrinsic}
\lean{Rd.BoundedLipschitzDomain.memSobolevOn_iff_exists_memApprox,
  Rd.BoundedLipschitzDomain.memSobolevOn_iff_modulusOn,
  Rd.BoundedLipschitzDomain.modNormOn_le_modNorm_of_ae_eq,
  Rd.BoundedLipschitzDomain.MemSobolevOn.exists_modNorm_lt_top}
\leanok
\uses{def:sobolev-lipschitz-domain, thm:sobolev-one-two-besov, thm:bounded-lipschitz-besov,
  lem:indicator-lp-eq-iff, lem:restriction-moduli}
For every integer $r \geq 2$ and every $f \in L_2(V)$ the following are equivalent:
\begin{enumerate}
  \item $f \in H^1(\Omega)$;
  \item there is $F$ in the approximation space $A^1_2$ of the band-limited scheme with $F = f$
    almost everywhere on $\Omega$;
  \item $\|f\|_{L_2(\Omega)} + \|(2^k \omega_r(f, 2^{-k}, \Omega)_2)_k\|_{\ell_2} < \infty$.
\end{enumerate}
For every integer $r \geq 2$ and every $f \in H^1(\Omega)$ there is $F \in L_2(V)$ with $F = f$
almost everywhere on $\Omega$ and $\|F\|_2 + \|(2^k \omega_r(F, 2^{-k})_2)_k\|_{\ell_2} < \infty$.
For every $r \in \N$ and all $F, f \in L_2(V)$ with $F = f$ almost everywhere on
$\Omega$,
\[
  \|f\|_{L_2(\Omega)} + \|(2^k \omega_r(f, 2^{-k}, \Omega)_2)_k\|_{\ell_2}
  \leq \|F\|_2 + \|(2^k \omega_r(F, 2^{-k})_2)_k\|_{\ell_2}.
\]
\end{theorem}

\begin{proof}
\leanok
(1) and (2) agree by \Cref{thm:sobolev-one-two-besov}, with $F = f$ on $\Omega$ written as
$1_\Omega F = 1_\Omega f$ (\Cref{lem:indicator-lp-eq-iff}). (2) and (3) agree by
\Cref{thm:bounded-lipschitz-besov} with $p = q = 2$ and $s = 1 < r$. The function $F$ of
\Cref{def:sobolev-lipschitz-domain} has the finite norm by \Cref{thm:sobolev-one-two-besov}. The
inequality is \Cref{lem:restriction-moduli}.
\end{proof}

The extension theorem for bounded Lipschitz domains (\Cref{thm:bounded-lipschitz-besov}) and the
trace theorem below are qualitative, so the statements on $H^1(\Omega)$ are too; the inequality
gives the restriction with constant $1$.

\begin{theorem}[The weak gradient]
\label{thm:sobolev-lipschitz-domain-weak}
\lean{Rd.BoundedLipschitzDomain.MemSobolevOn.memWeakSobolev}
\leanok
\uses{def:sobolev-lipschitz-domain, lem:weak-derivative-basic}
For every $f \in H^1(\Omega)$, $f \in \Sob{1}{2}(\Omega)$.
\end{theorem}

\begin{proof}
\leanok
The function $F \in \Sob{1}{2}(V)$ of \Cref{def:sobolev-lipschitz-domain} lies in
$\Sob{1}{2}(\Omega)$, and so does $f$, which agrees with $F$ almost everywhere on $\Omega$
(\Cref{lem:weak-derivative-basic}).
\end{proof}

In particular $f$ and its weak derivatives $\partial_v f$, $v \in V$, lie in $L_2(\Omega)$. This
inclusion holds by the definition of $H^1(\Omega)$ through extensions; the converse inclusion is
the extension theorem (\Cref{cor:sobolev-extension-two}).

For a resolution of unity $\psi$ on $\Gamma$ (\Cref{def:boundary-partition}), as in
\Cref{thm:boundary-trace-space} with $p = q = 2$, $L = 2$ and $s + 1/p = 1$, a function $g$ on
$V$ is a boundary trace of $f \in L_2(V)$ if $g = \operatorname{Tr}_\Gamma F$ almost everywhere
on $\Gamma$ for every $F \in L_2(V)$ with $F = f$ almost everywhere on $\Omega$ and
$\|F\|_2 + \|(2^k \omega_2(F, 2^{-k})_2)_k\|_{\ell_2} < \infty$. Both $B^{1/2}_{2,2}(\Gamma)$ and
the boundary trace are taken with respect to $\psi$. By \Cref{thm:sobolev-one-two-besov} with
$r = 2$, these $F$ are the $F \in \Sob{1}{2}(V)$ with $F = f$ almost everywhere on $\Omega$
(\Cref{lem:sobolev-lipschitz-domain-trace-extensions}).

\begin{lemma}[Boundary traces of $H^1$ functions]
\label{lem:sobolev-lipschitz-domain-trace-extensions}
\lean{Rd.BoundedLipschitzDomain.isBoundaryTrace_iff_forall_memSobolev,
  Rd.BoundedLipschitzDomain.hasBoundaryTrace_iff_memSobolevOn}
\leanok
\uses{thm:sobolev-one-two-besov, thm:boundary-trace-space, lem:indicator-lp-eq-iff,
  def:boundary-trace-finite-norm, thm:sobolev-lipschitz-domain-intrinsic}
For every resolution of unity $\psi$ on $\Gamma$, every $f \in L_2(V)$ and every function $g$ on
$V$, $g$ is a boundary trace of $f$ if and only if $g = \operatorname{Tr}_\Gamma F$ almost
everywhere on $\Gamma$ for every $F \in \Sob{1}{2}(V)$ with $F = f$ almost everywhere on
$\Omega$. For $L = 2$, $\sigma = 1$ and $q = 2$, $f$ has the boundary trace $g$ in the sense of
\Cref{def:boundary-trace-finite-norm} if and only if $f \in H^1(\Omega)$ and $g$ is a boundary
trace of $f$.
\end{lemma}

\begin{proof}
\leanok
By \Cref{thm:sobolev-one-two-besov} with $r = 2$, $F \in \Sob{1}{2}(V)$ if and only if
$\|F\|_2 + \|(2^k \omega_2(F, 2^{-k})_2)_k\|_{\ell_2} < \infty$, and $F = f$ almost everywhere on
$\Omega$ if and only if $1_\Omega F = 1_\Omega f$ (\Cref{lem:indicator-lp-eq-iff}). By
\Cref{thm:sobolev-lipschitz-domain-intrinsic} with $r = 2$, the finite intrinsic norm of
\Cref{def:boundary-trace-finite-norm} describes $H^1(\Omega)$.
\end{proof}

\begin{theorem}[The trace space of $H^1(\Omega)$]
\label{thm:sobolev-lipschitz-domain-trace}
\lean{Rd.BoundedLipschitzDomain.exists_isBoundaryTrace_of_memSobolevOn,
  Rd.BoundedLipschitzDomain.exists_memSobolevOn_of_memBoundaryBesov,
  Rd.BoundedLipschitzDomain.MemSobolevOn.isBoundaryTrace_unique,
  Rd.BoundedLipschitzDomain.memBoundaryBesov_iff_exists_memSobolevOn}
\leanok
\uses{thm:sobolev-lipschitz-domain-intrinsic, def:boundary-besov, def:boundary-trace,
  thm:boundary-trace-space}
For every resolution of unity $\psi$ on $\Gamma$ (\Cref{def:boundary-partition}), with
$B^{1/2}_{2,2}(\Gamma)$ and boundary traces taken with respect to $\psi$: for every
$f \in H^1(\Omega)$ there is a boundary trace $g \in B^{1/2}_{2,2}(\Gamma)$ of $f$, and every two
boundary traces of $f$ agree almost everywhere on $\Gamma$; for every function $g$ on $V$,
$g \in B^{1/2}_{2,2}(\Gamma)$ if and only if $g$ is the boundary trace of some
$f \in H^1(\Omega)$; in particular, for every $g \in B^{1/2}_{2,2}(\Gamma)$ there is
$f \in H^1(\Omega)$ with boundary trace $g$.
\end{theorem}

\begin{proof}
\leanok
Here $s = 1/2$ and $p = 2$, so $s + 1/p = 1$, and by
\Cref{thm:sobolev-lipschitz-domain-intrinsic} with $r = 2$ the functions of $H^1(\Omega)$ are those
with
$\|f\|_{L_2(\Omega)} + \|(2^k \omega_2(f, 2^{-k}, \Omega)_2)_k\|_{\ell_2} < \infty$. For them
\Cref{thm:boundary-trace-space} gives existence and uniqueness of the boundary trace and the trace
space.
\end{proof}

\section{Sobolev spaces on open sets as Hilbert spaces}
\label{sec:sobolev-restriction-spaces}

Throughout this section $\Omega \subseteq V$ is open and $s \in \R$. An element $U$ of Mathlib's
bundled space $H^s(V) = H^{s,2}(V, \C)$ (\texttt{BesselPotentialSpace}) is a tempered
distribution with $J_s U \in L_2(V)$, and the inner product
$\langle U, W \rangle = \langle J_s U, J_s W \rangle_{L_2}$ makes $H^s(V)$ a Hilbert space. For
a Schwartz function $u$, $U(u)$ is the value of the distribution $U$ at $u$, and a test function
$\varphi \in \mathcal D(\Omega)$ is read as the Schwartz function of \Cref{def:test-schwartz}.
The Sobolev spaces on $\Omega$ are built inside $H^s(V)$ from orthogonal complements, orthogonal
projections and closures, as in Chandler-Wilde, Hewett and Moiola
\cite[\S 3.1]{ChandlerWildeHewettMoiola} and McLean \cite[Chapter 3]{McLean}, so they are
Hilbert spaces as well, and Mathlib's Hilbert space theory (quotients, projections, duals,
compact operators) applies to them.

\begin{lemma}[Evaluation is continuous]
\label{lem:sobolev-evaluation-continuous}
\lean{Rd.besselEvalCLM, Rd.besselEvalCLM_apply, Rd.continuous_toDistr_apply}
\leanok
For every $s \in \R$ and every Schwartz function $u$, the map $U \mapsto U(u)$ is a continuous
linear functional on $H^s(V)$.
\end{lemma}

\begin{proof}
\leanok
By Mathlib's \texttt{BesselPotentialSpace.\allowbreak besselPotential\_\allowbreak
neg\_\allowbreak toLp\_\allowbreak eq}, $U = J_{-s}(J_s U)$, so $U \mapsto U(u)$ is the
composition of $U \mapsto J_s U \in L_2(V)$, the embedding of $L_2(V)$ into the tempered
distributions, $J_{-s}$, and evaluation at $u$, all continuous linear maps (Mathlib's
\texttt{PointwiseConvergenceCLM.\allowbreak evalCLM}).
\end{proof}

\begin{definition}[Distributions vanishing on $\Omega$]
\label{def:sobolev-vanishing}
\lean{Rd.vanishOn, Rd.mem_vanishOn_iff}
\leanok
\uses{lem:sobolev-evaluation-continuous, def:test-schwartz}
$N^s(\Omega)$ is the set of the $U \in H^s(V)$ with $U(\varphi) = 0$ for every
$\varphi \in \mathcal D(\Omega)$.
\end{definition}

\begin{lemma}[$N^s(\Omega)$ is closed]
\label{lem:sobolev-vanishing-closed}
\lean{Rd.isClosed_vanishOn, Rd.completeSpace_vanishOn}
\leanok
\uses{def:sobolev-vanishing}
For every $s \in \R$ and every open $\Omega$, $N^s(\Omega)$ is a closed subspace of $H^s(V)$,
hence complete.
\end{lemma}

\begin{proof}
\leanok
$N^s(\Omega)$ is the intersection of the kernels of the continuous linear functionals
$U \mapsto U(\varphi)$, $\varphi \in \mathcal D(\Omega)$ (\Cref{lem:sobolev-evaluation-continuous}).
\end{proof}

\begin{proposition}[Vanishing on $\Omega$ for $s = 1$]
\label{prop:sobolev-vanishing-one}
\lean{Rd.toDistr_testSchwartz_eq_integral, Rd.mem_vanishOn_one_iff,
  Rd.mem_vanishOn_one_iff_restrict_eq_zero, Rd.vanishOn_one_eq_ker}
\leanok
\uses{def:sobolev-vanishing, def:sobolev-one-to-l2, lem:sobolev-one-to-l2}
For every $U \in H^1(V)$ and every $\varphi \in \mathcal D(\Omega)$,
$U(\varphi) = \int \varphi \, \iota U$. For every $U \in H^1(V)$, $U \in N^1(\Omega)$ iff
$\iota U = 0$ almost everywhere on $\Omega$, iff $(\iota U)|_\Omega = 0$ in $L_2(\Omega)$, and
$N^1(\Omega)$ is the kernel of $U \mapsto (\iota U)|_\Omega$.
\end{proposition}

\begin{proof}
\leanok
The tempered distribution of $\iota U$ is $U$ (\Cref{lem:sobolev-one-to-l2}). If $\iota U = 0$
almost everywhere on $\Omega$, then $\varphi \, \iota U = 0$ almost everywhere, since
$\varphi = 0$ outside $\Omega$. Conversely, $\iota U$ is locally integrable on $\Omega$ and has
vanishing integrals against all smooth real functions with compact support in $\Omega$, so it
vanishes almost everywhere on $\Omega$ (Mathlib's \texttt{IsOpen.\allowbreak ae\_\allowbreak
eq\_\allowbreak zero\_\allowbreak of\_\allowbreak integral\_\allowbreak contDiff\_\allowbreak
smul\_\allowbreak eq\_\allowbreak zero}).
\end{proof}

\begin{definition}[The Sobolev space $H^s(\Omega)$]
\label{def:sobolev-restriction-space}
\lean{Rd.RestrictionSpace, Rd.restrictProj}
\leanok
\uses{lem:sobolev-vanishing-closed}
$H^s(\Omega) = N^s(\Omega)^\perp$ is the orthogonal complement of $N^s(\Omega)$ in $H^s(V)$, a
Hilbert space, and the restriction $r_\Omega : H^s(V) \to H^s(\Omega)$ is the orthogonal
projection onto it.
\end{definition}

\begin{lemma}[The restriction]
\label{lem:sobolev-restriction-basic}
\lean{Rd.norm_restrictProj_apply_le, Rd.restrictProj_eq_zero_iff, Rd.restrictProj_coe,
  Rd.restrictProj_surjective, Rd.sub_restrictProj_mem}
\leanok
\uses{def:sobolev-restriction-space}
For every $U \in H^s(V)$: $\|r_\Omega U\| \leq \|U\|$; $r_\Omega U = 0$ iff $U \in N^s(\Omega)$;
and $U - r_\Omega U \in N^s(\Omega)$. For every $u \in H^s(\Omega)$, $r_\Omega u = u$, and
$r_\Omega$ is surjective.
\end{lemma}

\begin{proof}
\leanok
Properties of orthogonal projections onto $K^\perp$ for the closed subspace $K = N^s(\Omega)$,
with $K^{\perp\perp} = K$ (Mathlib's \texttt{Submodule.\allowbreak orthogonal\_\allowbreak
orthogonal}).
\end{proof}

\begin{theorem}[$H^s(\Omega)$ as a quotient]
\label{thm:sobolev-restriction-quotient}
\lean{Rd.quotientEquivRestrictionSpace, Rd.quotientEquivRestrictionSpace_mk,
  Rd.norm_restrictProj_le_of_mem,
  Rd.exists_norm_restrictProj_eq}
\leanok
\uses{lem:sobolev-restriction-basic}
For every $s \in \R$ and every open $\Omega$, $[U] \mapsto r_\Omega U$ is an isometric
isomorphism from $H^s(V) / N^s(\Omega)$ onto $H^s(\Omega)$. For every $U \in H^s(V)$,
$\|r_\Omega U\| \leq \|U + W\|$ for every $W \in N^s(\Omega)$, with equality for some
$W \in N^s(\Omega)$.
\end{theorem}

\begin{proof}
\leanok
A Hilbert space modulo a closed subspace $K$ is isometric to $K^\perp$ (Mathlib's
\texttt{Submodule.\allowbreak quotientEquivOrthogonal}), and $[U] = [r_\Omega U]$ by
\Cref{lem:sobolev-restriction-basic}. The norm of $[U]$ is at most $\|U + W\|$ for
$W \in N^s(\Omega)$, and $W = r_\Omega U - U$ gives equality.
\end{proof}

So the norm of $H^s(\Omega)$ is the minimal norm of an extension, the norm of the restriction
space of Chandler-Wilde, Hewett and Moiola \cite[\S 3.1.4]{ChandlerWildeHewettMoiola}.

\begin{definition}[$H^1(\Omega)$ in $L_2(\Omega)$]
\label{def:sobolev-restriction-l2}
\lean{Rd.restrictL2, Rd.restrictL2_apply, Rd.restrictL2_eq_iff_ae_eq}
\leanok
\uses{def:sobolev-restriction-space, def:sobolev-one-to-l2}
For $u \in H^1(\Omega) \subseteq H^1(V)$, $R_\Omega u = (\iota u)|_\Omega \in L_2(\Omega)$;
$R_\Omega$ is a bounded linear operator. For every $u \in H^1(\Omega)$ and every $f \in L_2(V)$,
$R_\Omega u = f|_\Omega$ if and only if $\iota u = f$ almost everywhere on $\Omega$.
\end{definition}

\begin{proposition}[$R_\Omega$ is injective with norm at most $1$]
\label{prop:sobolev-restriction-l2}
\lean{Rd.restrictL2_injective, Rd.norm_restrictL2_apply_le, Rd.restrictL2_restrictProj,
  Rd.restrictProj_eq_of_ae_eq}
\leanok
\uses{def:sobolev-restriction-l2, prop:sobolev-vanishing-one, lem:sobolev-restriction-basic,
  lem:sobolev-one-to-l2}
For every open $\Omega$, $R_\Omega$ is injective and $\|R_\Omega u\|_{L_2(\Omega)} \leq \|u\|$
for every $u \in H^1(\Omega)$. For every $U \in H^1(V)$, $R_\Omega(r_\Omega U) =
(\iota U)|_\Omega$. For every $U \in H^1(V)$ and every $u \in H^1(\Omega)$ with
$\iota U = \iota u$ almost everywhere on $\Omega$, $r_\Omega U = u$.
\end{proposition}

\begin{proof}
\leanok
If $R_\Omega u = 0$, then $u \in N^1(\Omega) \cap N^1(\Omega)^\perp = \{0\}$
(\Cref{prop:sobolev-vanishing-one}). The bound is $\|\iota u\|_2 \leq \|u\|$
(\Cref{lem:sobolev-one-to-l2}). $U - r_\Omega U$ lies in $N^1(\Omega)$, the kernel of
$U \mapsto (\iota U)|_\Omega$. If $\iota U = \iota u$ on $\Omega$, then
$R_\Omega(r_\Omega U) = (\iota U)|_\Omega = R_\Omega u$, and $R_\Omega$ is injective.
\end{proof}

\begin{theorem}[Rellich's theorem for $H^1(\Omega)$]
\label{thm:sobolev-restriction-rellich}
\lean{Rd.isCompactOperator_restrictL2}
\leanok
\uses{def:sobolev-restriction-l2, thm:rellich-compactness}
For every bounded open set $\Omega \subseteq V$, $R_\Omega : H^1(\Omega) \to L_2(\Omega)$ is a
compact operator.
\end{theorem}

\begin{proof}
\leanok
$R_\Omega$ is the compact operator $U \mapsto (\iota U)|_\Omega$ of
\Cref{thm:rellich-compactness} composed with the bounded inclusion of $H^1(\Omega)$ into
$H^1(V)$ (Mathlib's \texttt{IsCompactOperator.\allowbreak comp\_\allowbreak clm}).
\end{proof}

For every bounded open set, $R_\Omega$ is compact; for the restriction space $H^1(\Omega)$ this
is the compactness that Hiptmair, Pauly and Schulz require (their Assumption F)
\cite{HiptmairPaulySchulz}.

\begin{lemma}[Weak derivatives of $H^1(V)$]
\label{lem:sobolev-one-weak-gradient-norm}
\lean{Rd.memSobolev_toL2, Rd.exists_toL2_eq_of_memSobolev, Rd.exists_hasWeakLineDeriv_toL2,
  Rd.sq_norm_eq_of_hasWeakLineDeriv, Rd.norm_le_norm_toL2_add_sum}
\leanok
\uses{lem:sobolev-one-to-l2, thm:sobolev-one-two, cor:sobolev-tempered,
  thm:weak-derivative-tempered}
For every $U \in H^1(V)$, $\iota U \in \Sob{1}{2}(V)$, and for every $v \in V$ the weak
derivative $\partial_v \iota U$ on $V$ exists in $L_2(V)$. For every $f \in L_2(V)$ with
$f \in \Sob{1}{2}(V)$ there is $U \in H^1(V)$ with $\iota U = f$. For every orthonormal basis
$(e_i)_{i \in I}$ of $V$, every $U \in H^1(V)$ and all $g_i \in L_2(V)$ such that $g_i$ is the
weak derivative $\partial_{e_i} \iota U$ on $V$ for every $i \in I$,
\[
  \|U\|^2 = \|\iota U\|_2^2 + (2 \pi)^{-2} \sum_i \|g_i\|_2^2 ,
\]
and in particular $\|U\| \leq \|\iota U\|_2 + \sum_i \|g_i\|_2$.
\end{lemma}

\begin{proof}
\leanok
The tempered distribution of $\iota U$ is $U \in H^1(V)$ (\Cref{lem:sobolev-one-to-l2}), so
$\iota U \in \Sob{1}{2}(V)$ by \Cref{thm:sobolev-one-two}, and its tempered derivatives are
$L_2$ functions (\Cref{cor:sobolev-tempered}), which are weak derivatives
(\Cref{thm:weak-derivative-tempered}). Conversely, for $f \in \Sob{1}{2}(V)$ the tempered
distribution of $f$ lies in $H^1(V)$ (\Cref{thm:sobolev-one-two}), and $f$ is its $L_2$ function
(\Cref{lem:sobolev-one-to-l2}). The identity is
$(2\pi)^2 \|J_1 U\|_2^2 = (2\pi)^2 \|\iota U\|_2^2 + \sum_i \|g_i\|_2^2$ from
\Cref{thm:sobolev-one-two}, with $\|U\| = \|J_1 U\|_2$. The bound follows from
$(2\pi)^{-2} \leq 1$ and $\sum_i a_i^2 \leq (\sum_i a_i)^2$ for $a_i \geq 0$.
\end{proof}

\begin{lemma}[The weak derivative as an operator]
\label{lem:sobolev-one-weak-derivative-operator}
\lean{Rd.norm_le_of_hasWeakLineDeriv_toL2, Rd.weakDerivL2Linear,
  Rd.hasWeakLineDeriv_weakDerivL2Linear, Rd.weakDerivL2, Rd.norm_weakDerivL2_le,
  Rd.hasWeakLineDeriv_weakDerivL2, Rd.weakDerivL2_eq}
\leanok
\uses{lem:sobolev-one-weak-gradient-norm, thm:sobolev-one-two, lem:weak-derivative-basic}
For every $v \in V$, every $U \in H^1(V)$ and every $g \in L_2(V)$ that is the weak derivative
$\partial_v \iota U$ on $V$, $\|g\|_2 \leq 2 \pi \|v\| \|U\|$. For every $v \in V$, the map
$\partial_v \colon H^1(V) \to L_2(V)$ that sends $U$ to the weak derivative $\partial_v \iota U$
on $V$ is a bounded linear operator of norm at most $2 \pi \|v\|$.
\end{lemma}

\begin{proof}
\leanok
The bound is $\|\partial_v f\|_2 \leq 2 \pi \|v\| \|J_1 f\|_2$ of \Cref{thm:sobolev-one-two} with
$J_1 \iota U = U$. The weak derivative exists (\Cref{lem:sobolev-one-weak-gradient-norm}) and is
unique in $L_2(V)$ (\Cref{lem:weak-derivative-basic}), hence linear in $U$.
\end{proof}

\begin{proposition}[The weak gradient on $\Omega$]
\label{prop:sobolev-restriction-weak-gradient}
\lean{Rd.memWeakSobolev_toL2, Rd.sq_norm_restr_add_sum_le, Rd.sq_norm_restrictL2_add_sum_le}
\leanok
\uses{lem:sobolev-one-weak-gradient-norm, def:sobolev-restriction-l2, lem:weak-derivative-basic}
For every open $\Omega$ and every $U \in H^1(V)$, $\iota U \in \Sob{1}{2}(\Omega)$. For every
orthonormal basis $(e_i)_{i \in I}$ of $V$, every $u \in H^1(\Omega)$ and all $g_i \in L_2(V)$
such that $g_i$ is the weak derivative $\partial_{e_i} \iota u$ on $V$ for every $i \in I$,
\[
  \|R_\Omega u\|_{L_2(\Omega)}^2 + (2 \pi)^{-2} \sum_i \|g_i|_\Omega\|_{L_2(\Omega)}^2
  \leq \|u\|^2 .
\]
For every orthonormal basis $(e_i)_{i \in I}$ of $V$, every $u \in H^1(\Omega)$ and all
$g_i \in L_2(\Omega)$ such that $g_i$ is the weak derivative $\partial_{e_i} \iota u$ on $\Omega$
for every $i \in I$,
\[
  \|R_\Omega u\|_{L_2(\Omega)}^2 + (2 \pi)^{-2} \sum_i \|g_i\|_{L_2(\Omega)}^2 \leq \|u\|^2 .
\]
\end{proposition}

\begin{proof}
\leanok
$\iota U \in \Sob{1}{2}(V)$ (\Cref{lem:sobolev-one-weak-gradient-norm}), and weak derivatives on
$V$ are weak derivatives on $\Omega$ (\Cref{lem:weak-derivative-basic}). Restriction to $\Omega$
has norm at most $1$ on $L_2$; apply \Cref{lem:sobolev-one-weak-gradient-norm} to
$u \in H^1(V)$. Weak derivatives on $\Omega$ agree almost everywhere on $\Omega$ with the
restrictions of the weak derivatives on $V$ (\Cref{lem:weak-derivative-basic}).
\end{proof}

Since $(2\pi)^2 \geq 1$, the norm
$\bigl(\|f\|_{L_2(\Omega)}^2 + \sum_i \|\partial_{e_i} f\|_{L_2(\Omega)}^2\bigr)^{1/2}$ of
$\Sob{1}{2}(\Omega)$ satisfies $\|R_\Omega u\|_{\Sob{1}{2}(\Omega)} \leq 2\pi \|u\|$ for every
$u \in H^1(\Omega)$.

\begin{lemma}[Weak derivatives vanish where the function vanishes]
\label{lem:weak-derivative-vanishing}
\lean{Rd.weakDeriv_ae_zero_of_mem_vanishOn}
\leanok
\uses{lem:weak-derivative-basic, prop:sobolev-vanishing-one}
For every $W \in N^1(\Omega)$, every $k \in \N$, every $v \in V^k$ and every weak derivative
$g = \partial^v \iota W$ on $V$, $g = 0$ almost everywhere on $\Omega$.
\end{lemma}

\begin{proof}
\leanok
$\iota W = 0$ almost everywhere on $\Omega$ (\Cref{prop:sobolev-vanishing-one}), so $g = 0$
almost everywhere on $\Omega$ (\Cref{lem:weak-derivative-basic}, item (7)).
\end{proof}

\begin{theorem}[$H^1(\Omega)$ of a bounded Lipschitz domain]
\label{thm:sobolev-restriction-lipschitz}
\lean{Rd.BoundedLipschitzDomain.memSobolevOn_iff_exists_restrictL2}
\leanok
\uses{def:sobolev-lipschitz-domain, prop:sobolev-restriction-l2,
  lem:sobolev-one-weak-gradient-norm, thm:sobolev-one-two, lem:sobolev-one-to-l2}
For every bounded Lipschitz domain $\Omega$ and every $f \in L_2(V)$, $f$ lies in $H^1(\Omega)$
in the sense of \Cref{def:sobolev-lipschitz-domain} iff $f|_\Omega = R_\Omega u$ for some $u$ in
the Hilbert space $H^1(\Omega)$ of \Cref{def:sobolev-restriction-space}.
\end{theorem}

\begin{proof}
\leanok
If $F \in \Sob{1}{2}(V)$ equals $f$ on $\Omega$, then $F = \iota U$ for some $U \in H^1(V)$
(\Cref{thm:sobolev-one-two}, \Cref{lem:sobolev-one-to-l2}), and $u = r_\Omega U$ has
$R_\Omega u = F|_\Omega = f|_\Omega$ (\Cref{prop:sobolev-restriction-l2}). Conversely,
$F = \iota u$ lies in $\Sob{1}{2}(V)$ (\Cref{lem:sobolev-one-weak-gradient-norm}) and equals $f$
almost everywhere on $\Omega$.
\end{proof}

\begin{definition}[$\tilde H^s(\Omega)$ and $H^s_0(\Omega)$]
\label{def:sobolev-tilde-zero-boundary}
\lean{Rd.ofTest, Rd.toDistr_ofTest, Rd.testSpan, Rd.tildeSpace, Rd.zeroBoundarySpace,
  Rd.tildeToRestrictionL}
\leanok
\uses{def:sobolev-restriction-space, def:test-schwartz}
Every $\varphi \in \mathcal D(\Omega)$, a Schwartz function, is an element of $H^s(V)$.
$T^s(\Omega)$ is the complex linear span of these elements in $H^s(V)$, $\tilde H^s(\Omega)$
its closure in $H^s(V)$, and $H^s_0(\Omega)$ the closure of $r_\Omega(T^s(\Omega))$ in
$H^s(\Omega)$. The map
$\tilde H^s(\Omega) \to H^s(\Omega)$, $U \mapsto r_\Omega U$, is a bounded linear operator.
\end{definition}

\begin{lemma}[$\tilde H^s(\Omega)$ and $H^s_0(\Omega)$ are Hilbert spaces]
\label{lem:sobolev-tilde-zero-boundary-closed}
\lean{Rd.isClosed_tildeSpace, Rd.completeSpace_tildeSpace, Rd.isClosed_zeroBoundarySpace,
  Rd.completeSpace_zeroBoundarySpace}
\leanok
\uses{def:sobolev-tilde-zero-boundary}
For every $s \in \R$ and every open $\Omega$, $\tilde H^s(\Omega)$ is closed in $H^s(V)$ and
$H^s_0(\Omega)$ is closed in $H^s(\Omega)$; both are complete.
\end{lemma}

\begin{proof}
\leanok
Closures are closed, and closed subspaces of complete spaces are complete.
\end{proof}

\begin{lemma}[Restriction maps $\tilde H^s(\Omega)$ into $H^s_0(\Omega)$]
\label{lem:sobolev-tilde-restriction-range}
\lean{Rd.range_tildeToRestrictionL_le, Rd.closure_range_tildeToRestrictionL}
\leanok
\uses{def:sobolev-tilde-zero-boundary, lem:sobolev-tilde-zero-boundary-closed}
For every $s \in \R$ and every open $\Omega$, $r_\Omega(\tilde H^s(\Omega)) \subseteq
H^s_0(\Omega)$, and $H^s_0(\Omega)$ is the closure of $r_\Omega(\tilde H^s(\Omega))$.
\end{lemma}

\begin{proof}
\leanok
$r_\Omega$ is continuous, so it maps the closure of $T^s(\Omega)$ into the closure of
$r_\Omega(T^s(\Omega))$; and $r_\Omega(T^s(\Omega)) \subseteq r_\Omega(\tilde H^s(\Omega))$.
\end{proof}

\begin{lemma}[Test functions are Schwartz functions supported in $\Omega$]
\label{lem:sobolev-test-span-schwartz}
\lean{Rd.toL2_ofTest_ae, Rd.exists_schwartz_of_mem_testSpan, Rd.exists_opens_toL2_ae_eq_zero}
\leanok
\uses{def:sobolev-tilde-zero-boundary, lem:sobolev-one-to-l2}
For every $\varphi \in \mathcal D(\Omega)$, $\iota \varphi = \varphi$ almost everywhere. For
every $U \in T^s(\Omega)$ there is a Schwartz function $w$ with closed support in $\Omega$ whose
tempered distribution is $U$. For every $U \in T^1(\Omega)$ there is an open set $O \supseteq
V \setminus \Omega$ with $\iota U = 0$ almost everywhere on $O$.
\end{lemma}

\begin{proof}
\leanok
An $L_2$ function is determined by its tempered distribution (\Cref{lem:sobolev-one-to-l2}).
Sums and multiples of Schwartz functions with closed support in $\Omega$ are such functions. For
$U \in T^1(\Omega)$, $\iota U = w$ almost everywhere, and $O$ is the complement of the closed
support of $w$.
\end{proof}

\begin{lemma}[Schwartz functions in $H^s(V)$]
\label{lem:sobolev-of-schwartz}
\lean{Rd.ofSchwartz, Rd.toDistr_ofSchwartz, Rd.ofSchwartz_eq, Rd.ofTest_eq_ofSchwartz,
  Rd.toL2_ofSchwartz, Rd.weakDerivL2_ofSchwartz}
\leanok
\uses{lem:bessel-schwartz-distribution, def:sobolev-tilde-zero-boundary, lem:sobolev-one-to-l2,
  lem:sobolev-one-weak-derivative-operator, cor:weak-derivative-schwartz}
For every $s \in \R$, the map that sends a Schwartz function $w$ to the element of $H^s(V)$ with
the tempered distribution $w$ is a continuous linear map from the Schwartz space to $H^s(V)$; on
$\mathcal D(\Omega)$ it is the inclusion of \Cref{def:sobolev-tilde-zero-boundary}. For every
Schwartz function $w$, regarded as an element of $H^1(V)$ ($s = 1$), and every $v \in V$,
$\iota w = w$ and $\partial_v \iota w = \partial_v w$ in $L_2(V)$.
\end{lemma}

\begin{proof}
\leanok
$w \mapsto J_s w$ is continuous on the Schwartz space, the inclusion of the Schwartz space into
$L_2(V)$ is continuous, and $J_s$ is an isometry of $H^s(V)$ onto $L_2(V)$; the tempered
distribution of the image is $J_{-s} J_s w = w$ (\Cref{lem:bessel-schwartz-distribution}). An
$L_2$ function is determined by its tempered distribution (\Cref{lem:sobolev-one-to-l2}), and
$\partial_v w$ is the weak derivative of $w$ (\Cref{cor:weak-derivative-schwartz}).
\end{proof}

\begin{theorem}[Schwartz functions are dense in $H^s(V)$]
\label{thm:sobolev-schwartz-dense}
\lean{Rd.denseRange_ofSchwartz, Rd.dense_ofSchwartz_hasCompactSupport,
  Rd.ofSchwartz_mem_testSpan, Rd.tildeSpace_top}
\leanok
\uses{lem:sobolev-of-schwartz, lem:bessel-schwartz-distribution, lem:schwartz-cutoff-tendsto,
  lem:weak-derivative-complex-test, def:sobolev-tilde-zero-boundary}
For every $s \in \R$, the Schwartz functions and the compactly supported Schwartz functions are
dense in $H^s(V)$. For every $s \in \R$, every open $\Omega \subseteq V$ and every Schwartz
function $w$ with compact support in $\Omega$, $w \in T^s(\Omega)$. For every $s \in \R$,
$\tilde H^s(V) = H^s(V)$.
\end{theorem}

\begin{proof}
\leanok
$J_s$ maps the Schwartz space onto itself, with inverse $J_{-s}$
(\Cref{lem:bessel-schwartz-distribution}), the Schwartz functions are dense in $L_2(V)$ (Mathlib's
\texttt{SchwartzMap.\allowbreak denseRange\_\allowbreak toLpCLM}), and $J_s$ is an isometry of
$H^s(V)$ onto $L_2(V)$. The cut-offs of a Schwartz function converge to it in the Schwartz space
(\Cref{lem:schwartz-cutoff-tendsto}), hence in $H^s(V)$ by \Cref{lem:sobolev-of-schwartz}. A
Schwartz function with compact support in $\Omega$ is $a + i b$ with real test functions
$a, b \in \mathcal D(\Omega)$ (\Cref{lem:weak-derivative-complex-test}), so it lies in
$T^s(\Omega)$; for $\Omega = V$ this gives $\tilde H^s(V) = H^s(V)$.
\end{proof}

\begin{lemma}[Functions with disjoint supports]
\label{lem:norm-disjoint-supports}
\lean{Rd.norm_le_norm_add_of_ae_eq_zero}
\leanok
For every measure space $(X, \mu)$, every normed group $E$, every $0 < p \leq \infty$, every
measurable $S \subseteq X$ and all $a, b \in L_p(\mu; E)$ such that $a = 0$ almost everywhere on
$X \setminus S$ and $b = 0$ almost everywhere on $S$, $\|a\|_p \leq \|a + b\|_p$.
\end{lemma}

\begin{proof}
\leanok
$|a + b| \geq |a|$ almost everywhere.
\end{proof}

\begin{proposition}[Adding a function vanishing on $\Omega$]
\label{prop:sobolev-test-norm-minimal}
\lean{Rd.norm_le_norm_add_of_mem_testSpan}
\leanok
\uses{lem:sobolev-one-weak-gradient-norm, lem:sobolev-test-span-schwartz,
  lem:weak-derivative-vanishing, lem:norm-disjoint-supports, prop:sobolev-vanishing-one}
For every $U \in T^1(\Omega)$ and every $W \in N^1(\Omega)$, $\|U\| \leq \|U + W\|$.
\end{proposition}

\begin{proof}
\leanok
Let $(e_i)$ be an orthonormal basis of $V$ and $g_i$, $h_i$ the weak derivatives
$\partial_{e_i}$ of $\iota U$ and $\iota W$ on $V$; $g_i + h_i$ is the one of $\iota(U + W)$.
By \Cref{lem:sobolev-one-weak-gradient-norm},
$\|U\|^2 = \|\iota U\|_2^2 + (2\pi)^{-2} \sum_i \|g_i\|_2^2$, and the same for $U + W$ with
$\iota U + \iota W$ and $g_i + h_i$. $\iota U$ and the $g_i$ vanish almost everywhere outside
$\Omega$ (\Cref{lem:sobolev-test-span-schwartz}, \Cref{lem:weak-derivative-vanishing}), and
$\iota W$ and the $h_i$ vanish almost everywhere on $\Omega$
(\Cref{prop:sobolev-vanishing-one}, \Cref{lem:weak-derivative-vanishing}), so every $L_2$ norm
grows by \Cref{lem:norm-disjoint-supports}.
\end{proof}

\begin{theorem}[Restriction is isometric on $\tilde H^1(\Omega)$]
\label{thm:sobolev-tilde-restriction-isometry}
\lean{Rd.norm_restrictProj_of_mem_testSpan, Rd.norm_restrictProj_of_mem_tildeSpace,
  Rd.tildeToRestriction, Rd.tildeToRestriction_apply}
\leanok
\uses{prop:sobolev-test-norm-minimal, thm:sobolev-restriction-quotient,
  lem:sobolev-restriction-basic}
For every open $\Omega \subseteq V$ and every $U \in \tilde H^1(\Omega)$,
$\|r_\Omega U\| = \|U\|$.
\end{theorem}

\begin{proof}
\leanok
For $U \in T^1(\Omega)$, $\|r_\Omega U\| = \|U + W\|$ for some $W \in N^1(\Omega)$
(\Cref{thm:sobolev-restriction-quotient}), which is at least $\|U\|$
(\Cref{prop:sobolev-test-norm-minimal}), and $\|r_\Omega U\| \leq \|U\|$
(\Cref{lem:sobolev-restriction-basic}). Both sides are continuous in $U$, so the identity
extends to the closure $\tilde H^1(\Omega)$.
\end{proof}

\begin{theorem}[$\tilde H^1(\Omega) = H^1_0(\Omega)$ isometrically]
\label{thm:sobolev-tilde-zero-boundary}
\lean{Rd.range_tildeToRestriction, Rd.tildeEquivZeroBoundary,
  Rd.coe_tildeEquivZeroBoundary_apply}
\leanok
\uses{thm:sobolev-tilde-restriction-isometry, lem:sobolev-tilde-restriction-range,
  lem:sobolev-tilde-zero-boundary-closed}
For every open set $\Omega \subseteq V$, $U \mapsto r_\Omega U$ is an isometric isomorphism from
$\tilde H^1(\Omega)$ onto $H^1_0(\Omega)$.
\end{theorem}

\begin{proof}
\leanok
By \Cref{thm:sobolev-tilde-restriction-isometry} the map is a linear isometry. Its range is
complete, since $\tilde H^1(\Omega)$ is complete, hence closed, and its closure is
$H^1_0(\Omega)$ (\Cref{lem:sobolev-tilde-restriction-range}).
\end{proof}

McLean proves $\tilde H^s(\Omega) = H^s_0(\Omega)$ on Lipschitz domains for $s \geq 0$ with
$s - \frac12 \notin \Z$, in particular for $s = 1$, up to equivalence of norms
\cite[Theorem 3.33]{McLean}. For $s = 1$, \Cref{thm:sobolev-tilde-zero-boundary} holds for every
open set and with equal norms.

\begin{lemma}[$\tilde H^1(\Omega)$ lies in $H^1(\Omega)$]
\label{lem:sobolev-tilde-in-restriction}
\lean{Rd.tildeSpace_one_le_orthogonal, Rd.coe_tildeEquivZeroBoundary_symm_apply}
\leanok
\uses{thm:sobolev-tilde-restriction-isometry, thm:sobolev-tilde-zero-boundary}
For every open set $\Omega \subseteq V$, $\tilde H^1(\Omega) \subseteq N^1(\Omega)^\perp =
H^1(\Omega)$. For every $u \in H^1_0(\Omega)$, the element $\tilde u \in \tilde H^1(\Omega)$ with
$r_\Omega \tilde u = u$ is $u$ itself, as an element of $H^1(V)$.
\end{lemma}

\begin{proof}
\leanok
$r_\Omega$ is the orthogonal projection onto $H^1(\Omega)$, and $\|r_\Omega U\| = \|U\|$ for
$U \in \tilde H^1(\Omega)$ (\Cref{thm:sobolev-tilde-restriction-isometry}). By Pythagoras,
$\|U\|^2 = \|r_\Omega U\|^2 + \|U - r_\Omega U\|^2$, so $U = r_\Omega U \in H^1(\Omega)$ (Mathlib's
\texttt{Submodule.\allowbreak mem\_\allowbreak iff\_\allowbreak norm\_\allowbreak
starProjection}). Hence $u = r_\Omega \tilde u = \tilde u$.
\end{proof}

So $\tilde H^1(\Omega)$ and $H^1_0(\Omega)$ are the same subspace of $H^1(V)$, and the isometry of
\Cref{thm:sobolev-tilde-zero-boundary} is the identity on it.

\section{Duality of Sobolev spaces on open sets}
\label{sec:sobolev-duality}

Throughout this section $\Omega \subseteq V$ is open and $s \in \R$; the spaces $N^s(\Omega)$,
$H^s(\Omega)$, $T^s(\Omega)$, $\tilde H^s(\Omega)$ and $H^s_0(\Omega)$ are those of
\Cref{sec:sobolev-restriction-spaces}. For a Schwartz function $u$ and $t \in \R$,
$J_t u = \mathcal F^{-1}(\langle \cdot \rangle^t \widehat u)$ is the Bessel potential on Schwartz
functions (\Cref{def:bessel-schwartz}); on tempered distributions $J_t$ is Mathlib's Bessel
potential, and $J_s U \in L_2(V)$ for $U \in H^s(V)$. The section proves that
$\tilde H^s(\Omega)$ and $H^{-s}(\Omega)$ are dual to each other, as in Chandler-Wilde, Hewett
and Moiola \cite[Theorem 3.3]{ChandlerWildeHewettMoiola}, for every open set and every real $s$.
The tool is the isometry $J_{2s} : H^s(V) \to H^{-s}(V)$: it carries the orthogonal complement
of $\tilde H^s(\Omega)$ onto $N^{-s}(\Omega)$, hence $\tilde H^s(\Omega)$ onto $H^{-s}(\Omega)$,
and Riesz's representation theorem does the rest.

\begin{lemma}[$J_t$ keeps real functions real]
\label{lem:bessel-schwartz-real}
\lean{Rd.conj_besselSchwartz_apply, Rd.conj_testSchwartz}
\leanok
\uses{lem:bessel-schwartz-distribution, lem:bessel-potential-transpose, lem:fourier-conjugation,
  def:test-schwartz}
For every $t \in \R$, every Schwartz function $u$ with real values and every $x \in V$,
$J_t u(x)$ is real. For every
$\varphi \in \mathcal D(\Omega)$, the Schwartz function of $\varphi$ (\Cref{def:test-schwartz})
has real values.
\end{lemma}

\begin{proof}
\leanok
By \Cref{lem:fourier-conjugation},
$\overline{J_t u(x)} = \mathcal F(\langle \cdot \rangle^t \, \overline{\widehat u})(x) =
\mathcal F(\langle \cdot \rangle^t \, \mathcal F^{-1} \overline u)(x)$, which is $J_t u(x)$ for
real $u$ (\Cref{lem:bessel-potential-transpose}).
\end{proof}

\begin{lemma}[The Bessel potential of an $L_2$ function at a Schwartz function]
\label{lem:bessel-potential-l2-schwartz}
\lean{Rd.besselPotential_toLp_apply, Rd.inner_toLp_besselSchwartz}
\leanok
\uses{lem:bessel-schwartz-real, lem:bessel-potential-transpose}
For every $t \in \R$, every $g \in L_2(V)$ and every Schwartz function $u$,
$(J_t g)(u) = \int J_t u \cdot g$. For every $t \in \R$, every Schwartz function $u$ with real
values and every $g \in L_2(V)$, $\langle J_t u, g \rangle_{L_2} = (J_t g)(u)$.
\end{lemma}

\begin{proof}
\leanok
$(J_t g)(u) = g(J_t u) = \int J_t u \cdot g$ (\Cref{lem:bessel-potential-transpose}). The inner
product is $\int \overline{J_t u} \, g$, and
$J_t u$ is real for real $u$.
\end{proof}

\begin{lemma}[Schwartz functions in $H^t(V)$]
\label{lem:sobolev-schwartz-elements}
\lean{Rd.inner_toLp_toLp, Rd.toLp_schwartz, Rd.toLp_ofTest, Rd.inner_ofTest_toLp,
  Rd.inner_ofTest_left}
\leanok
\uses{lem:bessel-schwartz-distribution, lem:bessel-potential-l2-schwartz,
  def:sobolev-tilde-zero-boundary, lem:lp-distribution-determined}
For every $t \in \R$ and all $U, W \in H^t(V)$,
$\langle U, W \rangle = \langle J_t U, J_t W \rangle_{L_2}$. For every $t \in \R$ and every
Schwartz function $u$, the $L_2$ function of $u \in H^t(V)$ is $J_t u$. For every $t \in \R$,
every $\varphi \in \mathcal D(\Omega)$ and every $g \in L_2(V)$,
$\langle J_t \varphi, g \rangle_{L_2} = (J_t g)(\varphi)$; for every $U \in H^t(V)$,
$\langle \varphi, U \rangle = (J_t (J_t U))(\varphi)$.
\end{lemma}

\begin{proof}
\leanok
$U \mapsto J_t U$ is a linear isometry onto $L_2(V)$, and linear isometries preserve inner
products. The tempered distributions of $J_t u$ as an element of $L_2(V)$ and of $J_t$ applied to
$u$ agree (\Cref{lem:bessel-schwartz-distribution}), and $L_2$ functions are determined by their
distributions (\Cref{lem:lp-distribution-determined}). The last two statements follow with
\Cref{lem:bessel-potential-l2-schwartz}, since
test functions are real (\Cref{lem:bessel-schwartz-real}).
\end{proof}

\begin{definition}[The isometry $J_{2s}$]
\label{def:bessel-shift}
\lean{Rd.besselShift}
\leanok
For $s \in \R$ and $U \in H^s(V)$, $\Theta_s U$ is the element of $H^{-s}(V)$ whose $L_2$
function is $J_s U$. $\Theta_s : H^s(V) \to H^{-s}(V)$ is an isometric isomorphism, the
composition of the isometries $H^s(V) \to L_2(V)$ and $L_2(V) \to H^{-s}(V)$.
\end{definition}

\begin{lemma}[$\Theta_s$ is $J_{2s}$]
\label{lem:bessel-shift}
\lean{Rd.toLp_toLpLinearIsometry_symm_apply, Rd.toLp_besselShift,
  Rd.toDistr_besselShift_eq_besselPotential_toLp, Rd.toDistr_besselShift,
  Rd.besselShift_symm_apply, Rd.toLp_besselShift_symm, Rd.toDistr_besselShift_symm,
  Rd.inner_ofTest_eq}
\leanok
\uses{def:bessel-shift, lem:sobolev-schwartz-elements}
For every $s \in \R$ and every $U \in H^s(V)$, $\Theta_s U$ is the tempered distribution
$J_s(J_s U) = J_{2s} U$. For every $s \in \R$ and every $W \in H^{-s}(V)$,
$\Theta_s^{-1} W = \Theta_{-s} W$, the element of $H^s(V)$ with $L_2$ function $J_{-s} W$, and its
tempered distribution is $J_{-s}(J_{-s} W)$. For every $s \in \R$, every $U \in H^s(V)$ and every
$\varphi \in \mathcal D(\Omega)$, $\langle U, \varphi \rangle = \overline{(\Theta_s U)(\varphi)}$.
\end{lemma}

\begin{proof}
\leanok
The tempered distribution of an element $W$ of $H^{-s}(V)$ is $J_s$ of its $L_2$ function, and
$J_s J_s = J_{2s}$. The last statement is \Cref{lem:sobolev-schwartz-elements}, conjugated.
\end{proof}

\begin{definition}[The pairing of $H^{-s}(V)$ and $H^s(V)$]
\label{def:sobolev-pairing}
\lean{Rd.pairing, Rd.pairing_add_left, Rd.pairing_sub_left, Rd.pairing_smul_left,
  Rd.pairing_add_right, Rd.pairing_sub_right, Rd.pairing_smul_right, Rd.pairingCLM,
  Rd.pairingCLM_apply}
\leanok
For $s \in \R$, $W \in H^{-s}(V)$ and $U \in H^s(V)$,
$[W, U] = \langle J_{-s} W, J_s U \rangle_{L_2}$. It is conjugate-linear in $W$ and linear in $U$,
and for every $W \in H^{-s}(V)$, $U \mapsto [W, U]$ is a continuous linear functional on
$H^s(V)$.
\end{definition}

\begin{proposition}[Properties of the pairing]
\label{prop:sobolev-pairing}
\lean{Rd.norm_pairing_le, Rd.pairing_besselShift, Rd.pairing_ofTest_left,
  Rd.pairing_ofTest_right, Rd.pairing_one_toLp}
\leanok
\uses{def:sobolev-pairing, lem:bessel-shift, lem:sobolev-schwartz-elements,
  lem:sobolev-one-to-l2, lem:plancherel-l2}
For every $s \in \R$, all $W \in H^{-s}(V)$ and $U \in H^s(V)$,
$|[W, U]| \leq \|W\| \, \|U\|$. For every $s \in \R$ and all $U', U \in H^s(V)$,
$[\Theta_s U', U] = \langle U', U \rangle$. For every $s \in \R$, every
$\varphi \in \mathcal D(\Omega)$, every $U \in H^s(V)$ and every $W \in H^{-s}(V)$,
$[\varphi, U] = U(\varphi)$ and $[W, \varphi] = \overline{W(\varphi)}$. For every Schwartz
function $\psi$, regarded as an element of $H^{-1}(V)$, and every $U \in H^1(V)$, the pairing of
$H^{-1}(V)$ and $H^1(V)$ satisfies $[\psi, U] = \langle \psi, \iota U \rangle_{L_2}$.
\end{proposition}

\begin{proof}
\leanok
The bound is the Cauchy--Schwarz inequality in $L_2(V)$, and $[\Theta_s U', U] =
\langle J_s U', J_s U \rangle_{L_2} = \langle U', U \rangle$
(\Cref{lem:sobolev-schwartz-elements}). By \Cref{lem:sobolev-schwartz-elements},
$[\varphi, U] = (J_{-s}(J_s U))(\varphi) = U(\varphi)$, and the second identity follows in the
same way after conjugation. For the last identity, by Plancherel's theorem
(\Cref{lem:plancherel-l2}) both sides are integrals over the frequencies,
$\int \overline{\langle \xi \rangle^{-1} \widehat \psi(\xi)} \, \widehat{J_1 U}(\xi) \, d\xi$ and
$\int \overline{\widehat \psi(\xi)} \, \langle \xi \rangle^{-1} \widehat{J_1 U}(\xi) \, d\xi$
(\Cref{lem:sobolev-one-to-l2}), and the weight is real.
\end{proof}

At $s = 1$ the pairing is the $L_2$ inner product on Schwartz functions, and for every $s$ it
extends the action of a distribution on test functions. So $[W, U]$ is the duality pairing of
Chandler-Wilde, Hewett and Moiola, in the sesquilinear form of Mathlib's inner products.

\begin{proposition}[The orthogonal complement of $\tilde H^t(\Omega)$]
\label{prop:orthogonal-tilde}
\lean{Rd.mem_orthogonal_tildeSpace_iff_toLp, Rd.orthogonal_tildeSpace_eq_comap,
  Rd.map_tildeSpace_eq_orthogonal}
\leanok
\uses{lem:sobolev-schwartz-elements, def:sobolev-tilde-zero-boundary,
  lem:sobolev-tilde-zero-boundary-closed, def:sobolev-vanishing, def:sobolev-restriction-space}
For every $t \in \R$, every open $\Omega$ and every $U \in H^t(V)$, $U$ is orthogonal to
$\tilde H^t(\Omega)$ iff $(J_t(J_t U))(\varphi) = 0$ for every $\varphi \in \mathcal D(\Omega)$.
For all $t, t' \in \R$, every open $\Omega$ and every isometric isomorphism
$e : H^t(V) \to H^{t'}(V)$ such that $e U$ is the tempered distribution $J_t(J_t U)$ for every
$U \in H^t(V)$, $\tilde H^t(\Omega)^\perp = e^{-1}(N^{t'}(\Omega))$ and
$e(\tilde H^t(\Omega)) = N^{t'}(\Omega)^\perp = H^{t'}(\Omega)$.
\end{proposition}

\begin{proof}
\leanok
The orthogonal complement of a closure is that of the set, and orthogonality to $T^t(\Omega)$ is
orthogonality to its generators $\varphi$, with
$\langle \varphi, U \rangle = (J_t(J_t U))(\varphi)$ (\Cref{lem:sobolev-schwartz-elements}). This
gives the first two statements. $\tilde H^t(\Omega)$ is closed
(\Cref{lem:sobolev-tilde-zero-boundary-closed}), hence its own double orthogonal complement, and
an isometric isomorphism maps orthogonal complements onto orthogonal complements (Mathlib's
\texttt{Submodule.\allowbreak map\_\allowbreak orthogonal\_\allowbreak equiv}), so
$e(\tilde H^t(\Omega)) = e(\tilde H^t(\Omega)^\perp)^\perp = N^{t'}(\Omega)^\perp$.
\end{proof}

The hypothesis on $e$ determines $e$, and such an $e$ exists only for $t' = -t$, namely
$e = \Theta_t$. \Cref{prop:orthogonal-tilde} keeps $t'$ free so that it applies to both
$\Theta_s$ and $\Theta_s^{-1}$, whose target $H^s(V)$ is $H^{-(-s)}(V)$ only up to the identity
$-(-s) = s$.

\begin{theorem}[$J_{2s}$ on $\tilde H^s(\Omega)$ and on $H^s(\Omega)$]
\label{thm:tilde-restriction-shift}
\lean{Rd.mem_orthogonal_tildeSpace_iff, Rd.map_tildeSpace_besselShift,
  Rd.map_tildeSpace_besselShift_symm, Rd.tildeEquivRestriction,
  Rd.coe_tildeEquivRestriction_apply, Rd.restrictionEquivTilde,
  Rd.coe_restrictionEquivTilde_apply}
\leanok
\uses{prop:orthogonal-tilde, lem:bessel-shift, def:sobolev-restriction-space}
For every open $\Omega$, every $s \in \R$ and every $U \in H^s(V)$, $U$ is orthogonal to
$\tilde H^s(\Omega)$ iff $\Theta_s U \in N^{-s}(\Omega)$. For every open $\Omega$ and every
$s \in \R$, $\Theta_s$ restricts to isometric isomorphisms of $\tilde H^s(\Omega)$ onto
$H^{-s}(\Omega)$ and of $H^s(\Omega)$ onto $\tilde H^{-s}(\Omega)$.
\end{theorem}

\begin{proof}
\leanok
Apply \Cref{prop:orthogonal-tilde} to $e = \Theta_s$ with $t = s$, $t' = -s$, and to
$e = \Theta_s^{-1}$ with $t = -s$, $t' = s$ (\Cref{lem:bessel-shift}): $\Theta_s$ maps
$\tilde H^s(\Omega)$ onto $H^{-s}(\Omega)$, and $\Theta_s^{-1}$ maps $\tilde H^{-s}(\Omega)$ onto
$H^s(\Omega)$.
\end{proof}

\begin{lemma}[Elements vanishing on $\Omega$ pair to zero]
\label{lem:sobolev-pairing-vanishing}
\lean{Rd.pairing_eq_zero_of_mem_vanishOn_left, Rd.pairing_eq_zero_of_mem_vanishOn_right}
\leanok
\uses{thm:tilde-restriction-shift, prop:sobolev-pairing}
For every open $\Omega$, every $s \in \R$, every $W \in N^{-s}(\Omega)$ and every
$U \in \tilde H^s(\Omega)$, $[W, U] = 0$. For every open $\Omega$, every $s \in \R$, every
$W \in \tilde H^{-s}(\Omega)$ and every $U \in N^s(\Omega)$, $[W, U] = 0$.
\end{lemma}

\begin{proof}
\leanok
$[W, U] = \langle \Theta_s^{-1} W, U \rangle$ (\Cref{prop:sobolev-pairing}). In the first case
$\Theta_s^{-1} W$ is orthogonal to $\tilde H^s(\Omega)$, since $\Theta_s \Theta_s^{-1} W = W$ lies
in $N^{-s}(\Omega)$; in the second case $\Theta_s^{-1} W$ lies in $H^s(\Omega)$, the orthogonal
complement of $N^s(\Omega)$ (\Cref{thm:tilde-restriction-shift}).
\end{proof}

\begin{theorem}[The dual of $\tilde H^s(\Omega)$ is $H^{-s}(\Omega)$]
\label{thm:sobolev-duality}
\lean{Rd.dualEquiv, Rd.dualEquiv_apply, Rd.dualEquiv_symm_apply, Rd.dualEquiv_eq_iff,
  Rd.norm_dualEquiv, Rd.dualEquiv_apply_of_restrictProj_eq}
\leanok
\uses{thm:tilde-restriction-shift, prop:sobolev-pairing, lem:sobolev-pairing-vanishing}
For every open $\Omega$, every $s \in \R$ and every bounded linear functional $\ell$ on
$\tilde H^s(\Omega)$ there is exactly one $W \in H^{-s}(\Omega)$ with $\ell(U) = [W, U]$ for every
$U \in \tilde H^s(\Omega)$, and $\|W\| = \|\ell\|$. Then $\ell(U) = [W', U]$ for every
$U \in \tilde H^s(\Omega)$ and every $W' \in H^{-s}(V)$ with $r_\Omega W' = W$. The map
$\ell \mapsto W$ is a conjugate-linear isometric isomorphism of the dual of $\tilde H^s(\Omega)$
onto $H^{-s}(\Omega)$, and for every $W \in H^{-s}(\Omega)$, $U \mapsto [W, U]$ is a bounded
linear functional on $\tilde H^s(\Omega)$.
\end{theorem}

\begin{proof}
\leanok
By Riesz's representation theorem (Mathlib's \texttt{InnerProductSpace.\allowbreak toDual}),
$\ell(U) = \langle y, U \rangle$ for exactly one $y \in \tilde H^s(\Omega)$, with
$\|y\| = \|\ell\|$. Let $W = \Theta_s y \in H^{-s}(\Omega)$ (\Cref{thm:tilde-restriction-shift});
then $[W, U] = \langle y, U \rangle$ (\Cref{prop:sobolev-pairing}) and $\|W\| = \|y\|$. If
$r_\Omega W' = W$, then $W' - W \in N^{-s}(\Omega)$, and $[W' - W, U] = 0$
(\Cref{lem:sobolev-pairing-vanishing}).
\end{proof}

Chandler-Wilde, Hewett and Moiola define the pairing of $H^{-s}(\Omega)$ and
$\tilde H^s(\Omega)$ through an arbitrary extension of the element of $H^{-s}(\Omega)$ to $V$;
the second statement of \Cref{thm:sobolev-duality} says that every extension gives the same
value.

\begin{corollary}[The dual of $H^s(\Omega)$ is $\tilde H^{-s}(\Omega)$]
\label{cor:sobolev-duality-restriction}
\lean{Rd.dualRestrictionEquiv, Rd.dualRestrictionEquiv_apply, Rd.dualRestrictionEquiv_symm_apply,
  Rd.dualRestrictionEquiv_eq_iff, Rd.dualRestrictionEquiv_apply_restrictProj}
\leanok
\uses{thm:tilde-restriction-shift, prop:sobolev-pairing, lem:sobolev-pairing-vanishing}
For every open $\Omega$, every $s \in \R$ and every bounded linear functional $\ell$ on
$H^s(\Omega)$ there is exactly one $W \in \tilde H^{-s}(\Omega)$ with $\ell(u) = [W, u]$ for every
$u \in H^s(\Omega)$, and $\|W\| = \|\ell\|$. Then $\ell(r_\Omega U) = [W, U]$ for every
$U \in H^s(V)$. The map $\ell \mapsto W$ is a conjugate-linear isometric isomorphism of the dual
of $H^s(\Omega)$ onto $\tilde H^{-s}(\Omega)$.
\end{corollary}

\begin{proof}
\leanok
Riesz's representation theorem on $H^s(\Omega)$ followed by $\Theta_s : H^s(\Omega) \to
\tilde H^{-s}(\Omega)$ (\Cref{thm:tilde-restriction-shift}), with
$[\Theta_s y, u] = \langle y, u \rangle$ (\Cref{prop:sobolev-pairing}). For $U \in H^s(V)$,
$U - r_\Omega U \in N^s(\Omega)$, and $[W, U - r_\Omega U] = 0$
(\Cref{lem:sobolev-pairing-vanishing}).
\end{proof}

\begin{corollary}[The dual of $H^1_0(\Omega)$ is $H^{-1}(\Omega)$]
\label{cor:sobolev-duality-zero-boundary}
\lean{Rd.dualZeroBoundaryEquiv, Rd.dualZeroBoundaryEquiv_apply,
  Rd.dualZeroBoundaryEquiv_symm_apply, Rd.dualZeroBoundaryEquiv_eq_iff}
\leanok
\uses{thm:sobolev-tilde-zero-boundary, lem:sobolev-tilde-in-restriction,
  thm:tilde-restriction-shift, prop:sobolev-pairing}
For every open $\Omega$ and every bounded linear functional $\ell$ on $H^1_0(\Omega)$ there is
exactly one $W \in H^{-1}(\Omega)$ with $\ell(u) = [W, u]$ for every $u \in H^1_0(\Omega)$, and
$\|W\| = \|\ell\|$; here $u$ is read as an element of $H^1(V)$. The map $\ell \mapsto W$ is a
conjugate-linear isometric isomorphism of the dual of $H^1_0(\Omega)$ onto $H^{-1}(\Omega)$.
\end{corollary}

\begin{proof}
\leanok
Riesz's representation theorem on $H^1_0(\Omega)$, the isometry $H^1_0(\Omega) \to
\tilde H^1(\Omega)$ inverse to the restriction (\Cref{thm:sobolev-tilde-zero-boundary}), and
$\Theta_1 : \tilde H^1(\Omega) \to H^{-1}(\Omega)$ (\Cref{thm:tilde-restriction-shift}); the
isometry preserves inner products, and inside $H^1(V)$ it is the identity
(\Cref{lem:sobolev-tilde-in-restriction}).
\end{proof}

$\tilde H^1(\Omega)$ and $H^1_0(\Omega)$ are the same subspace of $H^1(V)$
(\Cref{lem:sobolev-tilde-in-restriction}), so in \Cref{cor:sobolev-duality-zero-boundary} $W$ is
paired with $u$ itself.

\section{The trace space of a bounded Lipschitz domain}
\label{sec:boundary-trace-space}

Throughout this section $\Omega$ is a bounded Lipschitz domain with boundary $\Gamma$, boundary
balls $B_j$, graph domains $\Omega_j$ over the hyperplanes $H_j$ and graph maps $\gamma_j$
(\Cref{def:bounded-lipschitz}), and $\operatorname{Tr}_j$ is the trace of the graph domain
$\Omega_j$ (\Cref{thm:trace-graph}). $H^1(\Omega)$ is the Hilbert space of
\Cref{def:sobolev-restriction-space} for the open set $\Omega$, inside $H^1(V)$, and $\iota$ reads
an element of $H^1(V)$ as an $L_2$ function (\Cref{def:sobolev-one-to-l2}). Boundary traces are
taken with differences of order $L = 2$, $p = q = 2$ and $s = 1/2$, so that $s + 1/p = 1$. The
trace space $H^{1/2}(\Gamma)$ is defined as the orthogonal complement in $H^1(\Omega)$ of the kernel
of the trace, so that by this definition it is a Hilbert space with a trace and an extension of
norm at most $1$; by the trace theorem of $H^1(\Omega)$ (\Cref{thm:sobolev-lipschitz-domain-trace})
its elements correspond to the functions of $B^{1/2}_{2,2}(\Gamma)$ modulo null sets of $\Gamma$.
This correspondence is qualitative: a norm equivalence of $H^{1/2}(\Gamma)$ with
$B^{1/2}_{2,2}(\Gamma)$ is not formalised.

The chart traces are bounded by the norm of $B^{1/2}_{2,1}(V)$ (\Cref{lem:trace-linear}). The
following three results bound this norm on $H^1(V)$, with explicit constants.

\begin{lemma}[The $\ell_2$ norm of $(2^{-k/2})_k$]
\label{lem:lq-cauchy-schwarz}
\lean{ENNReal.lqNorm_two_dyadicWeight_neg_half}
\leanok
\uses{lem:hardy, lem:lq-cauchy-schwarz-sequences}
$\|(2^{-k/2})_{k \in \N}\|_{\ell_2} = \sqrt 2$. In particular, by the Cauchy--Schwarz
inequality of \Cref{lem:lq-cauchy-schwarz-sequences}, for every family $(b_k)_{k \in \N}$ in
$[0, \infty]$,
\[
  \sum_{k \in \N} 2^{-k/2} b_k \leq \sqrt 2 \, \|(b_k)_k\|_{\ell_2} .
\]
\end{lemma}

\begin{proof}
\leanok
$\|(2^{-k/2})_k\|_{\ell_2}^2 = \sum_k 2^{-k} = 2$ (\Cref{lem:hardy}), and the inequality is
\Cref{lem:lq-cauchy-schwarz-sequences} with $a_k = 2^{-k/2}$.
\end{proof}

\begin{lemma}[$\Besov{1}{2}{2} \subseteq \Besov{1/2}{2}{1}$ with constant $\sqrt 2$]
\label{lem:besov-half-one}
\lean{Rd.modNorm_half_one_le}
\leanok
\uses{lem:lq-cauchy-schwarz-sequences, lem:lq-cauchy-schwarz}
For every $F \in L_2(V)$,
\[
  \|F\|_2 + \sum_k 2^{k/2} \omega_2(F, 2^{-k})_2
  \leq \sqrt 2 \bigl(\|F\|_2 + \|(2^k \omega_2(F, 2^{-k})_2)_k\|_{\ell_2}\bigr).
\]
\end{lemma}

\begin{proof}
\leanok
$2^{k/2} \omega_k = 2^{-k/2} (2^k \omega_k)$, and \Cref{lem:lq-cauchy-schwarz} gives
$\sum_k 2^{k/2} \omega_k \leq \sqrt 2 \, \|(2^k \omega_k)_k\|_{\ell_2}$; and $1 \leq \sqrt 2$.
\end{proof}

\begin{theorem}[$H^1(V)$ in Besov spaces, with explicit constants]
\label{thm:sobolev-one-besov-constant}
\lean{Rd.sobolevOneBesovConst, Rd.sobolevOneBesovConst_ne_top, Rd.modNorm_two_le_norm,
  Rd.modNorm_half_one_le_norm, Rd.modNorm_toL2_lt_top}
\leanok
\uses{thm:modulus-coincidence, thm:rd-besov, thm:tl-besov-sandwich,
  thm:sobolev-two-triebel-lizorkin, lem:sobolev-one-to-l2, lem:besov-half-one,
  lem:bessel-space-enorm}
Let $K$ be the constant of the inverse direction of
\Cref{thm:modulus-coincidence} for the band-limited scheme with the moduli of order $2$
(\Cref{thm:rd-besov}) at smoothness $1$, and
$C_1 = K \cdot 2 (1 - 2^{-1})^{-1} \cdot (2 \cdot 4)^{1/2}$. Then $C_1 < \infty$, and for every
$U \in H^1(V)$,
\[
  \|\iota U\|_2 + \|(2^k \omega_2(\iota U, 2^{-k})_2)_k\|_{\ell_2} \leq C_1 \|U\|
\]
and
\[
  \|\iota U\|_2 + \sum_k 2^{k/2} \omega_2(\iota U, 2^{-k})_2 \leq \sqrt 2 \, C_1 \|U\|.
\]
In particular
$\|\iota U\|_2 + \|(2^{k(1/2 + 1/2)} \omega_2(\iota U, 2^{-k})_2)_k\|_{\ell_2} < \infty$ for every
$U \in H^1(V)$.
\end{theorem}

\begin{proof}
\leanok
The left side of the first bound is at most $K$ times the approximation norm of $\iota U$ in
$A^1_2$ of the band-limited scheme (\Cref{thm:modulus-coincidence}), which is at most
$2 (1 - 2^{-1})^{-1} \|\iota U\|_{F^1_{2,2}}$ (\Cref{thm:tl-besov-sandwich}), and
$\|\iota U\|_{F^1_{2,2}} \leq (2 \cdot 4)^{1/2} \|J_1 \iota U\|_2$
(\Cref{thm:sobolev-two-triebel-lizorkin}), where $J_1 \iota U$ is the $L_2$ function of $U$
(\Cref{lem:sobolev-one-to-l2}), and $\|J_1 \iota U\|_2 = \|U\|$ (\Cref{lem:bessel-space-enorm}).
The second bound follows with \Cref{lem:besov-half-one}.
\end{proof}

The finiteness in \Cref{thm:sobolev-one-besov-constant} is the hypothesis of the boundary trace
theorems with $s = 1/2$ and $p = q = 2$.

\begin{lemma}[The trace of $H^1(V)$ on a Lipschitz graph is bounded]
\label{lem:graph-trace-sobolev-one}
\lean{Rd.LipschitzGraph.exists_enorm_trace_toL2_le}
\leanok
\uses{lem:trace-linear, thm:sobolev-one-besov-constant}
For every Lipschitz graph domain with hyperplane $H$ and trace $\operatorname{Tr}$
(\Cref{thm:trace-graph}), taken with $L = 2$ and $p = 2$, there is $C < \infty$ with
$\|\operatorname{Tr} \iota U\|_{L_2(H)} \leq C \|U\|$ for every $U \in H^1(V)$.
\end{lemma}

\begin{proof}
\leanok
$\|\operatorname{Tr} \iota U\|_{L_2(H)}$ is at most $K$ times the $B^{1/2}_{2,1}$ quantity of
$\iota U$ (\Cref{lem:trace-linear}), which is at most $\sqrt 2 \, C_1 \|U\|$
(\Cref{thm:sobolev-one-besov-constant}); so $C = K \sqrt 2 \, C_1$.
\end{proof}

\begin{lemma}[The trace of $H^1(V)$ on a Lipschitz graph as an operator]
\label{lem:graph-trace-sobolev-one-operator}
\lean{Rd.LipschitzGraph.intrinsicNorm_toL2_lt_top, Rd.LipschitzGraph.traceH1Linear,
  Rd.LipschitzGraph.traceH1, Rd.LipschitzGraph.traceH1_apply}
\leanok
\uses{thm:trace-graph, lem:graph-trace-sobolev-one, lem:trace-linear,
  thm:sobolev-one-besov-constant}
For every Lipschitz graph domain with hyperplane $H$, $U \mapsto \operatorname{Tr} \iota U$
(\Cref{thm:trace-graph}, $L = 2$, $p = 2$) is a bounded linear operator $H^1(V) \to L_2(H)$.
\end{lemma}

\begin{proof}
\leanok
The trace is linear on $H^1(V)$ (\Cref{lem:trace-linear}, \Cref{thm:sobolev-one-besov-constant})
and bounded (\Cref{lem:graph-trace-sobolev-one}).
\end{proof}

\begin{lemma}[Traces of Schwartz functions]
\label{lem:graph-trace-schwartz}
\lean{Rd.LipschitzGraph.traceH1_ofSchwartz_ae}
\leanok
\uses{lem:graph-trace-sobolev-one-operator, thm:trace-continuous, lem:sobolev-of-schwartz}
For every Lipschitz graph domain with graph map $\gamma$ and every Schwartz function $w$, the
trace of $w \in H^1(V)$ is $w \circ \gamma$ almost everywhere on $H$.
\end{lemma}

\begin{proof}
\leanok
$\iota w = w$ is continuous (\Cref{lem:sobolev-of-schwartz}), and \Cref{thm:trace-continuous}
applies.
\end{proof}

\begin{lemma}[Nonnegative resolutions of unity]
\label{lem:resolution-of-unity-nonneg}
\lean{Rd.BoundedLipschitzDomain.exists_boundaryPartition_nonneg}
\leanok
\uses{def:boundary-partition, lem:resolution-of-unity}
There is a resolution of unity $\psi$ on $\Gamma$ with $\psi_j \geq 0$ for every $j$ and
$\sum_j \psi_j \leq 1$ on $V$.
\end{lemma}

\begin{proof}
\leanok
\Cref{lem:resolution-of-unity} for the closed set $\Gamma$ and its open cover by the balls $B_j$.
\end{proof}

\begin{definition}[The setting of the trace space]
\label{def:trace-space-setting}
\lean{Rd.BoundedLipschitzDomain.SobolevOne, Rd.BoundedLipschitzDomain.boundaryPartition,
  Rd.BoundedLipschitzDomain.boundaryPartition_nonneg,
  Rd.BoundedLipschitzDomain.sum_boundaryPartition_le_one,
  Rd.BoundedLipschitzDomain.BoundaryPartition.derivBound,
  Rd.BoundedLipschitzDomain.BoundaryPartition.derivBound_spec,
  Rd.BoundedLipschitzDomain.BoundaryPartition.localize}
\leanok
\uses{def:sobolev-restriction-space, def:boundary-partition, lem:product-modulus,
  lem:resolution-of-unity, lem:resolution-of-unity-nonneg}
$\psi^0$ is a fixed resolution of unity on $\Gamma$ with $\psi^0_j \geq 0$ for every $j$ and
$\sum_j \psi^0_j \leq 1$ on $V$. For every resolution of unity $\psi$ and every
$j$, $F \mapsto \psi_j F$ is the multiplication on $L_2(V)$, with a bound of the derivatives of
$\psi_j$ of order at most $2$.
\end{definition}

\begin{theorem}[The chart traces are bounded]
\label{thm:chart-traces-bounded}
\lean{Rd.BoundedLipschitzDomain.chartTraceLinear,
  Rd.BoundedLipschitzDomain.BoundaryPartition.intrinsicNorm_localize_lt_top,
  Rd.BoundedLipschitzDomain.chartTraceLinear_apply,
  Rd.BoundedLipschitzDomain.exists_enorm_trace_localize_le, Rd.BoundedLipschitzDomain.chartTrace,
  Rd.BoundedLipschitzDomain.chartTrace_apply, Rd.BoundedLipschitzDomain.chartTrace_ae_eq}
\leanok
\uses{def:trace-space-setting, thm:sobolev-one-besov-constant, lem:trace-linear,
  lem:product-modulus, lem:trace-partial-sums, thm:multipliers, thm:boundary-trace-bounded}
For every resolution of unity $\psi$ on $\Gamma$, all $j$, $k$ and every $U \in H^1(V)$,
$\psi_j \iota U$ has a finite intrinsic norm of smoothness $1 = 1/2 + 1/2$ on $\Omega_k$, with
$p = q = 2$. For every resolution of unity $\psi$ on $\Gamma$ and every $j$ there is $C < \infty$
with
$\|\operatorname{Tr}_j(\psi_j \iota U)\|_{L_2(H_j)} \leq C \|U\|$ for every $U \in H^1(V)$. For
every $j$, $T_j u = \operatorname{Tr}_j(\psi^0_j \iota u)$ defines a bounded linear operator
$T_j : H^1(\Omega) \to L_2(H_j)$, and for every $u \in H^1(\Omega)$,
$T_j u = (\psi^0_j \operatorname{Tr}_\Gamma \iota u) \circ \gamma_j$ almost everywhere on $H_j$,
with the boundary trace of $\psi^0$.
\end{theorem}

\begin{proof}
\leanok
The functions $\psi_j \iota U$ have finite intrinsic norms (\Cref{thm:multipliers},
\Cref{thm:sobolev-one-besov-constant}, \Cref{lem:trace-partial-sums}).
$\|\operatorname{Tr}_j(\psi_j \iota U)\|_2$ is at most $K$ times the $B^{1/2}_{2,1}$ quantity of
$\psi_j \iota U$ (\Cref{lem:trace-linear}), which is at most $M$ times that of $\iota U$, with the
multiplier constant $M$ of $\psi_j$ (\Cref{thm:multipliers}), and that is at most
$\sqrt 2 \, C_1 \|U\|$ (\Cref{thm:sobolev-one-besov-constant}); so $C = K M \sqrt 2 \, C_1$. The
trace is additive on the functions of finite intrinsic norm and commutes with complex scalars
(\Cref{lem:trace-linear}), and so does the multiplication by $\psi^0_j$
(\Cref{lem:product-modulus}). The chart identity is \Cref{thm:boundary-trace-bounded}.
\end{proof}

\begin{lemma}[Vanishing on the boundary, chart by chart]
\label{lem:boundary-vanishing-charts}
\lean{Rd.BoundedLipschitzDomain.eventuallyEq_zero_iff_forall_chart}
\leanok
\uses{def:ae-boundary, def:boundary-partition, lem:boundary-charts}
For every resolution of unity $\psi$ on $\Gamma$ and every function $T$ on $V$, $T = 0$ almost
everywhere on $\Gamma$ if and only if, for every $j$, $(\psi_j T) \circ \gamma_j = 0$ almost
everywhere on $H_j$.
\end{lemma}

\begin{proof}
\leanok
If $T = 0$ almost everywhere on $\Gamma$, then $(\psi_j T)(\gamma_j(x')) = 0$ for almost every
$x'$ with $\gamma_j(x') \in B_j$, and $\psi_j = 0$ outside $B_j$. Conversely, $\psi_j T = 0$ almost
everywhere on $\Gamma \cap B_j$ for every $j$ (\Cref{def:ae-boundary}), every point
$\gamma_j(x')$ in $B_j$ lies on $\Gamma$ (\Cref{lem:boundary-charts}), and on $\Gamma$
$T = \sum_j \psi_j T$, where each term vanishes: inside $B_j$ by assumption, outside $B_j$ since
$\psi_j = 0$ there.
\end{proof}

\begin{definition}[The trace kernel]
\label{def:trace-kernel}
\lean{Rd.BoundedLipschitzDomain.traceKer, Rd.BoundedLipschitzDomain.isClosed_traceKer,
  Rd.BoundedLipschitzDomain.completeSpace_traceKer}
\leanok
\uses{thm:chart-traces-bounded}
The trace kernel $N_\Gamma$ is the set of the $u \in H^1(\Omega)$ with $T_j u = 0$ for every $j$.
It is a closed subspace of $H^1(\Omega)$ and complete.
\end{definition}

$N_\Gamma$ is closed as an intersection of kernels of bounded operators, and complete as a closed
subspace of a Hilbert space.

\begin{theorem}[The trace kernel by boundary traces]
\label{thm:trace-kernel-boundary}
\lean{Rd.BoundedLipschitzDomain.mem_traceKer_iff, Rd.BoundedLipschitzDomain.boundaryTrace_toL2_sub}
\leanok
\uses{def:trace-kernel, lem:boundary-vanishing-charts, thm:boundary-trace-charts,
  lem:trace-linear, thm:sobolev-one-besov-constant, lem:trace-partial-sums,
  thm:chart-traces-bounded}
For every resolution of unity $\psi$ on $\Gamma$ and every $u \in H^1(\Omega)$, $u \in N_\Gamma$ if
and only if $\operatorname{Tr}_\Gamma \iota u = 0$ almost everywhere on $\Gamma$, with the boundary
trace of $\psi$. For every resolution of unity $\psi$ and all $U, W \in H^1(V)$,
$\operatorname{Tr}_\Gamma \iota(U - W) = \operatorname{Tr}_\Gamma \iota U -
\operatorname{Tr}_\Gamma \iota W$ almost everywhere on $\Gamma$.
\end{theorem}

\begin{proof}
\leanok
By \Cref{thm:chart-traces-bounded}, $T_j u = 0$ if and only if
$(\psi^0_j \operatorname{Tr}_\Gamma \iota u) \circ \gamma_j = 0$ almost everywhere on $H_j$, so
$u \in N_\Gamma$ if and only if the boundary trace of $\psi^0$ vanishes almost everywhere on
$\Gamma$ (\Cref{lem:boundary-vanishing-charts}), and the boundary traces of $\psi^0$ and $\psi$
agree almost everywhere on $\Gamma$ (\Cref{thm:boundary-trace-charts}). For the difference, in the
chart $j$ the boundary traces are the traces $\operatorname{Tr}_j$
(\Cref{thm:boundary-trace-charts}), which are linear on $H^1(V)$ (\Cref{lem:trace-linear}), since
the functions of $H^1(V)$ have finite intrinsic norms (\Cref{thm:sobolev-one-besov-constant},
\Cref{lem:trace-partial-sums}).
\end{proof}

So the trace kernel does not depend on the resolution of unity in its definition.

\begin{definition}[The trace space $H^{1/2}(\Gamma)$]
\label{def:trace-space}
\lean{Rd.BoundedLipschitzDomain.TraceSpace, Rd.BoundedLipschitzDomain.trace,
  Rd.BoundedLipschitzDomain.extend}
\leanok
\uses{def:trace-kernel}
$H^{1/2}(\Gamma) = N_\Gamma^\perp$ is the orthogonal complement of the trace kernel in
$H^1(\Omega)$, a Hilbert space. The trace $\gamma : H^1(\Omega) \to H^{1/2}(\Gamma)$ is the
orthogonal projection onto it, and the extension $E : H^{1/2}(\Gamma) \to H^1(\Omega)$ is the
inclusion.
\end{definition}

\begin{theorem}[Trace and extension]
\label{thm:trace-space-operators}
\lean{Rd.BoundedLipschitzDomain.trace_extend, Rd.BoundedLipschitzDomain.norm_trace_apply_le,
  Rd.BoundedLipschitzDomain.norm_extend_apply, Rd.BoundedLipschitzDomain.trace_eq_zero_iff,
  Rd.BoundedLipschitzDomain.ker_trace, Rd.BoundedLipschitzDomain.sub_extend_trace_mem,
  Rd.BoundedLipschitzDomain.quotientEquivTraceSpace,
  Rd.BoundedLipschitzDomain.quotientEquivTraceSpace_mk}
\leanok
\uses{def:trace-space}
For every $t \in H^{1/2}(\Gamma)$, $\gamma E t = t$ and $\|E t\| = \|t\|$. For every
$u \in H^1(\Omega)$, $\|\gamma u\| \leq \|u\|$, $\gamma u = 0$ if and only if $u \in N_\Gamma$, and
$u - E \gamma u \in N_\Gamma$. The map $[u] \mapsto \gamma u$ is an isometric isomorphism of
$H^1(\Omega) / N_\Gamma$ onto $H^{1/2}(\Gamma)$.
\end{theorem}

\begin{proof}
\leanok
Properties of the orthogonal projection onto $K^\perp$ for the closed subspace $K = N_\Gamma$, with
$K^{\perp\perp} = K$ (Mathlib's \texttt{Submodule.\allowbreak
ker\_\allowbreak orthogonalProjectionOnto} for the kernel), and Mathlib's
\texttt{Submodule.\allowbreak quotientEquivOrthogonal}, as in
\Cref{thm:sobolev-restriction-quotient}.
\end{proof}

\begin{theorem}[The functions of $H^{1/2}(\Gamma)$]
\label{thm:trace-space-functions}
\lean{Rd.BoundedLipschitzDomain.trace_eq_iff, Rd.BoundedLipschitzDomain.boundaryTrace_extend_trace,
  Rd.BoundedLipschitzDomain.boundaryTrace_extend_eventuallyEq_iff,
  Rd.BoundedLipschitzDomain.memBoundaryBesov_toL2,
  Rd.BoundedLipschitzDomain.memBoundaryBesov_extend,
  Rd.BoundedLipschitzDomain.exists_traceSpace_of_memBoundaryBesov}
\leanok
\uses{thm:trace-kernel-boundary, thm:trace-space-operators, thm:boundary-trace-bounded,
  thm:sobolev-one-besov-constant, thm:sobolev-lipschitz-domain-trace,
  thm:sobolev-restriction-lipschitz, lem:sobolev-lipschitz-domain-trace-extensions,
  lem:sobolev-one-weak-gradient-norm, def:sobolev-restriction-l2}
Let $\psi$ be a resolution of unity on $\Gamma$, with $B^{1/2}_{2,2}(\Gamma)$ and the boundary trace
taken with respect to $\psi$. For all $u, u' \in H^1(\Omega)$, $\gamma u = \gamma u'$ if and only
if $\operatorname{Tr}_\Gamma \iota u = \operatorname{Tr}_\Gamma \iota u'$ almost everywhere on
$\Gamma$, and $\operatorname{Tr}_\Gamma \iota E \gamma u = \operatorname{Tr}_\Gamma \iota u$ almost
everywhere on $\Gamma$. For all $t, t' \in H^{1/2}(\Gamma)$,
$\operatorname{Tr}_\Gamma \iota E t = \operatorname{Tr}_\Gamma \iota E t'$ almost everywhere on
$\Gamma$ if and only if $t = t'$. For every $U \in H^1(V)$,
$\operatorname{Tr}_\Gamma \iota U \in B^{1/2}_{2,2}(\Gamma)$; in particular
$\operatorname{Tr}_\Gamma \iota E t \in B^{1/2}_{2,2}(\Gamma)$ for every
$t \in H^{1/2}(\Gamma)$. For every $g \in B^{1/2}_{2,2}(\Gamma)$ there is
$t \in H^{1/2}(\Gamma)$ with $\operatorname{Tr}_\Gamma \iota E t = g$ almost everywhere on
$\Gamma$.
\end{theorem}

\begin{proof}
\leanok
$\gamma u = \gamma u'$ if and only if $u - u' \in N_\Gamma$ (\Cref{thm:trace-space-operators}), if
and only if $\operatorname{Tr}_\Gamma \iota (u - u') = 0$ almost everywhere on $\Gamma$, and this is
$\operatorname{Tr}_\Gamma \iota u - \operatorname{Tr}_\Gamma \iota u'$ almost everywhere on $\Gamma$
(\Cref{thm:trace-kernel-boundary}). The second and third statements follow with
$\gamma E t = t$. The boundary trace lies in $B^{1/2}_{2,2}(\Gamma)$ by
\Cref{thm:boundary-trace-bounded} and \Cref{thm:sobolev-one-besov-constant}. For
$g \in B^{1/2}_{2,2}(\Gamma)$ there is $f \in H^1(\Omega)$ with boundary trace $g$
(\Cref{thm:sobolev-lipschitz-domain-trace}), $f = \iota u$ almost everywhere on $\Omega$ for some
$u$ in the Hilbert space $H^1(\Omega)$ (\Cref{thm:sobolev-restriction-lipschitz},
\Cref{def:sobolev-restriction-l2}), $\iota u \in \Sob{1}{2}(V)$
(\Cref{lem:sobolev-one-weak-gradient-norm}), so
$\operatorname{Tr}_\Gamma \iota u = g$ almost everywhere on $\Gamma$
(\Cref{lem:sobolev-lipschitz-domain-trace-extensions}), and $t = \gamma u$.
\end{proof}

So $t \mapsto \operatorname{Tr}_\Gamma \iota E t$ is a bijection of $H^{1/2}(\Gamma)$ onto
$B^{1/2}_{2,2}(\Gamma)$ modulo null sets of $\Gamma$, without a comparison of the norms.

\begin{definition}[The space $H^{-1/2}(\Gamma)$]
\label{def:dual-trace-space}
\lean{Rd.BoundedLipschitzDomain.DualTraceSpace}
\leanok
\uses{def:trace-space}
$H^{-1/2}(\Gamma)$ is the space of the bounded linear functionals on $H^{1/2}(\Gamma)$.
\end{definition}

\begin{theorem}[Riesz for $H^{-1/2}(\Gamma)$ and the adjoint of the trace]
\label{thm:dual-trace-space}
\lean{Rd.BoundedLipschitzDomain.dualTraceEquiv, Rd.BoundedLipschitzDomain.dualTraceEquiv_apply,
  Rd.BoundedLipschitzDomain.eq_dualTraceEquiv_of_forall, Rd.BoundedLipschitzDomain.adjoint_trace}
\leanok
\uses{def:dual-trace-space, thm:trace-space-operators}
For every $\ell \in H^{-1/2}(\Gamma)$ there is exactly one $t_\ell \in H^{1/2}(\Gamma)$ with
$\ell(t) = \langle t_\ell, t \rangle$ for every $t \in H^{1/2}(\Gamma)$, and
$\ell \mapsto t_\ell$ is a conjugate-linear isometric isomorphism of $H^{-1/2}(\Gamma)$ onto
$H^{1/2}(\Gamma)$. The adjoint of $\gamma$ is $E$.
\end{theorem}

\begin{proof}
\leanok
Riesz's representation theorem (Mathlib's \texttt{InnerProductSpace.\allowbreak toDual});
uniqueness, since an element is determined by its inner products (Mathlib's
\texttt{ext\_\allowbreak inner\_\allowbreak right}). By
Mathlib's \texttt{Submodule.\allowbreak adjoint\_\allowbreak orthogonalProjectionOnto}, the adjoint
of the orthogonal projection onto a closed subspace is the inclusion.
\end{proof}

\section{Divergence kernels}
\label{sec:divergence-kernels}

This section writes the difference $m_2 - m$ of two dilates of a smooth function $m$ as a
divergence $\sum_i \partial_{e_i} K_i$ of smooth kernels that vanish where $m$ vanishes along the
rays from the origin, with mean zero when the first moments of $m$ vanish. Dilates of a
mean-zero Schwartz function form a Calder\'on--Zygmund family, so the divergence kernels carry
the strong type of \Cref{thm:vector-valued-strong-type}. Throughout, $(e_i)_{i \in I}$ is an
orthonormal basis of $V$ and $\Phi_c(x) = c^n \Phi(c x)$ (\Cref{lem:dilate}).

\paragraph{Parametric integrals.}

\begin{lemma}[Smoothness of parametric integrals]
\label{lem:parametric-integral-smooth}
\lean{intervalIntegral.hasFDerivAt_integral_of_contDiff,
  intervalIntegral.fderiv_integral_apply_of_contDiff,
  intervalIntegral.contDiff_parametric_intervalIntegral_of_contDiff_aux,
  intervalIntegral.contDiff_parametric_intervalIntegral_of_contDiff}
\leanok
Let $H$ be a finite-dimensional real normed space, $E$ a real normed space and $a, b \in \R$.
For every continuously differentiable $F \colon H \times \R \to E$ and all $x_0, v \in H$, the
function $x \mapsto \int_a^b F(x, t) \, dt$ is differentiable at $x_0$ with
\[
  D\Bigl(\int_a^b F(\cdot, t) \, dt\Bigr)(x_0) \, v = \int_a^b D_x F(x_0, t) \, v \, dt .
\]
For every smooth $F \colon H \times \R \to E$ the function $x \mapsto \int_a^b F(x, t) \, dt$
is smooth.
\end{lemma}

\begin{proof}
\leanok
The partial derivative $(x, t) \mapsto D_x F(x, t)$ is continuous, hence bounded on the compact
set $\bar B(x_0, 1) \times [a, b]$, and Mathlib's differentiation under the integral sign
with a dominating bound gives the derivative. For smooth $F$ the derivative is again a
parametric integral, of the smooth function
$D_x F$ with values in the bounded linear maps $H \to E$, so the parametric integral is
$k$ times continuously differentiable for every $k$, by induction on $k$. The induction runs
through the spaces $E$, $L(H, E)$, $L(H, L(H, E))$, \dots, which Lean needs in one universe; a
lift of $E$ to the universe of $H$ provides it.
\end{proof}

Mathlib has the derivative under the integral sign and the continuity of parametric integrals;
\Cref{lem:parametric-integral-smooth} adds their smoothness.

\paragraph{The dilation identity.}

\begin{lemma}[Exponential dilations]
\label{lem:exponential-dilation}
\lean{Rd.contDiff_dilate_exp, Rd.continuous_dilate_exp, Rd.fderiv_dilate_apply,
  Rd.hasDerivAt_dilate_exp}
\leanok
\uses{lem:dilate}
For every continuous $m \colon V \to \C$ the function $(x, s) \mapsto m_{e^s}(x)$ on
$V \times \R$ is continuous, and it is $k$ times continuously differentiable if $m$ is. For every
differentiable $m$, every $c \in \R$ and all $x, v \in V$,
$D(m_c)(x) \, v = c^n \, Dm(c x)(c v)$, and for every $x \in V$ and every $s \in \R$,
\[
  \frac{d}{ds} \, m_{e^s}(x) = n \, m_{e^s}(x) + D(m_{e^s})(x) \, x .
\]
\end{lemma}

\begin{proof}
\leanok
$m_{e^s}(x) = e^{ns} m(e^s x)$; the product and chain rules give
$\frac{d}{ds} m_{e^s}(x) = n e^{ns} m(e^s x) + e^{ns} Dm(e^s x)(e^s x)$, and the second term is
$D(m_{e^s})(x) \, x$.
\end{proof}

\begin{definition}[Divergence kernels]
\label{def:divergence-kernel}
\lean{Rd.dilationAverage, Rd.divergenceKernel}
\leanok
\uses{lem:exponential-dilation}
For $m \colon V \to \C$ and $w \in V$, the \emph{dilation average} and the \emph{divergence
kernel} in the direction $w$ are
\[
  A_m(x) = \int_0^{\log 2} m_{e^s}(x) \, ds , \qquad K^m_w(x) = \langle w, x \rangle \, A_m(x) .
\]
\end{definition}

The substitution $t = e^s$ gives the form $\int_1^2 t^{-1} m_t(x) \, dt$ of the dilation average.

\begin{lemma}[The dilation average after substitution]
\label{lem:dilation-average-substitution}
\lean{Rd.dilationAverage_eq_integral_inv_mul}
\leanok
\uses{def:divergence-kernel}
For every continuous $m \colon V \to \C$ and every $x \in V$,
$A_m(x) = \int_1^2 t^{-1} m_t(x) \, dt$.
\end{lemma}

\begin{proof}
\leanok
Substitute $t = e^s$, $dt = e^s \, ds$.
\end{proof}

\begin{lemma}[Smoothness of the divergence kernels]
\label{lem:divergence-kernel-smooth}
\lean{Rd.contDiff_dilationAverage, Rd.contDiff_divergenceKernel,
  Rd.differentiable_dilationAverage}
\leanok
\uses{def:divergence-kernel, lem:parametric-integral-smooth, lem:exponential-dilation}
For every smooth $m \colon V \to \C$ and every $w \in V$, $A_m$ and $K^m_w$ are smooth. For every
continuously differentiable $m$, $A_m$ is differentiable.
\end{lemma}

\begin{proof}
\leanok
$(x, s) \mapsto m_{e^s}(x)$ is as smooth as $m$ (\Cref{lem:exponential-dilation}), and
\Cref{lem:parametric-integral-smooth} applies.
\end{proof}

\begin{lemma}[Supports of the divergence kernels]
\label{lem:divergence-kernel-support}
\lean{Rd.dilationAverage_eq_zero, Rd.divergenceKernel_eq_zero,
  Rd.divergenceKernel_eq_zero_of_lt, Rd.hasCompactSupport_divergenceKernel}
\leanok
\uses{def:divergence-kernel}
For every $m \colon V \to \C$, every $w \in V$ and every $x \in V$ with $m(t x) = 0$ for every
$t \in [1, 2]$, $A_m(x) = 0$ and $K^m_w(x) = 0$. For every $R \in \R$ and every $m$ vanishing
outside the closed ball $\bar B(0, R)$, every $K^m_w$ vanishes outside $\bar B(0, R)$ and has
compact support.
\end{lemma}

\begin{proof}
\leanok
For $s \in [0, \log 2]$, $e^s \in [1, 2]$, so the integrand of $A_m(x)$ vanishes. For
$\|x\| > R$ and $t \geq 1$, $\|t x\| \geq \|x\| > R$.
\end{proof}

\begin{definition}[Divergence kernels as Schwartz functions]
\label{def:divergence-schwartz}
\lean{Rd.divergenceSchwartz, Rd.coe_divergenceSchwartz}
\leanok
\uses{lem:divergence-kernel-smooth, lem:divergence-kernel-support}
For every smooth $m \colon V \to \C$ that vanishes outside a ball and every $w \in V$, the
divergence kernel $K^m_w$, smooth with compact support, is a Schwartz function.
\end{definition}

\begin{theorem}[The dilation identity]
\label{thm:dilation-divergence}
\lean{Rd.fderiv_dilationAverage_apply, Rd.continuous_fderiv_dilate_exp_apply,
  Rd.dilate_two_sub_eq, Rd.sum_fderiv_divergenceKernel, Rd.sum_lineDerivOp_divergenceSchwartz,
  Rd.sum_lineDerivOp_schwartzDilate_divergenceSchwartz}
\leanok
\uses{def:divergence-kernel, lem:exponential-dilation, lem:parametric-integral-smooth,
  lem:divergence-kernel-smooth, def:divergence-schwartz, lem:derivative-dilate}
For every continuously differentiable $m \colon V \to \C$ and every $x \in V$,
\[
  m_2(x) - m(x) = n \, A_m(x) + D A_m(x) \, x = \sum_{i \in I} \partial_{e_i} K^m_{e_i}(x) .
\]
If $m$ is smooth and vanishes outside a ball, then for every $c \neq 0$ and every $x \in V$,
\[
  \sum_{i \in I} \partial_{e_i} \bigl(K^m_{e_i}\bigr)_c (x) = c \, (m_2 - m)_c(x) .
\]
\end{theorem}

\begin{proof}
\leanok
By \Cref{lem:parametric-integral-smooth},
$D A_m(x) \, x = \int_0^{\log 2} D(m_{e^s})(x) \, x \, ds$,
so by \Cref{lem:exponential-dilation} $n A_m(x) + D A_m(x) \, x$ is the integral of the
derivative of $s \mapsto m_{e^s}(x)$ over $[0, \log 2]$, which is
$m_2(x) - m_1(x) = m_2(x) - m(x)$ by the fundamental theorem of calculus. The product rule gives
$\partial_{e_i} K^m_{e_i}(x) = \langle e_i, e_i \rangle A_m(x) +
\langle e_i, x \rangle \, D A_m(x) \, e_i$, and the sum over $i$ is
$n A_m(x) + D A_m(x) \bigl(\sum_i \langle e_i, x \rangle e_i\bigr) = n A_m(x) + D A_m(x) \, x$.
The second identity follows with $\partial_w (G_c) = c \, (\partial_w G)_c$
(\Cref{lem:derivative-dilate}).
\end{proof}

\paragraph{The mean.}

\begin{lemma}[The divergence kernels have mean zero]
\label{lem:divergence-kernel-mean}
\lean{Rd.integrable_inner_mul_dilate_exp, Rd.integral_divergenceKernel}
\leanok
\uses{def:divergence-kernel, lem:first-moments-dilate}
For every continuous $m \colon V \to \C$ that vanishes outside a ball and every $w \in V$ with
$\int_V \langle w, x \rangle \, m(x) \, dx = 0$, $\int_V K^m_w = 0$.
\end{lemma}

\begin{proof}
\leanok
The function $(x, s) \mapsto \langle w, x \rangle m_{e^s}(x)$ is continuous and vanishes for
$s \in [0, \log 2]$ outside the compact set $\bar B(0, R) \times [0, \log 2]$, so it is
integrable on
$V \times (0, \log 2]$, and Fubini's theorem gives
$\int K^m_w = \int_0^{\log 2} \int_V \langle w, x \rangle m_{e^s}(x) \, dx \, ds
= \int_0^{\log 2} e^{-s} \cdot 0 \, ds = 0$ by \Cref{lem:first-moments-dilate}.
\end{proof}

\paragraph{Mean-zero kernels as Calder\'on--Zygmund families.}

\begin{lemma}[Fourier transforms of mean-zero kernels]
\label{lem:mean-zero-fourier}
\lean{Rd.norm_fourier_le_of_integral_eq_zero, Rd.ofReal_mul_enorm_mul_enorm_fourier_le,
  Rd.enorm_fourier_le_mul_min}
\leanok
\uses{lem:fourier-character-lipschitz, lem:fourier-integrable, def:gradient-mass,
  lem:gradient-mass}
For every Schwartz function $G$ and every $\xi \in V$,
$2 \pi \|\xi\| \, |\widehat G(\xi)| \leq C_G$ with the gradient mass $C_G$ of
\Cref{def:gradient-mass}. If moreover $\int G = 0$, then
$|\widehat G(\xi)| \leq 2 \pi \|\xi\| \int_V \|x\| \, |G(x)| \, dx$ and
\[
  |\widehat G(\xi)| \leq C \min\bigl(\|\xi\|, \|\xi\|^{-1}\bigr), \qquad
  C = 2 \pi \int_V \|x\| \, |G(x)| \, dx + \frac{C_G}{2 \pi} .
\]
\end{lemma}

\begin{proof}
\leanok
$\widehat{\partial_\xi G}(\xi) = 2 \pi i \|\xi\|^2 \widehat G(\xi)$ (Mathlib's
\texttt{SchwartzMap.\allowbreak fourier\_\allowbreak lineDerivOp\_\allowbreak eq}), and
$|\widehat{\partial_\xi G}(\xi)| \leq \|\partial_\xi G\|_1 \leq \|\xi\| C_G$
(\Cref{lem:fourier-integrable}, \Cref{lem:gradient-mass}). If $\int G = 0$,
$\widehat G(\xi) = \int (e^{-2 \pi i \langle x, \xi \rangle} - 1) \, G(x) \, dx$, and
$|e^{-2 \pi i a} - 1| \leq 2 \pi |a|$ (\Cref{lem:fourier-character-lipschitz}) with
$|\langle x, \xi \rangle| \leq \|x\| \|\xi\|$.
\end{proof}

\begin{theorem}[Dilates of mean-zero kernels]
\label{thm:mean-zero-cz-family}
\lean{Rd.meanZeroOverlap, Rd.meanZeroOverlap_ne_top, Rd.tsum_enorm_fourier_dilate_sq_le,
  Rd.isCZFamily_dilate_of_integral_eq_zero, Rd.isCZFamily_divergenceSchwartz}
\leanok
\uses{lem:mean-zero-fourier, lem:dyadic-min-sum, lem:fourier-dilate, thm:cz-family-dilation,
  lem:uniform-schwartz-constants, def:cz-family, def:hormander-const,
  lem:divergence-kernel-mean, def:divergence-schwartz}
For every Schwartz function $G$ with $\int G = 0$ and every $\xi \in V$,
\[
  \sum_{j \in \N} \bigl|\widehat{G_{2^j}}(\xi)\bigr|^2 \leq 4 C^2
\]
with the constant $C$ of \Cref{lem:mean-zero-fourier}, and the dilates $(G_{2^j})_{j \in \N}$
form a Calder\'on--Zygmund family with Fourier overlap $4 C^2$ and H\"ormander constant $4 H_A$
for the seminorms $A_{k,l}$ of $G$ (\Cref{def:hormander-const}). In particular, for every smooth
$m$ that vanishes outside a ball and every $w \in V$ with
$\int_V \langle w, x \rangle \, m(x) \, dx = 0$, the dilates
$\bigl((K^m_w)_{2^j}\bigr)_{j \in \N}$ form a Calder\'on--Zygmund family.
\end{theorem}

\begin{proof}
\leanok
$\widehat{G_{2^j}}(\xi) = \widehat G(2^{-j} \xi)$ (\Cref{lem:fourier-dilate}). For $\xi = 0$ every
term vanishes. Otherwise, with $a = \|\xi\|^{-1}$ and $t_j = 2^j a$, \Cref{lem:mean-zero-fourier}
gives $|\widehat G(2^{-j} \xi)|^2 \leq C^2 \min(t_j, t_j^{-1})^2 \leq C^2 \min(t_j, t_j^{-1})$,
and $\sum_j \min(t_j, t_j^{-1}) \leq 4$ (\Cref{lem:dyadic-min-sum}). The constant family $G$ is
uniformly Schwartz (\Cref{lem:uniform-schwartz-constants}), and \Cref{thm:cz-family-dilation}
applies. The divergence kernels have mean zero by \Cref{lem:divergence-kernel-mean}.
\end{proof}

\section{The extension theorem for Sobolev spaces}
\label{sec:sobolev-extension}

Let $\Omega$ be a bounded Lipschitz domain (\Cref{def:bounded-lipschitz}) in $V$ of dimension
$n \geq 1$, let $(e_i)_{i \in I}$ be the orthonormal basis of the dyadic grid of $V$ and
$1 < p < \infty$. The extension theorem of Calder\'on and Stein \cite[Chapter VI]{SteinSingular}
says that every $f \in \Sob{1}{p}(\Omega)$ is the restriction of some $F \in \Sob{1}{p}(V)$ with
$\|F\|_{\Sob{1}{p}(V)} \leq C \|f\|_{\Sob{1}{p}(\Omega)}$ (\Cref{thm:sobolev-extension}). The
proof is Rychkov's local extension \cite{Rychkov} of \Cref{thm:local-extension}, here with the
local means of the gradient: the local means kernels of order two are divergences of dilated
divergence kernels (\Cref{lem:bump-divergence-identity}), the derivatives move onto $f$, and the
strong type of the Calder\'on--Zygmund families of the divergence kernels
(\Cref{thm:vector-valued-strong-type}) bounds the local means of Triebel--Lizorkin type of $f$ on
a patch by $\|f\|_{\Sob{1}{p}(\Omega)}$. The local extensions lie in
$\TL{1}{p}{2} = \Sob{1}{p}(V)$ and are pasted in $\Sob{1}{p}(V)$.

\paragraph{Local means by the gradient.}

\begin{definition}[The divergence kernels of a bump]
\label{def:bump-divergence-kernels}
\lean{Rd.LocalMeansBump.momentKernel_two_eq_zero, Rd.LocalMeansBump.divKernel,
  Rd.LocalMeansBump.hasCompactSupport_divKernel, Rd.LocalMeansBump.divKernel_eq_zero}
\leanok
\uses{def:local-means, def:moment-kernel, def:divergence-schwartz,
  lem:divergence-kernel-support}
For a bump $g$ of radius $R$ (\Cref{def:local-means}) with the moment kernel $m$ of order two,
the \emph{divergence kernels} of $g$ are the Schwartz functions $K_i = K^m_{e_i}$, $i \in I$. The
moment kernel vanishes outside $\bar B(0, 2R)$, the $K_i$ have compact support, and
$K_i(y) \neq 0$ only if $m(t y) \neq 0$ for some $t \in [1, 2]$.
\end{definition}

\begin{lemma}[The local means kernels as divergences]
\label{lem:bump-divergence-identity}
\lean{Rd.LocalMeansBump.kernel_two_succ, Rd.LocalMeansBump.kernel_two_succ_apply,
  Rd.LocalMeansBump.integral_inner_mul_momentKernel, Rd.LocalMeansBump.isCZFamily_divKernel}
\leanok
\uses{def:bump-divergence-kernels, thm:dilation-divergence, lem:moment-kernel-first-moment,
  thm:mean-zero-cz-family}
Let $g$ be a bump with the moment kernel $m$ of order two, the local means kernels $k_j$ of order
two and the divergence kernels $K_i$. For every $j \in \N$ and every $y \in V$,
\[
  k_{j+1} = (m_2 - m)_{2^j} , \qquad
  k_{j+1}(y) = 2^{-j} \sum_{i \in I} \partial_{e_i} (K_i)_{2^j}(y) .
\]
For every $w \in V$, $\int_V \langle w, x \rangle \, m(x) \, dx = 0$, and for every $i$,
$((K_i)_{2^j})_{j \in \N}$ is a Calder\'on--Zygmund family with the constants $B_i$, $H_i$ of
\Cref{thm:mean-zero-cz-family}.
\end{lemma}

\begin{proof}
\leanok
$k_{j+1} = m_{2^{j+1}} - m_{2^j} = (m_2 - m)_{2^j}$, and \Cref{thm:dilation-divergence} with
$c = 2^j$ writes it as a divergence. The first moments vanish by
\Cref{lem:moment-kernel-first-moment}, and \Cref{thm:mean-zero-cz-family} gives the
Calder\'on--Zygmund families.
\end{proof}

\begin{lemma}[Local means by the gradient at a point]
\label{lem:local-means-gradient-point}
\lean{Rd.LocalMeansBump.conv_kernel_two_succ_eq,
  Rd.LocalMeansBump.dyadicWeight_mul_enorm_conv_kernel_le}
\leanok
\uses{def:bump-divergence-kernels, lem:bump-divergence-identity,
  lem:weak-derivative-convolution-local, lem:weak-derivative-basic, lem:dilate}
Let $g$ be a bump with the divergence kernels $K_i$, let $T \subseteq V$ be open,
$1 \leq p \leq \infty$, $f \in L_p(T)$ with weak derivatives $g_i = \partial_{e_i} f$ on $T$,
$j \in \N$ and $x \in V$ with $x - y \in T$ for every $i$ and every $y$ in the support of
$(K_i)_{2^j}$. Then
\[
  \bigl(k_{j+1} * \mathbf 1_T f\bigr)(x) = 2^{-j} \sum_{i \in I}
  \bigl((K_i)_{2^j} * \mathbf 1_T g_i\bigr)(x) ,
\]
and therefore
\[
  2^{j+1} \bigl|\bigl(k_{j+1} * \mathbf 1_T f\bigr)(x)\bigr| \leq 2 \sum_{i \in I}
  \bigl|\bigl((K_i)_{2^j} * \mathbf 1_T g_i\bigr)(x)\bigr| .
\]
\end{lemma}

\begin{proof}
\leanok
By \Cref{lem:bump-divergence-identity},
$(k_{j+1} * \mathbf 1_T f)(x) = 2^{-j} \sum_i ((\partial_{e_i} (K_i)_{2^j}) * \mathbf 1_T f)(x)$,
and $\mathbf 1_T g_i = \partial_{e_i} (\mathbf 1_T f)$ weakly on $T$
(\Cref{lem:weak-derivative-basic}), so \Cref{lem:weak-derivative-convolution-local} moves each
derivative to $\mathbf 1_T f$. The bound follows from $2^{j+1} 2^{-j} = 2$.
\end{proof}

\begin{definition}[Local means of Triebel--Lizorkin type on a set]
\label{def:tl-local-norm-on}
\lean{Rd.LocalMeansBump.tlLocalNormOn}
\leanok
\uses{def:local-means}
For a bump $g$, a measurable $W \subseteq V$, $L \in \N$, $s \in \R$, $q \in [0, \infty]$ and
$f \in L_p(V)$, $1 \leq p < \infty$, let
\[
  \|f\|^{W,L}_{s,q} = \Bigl\| x \mapsto \bigl\| \bigl(2^{js} \, |\mathbf 1_W (k_j * f)(x)|\bigr)_j
  \bigr\|_{\ell_q} \Bigr\|_p .
\]
\end{definition}

\begin{lemma}[The local extension in Triebel--Lizorkin spaces]
\label{lem:tl-local-extension}
\lean{Rd.LocalMeansBump.enorm_indicatorLp_localMean_le,
  Rd.LocalMeansBump.dyadicWeight_mul_enorm_indicatorLp_localMean_le,
  Rd.LocalMeansBump.enorm_localExtTerm_le, Rd.LocalMeansBump.summable_localExtTerm_of_tlLocalNormOn,
  Rd.LocalMeansBump.enorm_localExtension_le_of_tlLocalNormOn,
  Rd.LocalMeansBump.triebelLizorkinNorm_localExtension_le}
\leanok
\uses{def:tl-local-norm-on, thm:local-extension, lem:mixed-norm-member, thm:tl-synthesis-dual}
Let $g$ be a bump, $W \subseteq V$ measurable, $L \in \N$, $1 \leq p < \infty$, $s > 0$,
$q \neq 0$ and $f \in L_p(V)$ with $X = \|f\|^{W,L}_{s,q} < \infty$. Then
$2^{js} \|\mathbf 1_W (k_j * f)\|_p \leq X$ for every $j$, the terms of the local extension
satisfy $\|\eta_j * \mathbf 1_W (k_j * f)\|_p \leq 2 \|m\|_1 2^{-js} X$, the local extension
$F = \sum_j \eta_j * \mathbf 1_W (k_j * f)$ converges in $L_p$ with
$\|F\|_p \leq 2 \|m\|_1 (1 - 2^{-s})^{-1} X$, and for $1 < p < \infty$ and $1 < q \leq \infty$,
$\|F\|_{F^s_{p,q}} \leq C X$ with the constant $C$ of \Cref{thm:tl-local-means}.
\end{lemma}

\begin{proof}
\leanok
Each member of a sequence is at most its $\ell_q$ norm, so
$2^{js} \|\mathbf 1_W (k_j * f)\|_p \leq X$ (\Cref{lem:mixed-norm-member}); the bounds of the
terms and the convergence follow as in \Cref{thm:local-extension}, and
\Cref{thm:tl-synthesis-dual} with $h_j = \mathbf 1_W (k_j * f)$ bounds $\|F\|_{F^s_{p,q}}$.
\end{proof}

\begin{theorem}[The local means of Triebel--Lizorkin type by the gradient]
\label{thm:tl-local-means-gradient}
\lean{Rd.LocalMeansBump.gradientConst, Rd.LocalMeansBump.gradientConst_ne_top,
  Rd.LocalMeansBump.ae_lqNorm_le_of_hasWeakLineDeriv,
  Rd.LocalMeansBump.eLpNorm_lqNorm_conv_divKernel_le,
  Rd.LocalMeansBump.tlLocalNormOn_le_of_hasWeakLineDeriv, Rd.sum_mul_le_sum_mul_sum,
  Rd.LocalMeansBump.patchConst, Rd.LocalMeansBump.patchConst_ne_top,
  Rd.LocalMeansBump.tlLocalNormOn_le_patchConst}
\leanok
\uses{lem:local-means-gradient-point, def:tl-local-norm-on, lem:bump-divergence-identity,
  thm:vector-valued-strong-type, def:cz-const, lem:local-means-on-sets, lem:lq-basic,
  lem:countable-minkowski}
Let $n \geq 1$, $1 < p < \infty$, $g$ a bump with the moment kernel $m$ of order two and the
divergence kernels $K_i$, $W \subseteq V$ measurable, $T \subseteq V$ open, $f \in L_p(T)$ with
weak derivatives $g_i = \partial_{e_i} f \in L_p(T)$ on $T$, and assume $x - y \in T$ for every
$x \in W$, every $j \in \N$, every $i$ and every $y$ in the support of $(K_i)_{2^j}$. Then
$F = \mathbf 1_T f \in L_p(V)$ satisfies
\[
  \|F\|^{W,2}_{1,2} \leq \|m\|_1 \|f\|_{L_p(T)} + 2 \sum_{i \in I} S_{V,p}(B_i, H_i)
  \|g_i\|_{L_p(T)} ,
\]
hence
\[
  \|F\|^{W,2}_{1,2} \leq C_g \Bigl(\|f\|_{L_p(T)} + \sum_{i \in I} \|g_i\|_{L_p(T)}\Bigr) ,
\]
with the constant $C_g = \|m\|_1 + 2 \sum_i S_{V,p}(B_i, H_i)$, where $S_{V,p}$ is the strong
type constant of \Cref{def:cz-const}.
\end{theorem}

\begin{proof}
\leanok
Almost everywhere $\mathbf 1_W (k_j * F) = \mathbf 1_W (k_j * \mathbf 1_T f)$. Split off the term
$j = 0$ (\Cref{lem:lq-basic}); for $x \in W$, \Cref{lem:local-means-gradient-point} bounds the
other terms, and Minkowski's inequality in $\ell_2$ (\Cref{lem:countable-minkowski}) gives
\[
  \bigl\|(2^j |\mathbf 1_W (k_j * F)(x)|)_j\bigr\|_{\ell_2} \leq |\mathbf 1_W (k_0 * F)(x)|
  + 2 \sum_i \bigl\|\bigl(((K_i)_{2^j} * \mathbf 1_T g_i)(x)\bigr)_j\bigr\|_{\ell_2}
\]
almost everywhere. The first term has $L_p$ norm at most $\|m\|_1 \|F\|_p$
(\Cref{lem:local-means-on-sets}), and the analysis operators of the Calder\'on--Zygmund
families $((K_i)_{2^j})_j$ are bounded by $S_{V,p}(B_i, H_i)$
(\Cref{thm:vector-valued-strong-type}). Finally
$\sum_i a_i b_i \leq (\sum_i a_i)(\sum_i b_i)$ for $a_i, b_i \in [0, \infty]$.
\end{proof}

\paragraph{Extensions from patches.}

\begin{definition}[Bounded extensions]
\label{def:extends-sobolev}
\lean{Rd.ExtendsSobolev}
\leanok
\uses{def:weak-line-derivative}
Let $O \subseteq V$ be open, $U \subseteq V$, $1 \leq p \leq \infty$ and $C \in [0, \infty]$.
\emph{Sobolev functions on $O$ extend from $U$ with the constant $C$} if for every
$f \in L_p(O)$ with weak derivatives $g_i = \partial_{e_i} f \in L_p(O)$ on $O$ there are
$F \in \Sob{1}{p}(V)$ with $F = f$ almost everywhere on $U$ and weak derivatives
$G_i = \partial_{e_i} F \in L_p(V)$ on $V$ with
\[
  \|F\|_p + \sum_{i \in I} \|G_i\|_p \leq C \Bigl(\|f\|_{L_p(O)} + \sum_{i \in I}
  \|g_i\|_{L_p(O)}\Bigr) .
\]
\end{definition}

\begin{lemma}[Bounded extensions on larger sets]
\label{lem:extends-sobolev-mono}
\lean{Rd.ExtendsSobolev.of_le}
\leanok
\uses{def:extends-sobolev, lem:weak-derivative-basic}
For all open $T \subseteq O \subseteq V$, every $U \subseteq V$, every $1 \leq p \leq \infty$ and
every $C \in [0, \infty]$: if Sobolev functions on $T$ extend from $U$ with the constant $C$, then
Sobolev functions on $O$ extend from $U$ with the constant $C$.
\end{lemma}

\begin{proof}
\leanok
Weak derivatives on $O$ restrict to $T$ (\Cref{lem:weak-derivative-basic}), and
$\|h\|_{L_p(T)} \leq \|h\|_{L_p(O)}$.
\end{proof}

\begin{theorem}[The extension from one patch]
\label{thm:one-patch-extension}
\lean{Rd.LocalMeansBump.extensionConst, Rd.LocalMeansBump.extensionConst_ne_top,
  Rd.LocalMeansBump.extendsSobolev}
\leanok
\uses{thm:tl-local-means-gradient, lem:tl-local-extension, def:extends-sobolev,
  thm:triebel-lizorkin-to-sobolev, thm:local-extension, lem:support-principle-two}
Let $n \geq 1$, $1 < p < \infty$, $g$ a bump with the dual kernels $\eta_j$ of order two and the
divergence kernels $K_i$, let $U, W \subseteq V$ be measurable and $T \subseteq V$ open with
$U \subseteq T$, $x - y \in W$ for every $x \in U$, every $j$ and every $y$ with
$\eta_j(y) \neq 0$, and $x - y \in T$ for every $x \in W$, every $j$, every $i$ and every $y$ in
the support of $(K_i)_{2^j}$. Then Sobolev functions on $T$ extend from $U$ with the constant
$\bigl(2 \|m\|_1 (1 - 2^{-1})^{-1} + (\sum_i D_{p, e_i}) A\bigr) C_g$, where $A$ is the constant
$C$ of \Cref{thm:tl-local-means} for $L = 2$, $q = 2$, $s = 1$, $D_{p,e_i}$ the constants of
\Cref{thm:triebel-lizorkin-to-sobolev}, and $C_g$ that of \Cref{thm:tl-local-means-gradient}.
\end{theorem}

\begin{proof}
\leanok
Let $F_0 = \mathbf 1_T f \in L_p(V)$ and $F = \sum_j \eta_j * \mathbf 1_W (k_j * F_0)$. By
\Cref{thm:tl-local-means-gradient} and \Cref{lem:tl-local-extension}, $F$ converges, lies in
$\TL{1}{p}{2}$ and its $L_p$ and $F^1_{p,2}$ norms are bounded by the right side; $F$ lies in
$\Sob{1}{p}(V)$ with $\|\partial_{e_i} F\|_p \leq D_{p,e_i} \|F\|_{F^1_{p,2}}$
(\Cref{thm:triebel-lizorkin-to-sobolev}). On $U$ the terms are those of the reproducing formula
(\Cref{lem:support-principle-two}, as in \Cref{thm:local-extension}), so $F = F_0 = f$ almost
everywhere on $U$.
\end{proof}

\paragraph{Pasting in $\Sob{1}{p}(V)$.}

\begin{lemma}[The product rule in $L_p$]
\label{lem:product-rule-lp}
\lean{Rd.hasWeakLineDeriv_mulLp, Rd.MemSobolev.mulLp}
\leanok
\uses{lem:weak-derivative-product, lem:product-modulus, lem:weak-derivative-basic}
Let $1 \leq p \leq \infty$, $w \in V$, $F, G \in L_p(V)$ with $G = \partial_w F$ weakly on $V$,
and $\chi$ a smooth real function with $\chi$ and $\partial_w \chi$ bounded. Then
$\chi G + (\partial_w \chi) F = \partial_w (\chi F)$ weakly on $V$, with the products in
$L_p(V)$. For every $F \in \Sob{1}{p}(V)$ and every smooth real $\chi$ with compact support,
$\chi F \in \Sob{1}{p}(V)$.
\end{lemma}

\begin{proof}
\leanok
\Cref{lem:weak-derivative-product} for the representatives, and
\Cref{lem:weak-derivative-basic}. A smooth function with compact support is bounded with all its
derivatives.
\end{proof}

\begin{theorem}[Pasting bounded extensions]
\label{thm:pasting-sobolev}
\lean{Rd.pasting_bound, Rd.extendsSobolev_of_cover}
\leanok
\uses{def:extends-sobolev, lem:product-rule-lp, lem:resolution-of-unity, lem:product-modulus,
  lem:weak-derivative-basic}
Let $O \subseteq V$ be open, $1 \leq p \leq \infty$, and $(U_k)_{k \in \kappa}$ a finite family of
bounded open sets covering the closure of $O$. If for every $k$ Sobolev functions on $O$ extend
from $O \cap U_k$ with a finite constant $C_k$, then Sobolev functions on $O$ extend from $O$
with a finite constant.
\end{theorem}

\begin{proof}
\leanok
Let $(\psi_k)$ be a smooth resolution of unity on the closure of $O$ subordinate to $(U_k)$
(\Cref{lem:resolution-of-unity}), with $|\psi_k| \leq a_k$ and $|\partial_{e_i} \psi_k| \leq
d_{k,i}$. For $f$ and its extensions $F_k$ with weak derivatives $G_{k,i}$, the function
$F = \sum_k \psi_k F_k$ equals $\sum_k \psi_k f = f$ almost everywhere on $O$, since
$\psi_k = 0$ outside $U_k$; it lies in $\Sob{1}{p}(V)$ with the weak derivatives
$G_i = \sum_k (\psi_k G_{k,i} + (\partial_{e_i} \psi_k) F_k)$ (\Cref{lem:product-rule-lp}), and
\[
  \|F\|_p + \sum_i \|G_i\|_p \leq \sum_k \Bigl(a_k + \sum_i d_{k,i}\Bigr)
  \Bigl(\|F_k\|_p + \sum_i \|G_{k,i}\|_p\Bigr) ,
\]
which is at most $\sum_k (a_k + \sum_i d_{k,i}) C_k Y$ with
$Y = \|f\|_{L_p(O)} + \sum_i \|g_i\|_{L_p(O)}$.
\end{proof}

\paragraph{The extension theorem.}

\begin{lemma}[Supports of the dilated divergence kernels]
\label{lem:divergence-kernel-supports}
\lean{Rd.LocalMeansBump.exists_momentKernel_ne_zero_of_divKernel_ne_zero,
  Rd.LocalMeansBump.tsupport_dilate_divKernel_subset,
  Rd.LocalMeansBump.tsupport_dilate_divKernel_subset_closedBall,
  Rd.LipschitzGraph.tsupport_dilate_divKernel_subset, Rd.LipschitzGraph.patch_support_divKernel}
\leanok
\uses{def:bump-divergence-kernels, lem:graph-cone, def:cone-supported, lem:patches}
Let $g$ be a bump of radius $R$ with the divergence kernels $K_i$. For every $i$ and every
$j \in \N$ the support of $(K_i)_{2^j}$ lies in $\bar B(0, 2R)$. If $g$ is supported in $-K$ for a
Lipschitz graph domain $\Omega$ with the closed cone $\bar K$ (\Cref{lem:graph-cone}), the support
of $(K_i)_{2^j}$ lies in $\{y : -y \in \bar K, \|y\| \leq 2R\}$, and for every
$\varepsilon \geq 2R$,
every $\rho$ and every $c \in V$, $x - y \in \Omega \cap B(c, \rho)$ for every
$x \in \Omega \cap B(c, \rho - \varepsilon)$ and every $y$ in the support of $(K_i)_{2^j}$.
\end{lemma}

\begin{proof}
\leanok
$K_i(y) \neq 0$ only if $m(t y) \neq 0$ for some $t \in [1, 2]$, so $\|y\| \leq \|t y\| \leq 2R$,
and $-t y \in K$ if $g$ is supported in $-K$, so $-y \in \bar K$. These sets are closed and
contain $y$ whenever they contain $c y$ for some $c \geq 1$, so they contain the supports of the
dilates. Finally $\Omega + \bar K \subseteq \Omega$ (\Cref{lem:graph-cone}) and
$\|x - y - c\| < \rho - \varepsilon + 2R \leq \rho$.
\end{proof}

\begin{lemma}[The patches of a bounded Lipschitz domain]
\label{lem:sobolev-patches}
\lean{Rd.BoundedLipschitzDomain.exists_extendsSobolev_boundary,
  Rd.BoundedLipschitzDomain.exists_extendsSobolev_interior}
\leanok
\uses{thm:one-patch-extension, lem:divergence-kernel-supports, lem:patches,
  lem:bumps-every-radius, lem:cone-bump, def:bounded-lipschitz, lem:extends-sobolev-mono}
Let $\Omega$ be a bounded Lipschitz domain in dimension $n \geq 1$ with boundary balls
$B(c_j, \rho_j)$. For every $1 < p < \infty$, every $\varepsilon > 0$, every $\delta > 0$ and
every $j$, Sobolev functions on $\Omega$ extend with finite constants from
$\Omega \cap B(c_j, \rho_j - 2 \varepsilon)$ and from
$\{x \in \Omega : \dist(x, V \setminus \Omega) > 2 \delta\}$.
\end{lemma}

\begin{proof}
\leanok
For the boundary patch take the cone bump of the graph
domain $\Omega_j$ rescaled to a radius $R$ with $2R \leq \varepsilon$ (\Cref{lem:cone-bump},
\Cref{lem:bumps-every-radius}), $U = \Omega_j \cap B(c_j, \rho_j - 2\varepsilon)$,
$W = \Omega_j \cap B(c_j, \rho_j - \varepsilon)$ and $T = \Omega_j \cap B(c_j, \rho_j)
= \Omega \cap B(c_j, \rho_j)$; the dual kernels satisfy the support condition by
\Cref{lem:patches} and the divergence kernels by \Cref{lem:divergence-kernel-supports}. For the
inner patch take the centred bump of radius $\delta/2$ and $T = \Omega$. In both cases
\Cref{thm:one-patch-extension} applies, and \Cref{lem:extends-sobolev-mono} passes from $T$ to
$\Omega$.
\end{proof}

\begin{theorem}[The extension theorem]
\label{thm:sobolev-extension}
\lean{Rd.BoundedLipschitzDomain.exists_extendsSobolev,
  Rd.exists_memSobolev_of_memWeakSobolev_of_subsingleton,
  Rd.BoundedLipschitzDomain.exists_memSobolev_of_memWeakSobolev,
  Rd.BoundedLipschitzDomain.memWeakSobolev_iff_exists_memSobolev}
\leanok
\uses{lem:sobolev-patches, thm:pasting-sobolev, lem:bounded-lipschitz-geometry,
  def:extends-sobolev, lem:weak-derivative-basic}
For every bounded Lipschitz domain $\Omega$ in dimension $n \geq 1$ and every $1 < p < \infty$
there is $C < \infty$ such that for every $f \in L_p(\Omega)$ with weak derivatives
$g_i = \partial_{e_i} f \in L_p(\Omega)$ on $\Omega$ there are $F \in \Sob{1}{p}(V)$ with $F = f$
almost everywhere on $\Omega$ and weak derivatives $G_i = \partial_{e_i} F \in L_p(V)$ on $V$
with
\[
  \|F\|_p + \sum_{i \in I} \|G_i\|_p \leq C \Bigl(\|f\|_{L_p(\Omega)} + \sum_{i \in I}
  \|g_i\|_{L_p(\Omega)}\Bigr) .
\]
For every bounded Lipschitz domain $\Omega$ in $V$ of any dimension, every $1 < p < \infty$ and
every $f \colon V \to \C$, $f \in \Sob{1}{p}(\Omega)$ if and only if $f = F$ almost everywhere on
$\Omega$ for some $F \in \Sob{1}{p}(V)$.
\end{theorem}

\begin{proof}
\leanok
By \Cref{lem:bounded-lipschitz-geometry} the boundary is covered by the balls
$B(c_j, \rho_j - 2\varepsilon)$ for some $\varepsilon > 0$, and the points of $\Omega$ outside
them have distance more than $2\delta$ from $V \setminus \Omega$ for some $\delta > 0$. These sets
cover the closure of $\Omega$; \Cref{lem:sobolev-patches} and \Cref{thm:pasting-sobolev} give the
first claim. For the second, the weak derivatives of $F \in \Sob{1}{p}(V)$ restrict to $\Omega$
(\Cref{lem:weak-derivative-basic}); conversely the first claim gives $F$ for $n \geq 1$, and for
$n = 0$ the open set $\Omega$ is empty, where $F = 0$ serves, or all of $V$, where $F = f$
serves.
\end{proof}

The theorem is of first order and gives, for each $f$, an extension with this bound; a bounded
linear extension operator and the exponents $p \in \{1, \infty\}$ are not formalised.

\begin{corollary}[$H^1(\Omega) = \Sob{1}{2}(\Omega)$]
\label{cor:sobolev-extension-two}
\lean{Rd.BoundedLipschitzDomain.memSobolevOn_iff_memWeakSobolev,
  Rd.BoundedLipschitzDomain.memWeakSobolev_iff_exists_restrictL2}
\leanok
\uses{thm:sobolev-extension, def:sobolev-lipschitz-domain, thm:sobolev-restriction-lipschitz}
For every bounded Lipschitz domain $\Omega$ and every $f \in L_2(V)$, the following are
equivalent: $f$ lies in the restriction space $H^1(\Omega)$ of
\Cref{def:sobolev-lipschitz-domain}; $f \in \Sob{1}{2}(\Omega)$; the restriction of $f$ to
$\Omega$ is $u|_\Omega$ for some $u$ in the Hilbert space $H^1(\Omega)$.
\end{corollary}

\begin{proof}
\leanok
\Cref{thm:sobolev-extension} for $p = 2$ and \Cref{thm:sobolev-restriction-lipschitz}.
\end{proof}

\begin{corollary}[The boundary traces of $\Sob{1}{2}(\Omega)$]
\label{cor:weak-sobolev-trace}
\lean{Rd.BoundedLipschitzDomain.exists_isBoundaryTrace_of_memWeakSobolev,
  Rd.BoundedLipschitzDomain.isBoundaryTrace_unique_of_memWeakSobolev}
\leanok
\uses{cor:sobolev-extension-two, thm:sobolev-lipschitz-domain-trace}
For every bounded Lipschitz domain $\Omega$, every resolution of unity $\psi$ on $\Gamma$ and
every $f \in L_2(V)$ with $f \in \Sob{1}{2}(\Omega)$, there is a boundary trace
$g \in B^{1/2}_{2,2}(\Gamma)$ of $f$, and every two boundary traces of $f$ agree almost everywhere
on $\Gamma$.
\end{corollary}

\begin{proof}
\leanok
By \Cref{cor:sobolev-extension-two}, $f \in H^1(\Omega)$, and
\Cref{thm:sobolev-lipschitz-domain-trace} applies.
\end{proof}

\begin{corollary}[Rellich's theorem in $\Sob{1}{2}(\Omega)$]
\label{cor:rellich-weak-sobolev}
\lean{Rd.sq_norm_add_sum_sq_le, Rd.BoundedLipschitzDomain.exists_subseq_tendsto_of_memWeakSobolev}
\leanok
\uses{thm:sobolev-extension, cor:rellich-weak-gradient}
For every bounded Lipschitz domain $\Omega$ in dimension $n \geq 1$, every $M < \infty$ and all
$f_k \in L_2(\Omega)$ with weak derivatives $g_{k,i} = \partial_{e_i} f_k \in L_2(\Omega)$ on
$\Omega$ and $\|f_k\|_{L_2(\Omega)} + \sum_i \|g_{k,i}\|_{L_2(\Omega)} \leq M$ for every
$k \in \N$, some subsequence of $(f_k)$ converges in $L_2(\Omega)$.
\end{corollary}

\begin{proof}
\leanok
The extensions $F_k$ of \Cref{thm:sobolev-extension} satisfy
$\|F_k\|_2 + \sum_i \|G_{k,i}\|_2 \leq C M$, hence
$\|F_k\|_2^2 + \sum_i \|G_{k,i}\|_2^2 \leq (C M)^2$, and \Cref{cor:rellich-weak-gradient} gives
a subsequence whose restrictions to $\Omega$ converge in $L_2(\Omega)$; they are the $f_k$.
\end{proof}

\begin{corollary}[The quotient norm of $H^1(\Omega)$]
\label{cor:sobolev-quotient-norm}
\lean{Rd.BoundedLipschitzDomain.exists_norm_le_restrictL2}
\leanok
\uses{thm:sobolev-extension, lem:sobolev-one-weak-gradient-norm, prop:sobolev-restriction-l2,
  def:sobolev-restriction-space, thm:sobolev-one-two}
For every bounded Lipschitz domain $\Omega$ in dimension $n \geq 1$ there is $C$ such that for
every $u$ in the Hilbert space $H^1(\Omega)$ and all $g_i \in L_2(\Omega)$ that are weak
derivatives $\partial_{e_i}$ of the $L_2$ function of $u$ on $\Omega$,
\[
  \|u\|_{H^1(\Omega)} \leq C \Bigl(\|u|_\Omega\|_{L_2(\Omega)} + \sum_{i \in I}
  \|g_i\|_{L_2(\Omega)}\Bigr) .
\]
\end{corollary}

\begin{proof}
\leanok
The extension $F$ of $u|_\Omega$ (\Cref{thm:sobolev-extension}) is an element $U$ of $H^1(V)$
(\Cref{thm:sobolev-one-two}) whose restriction to $\Omega$ is $u$, since both have the same
restriction to $L_2(\Omega)$ and that map is injective (\Cref{prop:sobolev-restriction-l2}). The
norm of $u$ is at most $\|U\| \leq \|F\|_2 + \sum_i \|G_i\|_2$
(\Cref{lem:sobolev-one-weak-gradient-norm}).
\end{proof}

With \Cref{prop:sobolev-restriction-weak-gradient}, the quotient norm of $H^1(\Omega)$ and the
norm of $\Sob{1}{2}(\Omega)$ are equivalent on bounded Lipschitz domains.

\section{Gauss--Green on Lipschitz graph domains}
\label{sec:lipschitz-graph-green}

Throughout this section $\Omega = \{x : \varphi(x) < \langle x, e \rangle\}$ is a Lipschitz graph
domain (\Cref{def:lipschitz-graph}) in $V$, $\dim V = n$, with Lipschitz constant $M$, hyperplane
$H = e^\perp$ of dimension $n - 1$, projection $P$ and graph map $\gamma(x') = x' + \varphi(x') e$
(\Cref{def:graph-map}), with the gradient $\nabla\varphi$, the Gram density
$\sqrt{1 + \|\nabla\varphi\|^2}$ and the outward unit normal $\nu$ of
\Cref{sec:lipschitz-graph-area}; $\mathcal H^{n-1}$ is Mathlib's Euclidean Hausdorff measure. The
divergence theorem on $\Omega$ is proved
first for compactly supported $C^1$ functions (\Cref{thm:graph-gauss-green}), then for products
of $H^1(V)$ functions with their traces as boundary values (\Cref{thm:graph-green}), as in McLean
\cite[Chapter 3]{McLean}. The vertical direction $e$ is the fundamental theorem of calculus on the
vertical lines above the graph. For the horizontal directions the proof needs of $\varphi$ only
Rademacher's theorem: the integral $\Psi(y) = \int_0^\infty g(y + t e) \, dt$ is
differentiable with bounded derivative and $\partial_e \Psi = -g$, the function
$\Psi \circ \gamma$ on $H$ is Lipschitz with compact support, and the integral over $H$ of its
directional derivatives vanishes.

\paragraph{Product coordinates.}

\begin{lemma}[Product coordinates]
\label{lem:graph-product-coordinates}
\lean{Rd.LipschitzGraph.lineEquiv, Rd.LipschitzGraph.coe_lineEquiv, Rd.LipschitzGraph.prodEquiv,
  Rd.LipschitzGraph.prodEquiv_apply, Rd.LipschitzGraph.measurePreserving_prodEquiv,
  Rd.LipschitzGraph.graphShear, Rd.LipschitzGraph.measurePreserving_graphShear,
  Rd.LipschitzGraph.graphEquiv, Rd.LipschitzGraph.graphEquiv_apply,
  Rd.LipschitzGraph.measurePreserving_graphEquiv, Rd.LipschitzGraph.prodEquiv_mem_domain_iff,
  Rd.LipschitzGraph.graphMap_add_smul_mem_domain_iff}
\leanok
\uses{def:lipschitz-graph, def:graph-map}
The map $t \mapsto t e$ is a linear isometry of $\R$ onto $H^\perp = \R e$. The maps
$(x', t) \mapsto x' + t e$ and $(x', t) \mapsto \gamma(x') + t e$ are measurable
bijections of $H \times \R$ onto $V$ with measurable inverses that carry the product of the
Lebesgue measures of $H$ and $\R$ to the Lebesgue measure of $V$. For every $x' \in H$ and every
$t \in \R$, $x' + t e \in \Omega$ if and only if $\varphi(x') < t$, and
$\gamma(x') + t e \in \Omega$ if and only if $t > 0$.
\end{lemma}

\begin{proof}
\leanok
$V$ is the orthogonal sum of $H$ and $H^\perp = \R e$, and Mathlib's
\texttt{Submodule.\allowbreak measurePreserving\_\allowbreak measurableEquivProd} carries
$\mathcal H^n$, which is the Lebesgue measure of $V$, to the product of the Lebesgue measures of
$H$ and $H^\perp$; $t \mapsto t e$ is an isometry of $\R$ onto $H^\perp$ and preserves the
Lebesgue measures. The second map is the first after the shear
$(x', t) \mapsto (x', t + \varphi(x'))$, which translates each vertical line and so preserves the
product measure. Finally $\varphi(x' + t e) = \varphi(x')$, $\langle x' + t e, e \rangle = t$ and
$\gamma(x') + t e = x' + (\varphi(x') + t) e$.
\end{proof}

\begin{theorem}[Fubini on a Lipschitz graph domain]
\label{thm:graph-fubini}
\lean{Rd.LipschitzGraph.setIntegral_domain_eq_of_measurePreserving,
  Rd.LipschitzGraph.setIntegral_domain_eq, Rd.LipschitzGraph.setIntegral_domain_eq_graphMap}
\leanok
\uses{lem:graph-product-coordinates}
For every real Banach space $E$ and every $F \colon V \to E$ integrable on $\Omega$,
\[
  \int_\Omega F = \int_H \int_{\varphi(x')}^\infty F(x' + t e) \, dt \, dx'
  = \int_H \int_0^\infty F(\gamma(x') + t e) \, dt \, dx' .
\]
More generally, for every measure preserving measurable bijection $T$ of $H \times \R$ onto $V$
with measurable inverse and every $a \colon H \to \R$ with $T(x', t) \in \Omega$ if and only if
$a(x') < t$,
$\int_\Omega F = \int_H \int_{a(x')}^\infty F(T(x', t)) \, dt \, dx'$.
\end{theorem}

\begin{proof}
\leanok
For a measure preserving bijection $T$ of $H \times \R$ onto $V$ with $T(x', t) \in \Omega$ if
and only if $a(x') < t$, the change of variables and Fubini's theorem (Mathlib's
\texttt{MeasurePreserving.\allowbreak integral\_\allowbreak comp'} and
\texttt{integral\_\allowbreak prod}) applied to $\mathbf 1_\Omega F$ give
$\int_\Omega F = \int_H \int_{a(x')}^\infty F(T(x', t)) \, dt \, dx'$. The two maps of
\Cref{lem:graph-product-coordinates} have $a = \varphi$ and $a = 0$.
\end{proof}

\paragraph{Ray integrals.}

\begin{definition}[Ray integrals]
\label{def:ray-integral}
\lean{Rd.rayIntegral}
\leanok
For a function $g$ on $V$ with values in a real Banach space $E$, a vector $w$ and $y \in V$,
the ray integral is $\Psi(y) = \int_0^\infty g(y + t w) \, dt$.
\end{definition}

\begin{lemma}[Compactly supported functions along lines]
\label{lem:compact-support-lines}
\lean{Rd.abs_add_inner_le_of_norm_add_smul_le, Rd.apply_add_smul_eq_zero_of_notMem_Icc,
  Rd.hasCompactSupport_comp_add_smul,
  Rd.integrable_comp_add_smul, Rd.hasDerivAt_comp_add_smul}
\leanok
For every unit vector $w$, every $R \in \R$, every function $h$ on $V$ with closed support in
$\bar B(0, R)$ and every $y \in V$, $h(y + t w) = 0$ for every $t$ outside the interval
$[-\langle y, w \rangle - R, -\langle y, w \rangle + R]$; so $t \mapsto h(y + t w)$ has compact
support, and it is integrable on $\R$ when $h$ is continuous. For every $C^1$ function $g$ and all
$y, w \in V$, $t \mapsto g(y + t w)$ has the derivative $Dg(y + t w) w$ at every $t$.
\end{lemma}

\begin{proof}
\leanok
$t + \langle y, w \rangle = \langle y + t w, w \rangle$ has absolute value at most
$\|y + t w\|$, which is at most $R$ where $h(y + t w) \neq 0$. The derivative is the chain rule.
\end{proof}

\begin{lemma}[Derivatives of ray integrals]
\label{lem:ray-integral-derivative}
\lean{Rd.hasFDerivAt_rayIntegral, Rd.fderiv_rayIntegral_apply, Rd.fderiv_rayIntegral_self,
  Rd.exists_norm_fderiv_rayIntegral_le, Rd.exists_lipschitzWith_rayIntegral}
\leanok
\uses{def:ray-integral, lem:compact-support-lines}
For every $C^1$ function $g \colon V \to E$ with compact support and every unit vector $w$, the
ray integral $\Psi$ is differentiable at every $y \in V$ with
\[
  D\Psi(y) v = \int_0^\infty Dg(y + t w) v \, dt \quad \text{for every } v \in V ,
  \qquad D\Psi(y) w = -g(y) ,
\]
and there is $K < \infty$ with $\|D\Psi(y)\| \leq K$ for every $y$; in particular $\Psi$ is
Lipschitz.
\end{lemma}

\begin{proof}
\leanok
Let $\|Dg\| \leq C$ and the closed support of $g$ lie in $\bar B(0, R)$. For $y$ in the unit ball
about $y_0$, $Dg(y + t w) = 0$ unless $|t| \leq \|y_0\| + 1 + R$
(\Cref{lem:compact-support-lines}), so differentiation under the integral sign (Mathlib's
\texttt{hasFDerivAt\_\allowbreak integral\_\allowbreak of\_\allowbreak dominated\_\allowbreak
of\_\allowbreak fderiv\_\allowbreak le}) applies with the integrable bound $C$ times the
indicator of that interval. In the direction $w$ the integrand is the derivative of
$t \mapsto g(y + t w)$, which vanishes for large $t$, and the fundamental theorem of calculus on
$(0, \infty)$ gives $-g(y)$. Each line meets the support of $g$ in an interval of length at most
$2R$, so $\|D\Psi(y)\| \leq 2 R C$, and the mean value inequality gives the Lipschitz bound.
\end{proof}

\begin{lemma}[Line derivatives of Lipschitz functions integrate to zero]
\label{lem:integral-line-derivative-zero}
\lean{Rd.integral_lineDeriv_eq_zero, Rd.integrable_lineDeriv,
  Rd.integral_lineDeriv_eq_zero_complex}
\leanok
For every finite-dimensional real normed space $F$ with a Haar measure, every Lipschitz function
$f \colon F \to \R$ with compact support and every $v \in F$, $\int \partial_v f = 0$. For every
real Banach space $G$, every Lipschitz function $f \colon F \to G$ with compact support and every
$v \in F$, $\partial_v f$ is integrable. For every Lipschitz function $f \colon F \to \C$ with
compact support and every $v \in F$, $\int \partial_v f = 0$.
\end{lemma}

\begin{proof}
\leanok
Integration by parts for Lipschitz functions with the constant function $1$, whose line
derivatives vanish (Mathlib's \texttt{LipschitzWith.\allowbreak integral\_\allowbreak
lineDeriv\_\allowbreak mul\_\allowbreak eq}). $\partial_v f$ is measurable (Mathlib's
\texttt{aestronglyMeasurable\_\allowbreak lineDeriv}), bounded by $C \|v\|$ for the Lipschitz
constant $C$, and zero outside the closed support of $f$. For complex $f$, $f$ is
differentiable almost everywhere by Rademacher's theorem, so the real and imaginary parts of
$\partial_v f$ are almost everywhere the line derivatives of the Lipschitz functions
$\operatorname{Re} f$ and $\operatorname{Im} f$.
\end{proof}

\paragraph{The normal weight and the ray integral from the graph.}

\begin{definition}[The normal weight]
\label{def:normal-weight}
\lean{Rd.LipschitzGraph.normalWeight}
\leanok
\uses{def:graph-gradient}
For $v \in V$ and $x' \in H$, the normal weight is
$n_v(x') = \langle \nabla\varphi(x'), v \rangle - \langle v, e \rangle$.
\end{definition}

\begin{lemma}[Properties of the normal weight]
\label{lem:normal-weight}
\lean{Rd.LipschitzGraph.inner_normalAt_mul_gram, Rd.LipschitzGraph.abs_normalWeight_le,
  Rd.LipschitzGraph.measurable_normalWeight, Rd.LipschitzGraph.memLp_top_normalWeight}
\leanok
\uses{def:normal-weight, def:graph-normal, lem:graph-gradient}
For every $v \in V$ and every $x' \in H$,
$n_v(x') = \langle v, \nu(x') \rangle \sqrt{1 + \|\nabla\varphi(x')\|^2}$ and
$|n_v(x')| \leq (M + 1) \|v\|$; $n_v$ is measurable, hence in $L_\infty(H)$.
\end{lemma}

\begin{proof}
\leanok
$\nu = (\nabla\varphi - e)/\sqrt{1 + \|\nabla\varphi\|^2}$, $\|\nabla\varphi\| \leq M$
(\Cref{lem:graph-gradient}) and $\|e\| = 1$.
\end{proof}

\begin{definition}[The ray integral from the graph]
\label{def:graph-ray-integral}
\lean{Rd.LipschitzGraph.graphRayIntegral}
\leanok
\uses{def:ray-integral, def:graph-map}
For $g \colon V \to \C$ and $x' \in H$, the ray integral from the graph is the ray integral of
$g$ in the direction $e$ at $\gamma(x')$,
$\Psi(\gamma(x')) = \int_0^\infty g(\gamma(x') + t e) \, dt
= \int_{\varphi(x')}^\infty g(x' + t e) \, dt$.
\end{definition}

\begin{lemma}[Properties of the ray integral from the graph]
\label{lem:graph-ray-integral}
\lean{Rd.LipschitzGraph.norm_le_norm_graphMap_add_smul,
  Rd.LipschitzGraph.apply_graphMap_add_smul_eq_zero_of_norm_lt,
  Rd.LipschitzGraph.hasCompactSupport_graphRayIntegral,
  Rd.LipschitzGraph.exists_lipschitzWith_graphRayIntegral, Rd.LipschitzGraph.hasFDerivAt_graphMap,
  Rd.LipschitzGraph.lineDeriv_graphRayIntegral}
\leanok
\uses{def:graph-ray-integral, lem:ray-integral-derivative, lem:graph-gradient, lem:graph-over}
For every $x' \in H$ and every $t \in \R$, $\|x'\| \leq \|\gamma(x') + t e\|$; so if a function
$h$ has closed support in $\bar B(0, R)$ and $\|x'\| > R$, then $h(\gamma(x') + t e) = 0$ for
every $t$. For every $C^1$ function $g \colon V \to \C$ with compact support, $\Psi \circ \gamma$
has compact support and is Lipschitz, and at every $x' \in H$ where $\varphi|_H$ is
differentiable, $D\gamma(x') v' = v' + D\varphi(x') v' \, e$ and
\[
  \partial_{v'} (\Psi \circ \gamma)(x') = D\Psi(\gamma(x')) v'
    - \langle \nabla\varphi(x'), v' \rangle \, g(\gamma(x'))
  \quad \text{for every } v' \in H .
\]
\end{lemma}

\begin{proof}
\leanok
$x'$ is the projection of $\gamma(x') + t e = x' + (\varphi(x') + t) e$ onto $H$. The ray
integral is Lipschitz (\Cref{lem:ray-integral-derivative}) and so is $\gamma$. The derivative is
the chain rule with $D\Psi(\gamma(x')) e = -g(\gamma(x'))$
(\Cref{lem:ray-integral-derivative}) and
$D\varphi(x') v' = \langle \nabla\varphi(x'), v' \rangle$ (\Cref{lem:graph-gradient}).
\end{proof}

\begin{lemma}[Integrability on the hyperplane]
\label{lem:hyperplane-integrable}
\lean{Rd.integrable_of_norm_le_of_forall_norm_gt_eq_zero,
  Rd.LipschitzGraph.integrable_comp_graphMap,
  Rd.LipschitzGraph.integrable_fderiv_rayIntegral,
  Rd.LipschitzGraph.integrable_inner_gradPhi_mul}
\leanok
\uses{lem:graph-ray-integral, lem:ray-integral-derivative, lem:graph-gradient}
On every proper normed group $F$ with a measure that is finite on compact sets, every a.e.
strongly measurable $f$ with $\|f\| \leq C$ and $f(x) = 0$ for $\|x\| > R$ is integrable. For
every $C^1$ function $g \colon V \to \C$ with compact support and all $v \in V$, $v' \in H$, the
functions $x' \mapsto g(\gamma(x'))$, $x' \mapsto D\Psi(\gamma(x')) v$ and
$x' \mapsto \langle \nabla\varphi(x'), v' \rangle g(\gamma(x'))$ are integrable on $H$.
\end{lemma}

\begin{proof}
\leanok
A bounded function on a closed ball of finite measure is integrable there. The three functions
are measurable (Mathlib's \texttt{measurable\_\allowbreak fderiv} for $D\Psi$), bounded
(\Cref{lem:ray-integral-derivative}, \Cref{lem:graph-gradient}), and vanish for $\|x'\| > R$ when
the closed support of $g$ lies in $\bar B(0, R)$ (\Cref{lem:graph-ray-integral}).
\end{proof}

\begin{theorem}[The horizontal directions]
\label{thm:graph-horizontal}
\lean{Rd.LipschitzGraph.integral_fderiv_rayIntegral_graphMap}
\leanok
\uses{lem:graph-ray-integral, lem:integral-line-derivative-zero, lem:hyperplane-integrable,
  lem:graph-gradient}
For every $C^1$ function $g \colon V \to \C$ with compact support and every $v' \in H$,
\[
  \int_H D\Psi(\gamma(x')) v' \, dx'
  = \int_H \langle \nabla\varphi(x'), v' \rangle \, g(\gamma(x')) \, dx' .
\]
\end{theorem}

\begin{proof}
\leanok
$\varphi|_H$ is differentiable almost everywhere (Rademacher's theorem,
\Cref{lem:graph-gradient}), so almost everywhere $\partial_{v'} (\Psi \circ \gamma)$ is the
difference of the two integrands (\Cref{lem:graph-ray-integral}), which are integrable
(\Cref{lem:hyperplane-integrable}). $\Psi \circ \gamma$ is Lipschitz with compact support, so
the integral of $\partial_{v'} (\Psi \circ \gamma)$ vanishes
(\Cref{lem:integral-line-derivative-zero}).
\end{proof}

\paragraph{Gauss--Green for $C^1$ functions.}

\begin{theorem}[Gauss--Green on a Lipschitz graph domain]
\label{thm:graph-gauss-green}
\lean{Rd.LipschitzGraph.setIntegral_lineDeriv_eq, Rd.LipschitzGraph.setIntegral_lineDeriv_eq_normal,
  Rd.LipschitzGraph.setIntegral_lineDeriv_e, Rd.LipschitzGraph.setIntegral_lineDeriv_hyperplane,
  Rd.LipschitzGraph.setIntegral_lineDeriv_eq_surface}
\leanok
\uses{thm:graph-fubini, lem:ray-integral-derivative, thm:graph-horizontal, lem:normal-weight,
  lem:hyperplane-integrable, thm:lipschitz-graph-area, def:graph-map}
For every $C^1$ function $g \colon V \to \C$ with compact support and every $v \in V$,
\begin{align*}
  \int_\Omega \partial_v g &= \int_H g(\gamma(x')) \, n_v(x') \, dx' \\
  &= \int_H g(\gamma(x')) \langle v, \nu(x') \rangle \sqrt{1 + \|\nabla\varphi(x')\|^2} \, dx' \\
  &= \int_{\gamma(H)} g(y) \langle v, \nu(P y) \rangle \, d\mathcal H^{n-1}(y) .
\end{align*}
In particular, for every $v' \in H$,
\[
  \int_\Omega \partial_e g = -\int_H g(\gamma(x')) \, dx' ,
  \qquad
  \int_\Omega \partial_{v'} g = \int_H g(\gamma(x')) \langle \nabla\varphi(x'), v' \rangle \, dx' .
\]
\end{theorem}

\begin{proof}
\leanok
By \Cref{thm:graph-fubini} in graph coordinates and \Cref{lem:ray-integral-derivative},
$\int_\Omega \partial_v g = \int_H D\Psi(\gamma(x')) v \, dx'$. With
$v = Pv + \langle v, e \rangle e$ and $D\Psi e = -g$ the integrand is
$D\Psi(\gamma(x')) Pv - \langle v, e \rangle g(\gamma(x'))$, and \Cref{thm:graph-horizontal}
for $v' = Pv$ replaces the first term by $\langle \nabla\varphi(x'), Pv \rangle g(\gamma(x'))$,
where $\langle \nabla\varphi(x'), Pv \rangle = \langle \nabla\varphi(x'), v \rangle$
(\Cref{def:graph-map}). The second form is \Cref{lem:normal-weight}, the third the area
formula \Cref{thm:lipschitz-graph-area}, and $n_e = -1$ and
$n_{v'} = \langle \nabla\varphi, v' \rangle$ for $v' \in H$.
\end{proof}

With the surface measure $\sigma$ of $\Gamma = \gamma(H)$ the theorem reads
$\int_\Omega \partial_v g = \int_\Gamma g \langle v, \nu \rangle \, d\sigma$.

\paragraph{Green's formula for $H^1(V)$.}

\begin{lemma}[Continuity of both sides of Green's formula]
\label{lem:green-continuous}
\lean{Rd.setIntegral_weakDeriv_mul_add_eq, Rd.continuous_setIntegral_weakDeriv_mul_add,
  Rd.LipschitzGraph.integral_mul_mul_normalWeight_eq,
  Rd.LipschitzGraph.continuous_setIntegral_weakDeriv_mul_add,
  Rd.LipschitzGraph.continuous_integral_traceH1_mul_normalWeight}
\leanok
\uses{lem:sobolev-one-weak-derivative-operator, lem:graph-trace-sobolev-one-operator,
  lem:normal-weight, lem:indicator-lp}
For every measurable $s \subseteq V$, every $v \in V$ and all $U, W \in H^1(V)$,
\[
  \int_s (\partial_v U) \, \iota W + \iota U \, \partial_v W
  = \langle \mathbf 1_s \partial_v U, \iota W \rangle + \langle \mathbf 1_s \iota U,
    \partial_v W \rangle
\]
with the bilinear pairing of $L_2(V)$ with itself, and this integral is continuous in
$(U, W) \in H^1(V) \times H^1(V)$. For every $v \in V$, the maps
\[
  (U, W) \mapsto \int_\Omega (\partial_v U) \, W + U \, \partial_v W ,
  \qquad
  (U, W) \mapsto \int_H \operatorname{Tr} U \cdot \operatorname{Tr} W \cdot n_v
\]
are continuous on $H^1(V) \times H^1(V)$, with $\partial_v U \in L_2(V)$ the weak derivative of
$\iota U$. For all $a, b \in L_2(H)$,
$\int_H a b \, n_v = \langle a, n_v b \rangle$ with the bilinear pairing of $L_2(H)$ with itself.
\end{lemma}

\begin{proof}
\leanok
$\int_s a b$ for $a, b \in L_2(V)$ is the bounded bilinear pairing of $\mathbf 1_s a$
with $b$ (\Cref{lem:indicator-lp}), and multiplication by
$n_v \in L_\infty(H)$ is a bounded operator of $L_2(H)$ (Mathlib's
\texttt{ContinuousLinearMap.\allowbreak holderL}). The operators $\partial_v$, $\iota$, the
indicator and the trace are bounded (\Cref{lem:sobolev-one-weak-derivative-operator},
\Cref{lem:indicator-lp}, \Cref{lem:graph-trace-sobolev-one-operator}).
\end{proof}

\begin{lemma}[Green's formula for Schwartz functions]
\label{lem:green-schwartz}
\lean{Rd.lineDeriv_mul, Rd.lineDeriv_mul_schwartz, Rd.LipschitzGraph.setIntegral_green_ofSchwartz}
\leanok
\uses{thm:graph-gauss-green, lem:graph-trace-schwartz, lem:sobolev-of-schwartz}
For all $f, g \colon V \to \C$ differentiable at $x$ and every $v \in V$,
$\partial_v (f g)(x) = \partial_v f(x) \, g(x) + f(x) \, \partial_v g(x)$. For all Schwartz
functions $w_1, w_2$ and every $v \in V$,
$\partial_v (w_1 w_2) = (\partial_v w_1) w_2 + w_1 \partial_v w_2$, and if $w_1$ has compact
support,
\[
  \int_\Omega (\partial_v w_1) \, w_2 + w_1 \, \partial_v w_2
  = \int_H \operatorname{Tr} w_1 \cdot \operatorname{Tr} w_2 \cdot n_v .
\]
\end{lemma}

\begin{proof}
\leanok
The product rule is that of derivatives (Mathlib's \texttt{fderiv\_\allowbreak
fun\_\allowbreak mul}). Then \Cref{thm:graph-gauss-green} for the compactly supported $C^1$ function
$w_1 w_2$, with $\iota w_i = w_i$, $\partial_v \iota w_i = \partial_v w_i$
(\Cref{lem:sobolev-of-schwartz}) and $\operatorname{Tr} w_i = w_i \circ \gamma$
(\Cref{lem:graph-trace-schwartz}).
\end{proof}

\begin{theorem}[Green's formula on a Lipschitz graph domain]
\label{thm:graph-green}
\lean{Rd.LipschitzGraph.setIntegral_green_besselPotentialSpace, Rd.LipschitzGraph.setIntegral_green,
  Rd.LipschitzGraph.setIntegral_green_surface}
\leanok
\uses{lem:green-schwartz, lem:green-continuous, thm:sobolev-schwartz-dense,
  lem:sobolev-one-weak-derivative-operator, lem:sobolev-one-weak-gradient-norm,
  thm:lipschitz-graph-area, lem:normal-weight}
For all $U, W \in H^1(V)$ and every $v \in V$,
\begin{align*}
  \int_\Omega (\partial_v U) \, W + U \, \partial_v W
  &= \int_H \operatorname{Tr} U \cdot \operatorname{Tr} W \cdot n_v \\
  &= \int_{\gamma(H)} \operatorname{Tr} U(P y) \operatorname{Tr} W(P y)
    \langle v, \nu(P y) \rangle \, d\mathcal H^{n-1}(y) .
\end{align*}
For all $f, g \in L_2(V)$ with $f, g \in \Sob{1}{2}(V)$, every $v \in V$ and all
$f_v, g_v \in L_2(V)$ with $f_v = \partial_v f$ and $g_v = \partial_v g$ weakly on $V$,
\[
  \int_\Omega f_v \, g + f \, g_v
  = \int_H \operatorname{Tr} f \cdot \operatorname{Tr} g \cdot n_v .
\]
\end{theorem}

\begin{proof}
\leanok
Both sides of the first identity are continuous on $H^1(V) \times H^1(V)$
(\Cref{lem:green-continuous}) and agree on pairs of compactly supported Schwartz functions
(\Cref{lem:green-schwartz}), which are dense (\Cref{thm:sobolev-schwartz-dense}). The second
form is the area formula (\Cref{thm:lipschitz-graph-area}) with \Cref{lem:normal-weight}. For
$f, g \in \Sob{1}{2}(V)$, $f = \iota U$ and $g = \iota W$ for some $U, W \in H^1(V)$
(\Cref{lem:sobolev-one-weak-gradient-norm}), and $f_v = \partial_v U$, $g_v = \partial_v W$
by uniqueness of
weak derivatives (\Cref{lem:sobolev-one-weak-derivative-operator}).
\end{proof}

\section{Green's formula on bounded Lipschitz domains}
\label{sec:lipschitz-domain-green}

Throughout this section $\Omega$ is a bounded Lipschitz domain in $V$, $\dim V = n$, with boundary
$\Gamma$, boundary balls $B_j$ and graph domains $\Omega_j$ over the hyperplanes $H_j$, with graph
maps $\gamma_j$, projections $P_j$, traces $\operatorname{Tr}_j$, normal weights
$n^j_v = \langle \nabla\varphi_j, v \rangle - \langle v, e_j \rangle$ (\Cref{def:normal-weight})
and chart normals $\nu_j$ (\Cref{def:graph-normal}); $\sigma$ is the surface measure,
$T_\Gamma \colon H^1(V) \to L_2(\Gamma)$ and $\gamma_{L_2} \colon H^1(\Omega) \to L_2(\Gamma)$ are
the traces of \Cref{sec:surface-measure}, $\nu$ is the unit normal field
(\Cref{def:boundary-normal}) and $N_\Gamma$ the trace kernel (\Cref{def:trace-kernel}). For
$U \in H^1(V)$, $\partial_v U \in L_2(V)$ is the weak derivative of $\iota U$. Green's formula on
$\Omega$ (McLean \cite[Theorem 3.34]{McLean}) is pasted from Green's formula on the graph domains
(\Cref{thm:graph-green}) with a resolution of unity that equals one near $\Gamma$; the interior
piece integrates to zero by integration by parts on $V$. The chart normals agree almost
everywhere on the overlaps of the balls, by Gauss--Green in two charts, so the boundary term is
$\int_\Gamma \gamma_{L_2} u \, \gamma_{L_2} w \, \langle v, \nu \rangle \, d\sigma$. For $u$ in
the trace kernel, Green's formula has no boundary term and gives the weak derivatives
$\partial_v(\mathbf 1_\Omega \iota u) = \mathbf 1_\Omega \partial_v u$; the traces from both sides
of each graph give the converse.

\paragraph{Integration by parts in $H^1(V)$.}

\begin{lemma}[Weak derivatives in $H^1(V)$ are local]
\label{lem:sobolev-one-weak-derivative-local}
\lean{Rd.weakDerivL2_ae_eq_zero_of_ae_eq_zero, Rd.weakDerivL2_ae_eq_of_ae_eq}
\leanok
\uses{lem:sobolev-one-weak-derivative-operator, lem:weak-derivative-basic}
For every open $O \subseteq V$, every $v \in V$ and all $U, W \in H^1(V)$ with $\iota U = \iota W$
almost everywhere on $O$, $\partial_v U = \partial_v W$ almost everywhere on $O$; in particular
$\partial_v U = 0$ almost everywhere on $O$ if $\iota U = 0$ almost everywhere on $O$.
\end{lemma}

\begin{proof}
\leanok
A weak derivative on $V$ is one on $O$, and it vanishes almost everywhere on $O$ where the
function does (\Cref{lem:weak-derivative-basic}); apply this to $U - W$.
\end{proof}

\begin{lemma}[Green's integrand]
\label{lem:green-integrand}
\lean{Rd.integrable_weakDeriv_mul_add, Rd.setIntegral_weakDeriv_mul_add_sum,
  Rd.setIntegral_weakDeriv_mul_add_add, Rd.continuous_integral_weakDeriv_mul_add}
\leanok
\uses{lem:green-continuous, lem:sobolev-one-weak-derivative-operator}
For every $v \in V$ and all $U, W \in H^1(V)$, $(\partial_v U) \, \iota W + \iota U \,
\partial_v W$ is integrable on $V$. For every measurable $s \subseteq V$, every $v \in V$, every
$U \in H^1(V)$ and every finite family $(W_i)_{i \in I}$ in $H^1(V)$,
\[
  \int_s (\partial_v U) \, \iota \Bigl(\sum_i W_i\Bigr)
    + \iota U \, \partial_v \Bigl(\sum_i W_i\Bigr)
  = \sum_i \int_s (\partial_v U) \, \iota W_i + \iota U \, \partial_v W_i ,
\]
in particular for two summands, and for every $v \in V$,
$(U, W) \mapsto \int_V (\partial_v U) \, \iota W + \iota U \, \partial_v W$ is continuous on
$H^1(V) \times H^1(V)$.
\end{lemma}

\begin{proof}
\leanok
Products of two $L_2$ functions are integrable. By \Cref{lem:green-continuous} the integral over
$s$ is a sum of bilinear pairings, linear in $W$ since $\iota$ and $\partial_v$ are linear
(\Cref{lem:sobolev-one-weak-derivative-operator}), and continuous; for $s = V$ this is the
continuity.
\end{proof}

\begin{theorem}[Integration by parts in $H^1(V)$]
\label{thm:sobolev-one-integration-by-parts}
\lean{Rd.integral_lineDerivOp_mul_add_eq_zero, Rd.integral_weakDeriv_mul_add_eq_zero,
  Rd.integral_mul_add_eq_zero_of_hasWeakLineDeriv}
\leanok
\uses{lem:green-integrand, thm:sobolev-schwartz-dense, lem:sobolev-of-schwartz,
  lem:sobolev-one-weak-gradient-norm, lem:sobolev-one-weak-derivative-operator}
For every $v \in V$ and all Schwartz functions $w_1, w_2$,
$\int_V (\partial_v w_1) \, w_2 + w_1 \, \partial_v w_2 = 0$. For every $v \in V$ and all
$U, W \in H^1(V)$,
\[
  \int_V (\partial_v U) \, \iota W + \iota U \, \partial_v W = 0 .
\]
For all $f, g \in L_2(V)$ with $f, g \in \Sob{1}{2}(V)$, every $v \in V$ and all
$f_v, g_v \in L_2(V)$ with $f_v = \partial_v f$ and $g_v = \partial_v g$ weakly on $V$,
$\int_V f_v \, g + f \, g_v = 0$.
\end{theorem}

\begin{proof}
\leanok
For Schwartz functions this is Mathlib's integration by parts
\texttt{SchwartzMap.\allowbreak integral\_\allowbreak mul\_\allowbreak lineDerivOp\_\allowbreak
right\_\allowbreak eq\_\allowbreak neg\_\allowbreak left}. Both sides of the second identity are
continuous on $H^1(V) \times H^1(V)$ (\Cref{lem:green-integrand}), and they agree on pairs of
Schwartz functions (\Cref{lem:sobolev-of-schwartz}), which are dense
(\Cref{thm:sobolev-schwartz-dense}). For $f, g \in \Sob{1}{2}(V)$, $f = \iota U$ and
$g = \iota W$ (\Cref{lem:sobolev-one-weak-gradient-norm}), and $f_v = \partial_v U$,
$g_v = \partial_v W$ by uniqueness of weak derivatives
(\Cref{lem:sobolev-one-weak-derivative-operator}).
\end{proof}

\paragraph{Localisation.}

\begin{lemma}[Localisation of integrals]
\label{lem:localisation-integrals}
\lean{Rd.setIntegral_eq_of_inter_eq, Rd.integral_smul_indicatorLp_eq_setIntegral_mul,
  Rd.lineDeriv_eq_zero_of_notMem_tsupport}
\leanok
\uses{lem:indicator-lp}
For every measure space, all measurable sets $s, t$ and every set $B$ with
$s \cap B = t \cap B$, and every $F$ with values in a real normed space that vanishes almost
everywhere outside $B$, $\int_s F = \int_t F$. For every measurable $s \subseteq V$, every
$F \in L_p(V)$, every $\phi \colon V \to \R$ and every $h \colon V \to \C$ with $h = \phi$ almost
everywhere, $\int_V \phi \, \mathbf 1_s F = \int_s F h$. For every $g \colon V \to \C$, every
$x \notin \operatorname{supp} g$ and every $v \in V$, $\partial_v g(x) = 0$.
\end{lemma}

\begin{proof}
\leanok
Almost everywhere $\mathbf 1_s F = \mathbf 1_t F$: inside $B$ the indicators agree and outside
$B$ the function vanishes. The zero extension $\mathbf 1_s F$ in $L_p(V)$ is $F$ on $s$ and $0$
outside almost everywhere (\Cref{lem:indicator-lp}). A function that vanishes on a neighbourhood
of $x$ has zero directional derivatives at $x$.
\end{proof}

\begin{lemma}[Products with smooth functions]
\label{lem:sobolev-one-products}
\lean{Rd.exists_toL2_eq_mulLp, Rd.ae_weakDeriv_mul_add_eq_zero_of_mulLp,
  Rd.BoundedLipschitzDomain.BoundaryPartition.hasCompactSupport,
  Rd.BoundedLipschitzDomain.BoundaryPartition.exists_toL2_eq_localize}
\leanok
\uses{lem:product-rule-lp, lem:sobolev-one-weak-gradient-norm,
  lem:sobolev-one-weak-derivative-local, def:trace-space-setting}
For every smooth real $\chi$ with compact support and every $W \in H^1(V)$ there is
$W' \in H^1(V)$ with $\iota W' = \chi \, \iota W$. For every continuous bounded real $\chi$, all
$U, W, W' \in H^1(V)$ with $\iota W' = \chi \, \iota W$ and every $v \in V$,
$(\partial_v U) \, \iota W' + \iota U \, \partial_v W' = 0$ almost everywhere outside
$\operatorname{supp} \chi$. For every resolution of unity $\psi$ on $\Gamma$ and every $j$,
$\psi_j$ has compact support, and for every $W \in H^1(V)$ there is $W' \in H^1(V)$ with
$\iota W' = \psi_j \, \iota W$.
\end{lemma}

\begin{proof}
\leanok
$\iota W \in \Sob{1}{2}(V)$, so $\chi \, \iota W \in \Sob{1}{2}(V)$ (\Cref{lem:product-rule-lp}),
and it is $\iota W'$ (\Cref{lem:sobolev-one-weak-gradient-norm}). Outside
$\operatorname{supp} \chi$, an open set,
$\iota W' = 0$, hence $\partial_v W' = 0$ almost everywhere there
(\Cref{lem:sobolev-one-weak-derivative-local}). The $\psi_j$ are smooth with closed support in
the bounded ball $B_j$.
\end{proof}

\paragraph{Test functions on the hyperplane and the domain below the graph.}

\begin{lemma}[Test functions on the hyperplane are restrictions]
\label{lem:hyperplane-test-restriction}
\lean{Rd.LipschitzGraph.exists_contDiff_comp_graphMap_eq}
\leanok
\uses{def:graph-map, lem:resolution-of-unity}
For every Lipschitz graph domain with hyperplane $H$, graph map $\gamma$ and projection $P$, every
open $W \subseteq V$ and every smooth $\phi \colon H \to \R$ with compact support in
$\gamma^{-1}(W)$, there is a smooth $g \colon V \to \R$ with compact support in $W$ and
$g(\gamma(x')) = \phi(x')$ for every $x' \in H$.
\end{lemma}

\begin{proof}
\leanok
$K = \gamma(\operatorname{supp} \phi)$ is a compact subset of $W$, so there is a smooth $\chi$
with support in the bounded open set $W \cap \{y : \operatorname{dist}(y, K) < 1\}$ that equals
$1$ on a neighbourhood of $K$ (\Cref{lem:resolution-of-unity} with one open set); its support is
compact. Then $g(x) = \phi(P x) \, \chi(x)$, since $P \gamma(x') = x'$.
\end{proof}

\begin{definition}[The domain below the graph]
\label{def:lower-graph-domain}
\lean{Rd.LipschitzGraph.lower}
\leanok
\uses{def:lipschitz-graph}
For a Lipschitz graph domain $\Omega = \{x : \varphi(x) < \langle x, e \rangle\}$ with Lipschitz
constant $M$, the lower graph domain $\Omega_-$ is the Lipschitz graph domain with direction $-e$,
function $-\varphi$ and Lipschitz constant $M$.
\end{definition}

\begin{lemma}[The domain below the graph]
\label{lem:lower-graph-domain}
\lean{Rd.LipschitzGraph.mem_lower_domain_iff, Rd.LipschitzGraph.notMem_domain_of_mem_lower_domain,
  Rd.LipschitzGraph.lower_graphMap_proj_graphMap}
\leanok
\uses{def:lower-graph-domain, lem:boundary-charts}
For every Lipschitz graph domain $\Omega$ with graph map $\gamma$ and lower graph domain
$\Omega_-$ with graph map $\gamma_-$ and projection $P_-$:
$\Omega_- = \{x : \langle x, e \rangle < \varphi(x)\}$, $\Omega_- \cap \Omega = \emptyset$, and
$\gamma_-(P_- \gamma(x')) = \gamma(x')$ for every $x' \in H$.
\end{lemma}

\begin{proof}
\leanok
$-\varphi(x) < \langle x, -e \rangle$ means $\langle x, e \rangle < \varphi(x)$. A point
$y = \gamma(x')$ has $\langle y, e \rangle = \varphi(y)$, so it lies on the graph of $\Omega_-$
(\Cref{lem:boundary-charts}).
\end{proof}

\begin{theorem}[The traces from both sides of a graph]
\label{thm:traces-both-sides}
\lean{Rd.LipschitzGraph.traceH1_ae_eq_traceH1_lower, Rd.LipschitzGraph.traceH1_ae_eq_zero_of_lower}
\leanok
\uses{lem:lower-graph-domain, thm:change-of-chart, lem:trace-locality,
  lem:graph-trace-sobolev-one-operator}
For every Lipschitz graph domain $\Omega$ and every $U \in H^1(V)$,
$\operatorname{Tr} U(x') = \operatorname{Tr}_- U(P_- \gamma(x'))$ for almost every $x' \in H$,
with the trace $\operatorname{Tr}_-$ of the lower graph domain $\Omega_-$. For every open
$O \subseteq V$ and every $U \in H^1(V)$ with $\iota U = 0$ almost everywhere on $\Omega_- \cap O$,
$\operatorname{Tr} U(x') = 0$ for almost every $x' \in H$ with $\gamma(x') \in O$.
\end{theorem}

\begin{proof}
\leanok
The two graphs coincide (\Cref{lem:lower-graph-domain}), so the change of chart
(\Cref{thm:change-of-chart}) identifies the two traces. If $\iota U = 0$ on $\Omega_- \cap O$, the
trace from below vanishes on $\gamma_-^{-1}(O)$ (\Cref{lem:trace-locality}); the null sets of
$H_-$ correspond to null sets of $H$ (\Cref{thm:change-of-chart}).
\end{proof}

\paragraph{Green's formula in chart form.}

\begin{lemma}[A resolution of unity equal to one near the boundary]
\label{lem:boundary-partition-nhds}
\lean{Rd.BoundedLipschitzDomain.exists_boundaryPartition_nhds,
  Rd.BoundedLipschitzDomain.indicator_one_sub_sum_eventuallyEq_zero,
  Rd.BoundedLipschitzDomain.exists_interior_cutoff}
\leanok
\uses{lem:resolution-of-unity, def:boundary-partition, def:bounded-lipschitz}
There is a resolution of unity $\psi$ on $\Gamma$ with $\psi_j \geq 0$ for every $j$,
$\sum_j \psi_j \leq 1$ on $V$ and $\sum_j \psi_j = 1$ on a neighbourhood of $\Gamma$. For every
resolution of unity $\psi$ with $\sum_j \psi_j = 1$ on a neighbourhood of $\Gamma$, the function
$\mathbf 1_\Omega (1 - \sum_j \psi_j)$ vanishes on a neighbourhood of every point of
$V \setminus \Omega$, and there is a smooth real $\chi$ with compact support in $\Omega$ and
$\chi + \sum_j \psi_j = 1$ on $\Omega$.
\end{lemma}

\begin{proof}
\leanok
\Cref{lem:resolution-of-unity} for the compact set $\Gamma$ and the balls $B_j$. A point outside
$\Omega$ lies on $\Gamma$, where $\sum_j \psi_j = 1$ nearby, or outside
$\overline\Omega$. So $\chi = \mathbf 1_\Omega (1 - \sum_j \psi_j)$ agrees near every point with
a smooth function, $1 - \sum_j \psi_j$ inside $\Omega$ and $0$ elsewhere, and its closed support
lies in the bounded set $\Omega$.
\end{proof}

\begin{lemma}[The pieces of Green's formula]
\label{lem:green-pieces}
\lean{Rd.BoundedLipschitzDomain.setIntegral_weakDeriv_mul_add_localize,
  Rd.BoundedLipschitzDomain.setIntegral_weakDeriv_mul_add_mulLp_eq_zero,
  Rd.BoundedLipschitzDomain.toL2_ae_eq_sum_add}
\leanok
\uses{lem:sobolev-one-products, lem:localisation-integrals, thm:graph-green, lem:trace-products,
  thm:sobolev-one-integration-by-parts}
For every resolution of unity $\psi$ on $\Gamma$, every $j$, all $U, W, W' \in H^1(V)$ with
$\iota W' = \psi_j \, \iota W$ and every $v \in V$,
\[
  \int_\Omega (\partial_v U) \, \iota W' + \iota U \, \partial_v W'
  = \int_{H_j} \psi_j(\gamma_j(x')) \operatorname{Tr}_j U(x') \operatorname{Tr}_j W(x') \,
    n^j_v(x') \, dx' .
\]
For every continuous bounded real $\chi$ with closed support in $\Omega$, all
$U, W, W' \in H^1(V)$ with $\iota W' = \chi \, \iota W$ and every $v \in V$,
$\int_\Omega (\partial_v U) \, \iota W' + \iota U \, \partial_v W' = 0$. For every resolution of
unity $\psi$, every continuous bounded real $\chi$ with $\chi + \sum_j \psi_j = 1$ on $\Omega$, and
all $W, W_j, W_0 \in H^1(V)$ with $\iota W_j = \psi_j \, \iota W$ and $\iota W_0 = \chi \, \iota W$,
$\iota W = \iota (\sum_j W_j + W_0)$ almost everywhere on $\Omega$.
\end{lemma}

\begin{proof}
\leanok
The first integrand vanishes almost everywhere outside $B_j$ (\Cref{lem:sobolev-one-products}),
and $\Omega \cap B_j = \Omega_j \cap B_j$, so the integral over $\Omega$ is the integral over
$\Omega_j$ (\Cref{lem:localisation-integrals}); Green's formula on $\Omega_j$
(\Cref{thm:graph-green}) and $\operatorname{Tr}_j (\psi_j \iota W) = (\psi_j \circ \gamma_j)
\operatorname{Tr}_j W$ (\Cref{lem:trace-products}) give the right side. The second integrand
vanishes almost everywhere outside $\Omega$, so its integral over $\Omega$ is its integral over
$V$, which vanishes (\Cref{thm:sobolev-one-integration-by-parts}). The third statement is
$\sum_j \psi_j + \chi = 1$ on $\Omega$.
\end{proof}

\begin{theorem}[Green's formula in chart form]
\label{thm:domain-green-charts-nhds}
\lean{Rd.BoundedLipschitzDomain.setIntegral_green_of_nhds}
\leanok
\uses{lem:boundary-partition-nhds, lem:green-pieces, lem:green-integrand,
  lem:sobolev-one-weak-derivative-local}
For every resolution of unity $\psi$ on $\Gamma$ with $\sum_j \psi_j = 1$ on a neighbourhood of
$\Gamma$, all $U, W \in H^1(V)$ and every $v \in V$,
\[
  \int_\Omega (\partial_v U) \, \iota W + \iota U \, \partial_v W
  = \sum_j \int_{H_j} \psi_j(\gamma_j(x')) \operatorname{Tr}_j U(x') \operatorname{Tr}_j W(x') \,
    n^j_v(x') \, dx' .
\]
\end{theorem}

\begin{proof}
\leanok
Take $\chi$ from \Cref{lem:boundary-partition-nhds} and $W_j$, $W_0$ with
$\iota W_j = \psi_j \, \iota W$, $\iota W_0 = \chi \, \iota W$ (\Cref{lem:sobolev-one-products}).
On the open set $\Omega$, $\iota W = \iota (\sum_j W_j + W_0)$, hence also
$\partial_v W = \partial_v (\sum_j W_j + W_0)$ (\Cref{lem:green-pieces},
\Cref{lem:sobolev-one-weak-derivative-local}). The integral is linear in $W$
(\Cref{lem:green-integrand}), and \Cref{lem:green-pieces} evaluates the pieces.
\end{proof}

\paragraph{The unit normal.}

\begin{lemma}[Gauss--Green in one chart]
\label{lem:domain-gauss-green-chart}
\lean{Rd.BoundedLipschitzDomain.setIntegral_lineDeriv_eq_chart,
  Rd.BoundedLipschitzDomain.integrable_mul_inner_normalAt,
  Rd.BoundedLipschitzDomain.integral_mul_inner_normalAt_sub_eq_zero,
  Rd.BoundedLipschitzDomain.integral_graphMap_mul_inner_normalAt_sub_eq_zero,
  Rd.BoundedLipschitzDomain.locallyIntegrable_inner_normalAt_sub_mul_gram}
\leanok
\uses{thm:graph-gauss-green, lem:localisation-integrals, thm:surface-measure-charts-complex,
  thm:surface-measure-null-sets, lem:graph-normal, lem:graph-gradient}
For every $j$, every $C^1$ function $g \colon V \to \C$ with closed support in $B_j$ and every
$v \in V$,
\[
  \int_\Omega \partial_v g = \int_\Gamma g(y) \langle v, \nu_j(P_j y) \rangle \, d\sigma(y) .
\]
For every $j$, every continuous $g \colon V \to \C$ with compact support and every $v \in V$,
$y \mapsto g(y) \langle v, \nu_j(P_j y) \rangle$ is $\sigma$-integrable. For all $j, k$, every
$C^1$ function $g \colon V \to \C$ with closed support in $B_j \cap B_k$
and every $v \in V$,
\[
  \int_{H_j} g(\gamma_j(x')) \bigl\langle v, \nu_j(x') - \nu_k(P_k \gamma_j(x')) \bigr\rangle
    \sqrt{1 + \|\nabla\varphi_j(x')\|^2} \, dx' = 0 ,
\]
and the density in this integral is locally integrable on $H_j$.
\end{lemma}

\begin{proof}
\leanok
A closed subset of the ball $B_j$ is compact. Inside $B_j$, $\Omega$ agrees with $\Omega_j$ and
$\partial_v g$ vanishes outside the support of $g$ (\Cref{lem:localisation-integrals}), so the
left side is Gauss--Green on $\Omega_j$
(\Cref{thm:graph-gauss-green}), and the surface measure in the chart
(\Cref{thm:surface-measure-charts-complex}) gives the right side. For $g$ supported in
$B_j \cap B_k$, both charts give $\int_\Omega \partial_v g$, the boundary integrands are bounded
(\Cref{thm:surface-measure-null-sets}: $\sigma$ is finite), and the difference in the chart $j$ is
the second identity. The density is measurable and bounded by
$2 \|v\| \sqrt{1 + M_j^2}$ (\Cref{lem:graph-normal}, \Cref{lem:graph-gradient}).
\end{proof}

\begin{theorem}[The chart normals agree]
\label{thm:chart-normals-agree}
\lean{Rd.BoundedLipschitzDomain.normalAt_graphMap_ae_eq,
  Rd.BoundedLipschitzDomain.normalAt_proj_ae_eq}
\leanok
\uses{lem:domain-gauss-green-chart, lem:hyperplane-test-restriction,
  thm:surface-measure-null-sets, lem:graph-gradient}
For all $j, k$, $\nu_j(x') = \nu_k(P_k \gamma_j(x'))$ for almost every $x' \in H_j$ with
$\gamma_j(x') \in B_j \cap B_k$, and $\nu_j(P_j y) = \nu_k(P_k y)$ for $\sigma$-almost every
$y \in B_j \cap B_k$.
\end{theorem}

\begin{proof}
\leanok
For every vector $v$ of an orthonormal basis of $V$, the density of
\Cref{lem:domain-gauss-green-chart} integrates to zero against every smooth $\phi$ with compact
support in the open set $\gamma_j^{-1}(B_j \cap B_k)$, since $\phi = g \circ \gamma_j$ for a smooth
$g$ with compact support in $B_j \cap B_k$ (\Cref{lem:hyperplane-test-restriction}). So it vanishes
almost everywhere on that set (Mathlib's \texttt{IsOpen.\allowbreak ae\_\allowbreak eq\_\allowbreak
zero\_\allowbreak of\_\allowbreak integral\_\allowbreak contDiff\_\allowbreak smul\_\allowbreak
eq\_\allowbreak zero}), and the Gram density is positive (\Cref{lem:graph-gradient}). The second
statement is the first in the chart $j$ (\Cref{thm:surface-measure-null-sets}).
\end{proof}

\begin{theorem}[The unit normal]
\label{thm:boundary-normal-chart}
\lean{Rd.BoundedLipschitzDomain.normal_ae_eq_normalAt, Rd.BoundedLipschitzDomain.norm_normal_ae,
  Rd.BoundedLipschitzDomain.inner_normal_neg_ae}
\leanok
\uses{thm:chart-normals-agree, def:boundary-normal, lem:graph-normal, def:surface-measure}
For every $j$ and $\sigma$-almost every $y \in B_j$, $\nu(y) = \nu_j(P_j y)$ and
$\langle \nu(y), w \rangle < 0$ for every $w$ in the cone of $\Omega_j$. For $\sigma$-almost every
$y$, $\|\nu(y)\| = 1$.
\end{theorem}

\begin{proof}
\leanok
For $\sigma$-almost every $y \in \Gamma \cap B_j$, every chart normal $\nu_k(P_k y)$ with
$\psi^0_k(y) \neq 0$ equals $\nu_j(P_j y)$ (\Cref{thm:chart-normals-agree}), and
$\sum_k \psi^0_k(y) = 1$. The chart normals are unit vectors and point away from the cone
(\Cref{lem:graph-normal}), and the balls cover $\Gamma$, which carries $\sigma$.
\end{proof}

\paragraph{Green's formula.}

\begin{lemma}[The boundary integral in the charts]
\label{lem:boundary-integral-charts}
\lean{Rd.BoundedLipschitzDomain.ambientTraceL2_graphMap_ae,
  Rd.BoundedLipschitzDomain.normal_graphMap_ae,
  Rd.BoundedLipschitzDomain.integrable_ambientTraceL2_mul_inner_normal,
  Rd.BoundedLipschitzDomain.integral_ambientTraceL2_mul_inner_normal_eq}
\leanok
\uses{thm:surface-measure-null-sets, thm:ambient-trace-l2, thm:boundary-trace-charts,
  thm:boundary-normal-chart, lem:boundary-normal, thm:surface-measure-charts-complex,
  lem:normal-weight}
For every $j$ and every
$U \in H^1(V)$, $T_\Gamma U(\gamma_j(x')) = \operatorname{Tr}_j U(x')$ and
$\nu(\gamma_j(x')) = \nu_j(x')$ for almost every $x' \in H_j$ with $\gamma_j(x') \in B_j$. For all
$U, W \in H^1(V)$ and every $v \in V$, $T_\Gamma U \, T_\Gamma W \langle v, \nu \rangle$ is
$\sigma$-integrable, and for every resolution of unity $\psi$ on $\Gamma$,
\[
  \int_\Gamma T_\Gamma U \, T_\Gamma W \langle v, \nu \rangle \, d\sigma
  = \sum_j \int_{H_j} \psi_j(\gamma_j(x')) \operatorname{Tr}_j U(x') \operatorname{Tr}_j W(x') \,
    n^j_v(x') \, dx' .
\]
\end{lemma}

\begin{proof}
\leanok
The null sets of $\sigma$ on $B_j$ are those of the chart (\Cref{thm:surface-measure-null-sets}).
$T_\Gamma U$ is the boundary trace $\sigma$-almost everywhere (\Cref{thm:ambient-trace-l2}), which
is $\operatorname{Tr}_j U$ in the chart (\Cref{thm:boundary-trace-charts}), and $\nu = \nu_j \circ
P_j$ (\Cref{thm:boundary-normal-chart}). The integrand is a product of two $L_2(\Gamma)$
functions and a function bounded by $\|v\|$ (\Cref{lem:boundary-normal}). The surface measure in
the charts (\Cref{thm:surface-measure-charts-complex}) and
$\langle v, \nu_j \rangle \sqrt{1 + \|\nabla\varphi_j\|^2} = n^j_v$ (\Cref{lem:normal-weight})
give the identity.
\end{proof}

\begin{theorem}[Green's formula on a bounded Lipschitz domain]
\label{thm:domain-green}
\lean{Rd.BoundedLipschitzDomain.setIntegral_green_surface,
  Rd.BoundedLipschitzDomain.setIntegral_green_sobolevOne,
  Rd.BoundedLipschitzDomain.setIntegral_green_sobolevOne_of_hasWeakLineDeriv,
  Rd.BoundedLipschitzDomain.setIntegral_green_besselPotentialSpace,
  Rd.BoundedLipschitzDomain.setIntegral_green}
\leanok
\uses{thm:domain-green-charts-nhds, lem:boundary-integral-charts, lem:boundary-partition-nhds,
  thm:trace-l2, lem:sobolev-one-weak-gradient-norm, lem:sobolev-one-weak-derivative-operator,
  lem:weak-derivative-basic}
For every resolution of unity $\psi$ on $\Gamma$, all $U, W \in H^1(V)$ and every $v \in V$,
\begin{align*}
  \int_\Omega (\partial_v U) \, \iota W + \iota U \, \partial_v W
  &= \int_\Gamma T_\Gamma U \, T_\Gamma W \langle v, \nu \rangle \, d\sigma \\
  &= \sum_j \int_{H_j} \psi_j(\gamma_j(x')) \operatorname{Tr}_j U(x') \operatorname{Tr}_j W(x') \,
    n^j_v(x') \, dx' .
\end{align*}
For all $u, w \in H^1(\Omega)$ and every $v \in V$,
\[
  \int_\Omega (\partial_v u) \, w + u \, \partial_v w
  = \int_\Gamma \gamma_{L_2} u \, \gamma_{L_2} w \langle v, \nu \rangle \, d\sigma ,
\]
with $u, w$ read as elements of $H^1(V)$ on the left, and the same holds with $\partial_v u$,
$\partial_v w$ replaced by any weak derivatives $u_v$, $w_v$ of $u$, $w$ on $\Omega$. For every
resolution of unity $\psi$, all
$f, g \in L_2(V)$ with $f, g \in \Sob{1}{2}(V)$, every $v \in V$ and all $f_v, g_v \in L_2(V)$ with
$f_v = \partial_v f$ and $g_v = \partial_v g$ weakly on $V$,
\[
  \int_\Omega f_v \, g + f \, g_v
  = \sum_j \int_{H_j} \psi_j(\gamma_j(x')) \operatorname{Tr}_j f(x') \operatorname{Tr}_j g(x') \,
    n^j_v(x') \, dx' .
\]
\end{theorem}

\begin{proof}
\leanok
With a resolution of unity equal to one near $\Gamma$ (\Cref{lem:boundary-partition-nhds}), the
left side is the chart form of \Cref{thm:domain-green-charts-nhds}, which is the boundary integral
by \Cref{lem:boundary-integral-charts}; the same lemma for any $\psi$ gives the chart form for
$\psi$. For $u, w \in H^1(\Omega)$, $\gamma_{L_2} u = T_\Gamma u$ (\Cref{thm:trace-l2}), and
weak derivatives on $\Omega$ agree almost everywhere on $\Omega$ with the weak derivatives on $V$
(\Cref{lem:weak-derivative-basic}). For $f, g \in \Sob{1}{2}(V)$, $f = \iota U$ and
$g = \iota W$ (\Cref{lem:sobolev-one-weak-gradient-norm}), and $f_v = \partial_v U$,
$g_v = \partial_v W$
by uniqueness of weak derivatives (\Cref{lem:sobolev-one-weak-derivative-operator}).
\end{proof}

In chart form the formula holds for every resolution of unity, of any sign.

\paragraph{The zero extension.}

\begin{lemma}[Functions with vanishing trace]
\label{lem:green-trace-kernel}
\lean{Rd.BoundedLipschitzDomain.setIntegral_green_of_mem_traceKer,
  Rd.BoundedLipschitzDomain.hasWeakLineDeriv_indicatorLp_of_mem_traceKer}
\leanok
\uses{thm:domain-green, thm:trace-l2, lem:weak-derivative-integral, lem:sobolev-of-schwartz,
  def:test-schwartz, lem:directional-derivatives-commute, lem:localisation-integrals}
For every $u \in N_\Gamma$, every $W \in H^1(V)$ and every $v \in V$,
$\int_\Omega (\partial_v u) \, \iota W + \iota u \, \partial_v W = 0$, and
$\mathbf 1_\Omega \partial_v u$ is the weak derivative $\partial_v (\mathbf 1_\Omega \iota u)$
on $V$.
\end{lemma}

\begin{proof}
\leanok
$\gamma_{L_2} u = 0$ (\Cref{thm:trace-l2}), so the boundary term of \Cref{thm:domain-green}
vanishes. For a real test function $\phi$ on $V$, take $W$ the Schwartz function $\phi$
(\Cref{lem:sobolev-of-schwartz}); then
$\int_V \phi \, \mathbf 1_\Omega \partial_v u = -\int_V (\partial_v \phi) \mathbf 1_\Omega \iota u$,
which is the weak derivative (\Cref{lem:weak-derivative-integral}).
\end{proof}

\begin{theorem}[The zero extension and the trace kernel]
\label{thm:zero-extension-trace-kernel}
\lean{Rd.BoundedLipschitzDomain.memSobolev_indicatorLp_of_mem_traceKer,
  Rd.BoundedLipschitzDomain.mem_traceKer_of_memSobolev_indicatorLp,
  Rd.BoundedLipschitzDomain.memSobolev_indicatorLp_iff_mem_traceKer}
\leanok
\uses{lem:green-trace-kernel, lem:weak-derivative-basic, thm:traces-both-sides,
  lem:trace-locality, lem:lower-graph-domain, thm:trace-kernel-boundary, thm:boundary-trace-charts,
  lem:sobolev-one-weak-gradient-norm}
For every $u \in H^1(\Omega)$, $\mathbf 1_\Omega \iota u \in \Sob{1}{2}(V)$ if and only if
$u \in N_\Gamma$.
\end{theorem}

\begin{proof}
\leanok
If $u \in N_\Gamma$, the weak derivatives of $\mathbf 1_\Omega \iota u$ are
$\mathbf 1_\Omega \partial_v u \in L_2(V)$ (\Cref{lem:green-trace-kernel}). Conversely, let
$F = \mathbf 1_\Omega \iota u = \iota U$ with $U \in H^1(V)$ and fix $j$. On $\Omega_j \cap B_j =
\Omega \cap B_j$, $F = \iota u$, so $\operatorname{Tr}_j u = \operatorname{Tr}_j F$ almost
everywhere on $\gamma_j^{-1}(B_j)$ (\Cref{lem:trace-locality}); the lower graph domain of $\Omega_j$
meets $B_j$ outside $\Omega$ (\Cref{lem:lower-graph-domain}), where $F = 0$, so
$\operatorname{Tr}_j F = 0$ there (\Cref{thm:traces-both-sides}). Hence the boundary trace of $u$
vanishes almost everywhere on $\Gamma$ (\Cref{thm:boundary-trace-charts}), and
$u \in N_\Gamma$ (\Cref{thm:trace-kernel-boundary}).
\end{proof}

\section{The kernel of the trace}
\label{sec:trace-kernel}

Throughout this section $\Omega$ is a bounded Lipschitz domain in $V$, $\dim V = n$, with boundary
$\Gamma$, boundary balls $B_j$ and graph domains $\Omega_j$ with cones $K_j$
(\Cref{def:bounded-lipschitz}, \Cref{lem:graph-cone}); $\tilde H^1(\Omega)$ and $H^1_0(\Omega)$ are
the spaces of \Cref{def:sobolev-tilde-zero-boundary}, $r_\Omega$ is the restriction to
$H^1(\Omega)$,
$N_\Gamma$ is the trace kernel (\Cref{def:trace-kernel}), $H^{1/2}(\Gamma)$ the trace space
(\Cref{def:trace-space}) and $\Phi_c(x) = c^n \Phi(c x)$ (\Cref{lem:dilate}). The kernel of the
trace is $H^1_0(\Omega)$ (McLean \cite[Theorem 3.40]{McLean}; Hiptmair, Pauly and Schulz
\cite{HiptmairPaulySchulz} cite it in this form). By \Cref{thm:zero-extension-trace-kernel},
$u \in N_\Gamma$ if and only if the zero extension $\mathbf 1_\Omega \iota u$ lies in
$\Sob{1}{2}(V)$, so
the theorem says that the elements of $H^1(V)$ vanishing outside $\Omega$ are the elements of
$\tilde H^1(\Omega)$. Such a function is cut into pieces near the boundary and an interior piece;
each boundary piece vanishes outside the closure of a graph domain, and convolution with the dilates
of a smooth function supported in the cone of the graph moves this closure into the graph domain.
The mollified pieces are test functions of $\Omega$ and converge in $H^1(V)$.

\paragraph{Mollification in $H^1(V)$.}

\begin{lemma}[Convergence in $H^1(V)$]
\label{lem:sobolev-one-convergence}
\lean{Rd.tendsto_of_tendsto_toL2}
\leanok
\uses{lem:sobolev-one-weak-gradient-norm, lem:sobolev-one-weak-derivative-operator}
For every filter, every family $(W_i)$ in $H^1(V)$ and every $U \in H^1(V)$ with
$\iota W_i \to \iota U$ and $\partial_v W_i \to \partial_v U$ in $L_2(V)$ for every $v \in V$,
$W_i \to U$ in $H^1(V)$.
\end{lemma}

\begin{proof}
\leanok
For an orthonormal basis $(e_k)$, $\|W_i - U\| \leq \|\iota W_i - \iota U\|_2 + \sum_k
\|\partial_{e_k} W_i - \partial_{e_k} U\|_2$ (\Cref{lem:sobolev-one-weak-gradient-norm}), since
$\partial_{e_k}$ is linear (\Cref{lem:sobolev-one-weak-derivative-operator}).
\end{proof}

\begin{theorem}[Mollification in $H^1(V)$]
\label{thm:sobolev-one-mollification}
\lean{Rd.tendsto_of_toL2_eq_convLp}
\leanok
\uses{lem:weak-derivative-convolution-lp, lem:sobolev-one-convergence, thm:approx-identity,
  def:schwartz-dilate, lem:dilate, lem:sobolev-one-weak-derivative-operator}
For every Schwartz function $\eta$ with $\int \eta = 1$, every filter, every family $(c_i)$ of
nonzero reals with $c_i \to \infty$, every $U \in H^1(V)$ and every family $(W_i)$ in $H^1(V)$ with
$\iota W_i = \eta_{c_i} * \iota U$, $W_i \to U$ in $H^1(V)$.
\end{theorem}

\begin{proof}
\leanok
The dilates $\eta_{c_i}$ are Schwartz functions (\Cref{def:schwartz-dilate}), so
$\partial_v W_i = \eta_{c_i} * \partial_v U$ (\Cref{lem:weak-derivative-convolution-lp}). Both
$\eta_{c_i} * \iota U \to \iota U$ and $\eta_{c_i} * \partial_v U \to \partial_v U$ in $L_2(V)$
(\Cref{thm:approx-identity}), and \Cref{lem:sobolev-one-convergence} concludes.
\end{proof}

\paragraph{Smooth functions supported in $\Omega$.}

\begin{lemma}[Smoothing with compact support]
\label{lem:smoothing-compact-support}
\lean{Rd.exists_schwartz_eq_conv}
\leanok
\uses{lem:conv-derivatives, def:conv}
For every smooth $K \colon V \to \C$ with compact support, every compact $S \subseteq V$ and
every locally integrable $f$ vanishing outside $S$, $K * f$ is a Schwartz function with compact
support in $\operatorname{supp} K + S$.
\end{lemma}

\begin{proof}
\leanok
$K * f$ is smooth (\Cref{lem:conv-derivatives}), and its support lies in
$\operatorname{supp} K + S$, which is compact, hence closed.
\end{proof}

\begin{definition}[Bump mollifiers]
\label{def:bump-mollifier}
\lean{Rd.bumpMollifier, Rd.contDiff_bumpMollifier, Rd.tsupport_bumpMollifier_subset,
  Rd.hasCompactSupport_bumpMollifier, Rd.integral_bumpMollifier}
\leanok
For every smooth bump function $f$ around $c \in V$ with outer radius $r$ (Mathlib's
\texttt{ContDiffBump}), the bump mollifier of $f$ is $f / \int f$ with complex values. It is
smooth with support in $\bar B(c, r)$, compactly supported, and has integral $1$.
\end{definition}

\begin{theorem}[Mollified functions in $\tilde H^1$]
\label{thm:tilde-space-mollified}
\lean{Rd.mem_tildeSpace_of_dilate, Rd.mem_tildeSpace_of_isCompact}
\leanok
\uses{thm:sobolev-one-mollification, lem:smoothing-compact-support, def:bump-mollifier,
  def:sobolev-tilde-zero-boundary, thm:sobolev-schwartz-dense, lem:dilate,
  lem:dual-kernel-derivatives, lem:conv-op, def:conv, lem:sobolev-of-schwartz}
For every open $\Omega' \subseteq V$, every $U \in H^1(V)$, every compact $S \subseteq V$ with
$\iota U = 0$ almost everywhere outside $S$, and every smooth $\eta$ with compact support and
$\int \eta = 1$ such that $\operatorname{supp} \eta_{2^k} + S \subseteq \Omega'$ for all large
$k$, $U \in \tilde H^1(\Omega')$. For every open $\Omega' \subseteq V$, every compact
$S \subseteq \Omega'$ and every $U \in H^1(V)$ with $\iota U = 0$ almost everywhere outside $S$,
$U \in \tilde H^1(\Omega')$.
\end{theorem}

\begin{proof}
\leanok
Replace $\iota U$ by $\mathbf 1_S \iota U$, equal almost everywhere. For every $k$,
$\eta_{2^k} * \mathbf 1_S \iota U$ is a Schwartz function $w_k$ with compact support in
$\operatorname{supp} \eta_{2^k} + S$ (\Cref{lem:smoothing-compact-support}); as an element of
$H^1(V)$ it converges to $U$ (\Cref{thm:sobolev-one-mollification}), and for large $k$ it lies in
$T^1(\Omega')$ (\Cref{thm:sobolev-schwartz-dense}). For compact $S \subseteq \Omega'$ take
$\delta > 0$ such that the $\delta$-neighbourhood of $S$ lies in $\Omega'$, and the bump
mollifier $\eta$ of a bump around $0$ with outer radius $1$ (\Cref{def:bump-mollifier}): for
large $k$, $\operatorname{supp} \eta_{2^k} + S$ lies in the $\delta$-neighbourhood of $S$
(\Cref{lem:dilate}).
\end{proof}

\paragraph{Cone mollifiers.}

\begin{definition}[The cone mollifier]
\label{def:cone-mollifier}
\lean{Rd.LipschitzGraph.coneMollifierBase, Rd.LipschitzGraph.coneMollifier,
  Rd.LipschitzGraph.tsupport_coneMollifier_subset}
\leanok
\uses{def:bump-mollifier, lem:graph-cone}
For every Lipschitz graph domain with cone $K$, unit vector $e$ and $\rho = 1/(8(M + 1))$, the cone
mollifier is the bump mollifier of the smooth bump around $e/4$ with radii $\rho/4$ and $\rho/2$.
Its support lies in $K \cap \bar B(0, 1/2)$.
\end{definition}

The support is $\bar B(e/4, \rho/2) \subseteq B(e/4, \rho) \subseteq K$ (\Cref{lem:graph-cone}),
and $\|e/4\| + \rho/2 \leq 1/2$.

\paragraph{$\tilde H^1(\Omega)$ by supports.}

\begin{lemma}[$\tilde H^1$ and vanishing outside $\Omega$]
\label{lem:tilde-space-vanishes-outside}
\lean{Rd.toL2_ae_eq_zero_of_mem_tildeSpace, Rd.indicatorLp_toL2_restrictProj}
\leanok
\uses{def:sobolev-tilde-zero-boundary, lem:sobolev-test-span-schwartz,
  prop:sobolev-vanishing-one, lem:sobolev-restriction-basic, lem:indicator-lp}
For every open $\Omega' \subseteq V$ and every $U \in \tilde H^1(\Omega')$, $\iota U = 0$ almost
everywhere outside $\Omega'$. For every open $\Omega' \subseteq V$ and every $U \in H^1(V)$ with
$\iota U = 0$ almost everywhere outside $\Omega'$, $\mathbf 1_{\Omega'} \iota (r_{\Omega'} U) =
\iota U$.
\end{lemma}

\begin{proof}
\leanok
The $U$ with $\iota U = 0$ almost everywhere outside $\Omega'$ form the kernel of a bounded
operator into $L_2(V \setminus \Omega')$, a closed subspace that contains $T^1(\Omega')$
(\Cref{lem:sobolev-test-span-schwartz}), hence $\tilde H^1(\Omega')$. $U - r_{\Omega'} U$ vanishes
on $\Omega'$ (\Cref{lem:sobolev-restriction-basic}), so $\iota (r_{\Omega'} U) = \iota U$ almost
everywhere on $\Omega'$ (\Cref{prop:sobolev-vanishing-one}).
\end{proof}

\begin{lemma}[Chart pieces]
\label{lem:tilde-space-chart-pieces}
\lean{Rd.BoundedLipschitzDomain.mem_tildeSpace_of_mulLp_chart,
  Rd.BoundedLipschitzDomain.eq_sum_add_of_ae_eq_zero}
\leanok
\uses{thm:tilde-space-mollified, def:cone-mollifier, lem:graph-domain-closure-cone,
  lem:dilate, def:bounded-lipschitz, lem:green-pieces, lem:sobolev-one-to-l2,
  lem:product-modulus, def:trace-space-setting, def:boundary-partition}
For every $j$, every $U \in H^1(V)$ with $\iota U = 0$ almost everywhere outside $\Omega$, every
continuous bounded real $\zeta$ with $\operatorname{supp} \zeta \subseteq B_j$ and every
$W \in H^1(V)$ with $\iota W = \zeta \, \iota U$, $W \in \tilde H^1(\Omega)$. For every resolution
of unity $\psi$ on $\Gamma$, every continuous bounded real $\chi$ with $\chi + \sum_j \psi_j = 1$
on $\Omega$, every $U \in H^1(V)$ with $\iota U = 0$ almost everywhere outside $\Omega$ and all
$W_j, W_0 \in H^1(V)$ with $\iota W_j = \psi_j \, \iota U$ and $\iota W_0 = \chi \, \iota U$,
$U = \sum_j W_j + W_0$.
\end{lemma}

\begin{proof}
\leanok
Since $\Omega \cap B_j = \Omega_j \cap B_j$, $\iota W$ vanishes almost everywhere outside the
compact set $S = \overline{\Omega_j} \cap \operatorname{supp} \zeta$. Let $\eta$ be the cone
mollifier of $\Omega_j$ (\Cref{def:cone-mollifier}) and $\delta > 0$ such that
the $\delta$-neighbourhood of $\operatorname{supp} \zeta$ lies in $B_j$. For large $k$,
$z \in \operatorname{supp} \eta_{2^k}$ and $y \in S$, $z \in K_j$
(\Cref{lem:graph-domain-closure-cone}), so $z + y \in \Omega_j$, and $z + y$ lies in the
$\delta$-neighbourhood of $S$ (\Cref{lem:dilate}), so $z + y \in B_j$; hence
$z + y \in \Omega$, and \Cref{thm:tilde-space-mollified} applies. The decomposition holds almost
everywhere on $\Omega$ (\Cref{lem:green-pieces}), both sides vanish almost everywhere outside
$\Omega$, and $\iota$ is injective (\Cref{lem:sobolev-one-to-l2}).
\end{proof}

\begin{theorem}[$\tilde H^1(\Omega)$ by supports]
\label{thm:tilde-space-supports}
\lean{Rd.BoundedLipschitzDomain.mem_tildeSpace_of_ae_eq_zero,
  Rd.BoundedLipschitzDomain.mem_tildeSpace_iff}
\leanok
\uses{lem:tilde-space-chart-pieces, thm:tilde-space-mollified, lem:boundary-partition-nhds,
  lem:sobolev-one-products, lem:tilde-space-vanishes-outside, lem:product-modulus,
  def:boundary-partition}
For every $U \in H^1(V)$, $U \in \tilde H^1(\Omega)$ if and only if $\iota U = 0$ almost everywhere
outside $\Omega$.
\end{theorem}

\begin{proof}
\leanok
One direction is \Cref{lem:tilde-space-vanishes-outside}. Conversely, take a resolution of unity
$\psi$ on $\Gamma$ equal to one near $\Gamma$ and the interior cut-off $\chi$
(\Cref{lem:boundary-partition-nhds}), and $W_j, W_0 \in H^1(V)$ with $\iota W_j = \psi_j \iota U$
and $\iota W_0 = \chi \iota U$ (\Cref{lem:sobolev-one-products}). Then $U = \sum_j W_j + W_0$,
the $W_j$ lie in $\tilde H^1(\Omega)$ (\Cref{lem:tilde-space-chart-pieces}), and $W_0$ vanishes
outside the compact set $\operatorname{supp} \chi \subseteq \Omega$
(\Cref{thm:tilde-space-mollified}).
\end{proof}

In the language of McLean \cite[Theorem 3.29]{McLean}, $\tilde H^1(\Omega)$ is the space of
the elements of $H^1(V)$ supported in $\overline\Omega$, since $\Gamma$ is a null set
(\Cref{cor:boundary-null}).

\paragraph{The kernel of the trace.}

\begin{theorem}[The kernel of the trace is $H^1_0(\Omega)$]
\label{thm:trace-kernel-zero-boundary}
\lean{Rd.BoundedLipschitzDomain.zeroBoundarySpace_le_traceKer,
  Rd.BoundedLipschitzDomain.traceKer_le_zeroBoundarySpace,
  Rd.BoundedLipschitzDomain.traceKer_eq_zeroBoundarySpace}
\leanok
\uses{thm:tilde-space-supports, lem:tilde-space-vanishes-outside,
  thm:zero-extension-trace-kernel, thm:sobolev-tilde-zero-boundary,
  lem:sobolev-one-weak-gradient-norm, def:trace-kernel, lem:indicator-lp,
  prop:sobolev-restriction-l2}
$N_\Gamma = H^1_0(\Omega)$.
\end{theorem}

\begin{proof}
\leanok
Every element of $H^1_0(\Omega)$ is $r_\Omega U$ for some $U \in \tilde H^1(\Omega)$
(\Cref{thm:sobolev-tilde-zero-boundary}); $U$ vanishes outside $\Omega$, so
$\mathbf 1_\Omega \iota (r_\Omega U) = \iota U \in \Sob{1}{2}(V)$
(\Cref{lem:tilde-space-vanishes-outside}) and $r_\Omega U \in N_\Gamma$
(\Cref{thm:zero-extension-trace-kernel}). Conversely, for $u \in N_\Gamma$,
$\mathbf 1_\Omega \iota u = \iota U$ for some $U \in H^1(V)$
(\Cref{thm:zero-extension-trace-kernel}, \Cref{lem:sobolev-one-weak-gradient-norm}); $U$ vanishes
outside $\Omega$, so $U \in \tilde H^1(\Omega)$ (\Cref{thm:tilde-space-supports}), and
$u = r_\Omega U \in H^1_0(\Omega)$ (\Cref{prop:sobolev-restriction-l2},
\Cref{thm:sobolev-tilde-zero-boundary}).
\end{proof}

\begin{corollary}[The quotient by $H^1_0(\Omega)$ and the dual of the trace kernel]
\label{cor:trace-kernel-quotient-duals}
\lean{Rd.BoundedLipschitzDomain.quotientZeroBoundaryEquivTraceSpace,
  Rd.BoundedLipschitzDomain.quotientZeroBoundaryEquivTraceSpace_mk,
  Rd.BoundedLipschitzDomain.tildeEquivTraceKer,
  Rd.BoundedLipschitzDomain.coe_tildeEquivTraceKer_apply,
  Rd.BoundedLipschitzDomain.coe_tildeEquivTraceKer_symm_apply,
  Rd.BoundedLipschitzDomain.dualTraceKerEquiv, Rd.BoundedLipschitzDomain.dualTraceKerEquiv_apply,
  Rd.BoundedLipschitzDomain.dualTraceKerEquiv_symm_apply,
  Rd.BoundedLipschitzDomain.dualTraceKerEquiv_eq_iff}
\leanok
\uses{thm:trace-kernel-zero-boundary, thm:trace-space-operators, thm:sobolev-tilde-zero-boundary,
  lem:sobolev-tilde-in-restriction, thm:tilde-restriction-shift, prop:sobolev-pairing,
  cor:sobolev-duality-zero-boundary}
The map $[u] \mapsto \gamma u$ is an isometric isomorphism of $H^1(\Omega) / H^1_0(\Omega)$ onto
$H^{1/2}(\Gamma)$. The restriction $U \mapsto r_\Omega U$ is an isometric isomorphism of
$\tilde H^1(\Omega)$ onto $N_\Gamma$, and the element of $\tilde H^1(\Omega)$ with restriction
$u \in N_\Gamma$ is $u$ itself, read in $H^1(V)$. For every bounded linear functional $\ell$ on
$N_\Gamma$ there is exactly one $W \in H^{-1}(\Omega)$ with $\ell(u) = [W, u]$ for every
$u \in N_\Gamma$, and $\|W\| = \|\ell\|$; the map $\ell \mapsto W$ is a conjugate-linear isometric
isomorphism of the dual of $N_\Gamma$ onto $H^{-1}(\Omega)$.
\end{corollary}

\begin{proof}
\leanok
$H^1(\Omega) / K \cong K^\perp$ for the closed subspace $K = H^1_0(\Omega) = N_\Gamma$, whose
orthogonal complement is $H^{1/2}(\Gamma)$, and $u - E \gamma u \in N_\Gamma$
(\Cref{thm:trace-space-operators}). The isometry is that of
\Cref{thm:sobolev-tilde-zero-boundary} followed by $H^1_0(\Omega) = N_\Gamma$, with
\Cref{lem:sobolev-tilde-in-restriction} for the inverse. The dual is Riesz's representation
theorem on $N_\Gamma$, this isometry and $\Theta_1 \colon \tilde H^1(\Omega) \to H^{-1}(\Omega)$
(\Cref{thm:tilde-restriction-shift}), as in \Cref{cor:sobolev-duality-zero-boundary}.
\end{proof}

For the trace spaces of Hiptmair, Pauly and Schulz \cite{HiptmairPaulySchulz} at $k = 0$ the
quotient $H^1(\Omega) / H^1_0(\Omega)$ is the classical $H^{1/2}(\Gamma)$.

\chapter{Besov spaces on compact manifolds}
\label{ch:manifolds}

This chapter works on a compact manifold $M$ modelled on a finite-dimensional real inner product
space $E$ with corners $I$, in the sense of Mathlib: charts $e$ map open sets of $M$ to open sets
of a model space $H$, and the extended charts $\kappa = e.\mathrm{extend}\, I$ take values in the
closed set $\Omega = \operatorname{range} I \subseteq E$, the half-space for manifolds with
boundary and $E$ itself for manifolds without boundary. The results assume that $\Omega$ is a
Besov domain (\Cref{def:besov-domain}), which holds for these two models
(\Cref{thm:half-space-besov-domain}), so manifolds with corners are not covered; the boundary
trace is built for the half-space model.

\section{Smooth measures}

\begin{definition}[Smooth measures]
\label{def:smooth-measure}
\lean{Manifold.IsSmoothMeasure, Manifold.measurableSet_extend_target}
\leanok
A measure $\mu$ on $M$ is smooth if for every chart $e$ of the atlas there is a function $\rho_e$,
continuous on the target $\kappa(e.\mathrm{source})$ with $0 < \rho_e < \infty$ there, such that
$\kappa_* (\mu|_{e.\mathrm{source}}) = \rho_e \, \mathrm{vol}|_{\kappa(e.\mathrm{source})}$. The
target is measurable, the intersection of an open set with the closed range of $I$.
\end{definition}

The Riemannian volume and the Lebesgue measure of a smooth domain are smooth measures.

\begin{lemma}[Integrals in a chart]
\label{lem:chart-integrals}
\lean{Manifold.IsSmoothMeasure.lintegral_source_eq, Manifold.IsSmoothMeasure.exists_density_bounds,
  Manifold.IsSmoothMeasure.exists_lintegral_chart_le}
\leanok
\uses{def:smooth-measure}
For a chart $e$ of the atlas and a measurable $F \geq 0$ on $M$,
\[
  \int_{e.\mathrm{source}} F \, d\mu = \int_{\kappa(e.\mathrm{source})} \rho_e(y) F(\kappa^{-1} y) \, dy .
\]
For a compact $K \subseteq e.\mathrm{source}$ there are $c > 0$ and $C < \infty$ such that every
measurable $F \geq 0$ that vanishes outside $K$ satisfies
$\int_{\kappa(e.\mathrm{source})} F \circ \kappa^{-1} \leq c^{-1} \int_M F \, d\mu$ and
$\int_M F \, d\mu \leq C \int_{\kappa(e.\mathrm{source})} F \circ \kappa^{-1}$.
\end{lemma}

\begin{proof}
\leanok
The first identity is the definition of the image measure with $F = (F \circ \kappa^{-1}) \circ
\kappa$ on the source. The continuous density is bounded above and below by positive finite
constants on the compact $\kappa(K)$, which carries $F \circ \kappa^{-1}$.
\end{proof}

\section{The Besov atlas}

From now on $M$ is compact, smooth, and the range $\Omega$ of $I$ is a Besov domain
(\Cref{def:besov-domain}); $\mu$ is a smooth measure and $1 \leq p < \infty$.

\begin{definition}[Besov atlases]
\label{def:besov-atlas}
\lean{Manifold.BesovAtlas, Manifold.BesovAtlas.nonempty, Manifold.contDiffOn_comp_extend_symm,
  Manifold.BesovAtlas.ext, Manifold.BesovAtlas.psi_le_one, Manifold.BesovAtlas.theta_mul_psi}
\leanok
A Besov atlas consists of finitely many charts $e_i$ of the atlas with extended charts
$\kappa_i$, a smooth partition of unity $(\psi_i)$ with $\operatorname{tsupport} \psi_i \subseteq
e_i.\mathrm{source}$, and smooth cut-offs $\theta_i$ with $\theta_i = 1$ on
$\operatorname{tsupport} \psi_i$ and $\operatorname{tsupport} \theta_i \subseteq
e_i.\mathrm{source}$. Every compact smooth manifold has a Besov atlas. A smooth function $g$ on $M$
gives a function $g \circ \kappa^{-1}$ that is smooth within the target of $\kappa$.
\end{definition}

\begin{proof}
\leanok
Finitely many chart sources cover $M$; Mathlib's smooth partitions of unity are subordinate to
them, and a smooth Urysohn function between $\operatorname{tsupport} \psi_i$ and the complement of
an open $V_i$ with $\overline{V_i} \subseteq e_i.\mathrm{source}$ is the cut-off.
\end{proof}

\begin{definition}[Chart pieces]
\label{def:chart-pieces}
\lean{Manifold.BesovAtlas.piece, Manifold.BesovAtlas.aestronglyMeasurable_piece,
  Manifold.BesovAtlas.exists_eLpNorm_piece_le, Manifold.BesovAtlas.memLp_piece,
  Manifold.BesovAtlas.pieceLp, Manifold.BesovAtlas.pieceLp_ae, Manifold.BesovAtlas.piece_add,
  Manifold.BesovAtlas.piece_smul, Manifold.BesovAtlas.piece_ae_eq, Manifold.BesovAtlas.pieceLp_add,
  Manifold.BesovAtlas.pieceLp_smul, Manifold.BesovAtlas.pieceLp_finsetSum}
\leanok
\uses{def:besov-atlas, lem:chart-integrals}
The $i$-th piece of $g : M \to \C$ is $(\psi_i g) \circ \kappa_i^{-1}$ on the target of
$\kappa_i$ and $0$ elsewhere. It is linear in $g$, depends only on the class of $g$ modulo
$\mu$-null sets, and $\|(\psi_i g) \circ \kappa_i^{-1}\|_{L_p(E)} \leq C \|g\|_{L_p(\mu)}$; so the
pieces of $f \in L_p(\mu)$ lie in $L_p(E)$.
\end{definition}

\begin{lemma}[Changes of charts]
\label{lem:coord-change}
\lean{Manifold.range_inter_coordChange_source, Manifold.range_inter_coordChange_target,
  Manifold.mem_coordChange_source, Manifold.coordChange_target_eq_image,
  Manifold.coordChange_apply}
\leanok
For charts $e, e'$ the change of charts $\kappa' \circ \kappa^{-1}$ (Mathlib's
\texttt{ModelWithCorners.extendCoordChange}) has the source $\kappa(e.\mathrm{source} \cap
e'.\mathrm{source}) = \Omega \cap U$ and the target
$\kappa'(e.\mathrm{source} \cap e'.\mathrm{source}) = \Omega \cap U'$ with open $U, U'$.
\end{lemma}

\begin{proof}
\leanok
The image of an open subset $S$ of $H$ under $I$ is $I^{-1}$-preimage of $S$ intersected with
the range of $I$.
\end{proof}

\begin{definition}[The changes of charts with cut-offs]
\label{def:chart-transition}
\lean{Manifold.BesovAtlas.chartChange, Manifold.BesovAtlas.chartChange_apply,
  Manifold.BesovAtlas.tsupport_psi_mul_theta_subset, Manifold.BesovAtlas.transition,
  Manifold.BesovAtlas.transition_apply}
\leanok
\uses{def:cutoff-diffeo, lem:coord-change, def:besov-atlas}
For charts $e, e'$ of the atlas and a smooth $c : M \to \R$ with $\operatorname{tsupport} c
\subseteq e.\mathrm{source} \cap e'.\mathrm{source}$, the change of charts
$\Phi = \kappa' \circ \kappa^{-1}$ with the cut-off $\chi = c \circ \kappa'^{-1}$ and
$K_0 = \kappa'(\operatorname{tsupport} c)$ is a smooth change of variables with a cut-off on
$\Omega$. For two Besov atlases $A, A'$ and charts $i$ of $A$ and $j$ of $A'$, the cut-off
$c = \psi_i \theta'_j$ gives the change of variables $T_{ij}$, with
$\Phi_{ij} = \kappa'_j \circ \kappa_i^{-1}$ and $\chi_{ij} = (\psi_i \theta'_j) \circ
\kappa_j'^{-1}$.
\end{definition}

The smoothness of $\Phi_{ij}$ and of its inverse is the compatibility of the charts in the
maximal atlas, Mathlib's
\texttt{ModelWithCorners.\allowbreak contDiffOn\_\allowbreak extendCoordChange}.

\begin{theorem}[The pieces of two atlases]
\label{thm:pieces-two-atlases}
\lean{Manifold.BesovAtlas.piece_eq_sum_apply, Manifold.BesovAtlas.pieceLp_eq_sum}
\leanok
\uses{def:chart-pieces, def:chart-transition, lem:cutoff-diffeo-lp}
For every $g : M \to \C$, $(\psi_i g) \circ \kappa_i^{-1} = \sum_j T_{ij}\bigl((\psi'_j g) \circ
\kappa_j'^{-1}\bigr)$ at every point; for $f \in L_p(\mu)$ the same holds in $L_p(E)$.
\end{theorem}

\begin{proof}
\leanok
At $y = \kappa_i(x)$, $\psi_i(x) g(x) = \sum_j \psi_i(x) \psi'_j(x) g(x)$, and the $j$-th term is
$\chi_{ij}(\Phi_{ij} y) \cdot (\psi'_j g)(\kappa_j'^{-1} \Phi_{ij} y)$ when $x$ lies in the source
of $e'_j$, since $\theta'_j \psi'_j = \psi'_j$; otherwise $\psi'_j(x) = 0$ and $y$ lies outside the
source of $\Phi_{ij}$.
\end{proof}

\section{Besov spaces on compact manifolds}

\begin{definition}[The Besov quasi-norm of an atlas]
\label{def:manifold-besov-norm}
\lean{Manifold.BesovAtlas.besovNorm, Manifold.BesovAtlas.besovNorm_add_le,
  Manifold.BesovAtlas.besovNorm_const_smul, Manifold.BesovAtlas.exists_enorm_le_besovNorm}
\leanok
\uses{def:chart-pieces, lem:modnorm-on}
For a Besov atlas $A$, $r \in \N$, $s \in \R$ and $1 \leq q \leq \infty$,
\begin{gather*}
  \|f\|_{A, r, s, q} = \sum_i N_\Omega\bigl((\psi_i f) \circ \kappa_i^{-1}\bigr), \\
  N_\Omega(G) = \|G\|_{L_p(\Omega)} + \|(2^{ks} \omega_r(G, 2^{-k}, \Omega)_p)_k\|_{\ell_q} .
\end{gather*}
It is subadditive, homogeneous, and $\|f\|_{L_p(\mu)} \leq C \|f\|_{A, r, s, q}$.
\end{definition}

\begin{theorem}[Independence of the atlas]
\label{thm:atlas-independence}
\lean{Manifold.BesovAtlas.besovNorm_le_besovNorm}
\leanok
\uses{def:manifold-besov-norm, thm:pieces-two-atlases, thm:besov-domain-invariance}
For two Besov atlases $A, A'$, $0 < s < r$ and $1 \leq q \leq \infty$ there is $C < \infty$ with
$\|f\|_{A, r, s, q} \leq C \|f\|_{A', r, s, q}$ for all $f \in L_p(\mu)$.
\end{theorem}

\begin{proof}
\leanok
By \Cref{thm:pieces-two-atlases} and the subadditivity of $N_\Omega$,
$N_\Omega(\text{piece}_i f) \leq \sum_j N_\Omega(T_{ij}\, \text{piece}'_j f)$, and
$N_\Omega(T_{ij} G) \leq C_{ij} N_\Omega(G)$ on the Besov domain $\Omega$
(\Cref{thm:besov-domain-invariance}).
\end{proof}

\begin{definition}[Besov spaces on compact manifolds]
\label{def:manifold-besov}
\lean{Manifold.BesovAtlas.MemBesov, Manifold.BesovAtlas.memBesov_iff_besovNorm_lt_top,
  Manifold.BesovAtlas.memBesov_iff_memBesov, Manifold.MemBesov, Manifold.memBesov_iff,
  Manifold.BesovAtlas.MemBesov.add, Manifold.BesovAtlas.MemBesov.const_smul,
  Manifold.BesovAtlas.MemBesov.zero, Manifold.BesovAtlas.MemBesov.finsetSum,
  Manifold.BesovAtlas.MemBesov.modNormOn_pieceLp_lt_top}
\leanok
\uses{thm:atlas-independence, thm:besov-domain-invariance}
For $s > 0$, $f \in L_p(\mu)$ lies in $\Besov{s}{p}{q}(M)$ if $\|f\|_{A, \lfloor s \rfloor + 1, s,
q} < \infty$ for some Besov atlas $A$. This holds for one atlas if and only if it holds for every
atlas, and with any order $r > s$ in place of $\lfloor s \rfloor + 1$; then every piece has a
finite intrinsic quasi-norm of order $r$. $\Besov{s}{p}{q}(M)$ is a linear subspace of
$L_p(\mu)$.
\end{definition}

\begin{proof}
\leanok
\Cref{thm:atlas-independence} in both directions; for the order, the comparison of orders on a
Besov domain (\Cref{thm:besov-domain-invariance}) piece by piece.
\end{proof}

\begin{lemma}[Functions from a chart]
\label{lem:chart-lift}
\lean{Manifold.chartLift, Manifold.stronglyMeasurable_chartLift, Manifold.memLp_chartLift,
  Manifold.BesovAtlas.piece_chartLift, Manifold.BesovAtlas.chartLiftLp,
  Manifold.BesovAtlas.pieceLp_chartLiftLp, Manifold.BesovAtlas.memBesov_chartLiftLp}
\leanok
\uses{def:manifold-besov, def:chart-transition, thm:diffeo-invariance, lem:chart-integrals}
For a chart $e$ of the atlas, a continuous $\theta$ on $M$ with $\operatorname{tsupport} \theta
\subseteq e.\mathrm{source}$ and $F \in L_p(E)$, the lift $\theta \cdot F \circ \kappa$ on
$e.\mathrm{source}$, $0$ elsewhere, lies in $L_p(\mu)$. For a Besov atlas, the $i$-th piece of the
lift $\theta_k \cdot F \circ \kappa_k$ is $T_{ik} F$. If
$\|F\|_p + \|(2^{j\sigma} \omega_r(F, 2^{-j})_p)_j\|_{\ell_q} < \infty$ with $0 < \sigma < r$,
the lift lies in $\Besov{\sigma}{p}{q}(M)$.
\end{lemma}

\begin{proof}
\leanok
The chart $\kappa$ is continuous on the open source, so the lift is measurable (Mathlib's
\texttt{ContinuousOn.\allowbreak measurable\_\allowbreak piecewise}). Its integral is comparable to
the integral in the chart on the compact $\operatorname{tsupport} \theta$
(\Cref{lem:chart-integrals}), where $\theta$ is bounded. At $y = \kappa_i(x)$ with
$x \in e_k.\mathrm{source}$, $\psi_i(x) \theta_k(x) F(\kappa_k x) = \chi_{ik}(\Phi_{ik} y)
F(\Phi_{ik} y)$. The changes of charts are bounded into the intrinsic space on $\Omega$
(\Cref{thm:diffeo-invariance}).
\end{proof}

\section{Smooth bounded domains}

\begin{lemma}[Flattening a graph domain]
\label{lem:flattening}
\lean{Rd.LipschitzGraph.flatten, Rd.LipschitzGraph.unflatten, Rd.LipschitzGraph.phi_sub_smul,
  Rd.LipschitzGraph.unflatten_flatten, Rd.LipschitzGraph.flatten_unflatten,
  Rd.LipschitzGraph.inner_flatten, Rd.LipschitzGraph.flattenEquiv,
  Rd.LipschitzGraph.continuous_flatten, Rd.LipschitzGraph.continuous_unflatten,
  Rd.LipschitzGraph.contDiff_flatten, Rd.LipschitzGraph.contDiff_unflatten,
  Rd.LipschitzGraph.fderiv_phi_e, Rd.LipschitzGraph.hasFDerivAt_flatten,
  Rd.LipschitzGraph.det_fderiv_flatten, Rd.LipschitzGraph.measurePreserving_flatten}
\leanok
\uses{def:lipschitz-graph}
For a Lipschitz graph domain $\{x : \varphi(x) < \langle x, e \rangle\}$ the flattening
$F(x) = x - \varphi(x) e$ is a bijection of $V$ with inverse $y \mapsto y + \varphi(y) e$ and
$\langle F(x), e \rangle = \langle x, e \rangle - \varphi(x)$. For smooth $\varphi$ both maps are
smooth and $F$ preserves the Lebesgue measure.
\end{lemma}

\begin{proof}
\leanok
$\varphi$ is constant along $e$, so $\varphi(x - \varphi(x) e) = \varphi(x)$ and $D\varphi(x) e = 0$.
The derivative $\mathrm{id} - e \otimes D\varphi(x)$ is a transvection with determinant
$1 - D\varphi(x) e = 1$ (Mathlib's \texttt{LinearMap.transvection.det}), and the change of variables
formula gives $F_* \mathrm{vol} = \mathrm{vol}$.
\end{proof}

\begin{definition}[Smooth bounded domains]
\label{def:smooth-domain}
\lean{Rd.SmoothDomain, Rd.e0, Rd.norm_e0, Rd.coord_zero_eq_inner, Rd.alignTo,
  Rd.alignTo_apply_self, Rd.LipschitzGraph.closure_domain, Rd.mem_closure_iff_of_inter_eq}
\leanok
\uses{def:bounded-lipschitz}
A smooth bounded domain in $\R^n$ is a bounded Lipschitz domain whose graph functions are smooth.
The reflection $A_e$ in $(e - e_0)^\perp$ maps a unit vector $e$ to $e_0 = (1, 0, \dots, 0)$. The
closure of a graph domain is $\{x : \varphi(x) \leq \langle x, e \rangle\}$, and closures of sets
that agree on an open set $B$ agree on $B$.
\end{definition}

The carriers of smooth bounded domains are exactly the bounded sets with smooth local graphs
(\Cref{sec:smooth-local-graphs}).

\begin{definition}[Straighteners and their charts]
\label{def:straightener}
\lean{Rd.SmoothDomain.Straightener, Rd.SmoothDomain.Straightener.tau_mem_closure,
  Rd.SmoothDomain.Straightener.toMap, Rd.SmoothDomain.Straightener.coe_toMap,
  Rd.SmoothDomain.Straightener.invMap, Rd.SmoothDomain.Straightener.coe_invMap,
  Rd.SmoothDomain.Straightener.chart, Rd.SmoothDomain.Straightener.mem_chart_source,
  Rd.SmoothDomain.Straightener.mem_chart_target, Rd.SmoothDomain.Straightener.chart_apply_coe,
  Rd.SmoothDomain.Straightener.chart_symm_coe}
\leanok
\uses{def:smooth-domain}
A straightener of the closure of a smooth domain $\Omega$ is a smooth, measure-preserving bijection
$\sigma$ of $\R^n$ with smooth inverse $\tau$ and an open set $B$ such that for $x \in B$,
$x \in \overline{\Omega}$ if and only if $\sigma(x)_0 \geq 0$. Its chart maps
$\overline{\Omega} \cap B$ by $\sigma$ onto an open subset of the half-space
$\{y : y_0 \geq 0\}$, with inverse $\tau$.
\end{definition}

\begin{lemma}[Boundary and interior straighteners]
\label{lem:straighteners}
\lean{Rd.SmoothDomain.straighten, Rd.SmoothDomain.unstraighten, Rd.SmoothDomain.straighten_coord,
  Rd.SmoothDomain.boundaryStraightener, Rd.SmoothDomain.interiorRadius,
  Rd.SmoothDomain.interiorRadius_pos, Rd.SmoothDomain.ball_interiorRadius_subset,
  Rd.SmoothDomain.interior_coord_pos, Rd.SmoothDomain.interiorStraightener,
  Rd.SmoothDomain.boundaryIndex,
  Rd.SmoothDomain.mem_ball_boundaryIndex, Rd.SmoothDomain.straightenerAt,
  Rd.SmoothDomain.mem_straightenerAt}
\leanok
\uses{def:straightener, lem:flattening}
On the $j$-th boundary ball, $\sigma_j(x) = A_{e_j}(x - \varphi_j(x) e_j)$ is a straightener, with
$\sigma_j(x)_0 = \langle x, e_j \rangle - \varphi_j(x)$. Around $x_0 \in \Omega$ with
$B(x_0, \rho) \subseteq \Omega$ the translation $x \mapsto x - x_0 + 2\rho e_0$ is a straightener.
Every point of $\overline{\Omega}$ lies in the set of one of them.
\end{lemma}

\begin{proof}
\leanok
\Cref{lem:flattening} and the reflection give the properties of $\sigma_j$; on the ball,
$\overline{\Omega}$ agrees with the closure of the graph domain. For the translation,
$|(x - x_0)_0| \leq \|x - x_0\| < \rho$. A point of $\overline{\Omega} \setminus \Omega$ lies on the
boundary, which the boundary balls cover.
\end{proof}

\begin{theorem}[Smooth domains are manifolds with boundary]
\label{thm:smooth-domain-manifold}
\lean{Rd.SmoothDomain.chartedSpace, Rd.SmoothDomain.isManifold, Rd.SmoothDomain.compactSpace,
  Rd.SmoothDomain.volumeOn, Rd.SmoothDomain.Straightener.extend_target,
  Rd.SmoothDomain.isSmoothMeasure_volumeOn}
\leanok
\uses{lem:straighteners, def:smooth-measure}
The closure of a smooth bounded domain, with the charts of all straighteners, is a compact smooth
manifold with boundary modelled on the half-space $\{y \in \R^n : y_0 \geq 0\}$, and the Lebesgue
measure on it is a smooth measure with density $1$ in every chart.
\end{theorem}

\begin{proof}
\leanok
The change of charts from $(\sigma, \tau)$ to $(\sigma', \tau')$ is the restriction of the smooth
map $\sigma' \circ \tau$. The closure is closed and bounded. The image of the Lebesgue measure on
$\overline{\Omega} \cap B$ under the measure-preserving $\sigma$ is the Lebesgue measure on
$\sigma(\overline{\Omega} \cap B)$, the target of the chart.
\end{proof}

\section{Besov spaces of smooth domains by charts}

\begin{definition}[Global atlases]
\label{def:global-atlas}
\lean{Rd.SmoothDomain.GlobalAtlas, Rd.SmoothDomain.GlobalAtlas.nonempty,
  Rd.SmoothDomain.GlobalAtlas.toBesovAtlas, Rd.SmoothDomain.contMDiff_subtype_val,
  Rd.SmoothDomain.GlobalAtlas.hasCompactSupport_psi,
  Rd.SmoothDomain.GlobalAtlas.hasCompactSupport_theta, Rd.SmoothDomain.GlobalAtlas.psi_eq_zero,
  Rd.SmoothDomain.GlobalAtlas.theta_eq_zero, Rd.SmoothDomain.GlobalAtlas.theta_mul_psi}
\leanok
\uses{thm:smooth-domain-manifold, def:besov-atlas}
A global atlas of a smooth domain $\Omega$ consists of finitely many straighteners $(\sigma_i,
\tau_i, B_i)$ whose sets cover $\overline{\Omega}$, smooth functions $\psi_i \geq 0$ on $\R^n$ with
compact $\operatorname{tsupport} \psi_i \subseteq B_i$ and $\sum_i \psi_i = 1$ on
$\overline{\Omega}$, and smooth cut-offs $\theta_i = 1$ on $\operatorname{tsupport} \psi_i$ with
compact $\operatorname{tsupport} \theta_i \subseteq B_i$. Every smooth domain has one, and their
restrictions to $\overline{\Omega}$ form a Besov atlas, since the inclusion of $\overline{\Omega}$
into $\R^n$ is smooth.
\end{definition}

\begin{lemma}[Extension by zero and the pieces]
\label{lem:global-pieces}
\lean{Rd.SmoothDomain.extend, Rd.SmoothDomain.extend_coe, Rd.SmoothDomain.extend_eq_zero,
  Rd.SmoothDomain.stronglyMeasurable_extend, Rd.SmoothDomain.eLpNorm_extend,
  Rd.SmoothDomain.memLp_extend, Rd.SmoothDomain.extLp, Rd.SmoothDomain.extLp_ae,
  Rd.SmoothDomain.volume_coord_zero_eq_zero, Rd.SmoothDomain.Straightener.mem_image_iff,
  Rd.SmoothDomain.Straightener.homeomorph, Rd.SmoothDomain.Straightener.measurePreserving_tau,
  Rd.SmoothDomain.GlobalAtlas.piece_apply, Rd.SmoothDomain.GlobalAtlas.indicator_piece}
\leanok
\uses{def:global-atlas, def:chart-pieces}
The extension by $0$ is an isometry $L_p(\overline{\Omega}) \to L_p(\R^n)$. On the half-space
$\{y_0 \geq 0\}$, the $i$-th piece of $g$ is $(\psi_i\, \mathrm{ext}\, g) \circ \tau_i$. The
hyperplane $\{y_0 = 0\}$ is a null set, and $\tau_i$ preserves the Lebesgue measure.
\end{lemma}

\begin{proof}
\leanok
The image $\sigma_i(\overline{\Omega} \cap B_i)$ is $\{y : y_0 \geq 0,\ \tau_i(y) \in B_i\}$;
outside it on the half-space, $\psi_i \circ \tau_i = 0$.
\end{proof}

\begin{lemma}[The changes of variables between $\Omega$ and the charts]
\label{lem:global-changes}
\lean{Rd.CutoffDiffeo.global, Rd.CutoffDiffeo.global_apply, Rd.SmoothDomain.GlobalAtlas.toChart,
  Rd.SmoothDomain.GlobalAtlas.toChart_apply, Rd.SmoothDomain.GlobalAtlas.fromChart,
  Rd.SmoothDomain.GlobalAtlas.fromChart_apply, Rd.SmoothDomain.interior_univ_nonempty,
  Rd.modNormOn_congr}
\leanok
\uses{def:cutoff-diffeo, def:global-atlas, lem:modnorm-on}
A smooth bijection $\Phi$ of $V$ with smooth inverse and a smooth cut-off $\chi$ with compact
support form a change of variables with a cut-off on $V$, acting by $(\chi f) \circ \Phi$. In
particular $F \mapsto (\psi_i F) \circ \tau_i$ and $G \mapsto \theta_i (G \circ \sigma_i)$ are
bounded for the Besov quasi-norms of $\R^n$ (\Cref{thm:diffeo-invariance}). The intrinsic
quasi-norm on a set depends only on the restriction to the set.
\end{lemma}

\begin{proof}
\leanok
The charts of the atlas and their inverses are smooth bijections of $V$, and the functions of
the partition of unity are smooth with compact support, so each pair is a change of variables
with a cut-off; \Cref{thm:diffeo-invariance} gives the bounds. The intrinsic quasi-norm on a set
evaluates the function only at points of the set.
\end{proof}

\begin{theorem}[The Besov space of a smooth domain by charts]
\label{thm:smooth-domain-besov}
\lean{Rd.SmoothDomain.GlobalAtlas.modNormOn_pieceLp_lt_top,
  Rd.SmoothDomain.GlobalAtlas.modNormOn_extLp_lt_top, Rd.SmoothDomain.memBesov_iff_modNormOn_lt_top}
\leanok
\uses{lem:global-pieces, lem:global-changes, def:manifold-besov, thm:bounded-lipschitz-besov,
  thm:half-space-besov-domain}
For a smooth bounded domain $\Omega \subseteq \R^n$, $1 \leq p < \infty$, $0 < s < r$ and
$1 \leq q \leq \infty$, $f \in L_p(\overline{\Omega})$ lies in the Besov space
$\Besov{s}{p}{q}(\overline{\Omega})$ of the compact manifold $\overline{\Omega}$ with the Lebesgue
measure if and only if
$\|f\|_{L_p(\Omega)} + \|(2^{ks} \omega_r(f, 2^{-k}, \Omega)_p)_k\|_{\ell_q} < \infty$.
\end{theorem}

\begin{proof}
\leanok
If the intrinsic quasi-norm is finite, $f$ agrees on $\Omega$ with some $F \in \Besov{s}{p}{q}(\R^n)$
(\Cref{thm:bounded-lipschitz-besov}). On the half-space the piece $(\psi_i f) \circ \tau_i$ agrees
almost everywhere with $(\psi_i F) \circ \tau_i$: where $y_0 > 0$ and $\tau_i(y) \in B_i$, the
point $\tau_i(y)$ lies in $\Omega$, and $\tau_i$ preserves null sets. The latter function is in
$\Besov{s}{p}{q}(\R^n)$ by \Cref{lem:global-changes}. Conversely, the pieces extend from the
half-space (\Cref{thm:half-space-besov-domain}) to $G_i \in \Besov{s}{p}{q}(\R^n)$, and
$F = \sum_i \theta_i (G_i \circ \sigma_i) \in \Besov{s}{p}{q}(\R^n)$ agrees with $f$ on $\Omega$,
since $\theta_i \psi_i = \psi_i$ and $\sum_i \psi_i = 1$ on $\overline{\Omega}$.
\end{proof}

\section{Besov spaces on the boundary}

In this section $M$ is a compact smooth manifold with boundary modelled on the half-space
$\{y \in \R^n : y_0 \geq 0\}$. Its boundary hyperplane is $H_0 = e_0^\perp = \{y : y_0 = 0\}$, with
the orthogonal projection $P$ onto $H_0$ and the Lebesgue measure of $H_0$.

\begin{lemma}[Changes of charts preserve the boundary]
\label{lem:coord-change-boundary}
\lean{Manifold.flatDomain, Manifold.bdry, Manifold.mem_bdry_iff, Manifold.bdry_subset_range,
  Manifold.coe_proj_of_mem_bdry, Manifold.isClosed_bdry, Manifold.mem_interior_extend_target_iff,
  Manifold.coordChange_mem_bdry, Manifold.mem_coordChange_source_iff_of_bdry,
  Manifold.mem_coordChange_target_iff_of_bdry, Manifold.coe_proj_coordChange,
  Manifold.coe_proj_coordChange_symm}
\leanok
\uses{lem:coord-change}
$H_0$ is a closed subspace contained in the half-space. A point $y$ of the target of an extended
chart lies in the interior of the target if and only if $y_0 > 0$. For charts $e, e'$ of the
atlas, the change of charts maps every point $y$ of its source with $y_0 = 0$ to a point of $H_0$,
and so does its inverse.
\end{lemma}

\begin{proof}
\leanok
A point of $M$ is mapped into the interior of the half-space by one chart if and only if by every
chart (Mathlib's \texttt{ModelWithCorners.\allowbreak mem\_interior\_range\_iff\_of\_mem\_atlas}).
The change of charts maps into the half-space, so a point with $y_0 = 0$ goes to a point whose
first coordinate is $\geq 0$ and not $> 0$.
\end{proof}

\begin{definition}[Changes of charts on the boundary]
\label{def:boundary-chart-change}
\lean{Manifold.boundaryChartChange, Manifold.chartChange_apply_coe,
  Manifold.BesovAtlas.boundaryTransition}
\leanok
\uses{lem:coord-change-boundary, def:chart-transition, def:cutoff-diffeo}
Let $T$ be the change of charts with a cut-off of \Cref{def:chart-transition}, with $\Phi$, $\chi$
and $K_0$. Its restriction $T^\partial$ to $H_0$, with the change of charts $P \circ \Phi$ on
$H_0 \cap U$, its inverse $P \circ \Phi^{-1}$ on $H_0 \cap U'$, the cut-off $\chi|_{H_0}$ and the
compact set $K_0 \cap H_0$, is a smooth change of variables with a cut-off on the whole space
$H_0$, and $(T G)|_{H_0} = T^\partial(G|_{H_0})$ for every $G : \R^n \to \C$. For two Besov
atlases $A, A'$ the restriction of $T_{ij}$ is $T^\partial_{ij}$.
\end{definition}

\begin{definition}[Besov spaces on the boundary]
\label{def:manifold-boundary-besov}
\lean{Manifold.BesovAtlas.boundaryPiece, Manifold.BesovAtlas.MemBoundaryBesov}
\leanok
\uses{def:chart-pieces}
For a Besov atlas $A$, the $i$-th boundary piece of $g : M \to \C$ is the restriction
$\bigl((\psi_i g) \circ \kappa_i^{-1}\bigr)|_{H_0}$ of its $i$-th piece. For $1 \leq p < \infty$,
$s > 0$ and $1 \leq q \leq \infty$, $g$ lies in $\Besov{s}{p}{q}(\partial M)$ if every boundary
piece agrees almost everywhere on $H_0$ with some $G_i \in L_p(H_0)$ with
$\|G_i\|_p + \|(2^{ks} \omega_r(G_i, 2^{-k})_p)_k\|_{\ell_q} < \infty$, $r = \lfloor s \rfloor + 1$.
\end{definition}

\begin{theorem}[Independence of the atlas on the boundary]
\label{thm:boundary-atlas-independence}
\lean{Manifold.BesovAtlas.boundaryPiece_eq_sum, Manifold.BesovAtlas.interior_univ_bdry_nonempty,
  Manifold.BesovAtlas.MemBoundaryBesov.of_atlas,
  Manifold.BesovAtlas.memBoundaryBesov_iff_memBoundaryBesov}
\leanok
\uses{def:manifold-boundary-besov, def:boundary-chart-change, thm:pieces-two-atlases,
  thm:diffeo-invariance, lem:modnorm-on}
For two Besov atlases $A, A'$ and $g : M \to \C$, the $i$-th boundary piece of $g$ for $A$ is
$\sum_j T^\partial_{ij}$ applied to the $j$-th boundary piece for $A'$, at every point of $H_0$.
For $1 \leq p < \infty$, $s > 0$ and $1 \leq q \leq \infty$, $g \in \Besov{s}{p}{q}(\partial M)$
for $A$ if and only if for $A'$.
\end{theorem}

\begin{proof}
\leanok
The identity is \Cref{thm:pieces-two-atlases} restricted to $H_0$
(\Cref{def:boundary-chart-change}). The changes of charts $T^\partial_{ij}$ live on the whole
space $H_0$, so \Cref{thm:diffeo-invariance} bounds
$N(T^\partial_{ij} G'_j) \leq C_{ij} N(G'_j)$ for the Besov quasi-norm $N$ of $H_0$, and $N$ is
subadditive (\Cref{lem:modnorm-on}).
\end{proof}

\section{The trace theorem on compact manifolds}

In this section $M$ is a compact manifold with boundary as in the previous section, $\mu$ a smooth
measure, $1 \leq p < \infty$, and $\operatorname{Tr}$ the trace of order $L$ of the open
half-space $\Omega_0 = \{y : y_0 > 0\}$ (\Cref{thm:trace-every-order}).

\begin{lemma}[The closed and the open half-space]
\label{lem:closed-open-half-space}
\lean{Manifold.range_eq_closedHalfSpace, Manifold.range_ae_eq_domain,
  Manifold.cutoff_eq_indicatorLp, Manifold.modNormOn_range_eq}
\leanok
\uses{lem:null-set-domains}
The range of the model is the closed half-space $\{y : y_0 \geq 0\}$, which differs from
$\Omega_0$ by the null set $H_0$. So the cut-offs and the intrinsic quasi-norms on the two sets
agree.
\end{lemma}

\begin{proof}
\leanok
The range of the model is the closed half-space by its definition. The cut-offs and the
intrinsic quasi-norms are built from integrals over the set and over its translates, and these
change only by null sets when $\Omega_0$ is replaced by its closure
(\Cref{lem:null-set-domains}).
\end{proof}

\begin{theorem}[The trace commutes with the changes of charts]
\label{thm:trace-chart-change}
\lean{Manifold.trace_chartChange_applyLp_of_continuous,
  Manifold.trace_chartChange_applyLp_of_modNorm, Manifold.trace_chartChange_applyLp}
\leanok
\uses{def:boundary-chart-change, lem:trace-continuous-half-space, lem:partial-sums-approx,
  lem:trace-linear, thm:diffeo-invariance, lem:closed-open-half-space, def:besov-domain}
Let $T$ be a change of charts with a cut-off and $T^\partial$ its restriction to $H_0$
(\Cref{def:boundary-chart-change}), $s > 0$, $L > s + 1/p$, and let $G \in L_p(\R^n)$ satisfy
$\|G\|_{L_p(\Omega)} + \|(2^{k(s+1/p)} \omega_L(G, 2^{-k}, \Omega)_p)_k\|_{\ell_q} < \infty$ on the
closed half-space $\Omega$. Then $\operatorname{Tr}(T G) = T^\partial(\operatorname{Tr} G)$.
\end{theorem}

\begin{proof}
\leanok
Extend $G$ to $\R^n$, since the half-space is a Besov domain; neither side changes. For the
continuous partial sums $G_J$ of the reproducing formula, $T G_J$ is continuous on the closed
half-space and vanishes outside a compact set (\Cref{def:cutoff-diffeo}), so both sides are the
restriction of $T G_J$ to $H_0$ (\Cref{lem:trace-continuous-half-space}). As $J \to \infty$,
$G_J \to G$ in $\Besov{1/p}{p}{1}$ (\Cref{lem:partial-sums-approx}); $T$ maps this space boundedly
into the intrinsic space on $\Omega$ (\Cref{thm:diffeo-invariance}), where the trace is bounded
(\Cref{lem:trace-linear}), so $\operatorname{Tr}(T G_J) \to \operatorname{Tr}(T G)$; and
$\operatorname{Tr} G_J \to \operatorname{Tr} G$ in $L_p(H_0)$, where $T^\partial$ is bounded.
\end{proof}

\begin{definition}[The trace on the boundary]
\label{def:manifold-boundary-trace}
\lean{Manifold.BesovAtlas.boundaryTrace, Manifold.BesovAtlas.boundaryPiece_boundaryTrace_eq_sum}
\leanok
\uses{def:manifold-boundary-besov, def:boundary-chart-change}
For a Besov atlas $A$, an order $L$ and $f \in L_p(\mu)$, the trace on the boundary is
$\operatorname{Tr}_A f = \sum_i \theta_i \cdot \operatorname{Tr}(\text{piece}_i f) \circ P \circ
\kappa_i$ on the chart sources. At every point of $H_0$ its $k$-th boundary piece is
$\sum_i T^\partial_{ki} \bigl(\operatorname{Tr}(\text{piece}_i f)\bigr)$, since $\psi_k \theta_i$
is the cut-off of $T_{ki}$.
\end{definition}

\begin{theorem}[The boundary pieces of the trace]
\label{thm:boundary-trace-pieces}
\lean{Manifold.BesovAtlas.trace_pieceLp_eq_sum, Manifold.BesovAtlas.boundaryPiece_boundaryTrace,
  Manifold.BesovAtlas.boundaryPiece_boundaryTrace_of_atlas}
\leanok
\uses{def:manifold-boundary-trace, thm:trace-chart-change, thm:pieces-two-atlases, lem:trace-sums,
  thm:boundary-atlas-independence}
Let $s > 0$, $L > s + 1/p$ and $f \in \Besov{s+1/p}{p}{q}(M)$. For two Besov atlases $A, A'$,
$\operatorname{Tr}(\text{piece}_i f) = \sum_j T^\partial_{ij}
\bigl(\operatorname{Tr}(\text{piece}'_j f)\bigr)$, and the $i$-th boundary piece for $A$ of
$\operatorname{Tr}_{A'} f$ is $\operatorname{Tr}(\text{piece}_i f)$ almost everywhere on $H_0$. In
particular the trace does not depend on the atlas.
\end{theorem}

\begin{proof}
\leanok
Take the traces in $\text{piece}_i f = \sum_j T_{ij}(\text{piece}'_j f)$
(\Cref{thm:pieces-two-atlases}): the trace of the finite sum is the sum of the traces
(\Cref{lem:trace-sums}), and it commutes with the $T_{ij}$ (\Cref{thm:trace-chart-change}). For
$A' = A$ this is the pointwise formula of \Cref{def:manifold-boundary-trace}, almost everywhere; for
general $A'$ combine it with \Cref{thm:boundary-atlas-independence}.
\end{proof}

\begin{theorem}[The trace theorem on compact manifolds]
\label{thm:manifold-trace}
\lean{Manifold.exists_flatExtension, Manifold.BesovAtlas.memBoundaryBesov_boundaryTrace,
  Manifold.BesovAtlas.exists_boundaryTrace_eq,
  Manifold.BesovAtlas.memBoundaryBesov_iff_exists_boundaryTrace}
\leanok
\uses{thm:boundary-trace-pieces, thm:trace-every-order, thm:half-space-trace-space, lem:chart-lift}
Let $1 \leq p < \infty$, $s > 0$, $1 \leq q \leq \infty$ and $L \geq 2$ with $s + 1/p < L$. For
$f \in \Besov{s+1/p}{p}{q}(M)$, $\operatorname{Tr}_A f \in \Besov{s}{p}{q}(\partial M)$. Every
$g \in \Besov{s}{p}{q}(\partial M)$ is the trace of some $f \in \Besov{s+1/p}{p}{q}(M)$: the
boundary pieces of $\operatorname{Tr}_A f$ and of $g$ agree almost everywhere. So
$\Besov{s}{p}{q}(\partial M)$ is the trace space of $\Besov{s+1/p}{p}{q}(M)$ for every $s > 0$.
\end{theorem}

\begin{proof}
\leanok
The boundary pieces of $\operatorname{Tr}_A f$ are the traces of the pieces
(\Cref{thm:boundary-trace-pieces}), which lie in $\Besov{s}{p}{q}(H_0)$
(\Cref{thm:trace-every-order}). Conversely, the extensions $E G_k$ of moment order
$\lfloor s \rfloor + 1$ of the boundary pieces $G_k$ of $g$ have finite norms and traces $G_k$
(\Cref{thm:half-space-trace-space}). The pieces of
$f = \sum_k \theta_k \cdot (E G_k) \circ \kappa_k \in \Besov{s+1/p}{p}{q}(M)$
(\Cref{lem:chart-lift}) are $\sum_k T_{ik}(E G_k)$, with traces $\sum_k T^\partial_{ik} G_k$
(\Cref{thm:trace-chart-change}), the $i$-th boundary piece of $g$
(\Cref{thm:boundary-atlas-independence}).
\end{proof}

The closure of a smooth bounded domain is a compact manifold with boundary
(\Cref{thm:smooth-domain-manifold}), so this is the trace theorem on its boundary for every
$s > 0$, in charts.

\section{Smooth bounded domains from local graphs}
\label{sec:smooth-local-graphs}

The smooth bounded domains of \Cref{def:smooth-domain} carry their balls and graph functions as
data. This section starts from the usual definition by smooth local graphs and shows that it
describes exactly their carriers; the open balls are examples. In this section $V$ is a real
inner product space, and for a unit vector $\nu \in V$, $P_\nu$ is the orthogonal projection onto
the hyperplane $\nu^\perp$, as in \Cref{sec:local-lipschitz-graphs}.

\begin{definition}[Sets that are smooth near a point]
\label{def:smooth-near-point}
\lean{Rd.IsSmoothDomainAt, Rd.IsBoundedSmoothDomain}
\leanok
\uses{def:lipschitz-near-point}
Let $\Omega \subseteq V$ and $z \in V$. The set $\Omega$ is smooth near $z$ if there are $r > 0$,
a unit vector $\nu \in V$ and a function $\varphi : \nu^\perp \to \R$ that is smooth on
$\{y \in \nu^\perp : \|y\| < r\}$ with
\[
  \Omega \cap B(z, r)
  = \{x \in V : \varphi(P_\nu(x - z)) < \langle \nu, x - z \rangle\} \cap B(z, r) .
\]
The set $\Omega$ is a bounded smooth set if it is open and bounded and smooth near every
$z \in \partial\Omega$.
\end{definition}

This is the local condition of \cite[Definition 1.2.1.1]{Grisvard} with $\varphi$ of class
$C^\infty$, written as in \Cref{def:lipschitz-near-point}. The function $\varphi$ is asked to be
smooth only on the part of the hyperplane that the ball sees; a smooth cut-off turns it into a
smooth Lipschitz function on the whole hyperplane.

\begin{lemma}[The smooth cut-off]
\label{lem:smooth-cutoff}
\lean{Rd.contDiff_mul_of_tsupport_subset, Rd.exists_contDiff_lipschitzWith_eqOn_ball}
\leanok
Let $E$ be a real normed space. For every $n \in \N \cup \{\infty, \omega\}$, every
$b : E \to \R$ of class $C^n$, every open $U \subseteq E$ with
$\operatorname{tsupport} b \subseteq U$ and every $\varphi : E \to \R$ of class $C^n$ on $U$, the
product $b \varphi$ is of class $C^n$ on $E$. If $E$ is a finite-dimensional real inner product
space, then for every $r > 0$ and every $\varphi : E \to \R$ that is smooth on $B(0, r)$ there are
$L \geq 0$ and a smooth $L$-Lipschitz $\psi : E \to \R$ with $\psi = \varphi$ on $B(0, r/4)$.
\end{lemma}

\begin{proof}
\leanok
Near a point of $\operatorname{tsupport} b$ both factors are of class $C^n$; near any other point
the product vanishes. This is the flat analogue of Mathlib's
\texttt{ContMDiffOn.\allowbreak smul\_\allowbreak section\_\allowbreak of\_\allowbreak tsupport}.
For the second statement take a smooth bump $b$ equal to $1$ on
$\bar B(0, r/4)$ with $\operatorname{tsupport} b = \bar B(0, r/2) \subseteq B(0, r)$ (Mathlib's
\texttt{ContDiffBump}) and $\psi = b \varphi$. It is smooth with compact support, hence Lipschitz
(Mathlib's \texttt{ContDiff.\allowbreak lipschitzWith\_\allowbreak of\_\allowbreak
hasCompactSupport}).
\end{proof}

\begin{theorem}[Smooth near a point by smooth graph domains]
\label{thm:smooth-near-point}
\lean{Rd.IsSmoothDomainAt.exists_ball_inter_eq, Rd.isSmoothDomainAt_of_ball_inter_eq,
  Rd.isSmoothDomainAt_iff, Rd.IsSmoothDomainAt.isLipschitzDomainAt,
  Rd.IsBoundedSmoothDomain.isBoundedLipschitzDomain}
\leanok
\uses{def:smooth-near-point, lem:smooth-cutoff, lem:local-graph-domain, thm:lipschitz-near-point}
For every $\Omega \subseteq V$, every $z \in V$, every $r > 0$ and every Lipschitz graph domain
$\Omega'$ with smooth graph function and $B(z, r) \cap \Omega = B(z, r) \cap \Omega'$, the set
$\Omega$ is smooth near $z$. For every $\Omega \subseteq V$ and every $z \in V$, if $\Omega$ is
smooth near $z$, then $\Omega$ is Lipschitz near $z$; every bounded smooth set is a bounded
Lipschitz set. Let $V$ be finite-dimensional. For every $\Omega \subseteq V$ and every $z \in V$,
$\Omega$ is smooth near $z$ if and only if there are $r > 0$ and a Lipschitz graph domain
$\Omega'$ with smooth graph function and $B(z, r) \cap \Omega = B(z, r) \cap \Omega'$.
\end{theorem}

\begin{proof}
\leanok
For the first statement, the graph function seen from $z$ (\Cref{lem:local-graph-domain}) is
smooth on the whole hyperplane. For the second, let $\varphi$ be smooth on the ball of radius $r$
of $\nu^\perp$. It is of class $C^1$ at $0$, hence $K$-Lipschitz on a ball of some radius
$\varepsilon > 0$ (Mathlib's \texttt{ContDiffAt.\allowbreak exists\_\allowbreak lipschitzOnWith}),
and the last statement of \Cref{thm:lipschitz-near-point} applies with the radius
$\min(\varepsilon, r)$; the statement on bounded sets follows pointwise. For the equivalence,
\Cref{lem:smooth-cutoff} gives a smooth Lipschitz $\psi$ equal to $\varphi$ on the ball of radius
$r/4$, so $\Omega$ agrees on $B(z, r/4)$ with the region above the graph of $\psi$
(\Cref{thm:lipschitz-near-point}), the Lipschitz graph domain of the local graph $\psi$
(\Cref{lem:local-graph-domain}), whose graph function is smooth by the same lemma.
\end{proof}

\begin{theorem}[Bounded smooth sets and smooth bounded domains]
\label{thm:bounded-smooth-sets}
\lean{Rd.IsBoundedSmoothDomain.exists_chart, Rd.IsBoundedSmoothDomain.chart,
  Rd.IsBoundedSmoothDomain.chart_fst_pos, Rd.IsBoundedSmoothDomain.contDiff_chart_phi,
  Rd.IsBoundedSmoothDomain.ball_inter_chart, Rd.SmoothDomain.ofIsBoundedSmoothDomain,
  Rd.SmoothDomain.carrier_ofIsBoundedSmoothDomain, Rd.SmoothDomain.isBoundedSmoothDomain,
  Rd.isBoundedSmoothDomain_iff_exists}
\leanok
\uses{def:smooth-domain, thm:smooth-near-point, thm:bounded-lipschitz-sets}
For every $n \geq 1$ and every $\Omega \subseteq \R^n$, $\Omega$ is a bounded smooth set if and
only if it is the carrier of a smooth bounded domain.
\end{theorem}

\begin{proof}
\leanok
Let $\Omega$ be a bounded smooth set. \Cref{thm:smooth-near-point} gives for each
$z \in \partial\Omega$ a ball $B(z, r_z)$ and a Lipschitz graph domain with smooth graph function
that agrees with $\Omega$ on it. By \Cref{thm:bounded-lipschitz-sets} these charts give a bounded
Lipschitz domain with carrier $\Omega$ whose graph domains are among them, so its graph functions
are smooth. Conversely, near each boundary point a smooth bounded domain agrees with one of its
graph domains (\Cref{thm:bounded-lipschitz-sets}), whose graph function is smooth;
\Cref{thm:smooth-near-point} applies.
\end{proof}

\begin{theorem}[The ball as a smooth domain]
\label{thm:smooth-ball}
\lean{Rd.contDiffOn_sphereProfile_norm, Rd.isSmoothDomainAt_ball, Rd.isBoundedSmoothDomain_ball,
  Rd.SmoothDomain.ball, Rd.SmoothDomain.carrier_ball}
\leanok
\uses{def:sphere-profile, thm:ball-lipschitz, def:smooth-near-point, thm:bounded-smooth-sets}
For every real inner product space $E$ and every $R > 0$, the function $y \mapsto g_R(\|y\|)$ is
smooth on $B(0, R/2) \subseteq E$. Let $c \in V$.
\begin{enumerate}
  \item For every $R > 0$ and every $z \in V$ with $\|z - c\| = R$, the ball $B(c, R)$ is smooth
    near $z$.
  \item For every $R \in \R$, the ball $B(c, R)$ is a bounded smooth set; if $V = \R^n$ with
    $n \geq 1$, it is the carrier of a smooth bounded domain.
\end{enumerate}
\end{theorem}

\begin{proof}
\leanok
For $\|y\| < R/2$, $g_R(\|y\|) = R - \sqrt{R^2 - \|y\|^2}$ with $R^2 - \|y\|^2 > 0$; the function
$y \mapsto \|y\|^2$ is smooth (Mathlib's \texttt{contDiff\_\allowbreak norm\_\allowbreak sq}), and
so is the square root away from $0$. For $z$ on the sphere take radius $R/2$ and
$\nu = (c - z)/R$; \Cref{thm:ball-lipschitz} gives the set identity of
\Cref{def:smooth-near-point} with $\varphi(y) = g_R(\|y\|)$. For $R > 0$ the boundary of the ball
lies in the sphere $\|z - c\| = R$; for $R \leq 0$ the ball is empty. The last statement is
\Cref{thm:bounded-smooth-sets}.
\end{proof}

So the closure of every open ball of $\R^n$ is a compact manifold with boundary
(\Cref{thm:smooth-domain-manifold}), and the results of this chapter apply to it; for example its
Besov spaces by charts are the intrinsic ones.

\begin{corollary}[The Besov spaces of the closed ball]
\label{cor:ball-besov}
\lean{Rd.SmoothDomain.memBesov_ball_iff}
\leanok
\uses{thm:smooth-ball, thm:smooth-domain-besov}
For every $n \geq 1$, every $c \in \R^n$, every $R \in \R$, every $1 \leq p < \infty$, all
$0 < s < r$, every $1 \leq q \leq \infty$ and every $f \in L_p(M)$, where $M$ is the closure of
$B(c, R)$, $f$ lies in the Besov space $\Besov{s}{p}{q}(M)$ of the compact manifold $M$ with the
Lebesgue measure if and only if
$\|f\|_{L_p(B(c, R))} + \|(2^{ks} \omega_r(f, 2^{-k}, B(c, R))_p)_k\|_{\ell_q} < \infty$.
\end{corollary}

\begin{proof}
\leanok
\Cref{thm:smooth-domain-besov} for the smooth bounded domain of \Cref{thm:smooth-ball}.
\end{proof}
